\documentclass[11pt]{article}

\usepackage[margin=1.05in]{geometry}
\usepackage{amsmath,amssymb,amsthm,mathtools,amscd}
\usepackage{booktabs,array,longtable}
\usepackage{enumitem}
\usepackage{xcolor}
\usepackage{hyperref}
\IfFileExists{microtype.sty}{\usepackage{microtype}}{}

\hypersetup{colorlinks=true,linkcolor=blue!55!black,citecolor=blue!55!black,urlcolor=blue!55!black}
\setlist{nosep}
\allowdisplaybreaks
\makeatletter
\renewcommand*\l@subsection{\@dottedtocline{2}{1.5em}{3.0em}}
\makeatother

\newtheorem{theorem}{Theorem}[section]
\newtheorem{proposition}[theorem]{Proposition}
\newtheorem{lemma}[theorem]{Lemma}
\newtheorem{corollary}[theorem]{Corollary}
\newtheorem{conjecture}[theorem]{Conjecture}
\theoremstyle{definition}
\newtheorem{definition}[theorem]{Definition}
\newtheorem{example}[theorem]{Example}
\newtheorem{remark}[theorem]{Remark}
\newtheorem{warning}[theorem]{Warning}

\newcommand{\Qp}{\mathbf Q_p}
\newcommand{\Zp}{\mathbf Z_p}
\newcommand{\OK}{\mathcal O_K}
\newcommand{\OL}{\mathcal O_L}
\newcommand{\mK}{\mathfrak m_K}
\newcommand{\GK}{G_K}
\newcommand{\GL}{G_L}
\newcommand{\Rep}{\operatorname{Rep}}
\newcommand{\End}{\operatorname{End}}
\newcommand{\Hom}{\operatorname{Hom}}
\newcommand{\Ext}{\operatorname{Ext}}
\newcommand{\ad}{\operatorname{ad}}
\newcommand{\Cont}{\operatorname{Cont}}
\newcommand{\Sp}{\operatorname{Sp}}
\newcommand{\MC}{\operatorname{MC}}
\newcommand{\Tot}{\operatorname{Tot}}
\newcommand{\Fib}{\operatorname{Fib}}
\newcommand{\Gal}{\operatorname{Gal}}
\newcommand{\inv}{\operatorname{inv}}
\newcommand{\Fil}{\operatorname{Fil}}
\newcommand{\cris}{\mathrm{cris}}
\newcommand{\st}{\mathrm{st}}
\newcommand{\dR}{\mathrm{dR}}
\newcommand{\rig}{\mathrm{rig}}
\newcommand{\et}{\acute{\mathrm et}}
\newcommand{\cC}{\mathcal C}
\newcommand{\cJ}{\mathcal J}
\newcommand{\cR}{\mathcal R}
\newcommand{\cD}{\mathcal D}
\newcommand{\cA}{\mathcal A}
\newcommand{\cF}{\mathcal F}
\newcommand{\cP}{\mathcal P}
\newcommand{\cI}{\mathcal I}
\newcommand{\bR}{\mathbf R}
\newcommand{\arith}{\mathrm{arith}}
\newcommand{\bracket}[2]{\left[#1,#2\right]}
\newcommand{\Gm}{\mathbf G_m}
\newcommand{\Perf}{\operatorname{Perf}}

\title{\textbf{Profinite Arithmetic Cartan Geometry of $p$-adic Representations}\\[3pt]
\large Finite Covariant Atlases from Fox, Herr, and Cayley--Hamilton Carriers}
\author{}
\date{October 2026}

\begin{document}
\maketitle

\begin{abstract}
We study continuous $p$-adic Galois representations as profinite transport systems and organize their cohomology, deformation theory, and full transport in finite covariant arithmetic charts.  Continuous group cochains provide the intrinsic Spencer complex, Maurer--Cartan theory controls deformations, Herr complexes give finite cyclotomic cohomological charts, completed Fox calculus gives explicit transport resolutions in free and Demu\v{s}kin pro-$p$ sectors, and Cayley--Hamilton algebras retain full transport and nonsplit extension data.

The main finite results are tower-free.  Finite presentability of $G_K$ and $\operatorname{cd}_p(G_K)=2$ give bounded perfect coefficient carriers compatible with derived base change.  In the odd-prime Demu\v{s}kin sector, the quadratic Fox symbol recovers the cup form and the integral Fox--Bockstein data recover the power invariant.  At $p=2$, a separate orientation-enriched Fox--Bockstein certificate treats every dyadic Demu\v{s}kin normal form, including $q=2$ branches, with an explicit $G_{\mathbf Q_2}(2)$ anchor.  Rational cyclotomic Herr comparison also extends to $2$; integral cyclotomic torsion is instead retained by derived invariants.  The dyadic Fox--Wu identity identifies the mod-four cyclotomic class with the characteristic cup-square vector; a single twisted Fox row gives the exact lifting-depth obstruction for every unit-valued rank-one character and, at the canonical orientation, recovers integral Kummer--Tate cohomology and transfer-normalized Tate duality.  The fully marked Roe--Turturean reciprocity coordinates now give an exact, unsheared identification of the $G_{\mathbf Q_2}(2)$ Fox basis with the Kummer basis $(-1,5,2)$.  On every residual $2$-group-image sector, the full marked two-relation Jacobian splits integrally into an acyclic tame block and the Demu\v{s}kin Fox complex via an explicit Nielsen chain isomorphism.  Beyond pro-$2$ residual images, the tame $S_3\simeq\mathrm{GL}_2(\mathbf F_2)$ quotient yields an explicit nonnormal index-three Sylow--Schreier table, an integral Hecke projector with quadratic polynomial, and a complete coefficient-specific integral $8\times16$ adjoint marked-Jacobian Smith certificate.  For the distinct norm-zero natural lattice $T=\mathbf Z_2^2$ in this genuine non-pro-$2$ sector, the full marked coefficient complex contracts by explicit integral matrices to $T[-1]$, and an exact finite $16\times16$ Schreier-coordinate projector supplies a strict integral bar-cochain representative of the index-three Hecke/Sylow summand.  The full adjoint lattice now splits unimodularly via the tame $C_3$ projector as $\operatorname{ad}T\simeq\operatorname{Ind}_{G_{\mathbf Q_2(\zeta_3)}}^{G_{\mathbf Q_2}}\mathbf Z_2\oplus T$; its marked Jacobian admits a literal integral Fox diagonalization proving the seven-unit/one-two Smith certificate, and Shapiro identifies exactly which top torsion class the Sylow projector retains.  At the level of the entire framed $S_3$ residual deformation family, the marked tame relation forces $\rho(\tau)^3=I$ and yields a finite polynomial rank-two tame Reynolds projector; it becomes Galois-equivariant exactly on the closed wild-centralizer locus.  This locus has tangent codimension two, and a genuine dual-number Galois lift proves that no equivariant projector extending the fixed splitting exists on the unrestricted family.  Away from the split locus, exact off-diagonal block transport retains the missing wild coherence without asserting a false direct-sum decomposition.  For the trivial residual character this refines the integral rank-one deformation hypersurface $\mathcal O_E[[u,v,w]]/(u(u+2))$ to a complete full-group two-relation ideal.  Profinite idempotent powers in the marked full-$G_{\mathbf Q_2}$ presentation admit convergent residue-sector binomial matrix formulas and terminating Jacobian computations at each finite order.  Residual Sylow descent, local complete-intersection/Koszul models, and Cayley--Hamilton transport algebraization then assemble pointwise and sectorwise finite carriers, while wildness rules out a universal Jordan-type normal form.

The resulting derived faithful finite closure is unconditional: for fixed local field, coefficient field, and rank, finitely many residual sectors carry canonical two-layer derived charts consisting of a Cayley--Hamilton full-transport carrier and a perfect coefficient-cochain carrier, with Morita-valued descent.  The residual, Fitting, flag, and descent-exponent constructions upgrade these charts to a stratified recursively explicit native atlas: finite-stage descent, native presentation transitions, and every prescribed coherence level reduce to finite matrix operations and terminating chartwise procedures.  Raw Herr comparison is multiplicative at the intrinsic level and its direct Tate-dual linear representative is canonically normalized modulo continuous phantom maps; literal equality on one unreduced Robba representative is only an optional strict gauge.  Wildness therefore limits normal forms, not finite native globalization.  Period and prismatic descriptions are used only as comparison charts and are not part of the foundational construction.
\end{abstract}

\tableofcontents

\section{Introduction}

Let $K$ be a finite extension of $\Qp$, let $\overline K$ be an algebraic closure, and write
\[
\GK=\Gal(\overline K/K).
\]
A continuous $p$-adic representation is a finite-dimensional $E$-vector space $V$, for a finite extension $E/\Qp$, equipped with a continuous homomorphism
\[
\rho:\GK\longrightarrow \mathrm{GL}_E(V).
\]
Period functors isolate especially regular regions of this category:
\[
\Rep_{\cris}(\GK)\subset \Rep_{\st}(\GK)\subset \Rep_{\dR}(\GK)
\subset \Rep_{\mathrm{HT}}(\GK)\subset \Rep_E(\GK).
\]
Crystalline, semistable, de Rham, and Hodge--Tate representations admit finite-dimensional linear-algebraic invariants, and prismatic $F$-crystals give an integral classification of crystalline lattices \cite{BhattScholzePFC}.  Arbitrary continuous representations are also categorically encoded: Fontaine's local theory, the overconvergence theorem of Cherbonnier--Colmez, and the Herr complex identify them with etale $(\varphi,\Gamma)$-modules and compute their Galois cohomology \cite{FontaineLocal,CherbonnierColmez,HerrSurvey,LiuCohomology,BergerTrianguline,ColmezTrianguline}.

The gap addressed here is therefore not the absence of an equivalence of categories.  Standard theory does not by itself package the full representation, its intrinsic cochains, deformation structure, and descent data into one finite faithful geometric organization comparable in spirit to a finite atlas.  This paper first constructs the derived/Morita closure and then proves a stratified recursively explicit integral native globalization.  The only stronger questions left after that theorem concern optional choices of especially rigid closed-form representatives, not existence or globalization of the atlas.

The construction proved below is organized by the chain
\[
\boxed{
\begin{gathered}
\text{profinite arithmetic transport}\\
\longrightarrow \text{intrinsic Cartan--Spencer/cochain structure}\\
\longrightarrow \text{faithful finite covariant atlas}.
\end{gathered}}
\]
A \emph{finite covariant atlas}, or simply a \emph{finite atlas}, is a compatible system of finite projective arithmetic transport/cochain charts equipped with comparison, reconstruction, higher-coherence, and descent data.  The continuous bar complex remains the intrinsic Spencer complex, the deformation equation remains Maurer--Cartan, and Herr, Fox, Koszul, and Cayley--Hamilton constructions supply different finite gauges of the same representation-theoretic object.  Period and prismatic descriptions are comparison structures only and are not used to define the atlas.

\paragraph{Terminology.}
A \emph{carrier} retains specified transport, cohomological, or deformation data up to the stated equivalence; a \emph{chart} is an explicit presentation on a sector; an \emph{atlas} is a compatible family of such charts; and a \emph{model} is a chosen chain-level representative.

\subsection{Main results}

The main results have four layers.

\begin{enumerate}[label=\textbf{(\arabic*)}]
\item \textbf{Intrinsic transport and deformation.}
A representation is a flat multiplicative transport law.  Continuous Galois cochains form the intrinsic arithmetic Spencer complex, and the cup-commutator dg-Lie algebra gives the Maurer--Cartan deformation theory.  Rank-one, trianguline, and Herr calculations are retained only as low-rank consistency and reconstruction tests.

\item \textbf{Finite Fox--Spencer carriers.}
Finite profinite presentations give completed Fox complexes, with genuine finite resolutions in free and Demu\v{s}kin pro-$p$ sectors.  The quadratic Fox symbol recovers the cup form, while the integral linear Fox symbol and Bockstein lifting depth recover the Demu\v{s}kin power invariant in the odd-prime sector.  At $p=2$, a nonalternating cup square and the canonical dualizing orientation provide the extra dyadic data; Corollary~\ref{thm:dyadic-classification-certificate} gives the classification-faithful finite certificate.  Theorem~\ref{thm:dyadic-fox-wu-bockstein} identifies the mod-$4$ orientation with the cup-square/Bockstein characteristic class; Theorem~\ref{thm:dyadic-twisted-fox-lifting-depth} proves an exact all-level twisted-lifting criterion, and Theorem~\ref{thm:dyadic-fox-kummer-tate} reads out the entire integral local Kummer--Tate groups.  Theorem~\ref{thm:dyadic-fox-hilbert-marked} identifies the marked Fox quadratic matrix with Hilbert symbols over $\mathbf Q_2$, while Theorem~\ref{thm:dyadic-q2-rankone-integral-ring} makes its integral nonalternating square the equation of a rank-one deformation hypersurface.  Theorem~\ref{thm:dyadic-exact-hilbert-reciprocity-marking} fixes all three Kummer basis characters through local reciprocity, and Theorem~\ref{thm:dyadic-tame-jacobian-fox-splitting} supplies a strict integral full-marked/Fox chain comparison when the residual image is a $2$-group.  Theorems~\ref{thm:dyadic-s3-schreier-corestriction} and~\ref{thm:dyadic-s3-marked-adjoint-jacobian} give the first non-pro-$2$ residual test: finite integral Schreier/Hecke transport and the complete marked adjoint coefficient Smith normal form at the tame $S_3$ quotient of $G_{\mathbf Q_2}$.  Theorems~\ref{thm:dyadic-s3-natural-strict-retract} and~\ref{thm:dyadic-s3-strict-bar-hecke-projector} further construct the strict integral contraction of the natural-coefficient marked complex and its explicit rank-two Sylow bar-cochain summand.  Theorems~\ref{thm:dyadic-s3-adjoint-unimodular-splitting},~\ref{thm:dyadic-s3-adjoint-analytic-smith}, and~\ref{thm:dyadic-s3-adjoint-shapiro-hecke} extend this analysis to the fixed $S_3$ adjoint coefficient lattice by an explicit tame-isotypic decomposition, analytic integral Smith reduction, and a precise top-degree Hecke projection.  Theorems~\ref{thm:dyadic-f131-canonical-tame-extraction},~\ref{thm:dyadic-f131-universal-split-locus}, and~\ref{thm:dyadic-f131-first-order-nosplit} prove that tame $C_3$ rigidifies as an integral polynomial projector over the full framed residual ring, yet becomes a Galois-equivariant decomposition only on an explicit proper wild-centralizer subfunctor.  The exact two-block multiplication laws of Theorem~\ref{thm:dyadic-f131-curved-tame-transport} give the finite native replacement outside it.  Higher Magnus--Fox data and finite relation jets provide explicit finite deformation information when it is needed, without making a chosen presentation intrinsic.

\item \textbf{Full transport and residual descent.}
Finite presentability together with $\operatorname{cd}_p(G_K)=2$ yields bounded perfect coefficient carriers compatible with derived base change.  Sylow absorption and all-Sylow averaging provide integral residual descent, while the universal Cayley--Hamilton algebra retains completed transport and nonsplit extension data.  Local complete-intersection/Koszul carriers and finite-type Cayley--Hamilton charts meet on the same residual deformation sectors.  For the optional marked full-$G_{\mathbf Q_2}$ presentation, Theorems~\ref{thm:dyadic-omega2-power-formula} and~\ref{thm:dyadic-q2-marked-word-jet-solver} make every profinite $2$-idempotent power and relation Jacobian explicitly computable on any prescribed finite residual jet, without claiming a global finite polynomial relation.

\item \textbf{Derived faithful closure and stratified recursive native globalization.}
For fixed $K$, $E$, and rank, the Cayley--Hamilton transport layer together with the canonical all-Sylow coefficient carrier gives a finite residual derived atlas that is faithful for representation plus cohomology data, with Morita-valued descent on Azumaya strata.  The finite Fox--Cayley--Hamilton and lci/Koszul charts, Fitting and descent-exponent stratifications, finite-stage flag descent, and Cech--Stasheff/colored-transfer theorems further give a stratified recursively explicit native globalization: on every finite chart and at every prescribed coherence arity, all required matrices and fillers reduce to finite word evaluation, inverse-minor formulas, finite linear systems, and finite rooted- or leveled-tree expressions.  Rank jumps are handled by the unreduced perfect complexes and derived base change.  A global strict projection is neither natural nor necessary, and wildness excludes a universal normal form.  No further closure or globalization problem remains; only optional strict-gauge refinements may ask for one preferred closed raw-period representative of an already constructed intrinsic comparison.
\end{enumerate}

Together these results establish finite existence and substantial explicit reconstruction while keeping the distinction between the intrinsic representation and any chosen cohomological gauge.

\subsection{What is proved, interpreted, and proposed}

The paper keeps three levels separate.  \emph{Identities} are exact reformulations or consequences of standard group cohomology, deformation theory, Fox calculus, or period-module constructions.  \emph{Interpretations} recast these objects as Cartan transport, Spencer compatibility, contact equations, holonomy, or curvature; the novelty there is organizational rather than a new equivalence of categories.  \emph{Programmatic statements} now concern only stronger closed-form refinements of the finite native atlas already constructed theoremically.

Accordingly, no theorem below claims a new classification of all continuous $p$-adic representations by a finite list of invariants or by a universal normal form.  The theorem-level contribution has two successive strengths.  First, derived faithful finite closure identifies finite full-transport and coefficient carriers on every residual sector and proves reconstruction, perfectness, base-change, and Morita descent.  Second, Theorem~\ref{thm:recursive-native-globalization} globalizes these carriers after finite residual, Fitting, flag, and descent-exponent refinement into a stratified recursively explicit native atlas.  Recursive explicitness means that every requested finite coherence level is obtained from finite matrix operations and terminating chartwise descent procedures; it does not mean that one canonical closed formula or one global matrix frame exists.

No mainline comparison problem remains after passing to the finite native shadow.  The direct linear map $\Theta_{HF}^{\mathrm{raw}}$ exists, while Theorem~\ref{thm:raw-herr-abstract-ainfty-existence} gives an intrinsic representative-level $A_\infty$ comparison.  Theorem~\ref{thm:raw-herr-phantom-reduced-normalization} shows that their linear terms agree automatically modulo continuous phantom maps, i.e. after retaining exactly the finite Tate/Fox data seen by the native atlas.  Literal equality on one unreduced Robba-ring representative is only an optional strict-gauge normalization, summarized in Remark~\ref{rem:raw-herr-exact-gauge-boundary}; no such choice is needed for closure, globalization, multiplicativity, or finite higher readout.

The other former ``open'' items are now boundaries rather than problems.  The identity component of the linear presentation-comparison space is contractible by Theorem~\ref{thm:presentation-linear-comparison-contractible}, while the multiplicative and simplicial higher coherence required by the atlas is supplied independently by the common-source and Cech--Stasheff constructions.  Direct all-arity word-native formulas are only a stronger optional refinement, for which Theorem~\ref{thm:relative-decomposable-ainfty-lifting} gives a relative lifting criterion.  Non-rigid descent already has a terminating finite chartwise solver, and Theorem~\ref{thm:no-uniform-descent-exponent-bound} shows that no bound depending only on the fixed representation/flag combinatorics can exist under arbitrary nonreduced coefficient base change.

\subsection{Organization}

The paper first develops the intrinsic transport, Spencer, Maurer--Cartan, rank-one, trianguline, and Herr language.  It then constructs finite Fox--Spencer carriers, derived and lci/Kuranishi charts, Cayley--Hamilton transport, residual and flag atlases, and the homotopy-coherent finite-atlas formulation.  The appendices retain only the two low-level calculations used directly in the main text.


\section{Existing classifications and the gap}

For a period ring $B_*$, one forms
\[
D_*(V)=(B_*\otimes_{\Qp}V)^{\GK}.
\]
Crystalline, semistable, and de Rham representations are those for which the corresponding period readout has the expected dimension.  Prismatic $F$-crystals provide an integral classification of crystalline lattices, with logarithmic variants reaching semistable settings in suitable forms \cite{BhattScholzePFC,DuLiuPrismatic}.

Arbitrary continuous representations are already encoded by etale $(\varphi,\Gamma)$-modules, and their cohomology is computed by Herr complexes \cite{FontaineLocal,CherbonnierColmez,HerrSurvey,LiuCohomology,KPXFamilies}.  The problem addressed here is therefore not categorical existence.  Rather, the standard equivalences do not by themselves present the representation together with its cohomological and deformation data as one finite derived/Morita atlas.  The results below supply that closure and a stratified recursively explicit native chain atlas; they then distinguish this from the still stronger demand for uniform closed-form formulas in every gauge.

\begin{warning}[Encoding versus geometric classification]
The equivalence with etale $(\varphi,\Gamma)$-modules is not a finite geometric classification by slopes, filtrations, monodromy, and ramification data.  The present construction supplies both the derived faithful finite closure and its stratified recursively explicit native refinement; neither is intended as a replacement for Fontaine theory or as a universal normal-form classification.
\end{warning}

Cartan language is natural because a representation is a transport law, group cohomology controls its deformations and obstructions, the deformation equation is Maurer--Cartan, and the Herr complex is a two-directional compatibility complex.  The theorem-level issue is finite faithful derived organization together with stratified recursive native globalization.  After that globalization is established, any further demand for uniform closed-form representatives is an optional gauge refinement rather than part of closure.

\section{Intrinsic transport and cochains}

Let $G$ be a topological group and
$\rho:G\to\mathrm{GL}_E(V)$ a continuous representation.  Writing
$\Delta_g:=\rho(g)-1$, the transport law is
\[
 \Delta_{gh}=\Delta_g+\Delta_h+\Delta_g\Delta_h.
\]
Thus a representation is already a flat profinite transport system; no
commutativity of the matrices is implied or required.

Write
\[
 C^\bullet_{\mathrm{cont}}(G,V)
 =\bigl(C^n_{\mathrm{cont}}(G,V),d_G\bigr)
\]
for the normalized continuous bar cochain complex.  This is the intrinsic
coefficient complex used throughout the paper.  We occasionally call it the
\emph{arithmetic Spencer complex}: this is only an interpretation of the usual
bar compatibility operator, not a new cohomology theory.  Every finite Fox,
Herr, or Koszul cochain chart constructed below is compared with this same
intrinsic object.

\section{Maurer--Cartan deformation theory}\label{sec:deformation}

Fix $\rho:G\to\mathrm{GL}_E(V)$ and write $\mathfrak g_\rho=\ad(V)$ with adjoint action $g\cdot X=\rho(g)X\rho(g)^{-1}$.  Continuous cochains carry the cup product
\[
(a\smile b)(g_1,\dots,g_{r+s})
=a(g_1,\dots,g_r)(g_1\cdots g_r)\cdot b(g_{r+1},\dots,g_{r+s}),
\]
and the graded commutator makes the completed cochain complex a dg-Lie algebra for nilpotent coefficients.

Let $(A,\mathfrak m)$ be a local Artinian $E$-algebra and write a lift as
\[
\rho_A(g)=(I+a(g))\rho(g),
\qquad a(g)\in\mathfrak m\otimes_E\mathfrak g_\rho.
\]
The representation law is equivalent to
\begin{equation}\label{eq:nonlinear-cocycle}
a(gh)=a(g)+g\cdot a(h)+a(g)\,g\cdot a(h),
\end{equation}
and hence to
\begin{equation}\label{eq:MC}
d_Ga+a\smile a=0,
\end{equation}
or $d_Ga+\frac12[a,a]=0$ in characteristic zero.

\begin{theorem}[Maurer--Cartan integrability]\label{thm:MC}
The map $a\mapsto(I+a)\rho$ identifies lifts of $\rho$ to $A$ on the chosen module with Maurer--Cartan elements of
\[
\mathfrak m\widehat\otimes C^\bullet_{\mathrm{cont}}(G,\mathfrak g_\rho).
\]
Conjugation by $I+b$, with $b\in\mathfrak m\otimes C^0(G,\mathfrak g_\rho)$, is the usual gauge action.
\end{theorem}

\begin{proof}
Equation \eqref{eq:nonlinear-cocycle} is exactly the Maurer--Cartan equation, and nilpotence makes all products and inverses finite.
\end{proof}

Thus
\[
F(a):=d_Ga+a\smile a
\]
is the deformation curvature.  At first order, tangent classes lie in $H^1(G,\mathfrak g_\rho)$; at second order the obstruction is represented by $a_1\smile a_1$ in $H^2(G,\mathfrak g_\rho)$.  For local Galois groups, $\operatorname{cd}_p(G_K)=2$, so this is the terminal obstruction degree.  The dg-Lie identity gives the Bianchi relation
\[
(d_G+[a,-])F(a)=0.
\]

\section{Rank-one tests}

For rank one, $\ad V\cong E$ is commutative.  Rank-one examples therefore test holonomy, logarithmic depth, and coefficient cohomology, but not nonabelian Maurer--Cartan curvature.

For $V=\Qp(1)$, Kummer theory gives
\begin{equation}\label{eq:kummer}
H^1(\GK,\Qp(1))
\cong
\widehat{K^\times}^{\,p}\otimes_{\Zp}\Qp,
\qquad
\widehat{K^\times}^{\,p}:=\varprojlim_n K^\times/(K^\times)^{p^n}.
\end{equation}
Writing $K^\times\cong\pi_K^{\mathbf Z}\times\OK^\times$ separates a valuation component from principal units linearized by the $p$-adic logarithm.  A Kummer class is therefore a model of discrete holonomy together with infinitesimal logarithmic transport.  Local Tate duality gives $H^2(\GK,\Qp(1))\cong\Qp$, but this is not the rank-one deformation-obstruction group: since $\ad V\cong E$ with trivial action,
\[
H^2(\GK,\ad V)=H^2(\GK,E)=0.
\]
Rank-one deformations are unobstructed at the adjoint $H^2$ level.

Now choose a Lubin--Tate formal $\OK$-module for a uniformizer $\pi$, write $K_n=K(F[\pi^n])$, and let
\[
\chi_{\mathrm{LT}}:G_K\longrightarrow\OK^\times
\]
be the Lubin--Tate character.  One has
\begin{equation}\label{eq:LT-gal}
\Gal(K_n/K)\cong(\OK/\pi^n)^\times.
\end{equation}

\begin{proposition}[Connection depth recovers the Lubin--Tate level]\label{prop:LT-depth}
Let $1\le m<n$ and $m>e_K/(p-1)$, where $e_K=v_K(p)$.  Then
\[
g|_{K_n}\in\Gal(K_n/K_m)
\quad\Longleftrightarrow\quad
v_K\bigl(\chi_{\mathrm{LT}}(g)-1\bigr)\ge m.
\]
In this range,
\[
v_K\bigl(\log\chi_{\mathrm{LT}}(g)\bigr)
=v_K\bigl(\chi_{\mathrm{LT}}(g)-1\bigr).
\]
\end{proposition}

\begin{proof}
Under \eqref{eq:LT-gal}, $\Gal(K_n/K_m)$ is the principal-unit congruence subgroup.  The numerical bound places it in the valuation-preserving logarithmic range.
\end{proof}

Hence the Lubin--Tate level is read from logarithmic connection depth.  The following single-generator calculation shows why a scaling defect carries no further information.  If $J_n=\log(T^{p^n})$, then $J_n=p^nJ_0$.

\begin{proposition}[Single-generator logarithmic-scaling no-go]\label{prop:rank1-nogo}
For every such $T$,
\[
J_{n+1}-pJ_n=0.
\]
Hence a defect tower built only from powers of one transport generator contains no new information.
\end{proposition}

An informative higher defect must therefore involve several transport directions, nonlinear coefficient geometry, or genuine extension data.

\section{Trianguline extensions and Herr compatibility}

\subsection{Rank-one \texorpdfstring{$(\varphi,\Gamma)$}{(phi,Gamma)}-modules}

Let $\cR_E$ be the Robba ring over $E$.  A character
\[
\delta:\Qp^\times\longrightarrow E^\times
\]
defines $\cR_E(\delta)$.  A rank-two trianguline module is an extension \cite{ColmezTrianguline,BergerTrianguline}
\begin{equation}\label{eq:trianguline-extension}
0\longrightarrow D_1\longrightarrow D\longrightarrow D_2\longrightarrow0,
\end{equation}
where $D_i=\cR_E(\delta_i)$.

Choose a basis $e_1,e_2$ compatible with the filtration.  Write
\begin{align}
\varphi(e_1)&=\alpha_1e_1,
&\varphi(e_2)&=a_\varphi e_1+\alpha_2e_2,\label{eq:phi-matrix}\\
\gamma(e_1)&=\beta_1e_1,
&\gamma(e_2)&=a_\gamma e_1+\beta_2e_2,\label{eq:gamma-matrix}
\end{align}
where $\gamma$ is a topological generator of the procyclic part of $\Gamma$ after taking torsion invariants.

\subsection{Exact compatibility calculation}

The equality $\varphi\gamma(e_2)=\gamma\varphi(e_2)$ gives
\begin{align}
\alpha_1\varphi(a_\gamma)+\varphi(\beta_2)a_\varphi
=
\beta_1\gamma(a_\varphi)+\gamma(\alpha_2)a_\gamma.
\label{eq:extension-compatibility}
\end{align}
The diagonal equality
\[
\varphi(\beta_2)\alpha_2=\gamma(\alpha_2)\beta_2
\]
is the rank-one compatibility for $D_2$.

Now change the splitting by
\[
e_2'=e_2+b e_1.
\]
A direct computation gives
\begin{align}
 a_\varphi'&=a_\varphi+\alpha_1\varphi(b)-\alpha_2b,\label{eq:gauge-phi}\\
 a_\gamma'&=a_\gamma+\beta_1\gamma(b)-\beta_2b.\label{eq:gauge-gamma}
\end{align}
Thus $(a_\varphi,a_\gamma)$ is the off-diagonal connection datum and a change of splitting is gauge.

After twisting by $D_2^\vee$, choose a rank-one basis $f$ of
\[
M:=D_1\otimes D_2^\vee\cong \cR_E(\delta_1\delta_2^{-1})
\]
with
\[
\varphi_M(f)=\alpha_1\alpha_2^{-1}f,
\qquad
\gamma_M(f)=\beta_1\beta_2^{-1}f.
\]
Set
\begin{equation}\label{eq:herr-normalization}
y:=\frac{a_\varphi}{\alpha_2}f,
\qquad
z:=\frac{a_\gamma}{\beta_2}f,
\qquad
x:=bf.
\end{equation}
The diagonal rank-one compatibility gives the common clearing factor
\[
\Lambda:=\varphi(\beta_2)\alpha_2
=\gamma(\alpha_2)\beta_2.
\]
After multiplication by $\Lambda$, the Herr cocycle expression
\[
(\gamma_M-1)y-(\varphi_M-1)z
\]
is, up to sign, the difference of the two sides of \eqref{eq:extension-compatibility}; the only remainder is
\[
\frac{a_\gamma}{\beta_2}
\bigl(\alpha_2\varphi(\beta_2)-\beta_2\gamma(\alpha_2)\bigr),
\]
which vanishes by the diagonal relation.  Hence the Herr cocycle equation is exactly \eqref{eq:extension-compatibility}, while the Herr coboundary
\[
((\varphi_M-1)x,(\gamma_M-1)x)
\]
recovers \eqref{eq:gauge-phi}--\eqref{eq:gauge-gamma}.  The diagonal relation is therefore part of the comparison, not an auxiliary condition.

\subsection{The Herr complex as a finite arithmetic Spencer complex}

Let $\Delta\subset\Gamma$ be the finite torsion subgroup, put $D'=D^\Delta$, and choose a topological generator $\gamma$ of $\Gamma/\Delta$.  The Herr complex is
\begin{equation}\label{eq:herr-complex}
C^\bullet_{\varphi,\gamma}(D):
D'
\xrightarrow{d^0}
D'\oplus D'
\xrightarrow{d^1}
D',
\end{equation}
with
\begin{align}
 d^0(x)&=((\varphi-1)x,(\gamma-1)x),\label{eq:herr-d0}\\
 d^1(y,z)&=(\gamma-1)y-(\varphi-1)z.\label{eq:herr-d1}
\end{align}
Because $\varphi$ and $\gamma$ commute, $d^1d^0=0$.

\begin{theorem}[Herr]\label{thm:herr}
For an etale $(\varphi,\Gamma)$-module $D(V)$ associated with a $p$-adic representation $V$, the cohomology of \eqref{eq:herr-complex} canonically computes the continuous Galois cohomology of $V$:
\[
H^i_{\varphi,\gamma}(D(V))\cong H^i(\GK,V).
\]
\end{theorem}

The theorem is standard \cite{HerrSurvey,LiuCohomology}; here it fixes the finite chart used below.

\begin{proposition}[Finite Cartan--Spencer interpretation]\label{prop:herr-spencer}
The Herr complex is a two-directional arithmetic Spencer compatibility complex:
\begin{itemize}
\item $d^0$ assigns the formal first derivatives in the Frobenius and $\Gamma$ directions;
\item $d^1$ is the mixed contact defect comparing the two prolongations;
\item $H^0$ is the horizontal-section space;
\item $H^1$ classifies compatible off-diagonal connection data modulo gauge;
\item $H^2$ is the obstruction quotient.
\end{itemize}
\end{proposition}

\begin{proof}
Equations~\eqref{eq:herr-d0}--\eqref{eq:herr-d1} say that $(y,z)$ is compatible exactly when
\[
(\gamma-1)y=(\varphi-1)z,
\]
and that $x$ changes it by a coboundary.
\end{proof}

\subsection{Extension is not curvature}

For the extension \eqref{eq:trianguline-extension}, a nonzero class in
\[
H^1_{\varphi,\gamma}(D_1\otimes D_2^\vee)
\]
represents a non-split but still compatible extension.  Therefore:
\[
\boxed{\text{extension class} = \text{connection data modulo gauge},}
\]
not curvature.  Curvature is the failure of the Herr cocycle equation; such a pair does not define an extension.

\section{Fox--Herr comparison}\label{sec:fox-herr-bridge}

Herr and Fox--Spencer are finite charts of the same object, $R\Gamma(G_K,-)$.  Herr totalizes the cyclotomic Hochschild--Serre chart; Fox dualizes a finite resolution on an open pro-$p$ subgroup and then descends by transfer.

\subsection{Herr as a totalized normal chart}

Work in this subsection over the characteristic-zero coefficient field $E$, and let $p$ be arbitrary.  At $p=2$ the torsion subgroup of $\Gamma_K$ may have order $2$; it is harmless for $E$-linear comparison but must be retained derivedly for integral lattices (Proposition~\ref{prop:dyadic-integral-herr-derived}).  Put
\[
 K_\infty=K(\mu_{p^\infty}),
 \qquad
 H_K=G_{K_\infty},
 \qquad
 \Gamma_K=G_K/H_K.
\]
Let $\Delta_K\subset\Gamma_K$ be the finite torsion subgroup and choose a topological generator $\gamma$ of $\Gamma_K/\Delta_K$.  For a representation $V$, let $D(V)$ be its etale $(\varphi,\Gamma_K)$-module and write $D'=D(V)^{\Delta_K}$.

\begin{theorem}[Cyclotomic Hochschild--Serre factorization of the Herr complex]\label{thm:herr-hochschild-serre}
There is a natural factorization
\begin{equation}\label{eq:herr-hs-factorization}
 R\Gamma(G_K,V)
 \simeq
 R\Gamma\!\left(
 \Gamma_K,
 R\Gamma(H_K,V)
 \right)
 \simeq
 \Tot\!\left(
 D'
 \xrightarrow{\ (\varphi-1,\,\gamma-1)\ }
 D'\oplus D'
 \xrightarrow{\ (\gamma-1)\oplus(1-\varphi)\ }
 D'
 \right).
\end{equation}
The final total complex is the Herr complex of Equations~\eqref{eq:herr-complex}--\eqref{eq:herr-d1}.  The factorization is functorial in $V$, and replacing $\gamma$ by another generator changes the displayed complex by the standard explicit generator-change chain isomorphism.
\end{theorem}

\begin{proof}
Fontaine--Herr theory identifies the $H_K$-cohomology of $V$ with the two-term Frobenius complex
\[
 K_\varphi(D(V))
 :=
 [D(V)\xrightarrow{\varphi-1}D(V)]
\]
viewed as a complex with continuous $\Gamma_K$-action.  The Hochschild--Serre equivalence for
\[
1\longrightarrow H_K\longrightarrow G_K\longrightarrow\Gamma_K\longrightarrow1
\]
then gives the first equivalence in \eqref{eq:herr-hs-factorization}.  Since $|\Delta_K|$ is invertible in the characteristic-zero field $E$, taking $\Delta_K$-invariants on $E$-linear modules is exact, including at $p=2$.  No integral exactness is claimed here.  The quotient $\Gamma_K/\Delta_K\cong\mathbf Z_p$ has the standard two-term continuous resolution represented by $\gamma-1$.  Totalizing this two-by-two bicomplex gives exactly
\[
D'\xrightarrow{(\varphi-1,\gamma-1)}D'\oplus D'
\xrightarrow{(\gamma-1,1-\varphi)}D'.
\]
This is Herr's construction and computes continuous Galois cohomology \cite{HerrSurvey,LiuCohomology,KPXFamilies}.  The explicit independence of the generator is the standard Herr generator-change map.
\end{proof}

\begin{proposition}[Integral dyadic cyclotomic descent keeps the torsion direction]
\label{prop:dyadic-integral-herr-derived}
Let $p=2$, let $T$ be a $G_K$-stable $\mathcal O_E$-lattice,
let $\Gamma_K=\operatorname{Gal}(K(\mu_{2^\infty})/K)$,
and let $\Delta_K\subset\Gamma_K$ be its finite torsion subgroup.
Put $H_K=G_{K(\mu_{2^\infty})}$.  Then
\begin{equation}\label{eq:dyadic-integral-herr-hs}
R\Gamma(G_K,T)\simeq
R\Gamma\!\left(\Gamma_K/\Delta_K,\,
 R\Gamma\!\left(\Delta_K,R\Gamma(H_K,T)\right)\right).
\end{equation}
Whenever a compatible integral Fontaine--Herr Frobenius model
$K_\varphi(D^+(T))
=[D^+(T)\xrightarrow{\varphi-1}D^+(T)]$
is available for $R\Gamma(H_K,T)$, it may replace the
innermost factor in this derived expression.
If $\Delta_K=C_2$, replacing $R\Gamma(\Delta_K,-)$ by
ordinary termwise $\Delta_K$-invariants is generally invalid.
After inverting $\varpi_E$, the ordinary three-term
$D(V)^{\Delta_K}$ Herr complex does recover
$R\Gamma(G_K,V)$ for $V=T[1/\varpi_E]$.
\end{proposition}

\begin{proof}
Two applications of Hochschild--Serre, first to $H_K\triangleleft G_K$
and then to $\Delta_K\triangleleft\Gamma_K$, prove
\eqref{eq:dyadic-integral-herr-hs}; the derived quotient action
retains the extension's homotopy coherence.
For a trivial $C_2=\langle\delta\rangle$-action on $\mathbb Z_2$,
the periodic cochain resolution has alternating operators
$\delta-1=0$ and $1+\delta=2$, so
\[
H^{2j}(C_2,\mathbb Z_2)\cong\mathbb Z/2\mathbb Z,
\quad j\ge1.
\]
Thus integral invariants cannot be substituted for
derived invariants.  A compatible integral Frobenius model
can be substituted into the derived equivalence.
Over $E$, $2$ is invertible; finite-group averaging makes
$\Delta_K$-invariants exact and recovers the usual
$E$-linear Herr totalization.
\end{proof}

\begin{remark}[Herr as a compressed Hochschild--Serre chart]\label{rem:herr-wall-compression}
Herr compresses the kernel direction by the Frobenius complex and the quotient direction by the one-generator Koszul complex of $\Gamma_K/\Delta_K$.  This is a finite cohomological chart for the intrinsic Galois cochain object, not a replacement for full profinite transport.
\end{remark}

\subsection{Finite Fox--Herr comparison}

Let $\mathcal K_{\mathrm{SylFox}}^\bullet(V)$ denote any of the finite all-Sylow Fox retracts constructed in Theorem~\ref{thm:all-sylow-transfer}, or its family version in Theorem~\ref{thm:family-sylow-fox}.  Both the Fox carrier and the Herr complex represent the same intrinsic Galois cochain object.

\begin{theorem}[Fox--Herr comparison through the intrinsic cochain object]\label{thm:fox-herr-zigzag}
For every continuous $E$-representation $V$ and every cyclotomic generator $\gamma$, there is a functorial zigzag of quasi-isomorphisms
\begin{equation}\label{eq:fox-herr-zigzag}
 \mathcal K_{\mathrm{SylFox}}^\bullet(V)
 \xrightarrow{\ \sim\ }
 R\Gamma(G_K,V)
 \xleftarrow{\ \sim\ }
 C_{\varphi,\gamma}^\bullet(D(V)).
\end{equation}
The zigzag is compatible with derived coefficient base change on every residual framed family.  Choose a strict perfect representative on the Herr side,
\[
 P_{\mathrm{Herr}}^\bullet(V)
 \xrightarrow{\ \sim\ }
 C_{\varphi,\gamma}^\bullet(D(V)),
\]
and choose homotopy inverses in the perfect derived category.  The resulting finite chain quasi-isomorphism
\[
 \Theta_{\mathrm{FH}}:
 P_{\mathrm{Herr}}^\bullet(V)
 \xrightarrow{\ \sim\ }
 \mathcal K_{\mathrm{SylFox}}^\bullet(V)
\]
is well defined up to chain homotopy once the strict perfect models have been fixed.  A change of Sylow/Fox presentation is governed by the established Fox transition equivalences.  A change of $\gamma$ is governed by the explicit Herr generator-change chain isomorphism.
\end{theorem}

\begin{proof}
The two arrows in \eqref{eq:fox-herr-zigzag} are the all-Sylow transfer equivalence and Theorem~\ref{thm:herr-hochschild-serre}, and both commute with derived base change.  Choosing strict perfect representatives and a homotopy inverse therefore produces $\Theta_{\mathrm{FH}}$.  Its chain-homotopy class is independent of these auxiliary choices; no canonical continuous deformation retract of the raw Herr complex is claimed.
\end{proof}

The zigzag identifies the chain-homotopy class of the comparison.  Its degree-one Fox coordinates admit a direct formula.

\begin{proposition}[Degree-one Fox coordinates of a Herr class]\label{prop:herr-to-fox-h1}
Let
\[
 P=\langle x_1,\ldots,x_d\mid r_1,\ldots,r_m\rangle_{\mathrm{pro}\text{-}p}
\]
be a Fox chart occurring after restriction to an open subgroup, and let $(y,z)$ be a Herr $1$-cocycle.  Fix a cochain-level representative of the Herr-to-Galois quasi-isomorphism and let
\[
 c_{y,z}\in Z^1(G_K,V)
\]
be the resulting continuous Galois cocycle.  Changing the cochain-level comparison changes $c_{y,z}$ by a Galois coboundary.  Then the Fox coordinates of the restricted class are
\begin{equation}\label{eq:herr-to-fox-h1-evaluation}
 a_i:=c_{y,z}(x_i),
 \qquad 1\le i\le d,
\end{equation}
and they satisfy
\begin{equation}\label{eq:herr-to-fox-relation}
 \sum_{i=1}^d
 \rho\!\left(
 \overline{\frac{\partial r_j}{\partial x_i}}
 \right)a_i
 =0,
 \qquad 1\le j\le m.
\end{equation}
If $(y,z)$ is a Herr coboundary, then $(a_i)$ is the Fox coboundary of the corresponding degree-zero element.  Hence \eqref{eq:herr-to-fox-h1-evaluation} induces the degree-one component of the comparison in Theorem~\ref{thm:fox-herr-zigzag}.
\end{proposition}

\begin{proof}
For a continuous crossed homomorphism $c$ and a word $w$ in the free pro-$p$ generators, the twisted Fox formula gives
\[
 c(w)
 =
 \sum_i
 \rho\!\left(
 \overline{\frac{\partial w}{\partial x_i}}
 \right)c(x_i),
\]
with the involution convention used in Equation~\eqref{eq:fox-d1}.  Applying this to every relation $r_j=1$ gives \eqref{eq:herr-to-fox-relation}.  A Galois coboundary $c(g)=(g-1)v$ evaluates to $((x_i-1)v)_i$, which is exactly the Fox differential $d^0(v)$ of Equation~\eqref{eq:fox-d0}.
\end{proof}


\begin{theorem}[Raw cyclotomic Herr $H^1$ extraction by Tate duality]\label{thm:raw-herr-tate-extraction-core}
Let $U=\operatorname{Sp}A$ be a representation-induced affinoid residual chart and let $W$ be a finite projective $A$-module with continuous $G_K$-action.  Put $W^{\ast}:=W^{\vee}(1)$ and assume that
\[
 H^1(G_K,W),\qquad H^1(G_K,W^{\ast})
\]
are finite projective of the same rank $h$ on $U$, with the relative local Tate pairing perfect.  Choose raw Herr cocycles
\[
 e_a=(y_a,z_a)\in Z^1\!\left(C^\bullet_{\varphi,\gamma}(D(W))\right),
 \qquad
 \eta_b=(u_b,v_b)\in Z^1\!\left(C^\bullet_{\varphi,\gamma}(D(W^{\ast}))\right),
\]
whose cohomology classes are frames, for $1\le a,b\le h$.  For an arbitrary raw Herr $1$-cocycle $c=(y,z)$ define
\begin{equation}\label{eq:raw-herr-tate-functional-core}
 \mathcal P_b^{\gamma}(c)
 :=\operatorname{Ta}_{K}\!\left(
 \left[
 \operatorname{ev}\!\left(
 z\otimes\gamma(u_b)-y\otimes\varphi(v_b)
 \right)
 \right]
 \right)\in A,
\end{equation}
where $\operatorname{ev}:D(W)\otimes D(W^{\ast})\to D(A(1))$ is Cartier evaluation and $\operatorname{Ta}_{K}:H^2_{\varphi,\gamma}(D(A(1)))\to A$ is the normalized Herr--Tate map.  Put
\begin{equation}\label{eq:raw-herr-tate-matrix-core}
 \mathsf T^{\gamma}_{ba}:=\mathcal P_b^{\gamma}(e_a),
 \qquad
 \mathbf p^{\gamma}(c):=
 \bigl(\mathcal P_b^{\gamma}(c)\bigr)_{b=1}^{h}.
\end{equation}
Then:
\begin{enumerate}[label=\textnormal{(\roman*)}]
\item each $\mathcal P_b^{\gamma}$ depends only on the Herr cohomology class of $c$ and vanishes on Herr coboundaries;
\item $\mathsf T^{\gamma}$ is invertible, and if
\[
 [c]=\sum_{a=1}^{h}\lambda_a(c)[e_a],
\]
then the cohomology coordinates are obtained directly from the raw representative by
\begin{equation}\label{eq:raw-herr-tate-coordinate-core}
 \boxed{
 \boldsymbol\lambda_{\gamma}(c)
 =\bigl(\mathsf T^{\gamma}\bigr)^{-1}\mathbf p^{\gamma}(c).}
\end{equation}
In a Tate-dual choice of frames, $\mathsf T^{\gamma}=I_h$.
\item If $A^+\subset A$ is a $p$-adically complete integral model on which the chosen lattices, cocycles, and normalized Tate pairing are integral, then Equation~\eqref{eq:raw-herr-tate-coordinate-core} is integral on the locus $\det(\mathsf T^{\gamma})\in(A^+)^{\times}$.  Off that locus the only denominator is $\det(\mathsf T^{\gamma})$, via the adjugate formula.
\item If $\gamma'$ is another topological generator and the primal and dual frames are transported by the standard Herr generator-change chain isomorphism, then
\begin{equation}\label{eq:raw-herr-generator-invariance-core}
 \boldsymbol\lambda_{\gamma'}(U^1_{\gamma'\gamma}c)
 =\boldsymbol\lambda_{\gamma}(c).
\end{equation}
Thus the extraction is strictly compatible with change of cyclotomic generator on every common rank and Tate-pairing minor chart.
\end{enumerate}
In particular, extraction of a raw cyclotomic Herr $H^1$ class requires neither a solution of a $(\varphi-1)$ equation nor a raw-Herr-to-perfect homotopy inverse.
\end{theorem}

\begin{proof}
For degree-one Herr cocycles the relative Herr cup product is represented, in the conventions of Equations~\eqref{eq:herr-d0}--\eqref{eq:herr-d1}, by
\[
 (y,z)\smile(u_b,v_b)
 =z\otimes\gamma(u_b)-y\otimes\varphi(v_b)
\]
in degree two.  After Cartier evaluation and the normalized Tate map, Equation~\eqref{eq:raw-herr-tate-functional-core} is therefore exactly the relative local Tate pairing.  The relative cup product and Tate duality are compatible with families and base change \cite{KPXFamilies}; at a point this is Liu's local duality \cite{LiuCohomology}.  Hence the functionals vanish on coboundaries, and perfectness makes the matrix \eqref{eq:raw-herr-tate-matrix-core} invertible.  Bilinearity gives
\[
 \mathbf p^{\gamma}(c)=\mathsf T^{\gamma}\boldsymbol\lambda_{\gamma}(c),
\]
which proves~(ii).  Part~(iii) is finite linear algebra.  For~(iv), the standard generator-change chain isomorphism intertwines the Herr cup product, while the Tate pairing is intrinsic; transporting both frames therefore leaves the coordinate vector unchanged.
\end{proof}

\begin{corollary}[Direct raw Herr-to-Fox $H^1$ formula on a versal transport chart]\label{cor:raw-herr-direct-fox-core}
In the setting of Theorem~\ref{thm:raw-herr-tate-extraction-core}, suppose in addition that $W=\operatorname{Hom}_A(M_2,M_1)$ and that the chosen Herr frame $e_a$ comes from a cohomologically versal first-order Cayley--Hamilton extension family.  Thus, for a tangent basis $s_a$, the marked transport on a Fox generator $x_i$ has first-order block
\[
 \rho_{\mathcal E}(x_i)=
 \begin{pmatrix}
 R_{1,i}&\epsilon B_i(s_a)\\
 0&R_{2,i}
 \end{pmatrix},
 \qquad \epsilon^2=0,
\]
and the corresponding Fox cocycle column is
\begin{equation}\label{eq:versal-fox-column-core}
 J_F(s_a)_i:=B_i(s_a)R_{2,i}^{-1}.
\end{equation}
Let $K_F$ be the finite matrix whose $a$-th column is the Fox cohomology coordinate vector of $J_F(s_a)$, computed by any determinantal Fox cohomology projector on the chosen constant-rank minor chart.  Then every raw Herr $1$-cocycle $c$ has direct finite Fox cohomology coordinates
\begin{equation}\label{eq:raw-herr-direct-fox-coordinate-core}
 \boxed{
 [c]_{\mathrm{Fox}}
 =K_F\bigl(\mathsf T^{\gamma}\bigr)^{-1}
 \mathbf p^{\gamma}(c).}
\end{equation}
All finite matrices in the right-hand side are obtained from the marked Cayley--Hamilton transport Jacobian, the Fox differential, the Tate-pairing matrix, and inverse minors.  The formula is integral on the common unit-minor/Tate-determinant locus and is compatible with change of $\gamma$.
\end{corollary}

\begin{proof}
The normalized upper-right block in Equation~\eqref{eq:versal-fox-column-core} is the Galois/Fox cocycle of the same first-order extension whose relative $(\varphi,\gamma)$ derivative gives $e_a$.  Hence the canonical Fox--Herr cohomology comparison sends the $a$-th Herr frame vector to the $a$-th column of $K_F$.  Equation~\eqref{eq:raw-herr-tate-coordinate-core} gives the coordinates of an arbitrary raw class in that Herr frame, and composing the two finite coordinate maps gives \eqref{eq:raw-herr-direct-fox-coordinate-core}.
\end{proof}


\begin{theorem}[Tate-dual all-degree raw Herr perfectification]\label{thm:raw-herr-all-degree-perfectification}
Let $U=\operatorname{Sp}A$ be as in Theorem~\ref{thm:raw-herr-tate-extraction-core}, put $W^{\ast}=W^{\vee}(1)$, and assume for $i=0,1,2$ that
\[
 H^i(G_K,W),\qquad H^{2-i}(G_K,W^{\ast})
\]
are finite projective and that the relative local Tate pairing between them is perfect.  After a finite frame/minor refinement, choose raw Herr cocycle frames
\[
 e_{i,a}\in Z^i\!\left(C^\bullet_{\varphi,\gamma}(D(W))\right),
 \qquad
 \eta_{i,b}\in Z^{2-i}\!\left(C^\bullet_{\varphi,\gamma}(D(W^{\ast}))\right),
 \qquad 1\le a,b\le h_i.
\]
Let
\[
 \operatorname{ta}^{\mathrm{ch}}_{K,\gamma}:
 C^2_{\varphi,\gamma}(D(A(1)))\longrightarrow A,
 \qquad
 w\longmapsto \operatorname{Ta}_K([w]),
\]
be the chain-level top-degree Tate trace; it kills the image of the Herr degree-one differential.  Using the standard chain-level Herr cup product $\smile_\gamma$, define for an arbitrary raw degree-$i$ cochain $c$
\begin{align}
 \mathcal P^{\gamma}_{i,b}(c)
 &:=\operatorname{ta}^{\mathrm{ch}}_{K,\gamma}
 \left(
 \operatorname{ev}(c\smile_\gamma\eta_{i,b})
 \right),
 \label{eq:raw-herr-all-degree-functional}\\
 \mathsf T_i^{\gamma}
 &:=\bigl(\mathcal P^{\gamma}_{i,b}(e_{i,a})\bigr)_{b,a}.
 \label{eq:raw-herr-all-degree-tate-matrix}
\end{align}
Then every $\mathsf T_i^{\gamma}$ is invertible and the maps
\begin{equation}\label{eq:raw-herr-all-degree-projection}
 q_i^{\gamma}(c)
 :=(\mathsf T_i^{\gamma})^{-1}
   \bigl(\mathcal P^{\gamma}_{i,b}(c)\bigr)_b
 \in A^{h_i}
\end{equation}
assemble to a chain map
\begin{equation}\label{eq:raw-herr-finite-perfectification}
 \boxed{
 q_{\mathrm H}^{\gamma}:
 C^\bullet_{\varphi,\gamma}(D(W))
 \longrightarrow
 H^\bullet_{\mathrm{Tate}}(W)
 :=\left[A^{h_0}\xrightarrow{0}A^{h_1}\xrightarrow{0}A^{h_2}\right]
 }
\end{equation}
which is a quasi-isomorphism.  More precisely:
\begin{enumerate}[label=\textnormal{(\roman*)}]
\item $q_i^{\gamma}$ vanishes on degree-$i$ Herr boundaries, so $q_{\mathrm H}^{\gamma}$ is a chain map to the zero-differential finite complex;
\item on cohomology, $q_i^{\gamma}$ is the coordinate isomorphism in the frame $[e_{i,a}]$;
\item the construction commutes with coefficient base change preserving the chosen rank and Tate-pairing minor charts, and it is integral whenever the three pairing determinants are units in a chosen integral model;
\item after transporting the primal and dual frames by the standard Herr generator-change chain isomorphisms, the induced finite coordinates on raw cocycles are invariant under change of cyclotomic generator; equivalently the two perfectifications induce the same cohomology-coordinate map under the standard generator-change identification.
\end{enumerate}
The degree-one component is exactly Theorem~\ref{thm:raw-herr-tate-extraction-core}.  Thus the raw cyclotomic Herr complex admits a direct finite-output perfectification in all cohomological degrees without choosing a Frobenius parametrix, solving a $(\varphi-1)$ equation, or first choosing an abstract finite Herr representative.
\end{theorem}

\begin{proof}
The relative Herr cup product satisfies the Leibniz rule.  Since every $\eta_{i,b}$ is closed, if $c=d a$ is a boundary then
\[
 c\smile_\gamma\eta_{i,b}
 =d(a\smile_\gamma\eta_{i,b}).
\]
After Cartier evaluation this lies in the image of the degree-one differential of the Tate-twist Herr complex, and $\operatorname{ta}^{\mathrm{ch}}_{K,\gamma}$ kills that image.  Hence every functional in \eqref{eq:raw-herr-all-degree-functional} kills boundaries.  This proves the chain-map assertion because the target differential is zero.

On cycles, the same functionals are precisely the relative local Tate pairings with the dual frame classes.  Perfect Tate duality therefore makes each matrix \eqref{eq:raw-herr-all-degree-tate-matrix} invertible, and Equation~\eqref{eq:raw-herr-all-degree-projection} sends the class of $e_{i,a}$ to the $a$-th standard basis vector.  It is consequently an isomorphism on every cohomology module and hence a quasi-isomorphism.  Base change and integrality on the unit-determinant loci follow from the corresponding properties of the relative Herr cup product and Tate duality.  Generator change preserves the intrinsic Tate pairing; after transporting the raw cocycle frames, the resulting finite coordinates therefore agree on cocycles \cite{KPXFamilies,LiuCohomology}.
\end{proof}

\begin{corollary}[Direct all-degree raw Herr--Fox chain comparison]\label{cor:raw-herr-all-degree-fox}
Retain Theorem~\ref{thm:raw-herr-all-degree-perfectification}, and let $F^\bullet$ be a finite all-Sylow/Fox representative on the same native chart.  Refine to a constant-rank minor chart and choose a finite Fox contraction
\[
 \left(
 H_F^\bullet
 \mathop{\rightleftarrows}^{i_F}_{p_F}
 F^\bullet,h_F
 \right).
\]
Let
\[
 C_{HF}^{\mathrm{coh}}:
 H^\bullet_{\mathrm{Tate}}(W)
 \xrightarrow{\ \sim\ }
 H_F^\bullet
\]
be the finite block-diagonal matrix of the canonical Herr--Fox cohomology identification in the chosen frames.  Then
\begin{equation}\label{eq:raw-herr-all-degree-fox-map}
 \boxed{
 \Theta_{HF}^{\mathrm{raw}}
 :=i_F\,C_{HF}^{\mathrm{coh}}\,q_{\mathrm H}^{\gamma}:
 C^\bullet_{\varphi,\gamma}(D(W))
 \longrightarrow F^\bullet
 }
\end{equation}
is a direct chain quasi-isomorphism from the raw Herr complex to the finite native Fox core.

On a cohomologically versal Cayley--Hamilton extension chart, the degree-one block of $C_{HF}^{\mathrm{coh}}$ is the finite transport Jacobian $K_F$ of Corollary~\ref{cor:raw-herr-direct-fox-core}; the degree-zero block is the common invariant-coordinate map; and the degree-two block is determined by the finite Herr and Fox Tate-pairing matrices.  Thus all finite output coordinates of \eqref{eq:raw-herr-all-degree-fox-map} are obtained from raw Herr cup/Tate functionals, marked transport/Fox matrices, finite Tate pairings, and inverse minors.
\end{corollary}

\begin{proof}
Theorem~\ref{thm:raw-herr-all-degree-perfectification} makes $q_{\mathrm H}^{\gamma}$ a quasi-isomorphism onto the finite cohomology model.  The finite matrix $C_{HF}^{\mathrm{coh}}$ is an isomorphism and $i_F$ is the inclusion in a Fox contraction, hence a quasi-isomorphism.  Their composite is therefore a chain quasi-isomorphism.  The final coordinate description is Corollary~\ref{cor:raw-herr-direct-fox-core} in degree one and ordinary finite local Tate duality in degrees zero and two.
\end{proof}

\begin{warning}[Finite output is the correct raw-Herr notion]\label{warn:raw-herr-not-finite-a-matrix}
The terms of $C^\bullet_{\varphi,\gamma}(D(W))$ are finite projective over a relative Robba ring, but they are generally not finite projective as $A$-modules.  Therefore a request for a termwise finite $A$-matrix on the unreduced raw Herr terms is not the correct finite-atlas target.  The finite object is the output complex $H^\bullet_{\mathrm{Tate}}(W)$, and Equations~\eqref{eq:raw-herr-all-degree-functional}--\eqref{eq:raw-herr-finite-perfectification} give its finite family of native coordinate functionals.  Asking for a finite $A$-basis of the entire raw period module would reintroduce an artificial requirement unrelated to faithful finite closure.
\end{warning}

\begin{corollary}[Raw Herr classes feed the finite native higher operations]\label{cor:raw-herr-higher-readout-core}
On a chart satisfying Corollary~\ref{cor:raw-herr-direct-fox-core}, let $m_n^{F}$ be any prescribed finite transferred $A_\infty$ operation on the Fox cohomology gauge, or let $q_n^{F}$ denote the corresponding $L_\infty$ operation on adjoint cohomology.  For raw cyclotomic Herr $1$-cocycles $c_1,\ldots,c_n$, put
\[
 \Lambda_{HF}^{\gamma}(c_j)
 :=K_F(\mathsf T^{\gamma})^{-1}\mathbf p^{\gamma}(c_j).
\]
Then the cohomological higher readouts are computed directly by
\[
 m_n^{F}\bigl(\Lambda_{HF}^{\gamma}(c_1),\ldots,
                 \Lambda_{HF}^{\gamma}(c_n)\bigr),
 \qquad
 q_n^{F}\bigl(\Lambda_{HF}^{\gamma}(c_1),\ldots,
                 \Lambda_{HF}^{\gamma}(c_n)\bigr).
\]
For every fixed $n$ these are finite native formulas and require no period parametrix or $(\varphi-1)$ solve.  This closes raw cyclotomic Herr extraction and its finite cohomological higher readout.  It does not construct a direct representative-level $A_\infty$ morphism on the full raw Robba-ring Herr complex.
\end{corollary}


\begin{theorem}[Abstract representative-level raw-Herr multiplicative comparison exists]\label{thm:raw-herr-abstract-ainfty-existence}
Let $\mathcal A$ be a continuous coefficient dga (or an $E_1$ coefficient
algebra) in a fixed residual family sector.  Then the raw cyclotomic Herr gauge
and the finite all-Sylow/Fox gauge are equivalent as $E_1$ objects.  After
choosing ordinary dg/$A_\infty$ representatives there is a zigzag
\begin{equation}\label{eq:raw-herr-abstract-ainfty-zigzag}
 C^\bullet_{\varphi,\gamma}(D(\mathcal A))
 \simeq_{A_\infty}
 R\Gamma(G_K,\mathcal A)
 \simeq_{A_\infty}
 \mathcal K_{\bar\rho}^{\mathrm{allSyl}}(\mathcal A),
\end{equation}
and hence, on every finite Fox chart, some representative-level
$A_\infty$ quasi-isomorphism
\begin{equation}\label{eq:raw-herr-some-ainfty-map}
 G_{HF}^\infty:
 C^\bullet_{\varphi,\gamma}(D(\mathcal A))
 \rightsquigarrow F^\bullet.
\end{equation}
The construction is canonical as an equivalence of intrinsic $E_1$ objects and
is compatible with derived coefficient base change.  It does not assert that the
linear Taylor coefficient $(G_{HF}^\infty)_1$ is the direct Tate-dual map
$\Theta_{HF}^{\mathrm{raw}}$ of
Corollary~\ref{cor:raw-herr-all-degree-fox}, nor that its higher Taylor
coefficients are given by closed raw-period/native formulas.

Consequently there is no raw-Herr multiplicative existence problem.  The
stronger strict-representative question is whether one may prescribe the already
explicit linear term $\Theta_{HF}^{\mathrm{raw}}$ literally on the unreduced
Robba complex.  Theorem~\ref{thm:raw-herr-phantom-reduced-normalization} below
shows that this prescription is automatic modulo continuous phantoms, which is
the finite native information retained by the atlas.
\end{theorem}

\begin{proof}
The Herr comparison with continuous Galois cochains is compatible with cup
products and therefore identifies the two objects in the derived category of
$E_1$ coefficient algebras \cite{LiuCohomology,KPXFamilies}.  On the other side,
Theorem~\ref{thm:all-sylow-ainfty} transports the intrinsic $E_1$ structure of
continuous cochains to the all-Sylow perfect object and supplies bounded
$A_\infty$ representatives.  Composing these two equivalences gives
\eqref{eq:raw-herr-abstract-ainfty-zigzag}; choosing the finite Fox
representative gives~\eqref{eq:raw-herr-some-ainfty-map}.  Derived base-change
compatibility is inherited from the two comparison functors.  The final
qualification is immediate: an equivalence in the $E_1$ homotopy category fixes
neither a preferred linear chain representative nor closed formulas for its
higher Taylor components.
\end{proof}

\begin{definition}[Phantom-reduced raw-Herr linear class]\label{def:raw-herr-phantom-reduced-class}
Retain a Tate-duality chart of
Theorem~\ref{thm:raw-herr-all-degree-perfectification}.  Write
\[
 C_H:=C^\bullet_{\varphi,\gamma}(D(\mathcal A)),
 \qquad
 H_H:=H^\bullet_{\mathrm{Tate}}(\mathcal A),
\]
and let
\[
 i_H^\gamma:H_H\longrightarrow C_H
\]
be the chain map sending the chosen standard basis in degree $i$ to the raw
Herr cocycle frame $e_{i,a}$.  Then
$q_H^\gamma i_H^\gamma=\operatorname{id}_{H_H}$.  Let
$(H_F\mathop{\rightleftarrows}^{i_F}_{p_F}F,h_F)$ be the chosen finite Fox
contraction and let
$C_{HF}^{\mathrm{coh}}:H_H\xrightarrow{\sim}H_F$ be the canonical finite
Herr--Fox cohomology matrix.

For the continuous mapping complex put
\begin{equation}\label{eq:raw-herr-linear-shadow-map}
 \operatorname{sh}_1:
 H^0\underline{\operatorname{Hom}}_{A,\mathrm{cont}}(C_H,F)
 \longrightarrow
 \operatorname{Hom}^0_A\!\left(H^\bullet(C_H),H^\bullet(F)\right),
 \qquad
 [u]\longmapsto H^\bullet(u).
\end{equation}
This is intrinsic and well defined on continuous chain-homotopy classes.  In the
chosen Tate and Fox cohomology frames its matrix is exactly $p_Fu i_H^\gamma$.
Define the linear phantom submodule and the phantom-reduced native class module by
\begin{align}
 \operatorname{Ph}_1(C_H,F)
 &:=\ker(\operatorname{sh}_1),
 \label{eq:raw-herr-linear-phantom-submodule}\\
 \operatorname{Hom}^0_{\mathrm{nat}}(C_H,F)
 &:=
 H^0\underline{\operatorname{Hom}}_{A,\mathrm{cont}}(C_H,F)
 \big/\operatorname{Ph}_1(C_H,F).
 \label{eq:raw-herr-phantom-reduced-hom}
\end{align}
Thus two raw chain maps define the same native linear class exactly when they
induce the same cohomology morphism.  The quotient is therefore independent of
the chosen Tate frame, Fox contraction, and matrix coordinates; those choices
only display the intrinsic shadow in finite coordinates.  No continuous
splitting of the acyclic raw Herr complement is part of this definition.
\end{definition}

\begin{theorem}[Phantom-reduced native normalization is automatic]\label{thm:raw-herr-phantom-reduced-normalization}
Retain Definition~\ref{def:raw-herr-phantom-reduced-class}.  Choose the
representative-level multiplicative equivalence
\[
 G_{HF}^\infty:C_H\rightsquigarrow F
\]
of Theorem~\ref{thm:raw-herr-abstract-ainfty-existence} through the canonical
Herr--Galois--Fox comparison, so that its linear term induces
$C_{HF}^{\mathrm{coh}}$ on cohomology.  Then
\begin{equation}\label{eq:raw-herr-phantom-reduced-normalization}
 \boxed{
 [(G_{HF}^\infty)_1]
 =
 [\Theta_{HF}^{\mathrm{raw}}]
 \quad\text{in}\quad
 \operatorname{Hom}^0_{\mathrm{nat}}(C_H,F).}
\end{equation}
Equivalently,
\begin{equation}\label{eq:raw-herr-linear-difference-phantom}
 [(G_{HF}^\infty)_1-\Theta_{HF}^{\mathrm{raw}}]
 \in\operatorname{Ph}_1(C_H,F).
\end{equation}
Moreover, the class in
\eqref{eq:raw-herr-phantom-reduced-normalization} is compatible with change of
cyclotomic generator and with every coefficient base change on which the two
comparison maps are defined: after the standard identifications it is simply the
canonical Herr--Fox cohomology morphism.  After passage to the common finite
cohomology model, the $E_1/A_\infty$ products, Massey readouts, and, in
characteristic zero, the adjoint $L_\infty$ operations obtained from
$G_{HF}^\infty$ are the same intrinsic finite higher operations already used by
the native atlas.  Hence no phantom-vanishing assumption is required to
normalize the raw Herr comparison at the finite native level.

Exact equality $(G_{HF}^\infty)_1=\Theta_{HF}^{\mathrm{raw}}$ as maps on the unreduced Robba complex is a strictly stronger representative-gauge condition.  Remark~\ref{rem:raw-herr-exact-gauge-boundary} records this optional raw-topological refinement.  Its existence or failure changes only the chosen representative, not the intrinsic multiplicative equivalence or the finite native atlas.
\end{theorem}

\begin{proof}
By construction, the cohomology morphism of
$\Theta_{HF}^{\mathrm{raw}}$ is the canonical Herr--Fox comparison.  In the
chosen finite frames this reads
\[
 p_F\Theta_{HF}^{\mathrm{raw}}i_H^\gamma
 =p_Fi_FC_{HF}^{\mathrm{coh}}q_H^\gamma i_H^\gamma
 =C_{HF}^{\mathrm{coh}}.
\]
The chosen intrinsic comparison $G_{HF}^\infty$ induces the same canonical
Herr--Fox map on cohomology, hence its linear term has the same shadow.  Their
difference therefore lies in the intrinsic kernel of
\eqref{eq:raw-herr-linear-shadow-map}, proving
\eqref{eq:raw-herr-phantom-reduced-normalization}--
\eqref{eq:raw-herr-linear-difference-phantom}.  The shadow is the cohomology functor itself, so its quotient class is
preserved by the standard generator-change comparison and by derived base change
whenever the displayed chart survives.  Both multiplicative comparisons
represent the same intrinsic $E_1$ object $R\Gamma(G_K,\mathcal A)$.  Fixing the
common-source finite cohomology model therefore identifies their finite minimal
$A_\infty$ structures and all finite higher readouts up to the coherent
equivalence already built into the atlas.
The final statement is precisely the distinction between equality in the
phantom-reduced native class and equality of chosen continuous chain
representatives.
\end{proof}

\begin{remark}[Optional exact raw-Herr representative gauge]\label{rem:raw-herr-exact-gauge-boundary}
Theorem~\ref{thm:raw-herr-abstract-ainfty-existence} already gives a
representative-level multiplicative equivalence, while
Theorem~\ref{thm:raw-herr-phantom-reduced-normalization} identifies its linear
term with the direct Tate-dual map
$\Theta_{HF}^{\mathrm{raw}}$ after quotienting by continuous chain maps that are
invisible on cohomology.  Requiring literal equality of the two linear maps on
the unreduced Robba-ring complex is therefore a strictly stronger choice of raw
topological gauge.

Such a literal normalization would require additional continuous
chain-homotopy data on the acyclic raw-period directions.  Standard
$\varphi$--$\psi$ comparison removes an explicitly contractible
$\psi=0$ sector, and Iwasawa perfectness controls the remaining derived object
\cite{LiuCohomology,KPXFamilies}, but these facts do not by themselves choose a
global continuous strong deformation retract of the raw Robba complex.  No such
choice is needed below: the finite native atlas depends only on the canonical
phantom-reduced class and on the intrinsic $E_1/A_\infty$ equivalence.  We
therefore do not treat exact raw normalization as a further problem of this
paper.
\end{remark}


\begin{theorem}[Common-source chart-change strictification]\label{thm:common-source-chart-change}
Let $U$ be an affinoid chart, and let
\[
 \{C_\alpha^\bullet\}_{\alpha\in A}
\]
be a finite family of bounded complexes of finite projective $\mathcal O(U)$-modules.  Suppose that each $C_\alpha^\bullet$ is equipped with base-change-compatible isomorphisms
\[
 \lambda_{\alpha,i}:H^i(G_K,V)\xrightarrow{\ \sim\ }H^i(C_\alpha^\bullet),
\]
and that every $H^i(C_\alpha^\bullet)$ is finite locally free on $U$.  Put
\[
 H_\alpha:=\bigoplus_iH^i(C_\alpha^\bullet)[-i]
\]
and choose contractions
\begin{equation}\label{eq:tower-chart-contractions}
 \left(H_\alpha
 \mathop{\rightleftarrows}^{i_\alpha}_{p_\alpha}
 C_\alpha^\bullet,h_\alpha\right),
 \qquad
 p_\alpha i_\alpha=\operatorname{id}_{H_\alpha},
 \qquad
 \operatorname{id}_{C_\alpha}-i_\alpha p_\alpha
 =d h_\alpha+h_\alpha d.
\end{equation}
For $\alpha,\beta\in A$, define
\begin{align}
 \Psi_{\beta\alpha}
 &:=\bigoplus_i\lambda_{\beta,i}\lambda_{\alpha,i}^{-1}:
 H_\alpha\xrightarrow{\ \sim\ }H_\beta,\label{eq:tower-cohomology-transition}\\
 \Theta_{\beta\alpha}
 &:=i_\beta\Psi_{\beta\alpha}p_\alpha:
 C_\alpha^\bullet\longrightarrow C_\beta^\bullet.
 \label{eq:tower-chain-transition}
\end{align}
Then:
\begin{enumerate}[label=\textnormal{(\roman*)}]
\item $\Theta_{\beta\alpha}$ is a chain quasi-isomorphism, compatible with every base change on which the chosen contractions remain defined;
\item the transitions are strictly transitive:
\begin{equation}\label{eq:tower-strict-transitivity}
 \Theta_{\gamma\beta}\Theta_{\beta\alpha}
 =\Theta_{\gamma\alpha};
\end{equation}
\item $\Theta_{\alpha\alpha}=i_\alpha p_\alpha$ is chain-homotopic to the identity through $h_\alpha$, and $\Theta_{\alpha\beta}$ is a homotopy inverse of $\Theta_{\beta\alpha}$;
\item if a second collection of contractions gives transitions $\Theta'_{\beta\alpha}$, then the two strictifications are explicitly chain-homotopic.  Writing
\[
 \delta_{\beta\alpha}:=
 \Theta_{\beta\alpha}-\Theta'_{\beta\alpha},
\]
one may take
\begin{equation}\label{eq:chart-change-explicit-homotopy}
 \mathcal H_{\beta\alpha}
 :=h_\beta\delta_{\beta\alpha}
   +i_\beta p_\beta\delta_{\beta\alpha}h_\alpha,
 \qquad
 \delta_{\beta\alpha}
 =d\mathcal H_{\beta\alpha}
  +\mathcal H_{\beta\alpha}d.
\end{equation}
On determinantal minor charts, all entries of the chosen $\Theta_{\beta\alpha}$ and $\mathcal H_{\beta\alpha}$ are rational expressions in the comparison matrices and the inverted minors.
\end{enumerate}
\end{theorem}

\begin{proof}
Because $H_\alpha$ and $H_\beta$ have zero differential, Equation~\eqref{eq:tower-chain-transition} is a chain map and induces $\Psi_{\beta\alpha}$ on cohomology.  It is therefore a quasi-isomorphism.  The identity $p_\beta i_\beta=\operatorname{id}$ and the cocycle identity
$\Psi_{\gamma\beta}\Psi_{\beta\alpha}=\Psi_{\gamma\alpha}$ give
\[
 \Theta_{\gamma\beta}\Theta_{\beta\alpha}
 =i_\gamma\Psi_{\gamma\beta}p_\beta i_\beta
  \Psi_{\beta\alpha}p_\alpha
 =i_\gamma\Psi_{\gamma\alpha}p_\alpha.
\]
The contraction identity in \eqref{eq:tower-chart-contractions} proves (iii).

For (iv), the difference $\delta_{\beta\alpha}$ induces zero on cohomology, hence
$p_\beta\delta_{\beta\alpha}i_\alpha=0$.  Using that $\delta_{\beta\alpha}$ is a chain map and applying the two contraction identities gives
\begin{align*}
 d\mathcal H_{\beta\alpha}+\mathcal H_{\beta\alpha}d
 &=(\operatorname{id}-i_\beta p_\beta)\delta_{\beta\alpha}
   +i_\beta p_\beta\delta_{\beta\alpha}
    (\operatorname{id}-i_\alpha p_\alpha)\\
 &=\delta_{\beta\alpha}
   -i_\beta(p_\beta\delta_{\beta\alpha}i_\alpha)p_\alpha
 =\delta_{\beta\alpha}.
\end{align*}
The final assertion follows from Proposition~\ref{prop:determinantal-contraction}.
\end{proof}


\begin{theorem}[Common-source multiplicative chart atlas]\label{thm:common-source-multiplicative-tower}
Assume the setting of Theorem~\ref{thm:common-source-chart-change}, with $V$ replaced by a perfect coefficient dga $\mathcal A$ on which $G_K$ acts by dg-algebra automorphisms, and suppose that each $C_\alpha^\bullet$ represents the underlying perfect complex of
\[
 R\Gamma(G_K,\mathcal A).
\]
Write
\[
 H:=\bigoplus_i H^i(G_K,\mathcal A)[-i],
 \qquad
 \lambda_\alpha:=\bigoplus_i\lambda_{\alpha,i}:H\xrightarrow{\ \sim\ }H_\alpha.
\]
Choose a $K$-projective dga model $\mathcal B$ of the continuous cochain algebra and a contraction from $\mathcal B$ to $H$.  Homotopy transfer gives a minimal $A_\infty$ structure $m^H=(m_n^H)_{n\ge2}$ on $H$.  Transport it to $H_\alpha$ by
\begin{equation}\label{eq:tower-minimal-ainfty-transport}
 m_n^{(\alpha)}
 :=\lambda_\alpha\,m_n^H
   (\lambda_\alpha^{-1})^{\otimes n}.
\end{equation}
Then:
\begin{enumerate}[label=\textnormal{(\roman*)}]
\item every cohomology transition $\Psi_{\beta\alpha}$ of Equation~\eqref{eq:tower-cohomology-transition} is a strict $A_\infty$ isomorphism
\[
 (H_\alpha,m^{(\alpha)})
 \xrightarrow{\ \sim\ }
 (H_\beta,m^{(\beta)}),
\]
with no Taylor components above arity one, and these strict minimal transitions satisfy
\begin{equation}\label{eq:minimal-ainfty-tower-transitivity}
 \Psi_{\gamma\beta}\Psi_{\beta\alpha}
 =\Psi_{\gamma\alpha}
\end{equation}
as $A_\infty$ morphisms;
\item after the usual common mapping-cylinder stabilization, the contractions in Equation~\eqref{eq:tower-chart-contractions} extend to $A_\infty$ quasi-isomorphisms
\[
 I_\alpha^\infty:(H_\alpha,m^{(\alpha)})\longrightarrow C_\alpha^\bullet,
 \qquad
 P_\alpha^\infty:C_\alpha^\bullet\longrightarrow(H_\alpha,m^{(\alpha)}),
\]
whose first Taylor coefficients are $i_\alpha$ and $p_\alpha$.  Hence
\begin{equation}\label{eq:multiplicative-tower-transition}
 F_{\beta\alpha}
 :=I_\beta^\infty\circ\Psi_{\beta\alpha}\circ P_\alpha^\infty
 :C_\alpha^\bullet\longrightarrow C_\beta^\bullet
\end{equation}
is an $A_\infty$ quasi-isomorphism with
\begin{equation}\label{eq:multiplicative-tower-linear-term}
 (F_{\beta\alpha})_1
 =i_\beta\Psi_{\beta\alpha}p_\alpha
 =\Theta_{\beta\alpha};
\end{equation}
\item on every triple, $F_{\gamma\beta}\circ F_{\beta\alpha}$ and $F_{\gamma\alpha}$ are joined by an explicit $A_\infty$ homotopy obtained by inserting the transferred homotopy between
$P_\beta^\infty I_\beta^\infty$ and $\operatorname{id}_{H_\beta}$.  After choosing a coherent transfer convention, these homotopies extend to a homotopy-coherent diagram on the full tower-indexing nerve.  Its linear part is the strictly transitive atlas of Theorem~\ref{thm:common-source-chart-change};
\item every fixed Taylor component of $m^{(\alpha)}$, $I_\alpha^\infty$, $P_\alpha^\infty$, $F_{\beta\alpha}$, and every fixed coherence is a finite rooted- or leveled-tree expression in the common multiplication and the chosen contraction data.  On determinantal minor charts these formulas are rational in the inverted minors and commute with every base change preserving the chosen charts.  Different dg replacements, parametrices, contractions, and coherent transfer conventions give equivalent homotopy-coherent $A_\infty$ atlases.
\end{enumerate}
The same assertions hold for the dg-Lie algebra of adjoint cochains with $A_\infty$ replaced by $L_\infty$.  They also hold for perfect coefficient modules over the adjoint algebra, producing compatible $A_\infty$-module chart transitions.
\end{theorem}

\begin{proof}
Transfer the cup-product dga structure of $\mathcal B$ to the common cohomology complex $H$.  Equation~\eqref{eq:tower-minimal-ainfty-transport} transports one fixed minimal structure along the strict graded isomorphisms $\lambda_\alpha$.  Hence $\Psi_{\beta\alpha}=\lambda_\beta\lambda_\alpha^{-1}$ strictly intertwines every $m_n$, and \eqref{eq:minimal-ainfty-tower-transitivity} is the ordinary cocycle identity for the $\lambda_\alpha$.

After the stabilization used in Theorem~\ref{thm:all-sylow-ainfty}, homotopy transfer extends the linear contractions to the displayed $A_\infty$ quasi-isomorphisms.  Their arity-one composite is \eqref{eq:multiplicative-tower-linear-term}.  The comparison theorem also gives a homotopy from $P_\beta^\infty I_\beta^\infty$ to the identity; inserting it between two transition maps produces the triple-transition homotopy, and a cofibrant replacement of the indexing diagram supplies the higher coherences.  The rooted-tree formulas, base-change compatibility, uniqueness up to homotopy, and the dg-Lie and module variants are standard consequences of homotopy transfer and projective-resolution comparison \cite{BurkeTransfer,ManettiPerturbation,VokrinekHPT}.
\end{proof}


\begin{theorem}[Lax monoidal common-source chart atlas]\label{thm:common-source-lax-monoidal-tower}
Let $\mathcal M$ be a finite tensor-stable collection of perfect coefficient complexes with continuous $G_K$-action.  For every $M\in\mathcal M$ and every available chart $\alpha$, let $C_\alpha^\bullet(M)$ be a bounded finite-projective representative of $R\Gamma(G_K,M)$, equipped with common-source cohomology identifications and contractions as in Theorems~\ref{thm:common-source-chart-change} and~\ref{thm:common-source-multiplicative-tower}.  Assume that these choices are coefficient-functorial on the selected affinoid minor chart.

The external cup product gives
\begin{equation}\label{eq:common-source-external-cup}
 R\Gamma(G_K,M)\otimes^{\mathbf L}R\Gamma(G_K,N)
 \longrightarrow
 R\Gamma\!\left(G_K,M\otimes^{\mathbf L}N\right)
\end{equation}
and hence a graded map $\mu^H_{M,N}$ on cohomology.  For charts $\alpha,\beta,\gamma$ define
\begin{equation}\label{eq:minimal-lax-monoidal-chart-map}
 \mu^H_{\gamma;\alpha,\beta}
 :=\lambda_{\gamma,M\otimes N}\,\mu^H_{M,N}
 \bigl(\lambda_{\alpha,M}^{-1}\otimes
       \lambda_{\beta,N}^{-1}\bigr).
\end{equation}
Then:
\begin{enumerate}[label=\textnormal{(\roman*)}]
\item the minimal-model maps are strictly compatible with every chart transition:
\begin{equation}\label{eq:lax-monoidal-tower-naturality}
 \Psi^{M\otimes N}_{\gamma'\gamma}\,
 \mu^H_{\gamma;\alpha,\beta}
 =
 \mu^H_{\gamma';\alpha',\beta'}
 \bigl(\Psi^M_{\alpha'\alpha}\otimes
       \Psi^N_{\beta'\beta}\bigr).
\end{equation}
They satisfy the associativity, unit, and Koszul-symmetry identities inherited from the common-source cup product;
\item colored homotopy transfer produces an $A_\infty$-bilinear structure map
\[
 \mathfrak m^{\infty}_{\gamma;\alpha,\beta}:
 C_\alpha^\bullet(M)\otimes C_\beta^\bullet(N)
 \rightsquigarrow
 C_\gamma^\bullet(M\otimes N)
\]
whose binary linear component is
\begin{equation}\label{eq:lax-monoidal-chain-linear-term}
 \bigl(\mathfrak m^{\infty}_{\gamma;\alpha,\beta}\bigr)_{1,1}
 =i_{\gamma,M\otimes N}\,
   \mu^H_{\gamma;\alpha,\beta}
   \bigl(p_{\alpha,M}\otimes p_{\beta,N}\bigr);
\end{equation}
\item the squares obtained from Equation~\eqref{eq:lax-monoidal-tower-naturality} commute on the finite chain models up to explicit $A_\infty$ homotopies.  After choosing a coherent colored-transfer convention, the chart family becomes a homotopy-coherent lax symmetric-monoidal chart atlas.  Every fixed Taylor component and coherence is a finite rooted- or leveled-tree expression in the common cup product and the chosen contractions, and is rational in the inverse minors on a determinantal chart;
\item the same construction transports the coefficient evaluation maps
\[
 M^\vee\otimes M\longrightarrow\mathbf1
\]
and the composition maps for internal Homs.  On a Cayley--Hamilton splitting gerbe, scalar weights add under tensor product and change sign under duality.  Hence the transported structure maps are equivariant and descend to the derived Morita category; in particular the adjoint algebra, of scalar weight zero, carries a Morita-covariant lax symmetric-monoidal chart atlas.
\end{enumerate}
Changing parametrices, contractions, dg replacements, or coherent-transfer conventions changes the finite formulas by a coherent $A_\infty$ equivalence, not the transported lax monoidal object.
\end{theorem}

\begin{proof}
Continuous cochains are lax symmetric monoidal through the external cup product~\eqref{eq:common-source-external-cup}.  Transporting its cohomology maps along the strict graded isomorphisms $\lambda_{\alpha,M}$ gives Equation~\eqref{eq:minimal-lax-monoidal-chart-map}.  Equation~\eqref{eq:lax-monoidal-tower-naturality} is then the ordinary cancellation identity for the $\lambda$'s; the unit, associativity, and symmetry constraints are transported from the same common source.

Apply relative or colored homotopy transfer simultaneously to the coefficient objects, the external products, the adjoint algebra, and its modules.  The first transferred binary component is Equation~\eqref{eq:lax-monoidal-chain-linear-term}; the higher components and comparison homotopies are the standard colored rooted-tree formulas.  A cofibrant resolution of the chart and tensor indexing diagram supplies all higher coherences.  This proves (ii)--(iii) and the choice-independence statement \cite{BurkeTransfer,ManettiPerturbation,VokrinekHPT}.

Evaluation, duality, and internal-Hom composition are morphisms in the same symmetric monoidal coefficient category, so the identical transport argument proves (iv).  On the splitting gerbe, tensor products add scalar weights and duals negate them.  Proposition~\ref{prop:all-sylow-morita} and Proposition~\ref{prop:ch-lax-monoidal} therefore identify the transported maps with morphisms in the derived Morita category \cite{ToenAzumaya}.
\end{proof}

\section{Finite Fox--Spencer carriers}\label{sec:fox-spencer}

A finite profinite presentation turns transport into finite syntax, and completed Fox derivatives provide its noncommutative Jacobian.  The central question then shifts from finite existence to explicit and functorial comparison among the resulting charts.

\subsection{Why the tautological full-transport model is insufficient}

Retaining $V$ together with its complete $\Lambda_{K,E}$-action and bar construction is faithful but tautological: the transport map
\[
\Lambda_{K,E}\longrightarrow \End_E(V)
\]
has simply been retained as primitive input.

\begin{warning}[Degree-zero bar data are not a finite closure]\label{warn:no-degreezero-tautology}
A finite-dimensional representation together with its original $G_K$-action is not, by itself, a nontrivial finite-atlas closure.  A nontrivial closure must exhibit a finite \emph{transport resolution}: finitely many covariant transport cells and relations, independent of $V$, whose dual differential recovers the bar/Spencer differential and whose higher transferred operations recover the derived deformation theory.
\end{warning}

The nontrivial object is instead a finite transport resolution.  Write
\[
\Lambda_G^\circ:=\mathcal O_E[[G]],
\qquad
\Lambda_{G,E}:=\Lambda_G^\circ[1/\varpi_E].
\]
An \emph{integral finite transport resolution} is a bounded complex $(P_G^\circ)^\bullet$ of finite projective $\Lambda_G^\circ$-modules together with an augmentation
\[
(P_G^\circ)^\bullet\longrightarrow \mathcal O_E
\]
which is a quasi-isomorphism in the continuous derived category.  Its generic fibre
\[
P_{G,E}^\bullet:=(P_G^\circ)^\bullet[1/\varpi_E]
\longrightarrow E
\]
is then a finite transport resolution over $\Lambda_{G,E}$.  It produces, functorially in a continuous $E$-representation $V$,
\[
C_{P}^\bullet(G,V):=\operatorname{Hom}_{\Lambda_{G,E}}(P_{G,E}^\bullet,V)
\simeq R\Gamma(G,V).
\]
The finite cells of $(P_G^\circ)^\bullet$ are the finite-atlas data.  Fox complexes provide integral examples; generic fibres are obtained by inverting $\varpi_E$.

\subsection{Finite presentation of a local absolute Galois group}

If $n=[K:\Qp]$, then $G_K$ has minimal generator number $n+2$ and relation number $n+1$ \cite{JannsenWingberg,JardenShusterman,NeukirchLocal}; choose
\begin{equation}\label{eq:finite-profinite-presentation}
1\longrightarrow R_K\longrightarrow \widehat F_e
\longrightarrow G_K\longrightarrow1,
\qquad
e=n+2,
\end{equation}
where $R_K$ is the smallest closed normal subgroup containing finitely many relations
\[
r_1,\ldots,r_m,
\qquad m=n+1.
\]
Jannsen--Wingberg give more explicit marked tame--wild coordinates at odd residue characteristic.

\begin{remark}[Existence is not generally an explicit Fox Jacobian]\label{rem:nonconstructive-presentation}
For a general finite extension $K/\mathbf Q_p$, finite presentability is not a
canonical computable list of relation words.  The proof in
\cite{JardenShusterman} uses a cohomological criterion rather than producing
distinguished generators and relations.  The Jannsen--Wingberg description in
odd residue characteristic gives more explicit marked tame--wild coordinates,
but its marked pro-$p$ condition is not itself an ordinary finite word relation.
The residual strictification theorem below explains why such marked conditions
become automatic on an admissible residual deformation sector.

No explicit full-$G_K$ dyadic word presentation is needed for the finite-atlas
theorems.  At $p=2$ the all-Sylow construction restricts to open subgroups whose
maximal pro-$2$ quotients are the standard free or dyadic Demu\v{s}kin groups;
their finite one-relator Fox resolutions supply the coefficient charts used in
the closure and globalization arguments.  Thus the mainline uses finite
presentability for full transport and standard pro-$p$ structure after Sylow
reduction, rather than a special marked word normal form for $G_{\mathbf Q_2}$.
\end{remark}

\subsection{Residual automaticity}\label{subsec:marked-residual-automaticity}

Although the marked pro-$p$ condition is not a word relation, it becomes automatic on a residual sector whose wild image is a $p$-group.

\begin{theorem}[Residual strictification of a marked pro-$p$ presentation]\label{thm:marked-residual-strictification}
Let $\Gamma^{\mathrm{mark}}$ be a profinite group described by bounded generators
\[
 y_1,\ldots,y_e
\]
subject to finitely many word relations $r_1,\ldots,r_m$.  In addition, require the closed normal subgroup generated by a specified subset
\[
 y_{i_1},\ldots,y_{i_s}
\]
to be pro-$p$.  Let
\[
 \bar\rho:\Gamma^{\mathrm{mark}}\longrightarrow \mathrm{GL}_r(k_E)
\]
be a continuous residual representation.  Let $A$ be a complete Noetherian local $\mathcal O_E$-algebra with residue field $k_E$.

For every tuple of lifts
\[
 Y_1,\ldots,Y_e\in\mathrm{GL}_r(A)
\]
of the residual generator matrices which satisfies the word relations $r_j(Y_1,\ldots,Y_e)=I$, put
\[
 \Gamma_A:=\overline{\langle Y_1,\ldots,Y_e\rangle}
 \subseteq\mathrm{GL}_r(A).
\]
Then the closed normal subgroup of $\Gamma_A$ generated by $Y_{i_1},\ldots,Y_{i_s}$ is pro-$p$.  Consequently the framed deformation functor of $\bar\rho$ as a representation of $\Gamma^{\mathrm{mark}}$ is equal to the formal completion, at the residual tuple, of the ordinary finite word-relation locus.
\end{theorem}

\begin{proof}
Let
\[
 \Gamma_A=\overline{\langle Y_1,\ldots,Y_e\rangle}
 \subseteq\mathrm{GL}_r(A)
\]
and let $W_A\triangleleft\Gamma_A$ be the closed normal subgroup generated by the lifted marked matrices.  Reduction identifies the residual image of $\Gamma_A$ with $\bar\rho(\Gamma^{\mathrm{mark}})$.  Hence the reduction $\overline W$ is contained in the normal closure, inside $\bar\rho(\Gamma^{\mathrm{mark}})$, of the marked residual matrices.  This normal closure is exactly the image of the marked pro-$p$ normal subgroup, and is therefore a finite $p$-group.  The kernel of
\[
 \mathrm{GL}_r(A)\longrightarrow\mathrm{GL}_r(k_E)
\]
is pro-$p$: it is the inverse limit of the congruence kernels over the Artinian quotients $A/\mathfrak m_A^n$, whose successive congruence quotients are finite additive groups of characteristic $p$.  Therefore
\[
 1\longrightarrow
 W_A\cap\ker\bigl(\mathrm{GL}_r(A)\to\mathrm{GL}_r(k_E)\bigr)
 \longrightarrow W_A\longrightarrow\overline W\longrightarrow1
\]
exhibits $W_A$ as an extension of pro-$p$ groups, hence as pro-$p$.  Every finite quotient of the compact image generated by the $Y_i$ is therefore an admissible marked quotient: it satisfies the word relations and has $p$-group wild normal closure.  By the defining inverse-limit universal property of the marked quotient, the resulting homomorphism from the free profinite group factors continuously through $\Gamma^{\mathrm{mark}}$.  Thus the marked condition imposes no additional requirement on a lift of the fixed residual tuple beyond the displayed word relations.  Equality of the represented framed deformation functors follows.
\end{proof}

\begin{corollary}[Marked full-group word charts on odd residual sectors]\label{cor:marked-full-group-residual-word-chart}
In odd residue characteristic, the marked Jannsen--Wingberg presentation gives,
on every genuine residual representation sector to which the marked chart
applies, an ordinary finite word-equation formal chart for framed $p$-adic
deformations.  The marked pro-$p$ normal-subgroup condition contributes no
additional equation after completion at the residual tuple.

In these charts the differential of the relation map is the evaluated Fox
Jacobian.  The dyadic branch is not asserted to admit the same full-group marked
word coordinates; it is handled instead by the all-Sylow/free--Demu\v{s}kin
coefficient charts used later.
\end{corollary}

\begin{proof}
For a genuine residual representation in the marked odd-prime chart, the image
of the marked wild normal subgroup is a finite $p$-group because the source is
pro-$p$.  Apply Theorem~\ref{thm:marked-residual-strictification}.  Linearizing
the surviving finite word relations is exactly completed Fox differentiation.
\end{proof}

\begin{remark}[An optional marked full-$G_{\mathbf Q_2}$ residual chart]
\label{rem:q2-marked-full-group-source}
Roe--Turturean give a marked profinite presentation of
$G_{\mathbf Q_2}$ with four generators, two completed-word
relations, and the requirement that the normal closure
of its designated wild generators be pro-$2$
\cite{RoeTurtureanGQ2}.
Taking that marked presentation as source,
Theorem~\ref{thm:marked-residual-strictification} makes
the pro-$2$ condition automatic on each genuine residual
framed sector.  It consequently supplies a finite list
of formal completed-word equations, and
Theorem~\ref{thm:word-lci-reduction} gives their
local syzygy-complete lci/Koszul replacement.
The completed-word qualifier matters: profinite powers
need not be ordinary polynomial words on a global
generic-matrix chart.  This optional $K=\mathbf Q_2$
full-group source becomes finite-order matrix-computable by
Theorem~\ref{thm:dyadic-q2-marked-word-jet-solver} and is not used in proving the
all-Sylow dyadic coefficient atlas and does not
assert a uniform marked presentation for every
finite dyadic extension.
\end{remark}

\begin{warning}[Residual strictification is not derived Fox completeness]\label{warn:q2-marked-not-resolution}
Theorem~\ref{thm:marked-residual-strictification} is a statement about the classical framed deformation functor.  It does not turn the displayed word relators into a projective resolution of the trivial completed-group-algebra module, and it does not identify their naive Koszul derived zero locus with the intrinsic derived representation stack.  Redundant relations may create spurious derived directions.  Theorem~\ref{thm:word-lci-reduction} removes this formal syzygy-existence gap by replacing the redundant word list with a minimal regular sequence, but it still does not produce a full completed-group-algebra transport resolution.
\end{warning}

\subsection{Formal lci word charts}\label{subsec:word-lci-reduction}

The naive derived zero locus is inadequate, but after residual strictification the lci theorem supplies a finite syzygy-complete formal replacement.

\begin{theorem}[Lci reduction of a residual full-group word chart]\label{thm:word-lci-reduction}
Let $\bar\rho:G_K\to\mathrm{GL}_r(k_E)$ be a genuine residual representation lying in a marked full-group chart for which Theorem~\ref{thm:marked-residual-strictification} applies.  Let
\[
 A_{\mathrm{word}}
 =\mathcal O_E[[z_1,\ldots,z_N]]
\]
be the formally smooth framed matrix-coordinate ring, let
\[
 F_1,\ldots,F_M\in A_{\mathrm{word}}
\]
be the scalar matrix entries of the finitely many word relations, and put
\[
 I_{\mathrm{word}}=(F_1,\ldots,F_M),
 \qquad
 R_{\mathrm{word}}=A_{\mathrm{word}}/I_{\mathrm{word}}.
\]
Assume, as supplied by Corollary~\ref{cor:marked-full-group-residual-word-chart}, that $R_{\mathrm{word}}$ represents the framed deformation functor of $\bar\rho$.  Then:
\begin{enumerate}[label=\textnormal{(\roman*)}]
\item $I_{\mathrm{word}}$ is a complete-intersection ideal;
\item if
\[
 h=\dim_{k_E}
 \frac{I_{\mathrm{word}}}
      {\mathfrak m_{A_{\mathrm{word}}}I_{\mathrm{word}}},
\]
there is a matrix $U\in M_{h\times M}(A_{\mathrm{word}})$ such that, writing
\[
 (g_1,\ldots,g_h)^{\mathrm t}
 =U(F_1,\ldots,F_M)^{\mathrm t},
\]
the elements $g_1,\ldots,g_h$ generate $I_{\mathrm{word}}$ and form an $A_{\mathrm{word}}$-regular sequence;
\item the Koszul dg algebra
\[
 K_{\mathrm{word}}^{\mathrm{lci}}
 :=K_{A_{\mathrm{word}}}(g_1,\ldots,g_h)
\]
is a finite free syzygy-complete resolution of $R_{\mathrm{word}}$;
\item let
\[
 q_F:R_{\mathrm{word}}^{M}
 \longrightarrow
 I_{\mathrm{word}}/I_{\mathrm{word}}^2
\]
be the conormal map sending the $i$th basis vector to the class of $F_i$.  The regular sequence identifies
\[
 R_{\mathrm{word}}^{h}
 \xrightarrow{\sim}
 I_{\mathrm{word}}/I_{\mathrm{word}}^2,
 \qquad e_j\longmapsto[g_j].
\]
Modulo $I_{\mathrm{word}}$, its Jacobian is obtained from the evaluated word/Fox Jacobian by finite row reduction:
\begin{equation}\label{eq:word-lci-jacobian-reduction}
 J_g=\overline U\,J_F.
\end{equation}
After choosing a splitting of $q_F$, one has
\[
 R_{\mathrm{word}}^{M}
 \cong
 K_{\mathrm{rel}}\oplus R_{\mathrm{word}}^{h},
 \qquad K_{\mathrm{rel}}:=\ker(q_F),
\]
and the word Jacobian becomes
\begin{equation}\label{eq:word-jacobian-splitting}
 J_F=(0,J_g).
\end{equation}
Thus the naive relation-list cotangent complex is the intrinsic lci cotangent complex plus the redundant summand $K_{\mathrm{rel}}[1]$.  The latter is exactly the spurious derived direction created by the nonminimal word list; deleting it gives the syzygy-complete Koszul replacement.
\end{enumerate}
\end{theorem}

\begin{proof}
Theorem~\ref{thm:marked-residual-strictification} and its corollary identify $R_{\mathrm{word}}$ with the framed local Galois deformation ring.  The local complete-intersection theorem of B\"ockle--Iyengar--Pa\v{s}k\=unas, recalled in Theorem~\ref{thm:local-lci}, therefore implies that the quotient of the regular local ring $A_{\mathrm{word}}$ by $I_{\mathrm{word}}$ is a complete intersection.  In a Noetherian local ring, an ideal generated by a regular sequence has the property that every minimal generating set is a regular sequence \cite{StacksLCI}.  Choose $h$ $A_{\mathrm{word}}$-linear combinations of the $F_i$ whose classes form a basis of
\[
 I_{\mathrm{word}}/
 \mathfrak m_{A_{\mathrm{word}}}I_{\mathrm{word}}.
\]
Nakayama's lemma shows that the resulting $g_1,\ldots,g_h$ generate $I_{\mathrm{word}}$, and the complete-intersection criterion makes them a regular sequence.  Their Koszul complex is therefore a finite free resolution of the quotient \cite{StacksKoszul}.

Write $g=UF$.  Differentiating gives
\[
 dg=U\,dF+(dU)F,
\]
and the second term vanishes after passage to $R_{\mathrm{word}}$, giving \eqref{eq:word-lci-jacobian-reduction}.  The classes of the $g_j$ form a basis of the conormal module, so the displayed map $R_{\mathrm{word}}^h\to I_{\mathrm{word}}/I_{\mathrm{word}}^2$ is an isomorphism.  Since its target is free, the surjection $q_F$ splits.  If a coefficient vector lies in $K_{\mathrm{rel}}$, its linear combination of the $F_i$ lies in $I_{\mathrm{word}}^2$; differentiating shows that the corresponding row combination of $J_F$ vanishes modulo $I_{\mathrm{word}}$.  This proves \eqref{eq:word-jacobian-splitting}.  The extra summand is therefore precisely the redundant relation module, not part of the intrinsic cotangent complex.
\end{proof}

\begin{corollary}[Single-chart transport-explicit lci row reduction]\label{cor:local-transport-explicit-lci-row-reduction}
In the complete local setting of Theorem~\ref{thm:word-lci-reduction}, no
conormal-minor stratification is needed.  There is an $h$-element subset
\[
 J=\{j_1,\ldots,j_h\}\subset\{1,\ldots,M\}
\]
such that the original word equations
\[
 g_a:=F_{j_a},\qquad 1\le a\le h,
\]
generate $I_{\mathrm{word}}$ and form an $A_{\mathrm{word}}$-regular
sequence.  Every redundant equation has an exact expression
\begin{equation}\label{eq:local-redundant-word-row-expression}
 F_k=\sum_{a=1}^{h}c_{ka}g_a,
 \qquad c_{ka}\in A_{\mathrm{word}},
\end{equation}
and modulo $I_{\mathrm{word}}$ its Jacobian row satisfies
\begin{equation}\label{eq:local-transport-explicit-lci-row-reduction}
 dF_k=\sum_{a=1}^{h}\overline c_{ka}\,dg_a.
\end{equation}
Thus at a fixed residual point the intrinsic lci Jacobian is obtained from
one globally selected set of actual word/Fox rows; there is no local
\v{C}ech--Stasheff issue attached merely to choosing the regular sequence.
\end{corollary}

\begin{proof}
The classes of the original generators $F_1,\ldots,F_M$ span
\[
 I_{\mathrm{word}}/
 \mathfrak m_{A_{\mathrm{word}}}I_{\mathrm{word}}.
\]
Choose $h$ of them whose classes form a $k_E$-basis.  Nakayama's lemma
shows that the selected $g_a$ generate $I_{\mathrm{word}}$ itself.  Since
$I_{\mathrm{word}}$ is a complete-intersection ideal of height $h$, this
minimal generating set is a regular sequence.  Exact coefficients in
\eqref{eq:local-redundant-word-row-expression} therefore exist already over
$A_{\mathrm{word}}$, with no localization or refinement.  Differentiating
that identity and reducing modulo $I_{\mathrm{word}}$ gives
\eqref{eq:local-transport-explicit-lci-row-reduction}.
\end{proof}

\begin{corollary}[Formal syzygy-complete marked full-group charts]\label{cor:formal-full-group-syzygy-complete}
For every genuine residual sector in an odd-residue-characteristic
Jannsen--Wingberg marked word chart, the explicit full-group word equations
admit a finite syzygy-complete formal lci/Koszul enhancement.

Thus the marked full-group word branch contains no existence question for
higher formal syzygies.  The comparison with the intrinsic full-transport,
Fox, Herr, and Cayley--Hamilton carriers is supplied later at the finite-native
level.  The dyadic branch does not require such a full-group word chart: its
coefficient carrier is obtained after all-Sylow reduction from the standard
free/dyadic-Demu\v{s}kin pro-$2$ resolutions.
\end{corollary}

\begin{warning}[Lci reduction is not a full transport resolution]\label{warn:lci-word-not-group-resolution}
The Koszul complex $K_{\mathrm{word}}^{\mathrm{lci}}$ resolves the \emph{framed deformation ring} inside its formally smooth matrix-coordinate ring.  It does not resolve the trivial module over $\mathcal O_E[[G_K]]$, and it does not by itself reconstruct the full continuous transport of an arbitrary coefficient module.  The theorem removes the formal syzygy-existence gap of the word chart; it does not remove the need for an explicit comparison with a full-transport coefficient resolution.
\end{warning}

\begin{corollary}[Scope of marked full-$G_K$ word explicitization]\label{cor:q2-revised-full-gk-boundary}
In odd residue characteristic, Jannsen--Wingberg marked tame--wild coordinates
provide the full-group word branch described above.  On each genuine residual
sector the marked pro-$p$ condition is automatic, Theorem~\ref{thm:word-lci-reduction}
supplies a finite syzygy-complete lci/Koszul enhancement, and
Corollary~\ref{cor:local-transport-explicit-lci-row-reduction} selects actual
word/Fox rows for the regular sequence at the complete local point.  The
nonlocal conormal-minor stratification belongs to the finite-type
Cayley--Hamilton algebraization and is supplied later by
Theorem~\ref{thm:transport-explicit-lci-row-reduction}.

The pro-$p$ free-frame, blockwise/Morita descent, and recursive-globalization
theorems later carry out the invariant comparison tasks without requiring a
uniform marked full-group word chart.  In particular no uniform dyadic full-$G_K$ marked word normal form for
arbitrary finite $K/\mathbf Q_2$ is assumed; the special
$K=\mathbf Q_2$ input of Remark~\ref{rem:q2-marked-full-group-source}
is optional.  This is compatible with finite-atlas
closure, which uses finite presentability for the full transport layer and
all-Sylow descent to free or dyadic Demu\v{s}kin coefficient charts.
\end{corollary}

After conjugation, a framed rank-$r$ representation preserves an $\mathcal O_E$-lattice and is represented by
\[
(X_1,\ldots,X_e)\in \mathrm{GL}_r(\mathcal O_E)^e
\]
satisfying the finitely many profinite relation equations
\[
r_j(X_1,\ldots,X_e)=I_r,
\qquad 1\leq j\leq m.
\]
The integral target is required for profinite word evaluation; isomorphisms remain simultaneous $\mathrm{GL}_r(E)$-conjugacies.

\begin{proposition}[Bounded finite transport coordinates]\label{prop:finite-transport-coordinates}
For fixed $K$, $E$, and rank $r$, the groupoid of continuous $E$-representations of $G_K$ is equivalent to the following groupoid:
\begin{itemize}
\item objects are tuples $(X_1,\ldots,X_e)\in\mathrm{GL}_r(\mathcal O_E)^e$ satisfying the relations in \eqref{eq:finite-profinite-presentation};
\item a morphism from $(X_i)$ to $(X_i')$ is an $S\in\mathrm{GL}_r(E)$ such that $SX_i=X_i'S$ for every $i$.
\end{itemize}
Equivalently, every continuous representation is $E$-conjugate to a bounded integral solution, and two such solutions define isomorphic $E$-representations exactly when they are simultaneously $\mathrm{GL}_r(E)$-conjugate.
\end{proposition}

\begin{proof}
The image of a continuous representation is compact, so it preserves an $\mathcal O_E$-lattice after a change of basis.  Thus the representation is conjugate to one with image in the profinite group $\mathrm{GL}_r(\mathcal O_E)$.  The universal property of $\widehat F_e$ now gives a unique continuous homomorphism from the chosen integral generator images.  It factors through $G_K$ exactly when the closed normal subgroup generated by the $r_j$ lies in its kernel.  Finally, an isomorphism of the underlying $E$-representations is precisely an arbitrary $S\in\mathrm{GL}_r(E)$ intertwining all generator matrices.
\end{proof}

\begin{remark}[Why boundedness is necessary]\label{rem:bounded-profinite-coordinates}
The same statement with objects in all of $\mathrm{GL}_r(E)^e$ is false.  Already for $\widehat F_1=\widehat{\mathbf Z}$, the matrix $pI_r$ defines a homomorphism from the discrete group $\mathbf Z$ but its powers are unbounded, so it cannot be the image of a continuous homomorphism from the compact group $\widehat{\mathbf Z}$.  Passing to a stable lattice is therefore part of the mathematical statement, not a harmless choice of coordinates.
\end{remark}

Bounded coordinates therefore already exist.  The substantive issue is to equip them with a covariant differential and derived calculus.

\subsection{Completed Fox calculus}

Let $F$ be free pro-$p$ on $x_1,\ldots,x_d$.  Continuous Fox derivatives
\[
\frac{\partial}{\partial x_i}:\mathcal O_E[[F]]
\longrightarrow \mathcal O_E[[F]]
\]
are characterized by
\[
\frac{\partial x_j}{\partial x_i}=\delta_{ij},
\qquad
\frac{\partial(uv)}{\partial x_i}
=\frac{\partial u}{\partial x_i}
+u\frac{\partial v}{\partial x_i},
\]
and continuous extension from the dense group algebra.  Their fundamental identity is
\begin{equation}\label{eq:fox-fundamental}
w-1
=\sum_{i=1}^{d}
\frac{\partial w}{\partial x_i}(x_i-1).
\end{equation}
This is the noncommutative first-order Taylor identity \cite{Fox}.

Let
\[
P=F/\overline{\langle\!\langle r\rangle\!\rangle}
\]
be a one-relator pro-$p$ group, let $\Lambda_P^\circ=\mathcal O_E[[P]]$, and write a bar over a Fox derivative for its image in $\Lambda_P^\circ$.  Define
\begin{align}
\partial_1:(\Lambda_P^\circ)^d&\longrightarrow\Lambda_P^\circ,
&
\partial_1(a_1,\ldots,a_d)
&=\sum_{i=1}^{d}a_i(x_i-1),\label{eq:fox-boundary1}\\
\partial_2:\Lambda_P^\circ&\longrightarrow(\Lambda_P^\circ)^d,
&
\partial_2(a)
&=\left(a\,\overline{\frac{\partial r}{\partial x_i}}\right)_{i=1}^{d}.
\label{eq:fox-boundary2}
\end{align}
Equation \eqref{eq:fox-fundamental} gives
\[
\partial_1\partial_2(a)=a(r-1)=0.
\]
Thus the defining relation supplies a finite chain complex
\begin{equation}\label{eq:fox-chain-complex}
0\longrightarrow\Lambda_P^\circ
\xrightarrow{\partial_2}
(\Lambda_P^\circ)^d
\xrightarrow{\partial_1}
\Lambda_P^\circ
\xrightarrow{\epsilon}
\mathcal O_E
\longrightarrow0.
\end{equation}

\begin{proposition}[Fox identity as discrete Bianchi identity]\label{prop:fox-bianchi}
The equality $\partial_1\partial_2=0$ in \eqref{eq:fox-chain-complex} is precisely the linearized compatibility of the relation $r=1$.  After dualizing against a representation, it becomes $d^1d^0=0$ in the finite arithmetic Spencer complex.
\end{proposition}

\begin{proof}
The first assertion is the fundamental formula \eqref{eq:fox-fundamental} evaluated at $r$.  The second follows by applying $\operatorname{Hom}_{\Lambda_P^\circ}(-,T)$ to \eqref{eq:fox-chain-complex}.
\end{proof}

\subsection{The free and Demu\v{s}kin sectors}

For a $p$-adic local field $K$, the maximal pro-$p$ quotient
\[
P_K:=\Gal(K(p)/K)
\]
is free pro-$p$ of rank $[K:\Qp]+1$ when $\mu_p\not\subset K$, and is a Demu\v{s}kin group of rank $[K:\Qp]+2$ when $\mu_p\subset K$ \cite{JardenShusterman,NSW}.  A finite-rank Demu\v{s}kin group is a one-relator pro-$p$ Poincar\'e duality group of dimension two.  In that case \eqref{eq:fox-chain-complex} is the standard finite free resolution of the trivial module; in the free case the leftmost term is absent.

Let $T$ be a finite free $\mathcal O_E$-module with continuous $P_K$-action.  Applying $\operatorname{Hom}_{\Lambda_{P_K}^\circ}(-,T)$ gives
\begin{equation}\label{eq:fox-spencer-cochains}
C_{\mathrm{Fox}}^\bullet(P_K,T):
\qquad
T\xrightarrow{d^0}T^d\xrightarrow{d^1}T,
\end{equation}
where
\begin{align}
 d^0(t)&=((x_i-1)t)_{i=1}^{d},\label{eq:fox-d0}\\
 d^1(t_1,\ldots,t_d)
 &=\sum_{i=1}^{d}
 \overline{\frac{\partial r}{\partial x_i}}\,t_i.
\label{eq:fox-d1}
\end{align}
In the free case, $d^1$ and the final copy of $T$ are omitted.

\begin{theorem}[Fox--Spencer closure for the maximal pro-$p$ sector]\label{thm:fox-spencer-demushkin}
Let $P_K$ be the maximal pro-$p$ quotient of a $p$-adic local Galois group and let $T$ be a finite free $\mathcal O_E$-representation of $P_K$.
\begin{enumerate}
\item If $P_K$ is free pro-$p$, the two-term Fox complex computes $R\Gamma(P_K,T)$.
\item If $P_K$ is Demu\v{s}kin, the three-term complex \eqref{eq:fox-spencer-cochains} computes $R\Gamma(P_K,T)$.
\item After inverting $\varpi_E$, the same statements hold for $V=T[1/\varpi_E]$.
\end{enumerate}
The resulting complex is finite, tower-free, and constructed from the full pro-$p$ transport presentation rather than from a period-ring chart.
\end{theorem}

\begin{proof}
In the free case, the augmentation ideal has the standard length-one free resolution.  In the Demu\v{s}kin case, the minimal one-relator presentation and Poincar\'e duality dimension two give the exact finite free resolution \eqref{eq:fox-chain-complex}.  Applying $\operatorname{Hom}_{\Lambda_{P_K}^\circ}(-,T)$ computes continuous group cohomology.  Localization at $\varpi_E$ is exact.
\end{proof}

\begin{corollary}[Finite tangent--obstruction complex]\label{cor:fox-deformation}
For a representation $\rho:P_K\to\mathrm{GL}_E(V)$ in the Demu\v{s}kin sector, the tangent and obstruction spaces are computed by
\begin{equation}\label{eq:fox-adjoint-complex}
\ad V
\xrightarrow{d^0_\rho}
(\ad V)^d
\xrightarrow{d^1_\rho}
\ad V,
\end{equation}
where $x_i$ acts on $\ad V$ by $\operatorname{Ad}(\rho(x_i))$ and $d^1_\rho$ is the evaluated Fox Jacobian of $r$.  Hence
\[
H^1(P_K,\ad V)=\frac{\ker d^1_\rho}{\operatorname{im}d^0_\rho},
\qquad
H^2(P_K,\ad V)=\operatorname{coker}d^1_\rho.
\]
\end{corollary}

Here $(\ad V)^d$ is the space of first-order generator variations, $\operatorname{im}d^0_\rho$ is infinitesimal conjugation, and $d^1_\rho$ is the relation symbol.

\subsection{The quadratic Fox symbol and the explicit cup product}

The preceding complex gives an explicit differential.  In the one-relator pro-$p$ sector it also gives the first nontrivial multiplicative operation.  Let $F$ be the free pro-$p$ group on $x_1,\ldots,x_d$, let
\[
\varepsilon:\mathbf F_p[[F]]\longrightarrow\mathbf F_p
\]
be augmentation, and let $\partial_i$ denote the continuous Fox derivative with respect to $x_i$.  Normalize the completed Magnus expansion by
\[
\mathcal M(x_i)=1+X_i.
\]
For $w\in F$ put
\begin{equation}\label{eq:second-fox-coeff}
\varepsilon_{ij}(w)
:=
[X_iX_j]\,\mathcal M(w)
=
\varepsilon\!\left(\partial_{ij}w\right)
\in\mathbf F_p,
\end{equation}
where $\partial_{ij}$ denotes the standard ordered second Fox derivative corresponding to the monomial $X_iX_j$.  Thus \eqref{eq:second-fox-coeff} fixes the derivative-order convention and identifies the coefficient intrinsically.


\begin{theorem}[Twisted tensor Fox formula for the binary coefficient product]\label{thm:twisted-tensor-fox-cup}
Let
\[
 P=\langle x_1,\ldots,x_d\mid r_1,\ldots,r_m\rangle_{\mathrm{pro}\text{-}p}
\]
be a pro-$p$ presentation whose completed Fox $2$-complex computes continuous
cohomology through degree two; in particular this applies to the free and
Demu\v{s}kin charts used above.  Let $A$ be a complete coefficient algebra,
let $M,N,Q$ be finite projective $A$-modules with continuous $P$-actions, and
let
\[
 \mu:M\widehat\otimes_A N\longrightarrow Q
\]
be a continuous $P$-equivariant pairing.  Put $\Lambda_F=A[[F]]$ for the free
pro-$p$ group $F$ on the $x_i$, let
\[
 \Delta_F:\Lambda_F\longrightarrow
 \Lambda_F\widehat\otimes_A\Lambda_F,
 \qquad g\longmapsto g\otimes g,
\]
be the completed Hopf diagonal, and let $\partial_i$ be the completed Fox
derivatives.  Define the tensor Fox symbols
\begin{equation}\label{eq:tensor-fox-symbol}
 \Xi_{ij}(w)
 :=
 (\partial_i\widehat\otimes 1)\,\Delta_F(\partial_j w)
 \in
 \Lambda_F\widehat\otimes_A\Lambda_F.
\end{equation}
Write $\overline{\Xi}_{ij}(w)$ for their images in
$A[[P]]\widehat\otimes_A A[[P]]$.  If
$a\in Z^1_{\mathrm{cont}}(P,M)$ and
$b\in Z^1_{\mathrm{cont}}(P,N)$, put
\[
 A_i:=a(x_i),\qquad B_j:=b(x_j).
\]
Then the coefficient cup product, evaluated on the relation $2$-cell
$e^2_\alpha$ corresponding to $r_\alpha$, is
\begin{equation}\label{eq:twisted-tensor-fox-cup}
 \left\langle a\smile_\mu b,e^2_\alpha\right\rangle
 =
 \mu\!\left(
 \sum_{i,j=1}^{d}
 (\rho_M\widehat\otimes\rho_N)
 \bigl(\overline{\Xi}_{ij}(r_\alpha)\bigr)
 (A_i\otimes B_j)
 \right).
\end{equation}
Here $A[[P]]\widehat\otimes_AA[[P]]$ acts on
$M\widehat\otimes_A N$ factorwise.

For trivial coefficients, applying augmentation in both tensor factors gives
\begin{equation}\label{eq:tensor-fox-augmentation}
 (\varepsilon\otimes\varepsilon)\Xi_{ij}(w)
 =\varepsilon(\partial_i\partial_j w)
 =\varepsilon_{ij}(w),
\end{equation}
with the ordering convention of Equation~\eqref{eq:second-fox-coeff}.
Consequently Equation~\eqref{eq:twisted-tensor-fox-cup} specializes to
Theorem~\ref{thm:quadratic-fox-cup}.

If $M$ and $N$ lie in fixed residual determinant sectors, then
Theorem~\ref{thm:ch-tensor-evaluation-compression} sends every
$\overline{\Xi}_{ij}(r_\alpha)$ to the finite Cayley--Hamilton tensor algebra.
Hence Equation~\eqref{eq:twisted-tensor-fox-cup} is a finite $A$-matrix formula
after sector evaluation.  When the relators are finite words---as in the
standard Demu\v{s}kin charts---these matrices are regular finite expressions
in the generator transport matrices and their inverses on the corresponding
generic-matrix chart.
\end{theorem}

\begin{proof}
The completed Fox differential of the relation cell is
\[
 d_2(e^2_\alpha)=
 \sum_j \overline{\partial_j r_\alpha}\,e^1_j.
\]
Modulo any open ideal, write
$\overline{\partial_j r_\alpha}=\sum_g c_{g,j}g$.
The standard Fox-to-bar chain transformation sends
$e^1_j$ to $[x_j]$ and the relation cell to
\[
 \sum_{j,g}c_{g,j}[g\mid x_j];
\]
this is the Fenn--Sjerve comparison used in the discrete Fox--Magnus cup
formula \cite{SuciuWangCup}.  For left coefficient modules the normalized bar
cup product satisfies
\[
 (a\smile b)[g\mid h]
 =a(g)\otimes \rho_N(g)b(h).
\]
A crossed homomorphism obeys the universal Fox identity
\[
 a(g)=\sum_i\rho_M(\overline{\partial_i\widetilde g})A_i,
\]
for any lift $\widetilde g\in F$.  Substitution and regrouping give
\[
 \sum_{i,j}
 (\rho_M\widehat\otimes\rho_N)
 \left(
 \sum_g c_{g,j}\,
 \partial_i\widetilde g\otimes g
 \right)
 (A_i\otimes B_j).
\]
The parenthesized tensor is exactly
$(\partial_i\widehat\otimes1)\Delta_F(\partial_jr_\alpha)$.
The formula is compatible modulo every open ideal, so continuity passes it to
the completed inverse limit and proves
Equation~\eqref{eq:twisted-tensor-fox-cup}.
Applying augmentation gives
Equation~\eqref{eq:tensor-fox-augmentation}.  The final assertion follows by
applying the finite sector evaluation of
Theorem~\ref{thm:ch-tensor-evaluation-compression}; for finite-word relators,
Fox differentiation and the Hopf diagonal produce finite sums of transport
words, hence regular matrix expressions on a generic-matrix chart.
\end{proof}

\begin{corollary}[Finite Fox--Cayley--Hamilton binary cup matrices]\label{cor:fox-ch-binary-cup}
On every finite-word pro-$p$ factorization locus of
Theorem~\ref{thm:ch-fox-effectivity}, the binary coefficient cup product on the
universal Fox complex is represented by finite algebraic matrices obtained
from the evaluated tensor Fox symbols
\[
 (q_M\widehat\otimes q_N)
 \bigl(\overline{\Xi}_{ij}(r_\alpha)\bigr),
\]
where $q_M$ and $q_N$ are the two residual-sector transport maps of
Theorem~\ref{thm:ch-tensor-evaluation-compression}.  Thus this gives a
parametrix-free finite binary cochain formula on the Fox--Cayley--Hamilton
branch.  For trivial $\mathbf F_p$ coefficients in minimal pro-$p$ presentations,
Theorem~\ref{thm:higher-magnus-massey} further gives the first nonzero higher
Massey arity directly from the Magnus symbol.  Higher arity for genuinely nontrivial coefficient actions, comparison with
non-Fox gauges, and tower-native coherence are supplied later at the recursively
explicit atlas level by the all-Sylow, common-source, and Stasheff comparison
theorems.  At scalar residual pro-$p$ points,
Theorem~\ref{thm:scalar-magnus-deformation} also identifies the first nonzero
higher adjoint obstruction directly.  Thus the present corollary isolates the
binary native formula; the later theory closes its finite higher-coherence
assembly.
\end{corollary}

\begin{theorem}[Quadratic Fox symbol as the cup-product matrix]\label{thm:quadratic-fox-cup}
Let
\[
P=\langle x_1,\ldots,x_d\mid r\rangle_{\mathrm{pro}-p}
\]
be a minimal one-relator pro-$p$ group of cohomological dimension two.  Let $\chi_1,\ldots,\chi_d$ be the basis of $H^1(P,\mathbf F_p)$ dual to the images of the generators, and choose the generator $\omega\in H^2(P,\mathbf F_p)$ corresponding under transgression to the functional which sends the chosen relation class of $r$ to $1$.  Then
\begin{equation}\label{eq:fox-cup-matrix}
\chi_i\smile\chi_j
=
\varepsilon_{ij}(r)\,\omega
\qquad(1\le i,j\le d).
\end{equation}
Hence the binary multiplication on the mod-$p$ Fox--Spencer chart is computable directly from the degree-two Magnus--Fox symbol
\[
\sigma_2(r)
=
\sum_{i,j=1}^{d}\varepsilon_{ij}(r)X_iX_j.
\]
Under a change of minimal generators the coefficient matrix changes covariantly, while the resulting bilinear form on $H^1(P,\mathbf F_p)$ is intrinsic.
\end{theorem}

\begin{proof}
Minimality puts $r$ in the Frattini subgroup of $F$.  Consequently, after applying $\operatorname{Hom}(-,\mathbf F_p)$ to the completed Fox resolution, both displayed differentials vanish: $x_i-1$ acts trivially and the augmented first Fox derivatives of $r$ vanish modulo $p$.  Thus the dual generator classes give bases of $H^1$ and $H^2$.

The comparison from the completed cellular/Fox resolution to the normalized continuous bar resolution evaluates the cup product of the two generator-dual $1$-cochains on the relation $2$-cell by the degree-two Magnus coefficient of the relator.  Equivalently this coefficient is the iterated augmentation $\varepsilon(\partial_{ij}r)$ with the convention of \eqref{eq:second-fox-coeff}.  This is the completed pro-$p$ analogue of the Fox--Magnus cup-product formula for presentation $2$-complexes; the calculation is finite modulo every open ideal and therefore passes to the completed inverse limit.  See \cite{SuciuWangCup} for the discrete Fox--Magnus formula and \cite{LabuteClassification,NSW} for the minimal one-relator and duality structure of Demu\v{s}kin groups.  The normalization of $\omega$ absorbs the global orientation sign.
\end{proof}


\subsection{Higher Magnus--Fox symbols and the first nonzero Massey arity}

The quadratic symbol above is the arity-two case of a standard higher relation
between Magnus coefficients and Massey products in minimal pro-$p$
presentations.  We record it because it gives a genuinely closed higher-arity
Fox formula on a useful class of test sectors, while making clear that it does
not by itself solve the general twisted-coefficient $A_\infty$ problem.

Let $F$ be the free pro-$p$ group on $x_1,\ldots,x_d$ and write the mod-$p$
Magnus expansion
\[
 \mathcal M(f)
 =1+\sum_{k\ge1}\ \sum_{|I|=k}
 \varepsilon_I(f)X_I,
 \qquad
 X_I=X_{i_1}\cdots X_{i_k}.
\]
Write $F_{(n)}$ for the $n$th Zassenhaus step, equivalently the subgroup of
those $f$ for which every Magnus coefficient of length $<n$ vanishes.

\begin{theorem}[Higher Magnus--Fox symbol as the first nonzero Massey tensor]\label{thm:higher-magnus-massey}
Let
\[
 1\longrightarrow R\longrightarrow F\longrightarrow G\longrightarrow1
\]
be a minimal finite presentation of a finitely presented pro-$p$ group, and let
$r_1,\ldots,r_m$ be a minimal system of defining relations.  Let
$\chi_1,\ldots,\chi_d$ be the basis of $H^1(G,\mathbf F_p)$ dual to the
chosen generators.  Via transgression, choose the basis
$\omega_1,\ldots,\omega_m$ of $H^2(G,\mathbf F_p)$ dual to the relation
classes, with the global sign convention normalized so that the formulas below
hold.

Assume that
\[
 R\subseteq F_{(n)}
\]
for some $n\ge2$.  Then:
\begin{enumerate}[label=\textnormal{(\roman*)}]
\item For every $2\le k<n$, every $k$-fold Massey product on
$H^1(G,\mathbf F_p)$ is uniquely defined and identically zero.

\item The $n$-fold Massey product is uniquely defined and multilinear.  For
any multi-index $I=(i_1,\ldots,i_n)$,
\begin{equation}\label{eq:higher-magnus-massey}
 \left\langle
  \chi_{i_1},\ldots,\chi_{i_n}
 \right\rangle_n
 =
 \sum_{a=1}^{m}
 \varepsilon_I(r_a)\,\omega_a.
\end{equation}
Equivalently, the complete $n$-linear Massey tensor is the degree-$n$ part of
the Magnus expansions of the minimal relation system.

\item If $n$ is the exact Zassenhaus invariant of the presentation, meaning
that some $\varepsilon_I(r_a)$ of length $n$ is nonzero, then $n$ is the first
arity at which the uniquely defined Massey tensor can be nonzero.  Thus the
first nonzero higher cohomological symbol is read directly from the first
nonzero homogeneous Magnus--Fox relation symbol.
\end{enumerate}
\end{theorem}

\begin{proof}
The hypothesis $R\subseteq F_{(n)}$ is equivalent to vanishing of all Magnus
coefficients of the defining relations in lengths $<n$.  The pro-$p$
Magnus--Massey theorem of Vogel and Morishita, in the formulation
\cite[Theorem~4.5]{GartnerMassey}, identifies the trace of every uniquely
defined Massey product on a relation class with the corresponding Magnus
coefficient.  It also implies that the lower products of arity $<n$ are
uniquely defined and zero.  Once all lower products vanish uniquely, the
$n$-fold product is itself uniquely defined and multilinear.  Applying the
relation trace basis gives Equation~\eqref{eq:higher-magnus-massey}.  The last
claim is immediate from the definition of the exact Zassenhaus depth.
\end{proof}

\begin{corollary}[One-relator higher Fox formula]\label{cor:one-relator-higher-massey}
In the one-relator case $G=\langle x_1,\ldots,x_d\mid r\rangle_{\mathrm{pro}-p}$,
if $r\in F_{(n)}$ then, after the same orientation normalization,
\[
 \left\langle
 \chi_{i_1},\ldots,\chi_{i_n}
 \right\rangle_n
 =\varepsilon_{i_1\cdots i_n}(r)\,\omega.
\]
If $r\in F_{(n)}\setminus F_{(n+1)}$, the degree-$n$ Magnus symbol
\[
 \sigma_n(r)=\sum_{|I|=n}\varepsilon_I(r)X_I
\]
is therefore exactly the first nonzero higher Massey tensor.  For $n=2$ this
recovers Theorem~\ref{thm:quadratic-fox-cup}.
\end{corollary}

\begin{warning}[Higher Magnus symbols are not the full twisted $A_\infty$ chart]\label{warn:higher-magnus-not-twisted}
Theorem~\ref{thm:higher-magnus-massey} concerns trivial $\mathbf F_p$
coefficients and the first uniquely defined higher Massey layer forced by a
minimal pro-$p$ presentation.  Theorem~\ref{thm:scalar-magnus-deformation}
below extends this first nonzero layer to adjoint coefficients at a scalar
residual point, where the pro-$p$ action on the adjoint is trivial.  Neither
statement gives the chain-level operations $m_n$ for a genuinely nontrivial
coefficient action, nor a tower-native Fox--Herr comparison.  Those remain
higher-coherent coefficient problems.
\end{warning}

\begin{theorem}[Matrix-colored higher Fox obstruction at a scalar residual point]\label{thm:scalar-magnus-deformation}
Retain the minimal presentation and the hypothesis $R\subseteq F_{(n)}$ of
Theorem~\ref{thm:higher-magnus-massey}, and let $k$ be a finite field of
characteristic $p$.  Let
\[
 \bar\rho:G\longrightarrow\mathrm{GL}_s(k)
\]
be a scalar residual representation.  Then $\bar\rho$ is trivial, because
$G$ is pro-$p$ while $k^\times$ has order prime to $p$.  Consequently
\[
 \operatorname{ad}\bar\rho\cong M_s(k)
\]
with trivial $G$-action and
\begin{equation}\label{eq:scalar-adjoint-tensor-dga}
 C^\bullet_{\mathrm{cont}}(G,\operatorname{ad}\bar\rho)
 \cong
 C^\bullet_{\mathrm{cont}}(G,k)\otimes_k M_s(k)
\end{equation}
as dg algebras, where the product on the right is the scalar cup product
together with matrix multiplication.

Choose a minimal $A_\infty$ model on $H^\ast(G,k)$ whose restriction to
$H^1(G,k)$ realizes the uniquely defined Massey products of
Theorem~\ref{thm:higher-magnus-massey}; normalize the global sign by the same
relation-trace convention as Equation~\eqref{eq:higher-magnus-massey}.  The
tensor product contraction with $M_s(k)$ gives a minimal $A_\infty$ model on
$H^\ast(G,\operatorname{ad}\bar\rho)$ with the following properties.
\begin{enumerate}[label=\textnormal{(\roman*)}]
\item For $2\le q<n$,
\[
 m_q^{\operatorname{ad}}
 \big|_{(H^1(G,k)\otimes M_s(k))^{\otimes q}}=0.
\]

\item For $I=(i_1,\ldots,i_n)$ and matrices $X_1,\ldots,X_n\in M_s(k)$,
\begin{equation}\label{eq:matrix-colored-higher-fox}
 m_n^{\operatorname{ad}}
 \bigl(\chi_{i_1}\otimes X_1,\ldots,\chi_{i_n}\otimes X_n\bigr)
 =
 \sum_{a=1}^{m}
 \varepsilon_I(r_a)\,\omega_a\otimes X_1\cdots X_n.
\end{equation}
Thus the first higher operation is obtained from the Magnus--Fox tensor by
coloring its ordered letters with matrix multiplication.

\item For a universal tangent vector
\[
 u=\sum_{i=1}^{d}\chi_i\otimes X_i
 \in H^1(G,k)\otimes M_s(k),
\]
put
\begin{equation}\label{eq:matrix-magnus-relation-polynomial}
 R_a^{(n)}(X_1,\ldots,X_d)
 :=
 \sum_{i_1,\ldots,i_n}
 \varepsilon_{i_1\cdots i_n}(r_a)
 X_{i_1}\cdots X_{i_n}.
\end{equation}
Then the first possible nonzero homogeneous term of the minimal
$A_\infty$ Kuranishi map (equivalently, of its Maurer--Cartan series) is
\begin{equation}\label{eq:scalar-kuranishi-magnus}
 \kappa_n(u)
 :=m_n^{\operatorname{ad}}(u,\ldots,u)
 =\sum_{a=1}^{m}\omega_a\otimes
 R_a^{(n)}(X_1,\ldots,X_d).
\end{equation}
This is exactly the degree-$n$ initial relation obtained by substituting
$x_i\mapsto 1+X_i$ in the Magnus expansion of each $r_a$.  Hence, at the
scalar residual point, the first nonlinear homogeneous representation-level
obstruction is the matrix substitution of the first nonzero homogeneous
Magnus--Fox relation symbol.
\end{enumerate}
\end{theorem}

\begin{proof}
The triviality of the scalar character follows from the pro-$p$ hypothesis.
Equation~\eqref{eq:scalar-adjoint-tensor-dga} is then immediate from the
trivial conjugation action.  Tensor the chosen contraction defining the
minimal $A_\infty$ model of $C^\bullet_{\mathrm{cont}}(G,k)$ with the
identity contraction of the degree-zero algebra $M_s(k)$.  Every rooted-tree
term acquires the ordered product of the matrix labels, so the vanishing of
the lower scalar operations and Equation~\eqref{eq:higher-magnus-massey}
give~(i) and~\eqref{eq:matrix-colored-higher-fox}.  Expanding the latter on
$u^{\otimes n}$ gives~\eqref{eq:scalar-kuranishi-magnus}.

Finally, the Magnus expansion is characterized by $x_i\mapsto1+X_i$ in the
completed noncommutative power-series algebra.  Since all homogeneous terms of
the relators below degree $n$ vanish, the degree-$n$ term of the framed
relation equation is precisely
$R_a^{(n)}(X_1,\ldots,X_d)$.  This agrees with the general principle that the
higher multiplications in a minimal $A_\infty$ model encode the initial
relations and with the $A_\infty$ presentation of Galois deformation rings
\cite{WangErickson}; here the equality is already visible directly from the
Magnus expansion and Theorem~\ref{thm:higher-magnus-massey}.
\end{proof}

\begin{corollary}[First nonlinear scalar-residual deformation equation]\label{cor:scalar-magnus-kuranishi}
If the Zassenhaus invariant is exactly $n$, the framed deformation equations
at the trivial $s$-dimensional residual representation have no homogeneous
relation term in degrees $2,\ldots,n-1$, while their degree-$n$ initial term
is the matrix-valued Magnus polynomial
$R_a^{(n)}(X_1,\ldots,X_d)$ of
Equation~\eqref{eq:matrix-magnus-relation-polynomial}.  Equivalently, the
first nonzero associative $A_\infty$ obstruction tensor on adjoint
cohomology is obtained by matrix-coloring the first nonzero higher
Magnus--Fox symbol.  Antisymmetrization gives the corresponding leading
$L_\infty$ bracket, but that antisymmetrized bracket may vanish on special
matrix directions even when the associative tensor is nonzero.
\end{corollary}


\subsection{Noncommutative relation jets at arbitrary residual points}\label{subsec:arbitrary-residual-word-jets}

The scalar calculation above has a direct classical extension to arbitrary
residual transport.  What remains special to the scalar case is the
identification with the first intrinsic minimal $A_\infty$ operation, not the
existence of explicit nonlinear relation jets.

\begin{theorem}[Finite noncommutative word jets at an arbitrary residual point]\label{thm:arbitrary-residual-word-jets}
Let $F$ be the free pro-$p$ group on $x_1,\ldots,x_d$, let $k$ be a finite
field of characteristic $p$, and let
\[
 \bar\rho:F\longrightarrow \mathrm{GL}_s(k)
\]
be a continuous representation.  Put $A_i:=\bar\rho(x_i)$.  Let
\[
 S:=k[[X_{i,ab}:1\le i\le d,\ 1\le a,b\le s]]
\]
and write $X_i=(X_{i,ab})\in M_s(\mathfrak m_S)$.  Use the left-trivialized
formal generator deformation
\begin{equation}\label{eq:left-trivialized-generator-deformation}
 \rho_X(x_i):=(1+X_i)A_i.
\end{equation}
For a finite word
\[
 w=a_1\cdots a_L,
 \qquad
 a_\ell=x_{j_\ell}^{\epsilon_\ell},\quad \epsilon_\ell\in\{\pm1\},
\]
define homogeneous local letter terms $C_{\ell,q}$ by
\begin{equation}\label{eq:letter-jet-coefficients}
 C_{\ell,q}:=
 \begin{cases}
 A_{j_\ell},&\epsilon_\ell=+1,\ q=0,\\
 X_{j_\ell}A_{j_\ell},&\epsilon_\ell=+1,\ q=1,\\
 0,&\epsilon_\ell=+1,\ q\ge2,\\
 (-1)^qA_{j_\ell}^{-1}X_{j_\ell}^{q},
   &\epsilon_\ell=-1,\ q\ge0.
 \end{cases}
\end{equation}
Then:
\begin{enumerate}[label=\textnormal{(\roman*)}]
\item The normalized word map has the convergent formal expansion
\begin{equation}\label{eq:arbitrary-residual-word-jet-expansion}
 \rho_X(w)\bar\rho(w)^{-1}
 =1+\sum_{q\ge1}J_q^{\bar\rho}(w;X),
\end{equation}
where
\begin{equation}\label{eq:arbitrary-residual-word-jet}
 J_q^{\bar\rho}(w;X)
 =
 \left(
 \sum_{k_1+\cdots+k_L=q}
 C_{1,k_1}\cdots C_{L,k_L}
 \right)\bar\rho(w)^{-1}.
\end{equation}
For each fixed $q$, this is a finite noncommutative matrix polynomial.  A
crude bound for the number of displayed summands is
$\binom{q+L-1}{L-1}$.

\item Let the completed group ring act on $M_s(k)$ through the adjoint residual
transport, and let
\[
 \operatorname{Ad}_{\bar\rho}:k[[F]]
 \longrightarrow \operatorname{End}_k(M_s(k))
\]
be the continuous extension of the adjoint residual action
$g\cdot X=\bar\rho(g)X\bar\rho(g)^{-1}$.  Then the linear term is exactly
the evaluated Fox Jacobian:
\begin{equation}\label{eq:word-jet-linear-fox}
 J_1^{\bar\rho}(w;X)
 =\sum_{i=1}^{d}
 \operatorname{Ad}_{\bar\rho}(\partial_i w)(X_i).
\end{equation}
Thus the nonlinear jet tower extends, rather than replaces, the Fox
linearization already used in the framed tangent complex.

\item If $w=r$ is a defining relation for a quotient through which
$\bar\rho$ factors, then $\bar\rho(r)=1$ and
$J_q^{\bar\rho}(r;X)$ is the degree-$q$ homogeneous relation matrix in the
framed formal deformation chart.  Hence the complete classical relation jet
of every finite word relator is computable by the finite recursion
\eqref{eq:arbitrary-residual-word-jet} in every fixed degree.

\item If $\bar\rho$ is trivial, then $A_i=1$ and
\eqref{eq:arbitrary-residual-word-jet-expansion} is the ordinary matrix
Magnus substitution $x_i\mapsto1+X_i$.  In particular, when all relators have
Zassenhaus depth at least $n$, the terms of degree $<n$ vanish and the
first nonzero term is exactly the matrix-colored Magnus polynomial of
Theorem~\ref{thm:scalar-magnus-deformation}.
\end{enumerate}
\end{theorem}

\begin{proof}
Equation~\eqref{eq:left-trivialized-generator-deformation} gives
\[
 \rho_X(x_i)=(1+X_i)A_i,
 \qquad
 \rho_X(x_i^{-1})=A_i^{-1}(1+X_i)^{-1}
 =\sum_{q\ge0}(-1)^qA_i^{-1}X_i^q.
\]
Multiplying these letter expansions and collecting homogeneous degree gives
\eqref{eq:arbitrary-residual-word-jet}; for fixed degree only finitely many
compositions $k_1+\cdots+k_L=q$ occur.  This proves~(i) and~(iii).

For~(ii), the first-order normalized variation
\[
 c_X(w):=[\rho_X(w)\bar\rho(w)^{-1}]_{\deg1}
\]
is the crossed homomorphism with $c_X(x_i)=X_i$ for the adjoint
$F$-module $M_s(k)$.  The universal Fox identity for crossed homomorphisms \cite{Fox}
therefore gives
\[
 c_X(w)=\sum_i(\partial_iw)\cdot X_i,
\]
which is precisely \eqref{eq:word-jet-linear-fox} under the displayed
adjoint action.  Finally, when all $A_i=1$, the letter expansions are exactly
those obtained from $x_i\mapsto1+X_i$ and
$x_i^{-1}\mapsto(1+X_i)^{-1}$, proving~(iv).
\end{proof}

\begin{corollary}[Finite arbitrary-residual classical obstruction jets]\label{cor:arbitrary-residual-classical-obstruction}
Let $r_1,\ldots,r_m$ be finite relator words defining a framed residual word
chart at $\bar\rho$.  Then every finite jet of its relation map is an explicit
finite matrix polynomial in the residual transport matrices
$A_i^{\pm1}$ and the tangent variables $X_i$.  Restricting to the linear
tangent locus
\[
 J_1^{\bar\rho}(r_a;X)=0\qquad(1\le a\le m)
\]
and inspecting the first surviving homogeneous degree therefore gives a
finite classical obstruction polynomial at an arbitrary residual point.

If Theorem~\ref{thm:word-lci-reduction} replaces the word equations
$F=(F_1,\ldots,F_M)$ by a regular sequence $g=UF$, then the jets of $g$ are
obtained from the finite Cauchy products of the jets of $U$ with the word jets
\eqref{eq:arbitrary-residual-word-jet}.  Hence the syzygy-complete lci formal
chart also has computable finite jets in every prescribed degree.
\end{corollary}


\begin{theorem}[Finite Lyapunov--Schmidt reduction of arbitrary-residual word jets]\label{thm:word-jet-lyapunov-schmidt}
Let $F:\widehat V\to\widehat W$ be the formal framed relation map of a
finite residual word chart, or of the lci regular sequence obtained from it
by Theorem~\ref{thm:word-lci-reduction}.  Thus $F(0)=0$, and in chosen
finite coordinates every homogeneous Taylor coefficient of $F$ is explicitly
computable from Theorem~\ref{thm:arbitrary-residual-word-jets} and finite
Cauchy products.  Put
\[
 L:=dF_0:V\longrightarrow W.
\]
Work first over a residual field, or more generally on a constant-rank
minor chart on which the relevant finite projective modules split.  Choose
splittings
\[
 V=T\oplus N,
 \qquad
 W=I\oplus O
\]
such that $T=\ker L$, $I=\operatorname{im}L$, and
\[
 A:=\operatorname{pr}_I\circ L|_N:N\xrightarrow{\sim}I
\]
is an isomorphism.  After the corresponding linear row operation on the
target one has
\[
 dF_0(t,n)=(An,0).
\]
Write $F=(F_I,F_O)$ in these coordinates.  Then:
\begin{enumerate}[label=\textnormal{(\roman*)}]
\item There is a unique formal map
\[
 h:T\longrightarrow N,
 \qquad h(0)=0,
 \qquad dh_0=0,
\]
satisfying
\begin{equation}\label{eq:word-ls-solved-equations}
 F_I(t,h(t))=0.
\end{equation}
If $h=\sum_{q\ge2}h_q$ is decomposed into homogeneous parts and
$h_{<q}:=\sum_{2\le j<q}h_j$, then every coefficient is obtained from the
finite recursion
\begin{equation}\label{eq:word-ls-recursion}
 h_q(t)
 =-A^{-1}
 \left[F_I\bigl(t,h_{<q}(t)\bigr)\right]_q,
 \qquad q\ge2.
\end{equation}
Hence the $q$-jet of $h$ uses only the word/lci jets of degrees at most $q$
and one fixed inverse Jacobian minor.

\item Define the reduced obstruction map
\begin{equation}\label{eq:word-ls-kuranishi}
 \kappa_{\mathrm{word}}:T\longrightarrow O,
 \qquad
 \kappa_{\mathrm{word}}(t):=F_O(t,h(t)).
\end{equation}
It has no constant or linear term.  Its degree-$q$ homogeneous coefficient is
\begin{equation}\label{eq:word-ls-kuranishi-jet}
 \kappa_{\mathrm{word},q}(t)
 =\left[F_O\bigl(t,h_{<q}(t)\bigr)\right]_q,
 \qquad q\ge2.
\end{equation}
Thus every prescribed finite Taylor jet of $\kappa_{\mathrm{word}}$ is a
finite matrix expression in the residual transport matrices, the relation
word jets, the lci row-reduction matrices when present, and the inverse
minor $A^{-1}$.

\item The formal zero locus of the original framed relation map is
canonically identified, relative to the chosen splittings, with the zero
locus of the reduced obstruction map:
\begin{equation}\label{eq:word-ls-zero-locus}
 \{F=0\}
 \xrightarrow{\sim}
 \{\kappa_{\mathrm{word}}=0\},
 \qquad
 (t,n)\longleftrightarrow t,
 \qquad n=h(t).
\end{equation}
Consequently the classical obstruction problem at an arbitrary residual
point admits an explicit representative in tangent and obstruction
coordinates, not merely explicit ambient word equations.

\item In a family, after a finite determinantal refinement on which the rank
of $L$ is constant and the chosen minor is invertible, the same recursion is
regular on that chart and commutes with flat base change.  Under arbitrary
specialization the unreduced word/lci jets still specialize term by term; if
the chosen minor vanishes, one passes to another determinantal chart rather
than introducing a new obstruction differential.
\end{enumerate}
\end{theorem}

\begin{proof}
After the chosen splittings, the first component has the form
\[
 F_I(t,n)=An+F_I^{\ge2}(t,n).
\]
The formal implicit-function theorem applies because $A$ is invertible.
Collecting homogeneous degree $q$ after the coefficients of degrees $<q$
have already been determined gives exactly
\eqref{eq:word-ls-recursion}; this also proves uniqueness recursively.
Because the linear part of $F_O$ is zero, its degree-$q$ value after
substitution depends only on $h_j$ with $j<q$, which gives
\eqref{eq:word-ls-kuranishi-jet}.  Equation
\eqref{eq:word-ls-zero-locus} is then immediate from the unique solution of
\eqref{eq:word-ls-solved-equations}.  On a constant-rank family chart all
operations are finite matrix operations and inversion of the one chosen
minor, so they commute with flat base change.  The final specialization
statement is the usual determinantal-chart transition already used in
Theorem~\ref{thm:transport-explicit-lci-row-reduction}.
\end{proof}

\begin{corollary}[Finite classical Kuranishi jets at every residual point]\label{cor:finite-classical-kuranishi-jets}
On every residual finite-word/lci chart, and after a finite constant-rank
minor refinement, the classical framed deformation equations can be reduced
to a formal obstruction map on the tangent kernel with values in a chosen
obstruction complement.  Every fixed homogeneous component of that map is
computable by the finite recursion
\eqref{eq:word-ls-recursion}--\eqref{eq:word-ls-kuranishi-jet}.  Thus extraction of a classical Kuranishi-type obstruction series from the
transport words is already finite and recursive.  On sectors where an
intrinsic Fox contraction is fixed, Theorem~\ref{thm:kuranishi-jet-coordinate-recursion}
below likewise makes the coordinate comparison with the transferred Kuranishi
series finite and recursive.  The later native ordered, common-source, and
Stasheff theorems close presentation and tower coherence at the recursively
explicit finite-atlas level; only a preferred strict closed representative is
noncanonical.
\end{corollary}

\begin{corollary}[Word jets versus the intrinsic Kuranishi gauge]\label{cor:word-jets-kuranishi-boundary}
On a residual sector where Theorem~\ref{thm:fox-kuranishi-comparison} applies,
Theorem~\ref{thm:word-jet-lyapunov-schmidt} first turns the explicit word/lci
jets into an explicit reduced obstruction series
$\kappa_{\mathrm{word}}:T\to O$.  After splitting off the smooth framed/gauge
directions and identifying the tangent and obstruction fibers with
$H^1(G_K,\operatorname{ad}\rho)$ and $H^2(G_K,\operatorname{ad}\rho)$, this
series defines the same formal deformation zero locus as the intrinsic
transferred Kuranishi map, after formal source coordinates and an invertible
formal change of the obstruction frame.  They are not canonically equal term
by term.  Theorem~\ref{thm:kuranishi-jet-coordinate-recursion} makes the jets
of those coordinate changes computable by finite linear systems.  Thus
arbitrary residual \emph{classical} nonlinear transport equations, their
tangent/obstruction reduction, and the jetwise comparison with one chosen
intrinsic minimal gauge are now effective.  At this stage a direct Fox formula
for the genuinely twisted coefficient $A_\infty/L_\infty$ operations and their
presentation- and tower-native comparisons has not yet been produced.  The later
native higher-operation and globalization theorems close this at the recursively
explicit atlas level; only optional strict closed representatives remain noncanonical.
\end{corollary}

\begin{theorem}[Finite recursive comparison of classical and intrinsic Kuranishi jets]\label{thm:kuranishi-jet-coordinate-recursion}
Work on a constant-rank residual sector on which
Theorems~\ref{thm:word-jet-lyapunov-schmidt} and
\ref{thm:fox-kuranishi-comparison} both apply.  Choose the required minor
charts, splittings, and one contraction of the finite adjoint Fox complex.
After linear identifications of the tangent and obstruction modules, write
\[
 \kappa_{\mathrm{word}},\ \kappa_{\mathrm{Fox}}:T\longrightarrow O
\]
for the reduced word/lci obstruction series and the transferred intrinsic
Kuranishi series.  Suppose they are not identically zero.  The formal
equivalence in Theorem~\ref{thm:fox-kuranishi-comparison} gives a formal source
automorphism $\Phi$ and an invertible formal obstruction-frame matrix $B(t)$
such that
\begin{equation}\label{eq:kuranishi-formal-gauge-equivalence}
 \kappa_{\mathrm{Fox}}\bigl(\Phi(t)\bigr)
 =B(t)\,\kappa_{\mathrm{word}}(t).
\end{equation}
Absorb the invertible linear parts into the chosen linear identifications, so
that
\[
 \Phi(t)=t+\sum_{r\ge2}\Phi_r(t),
 \qquad
 B(t)=1+\sum_{r\ge1}B_r(t).
\]
Let $m\ge2$ be the common order of vanishing.  Then the leading homogeneous
terms agree; denote them by $\kappa_m$.  For every $s\ge1$, put $q=m+s$ and
assume that $\Phi_2,\ldots,\Phi_s$ and
$B_1,\ldots,B_{s-1}$ have already been chosen so that
\eqref{eq:kuranishi-formal-gauge-equivalence} holds through degree $q-1$.
Define the known degree-$q$ defect
\begin{equation}\label{eq:kuranishi-coordinate-defect}
 E_q(t):=
 \left[
 \kappa_{\mathrm{Fox}}\!\left(
 t+\sum_{2\le r\le s}\Phi_r(t)
 \right)
 -\left(1+\sum_{1\le r\le s-1}B_r(t)\right)
 \kappa_{\mathrm{word}}(t)
 \right]_q.
\end{equation}
Then:
\begin{enumerate}[label=\textnormal{(\roman*)}]
\item The new coefficients occur linearly and satisfy
\begin{equation}\label{eq:kuranishi-coordinate-linear-recursion}
 D\kappa_m(t)\bigl[\Phi_{s+1}(t)\bigr]
 -B_s(t)\kappa_m(t)
 =-E_q(t).
\end{equation}
No coefficient of degree $>q$ enters this equation.

\item In chosen finite bases, Equation~\eqref{eq:kuranishi-coordinate-linear-recursion}
is a finite linear system.  More intrinsically, its coefficient map is
\begin{equation}\label{eq:kuranishi-coordinate-linear-operator}
 \mathcal L_s:
 \left(T\otimes\operatorname{Sym}^{s+1}T^\vee\right)
 \oplus
 \left(\operatorname{End}(O)\otimes\operatorname{Sym}^{s}T^\vee\right)
 \longrightarrow
 O\otimes\operatorname{Sym}^{m+s}T^\vee,
\end{equation}
\[
 \mathcal L_s(\phi,b)
 =D\kappa_m(t)[\phi(t)]-b(t)\kappa_m(t).
\]
The existence of the full formal equivalence implies
$-E_q\in\operatorname{im}(\mathcal L_s)$.  Hence the solutions form a nonempty
affine torsor under $\ker(\mathcal L_s)$.

\item After a finite Fitting/minor refinement for the matrices
$\mathcal L_s$, a chosen solution of every prescribed finite stage is obtained
by ordinary row reduction.  Consequently, for every $N$, the $N$-jets of
$\Phi$ and $B$ are finite matrix expressions in the word/lci jets, the chosen
Fox contraction matrices, and finitely many inverse minors.  The word-side
jets are supplied by Theorem~\ref{thm:word-jet-lyapunov-schmidt}; the
Fox-side jets are the finite rooted-tree coefficients of the chosen
homotopy transfer.

\item On a family chart on which the selected minors remain invertible, this
finite recursion commutes with flat base change.  Under a rank jump the full
word and Fox perfect objects still specialize; only the chosen coordinate
solution may require passage to another determinantal chart.
\end{enumerate}
If both Kuranishi maps vanish identically, the deformation problem is
unobstructed and no comparison recursion is needed.
\end{theorem}

\begin{proof}
Theorem~\ref{thm:fox-kuranishi-comparison} gives a formal source change and
an invertible formal change of obstruction coordinates.  In finite
coordinates, any target change $\Psi(y)$ with $\Psi(0)=0$ can be written
$\Psi(y)=C(y)y$ for a formal matrix $C$ with $C(0)=d\Psi_0$; after substituting
$y=\kappa_{\mathrm{word}}(t)$ this yields
\eqref{eq:kuranishi-formal-gauge-equivalence} with
$B(t)=C(\kappa_{\mathrm{word}}(t))$.  After the linear parts are absorbed,
invertibility forces the two maps to have the same order of vanishing and the
same leading homogeneous tensor.  Assume the identity has been solved below
degree $q=m+s$.  Substituting a new homogeneous
source coefficient $\Phi_{s+1}$ into the leading term $\kappa_m$ contributes
\[
 D\kappa_m(t)[\Phi_{s+1}(t)]
\]
in degree $m+s=q$.  Inserting $\Phi_{s+1}$ into any higher
$\kappa_{\mathrm{Fox},j}$ with $j>m$ has degree at least $q+1$.  Likewise, the
only new degree-$q$ contribution from the obstruction-frame change is
$B_s\kappa_m$.  All other terms are already contained in the known defect
\eqref{eq:kuranishi-coordinate-defect}.  This proves
\eqref{eq:kuranishi-coordinate-linear-recursion} and
\eqref{eq:kuranishi-coordinate-linear-operator}.

A full formal equivalence exists, so the displayed finite system is solvable
at every stage.  Finite projectivity turns it into ordinary matrix algebra;
Fitting stratification and an invertible minor give a concrete row-reduced
solution.  The fixed-degree transferred Fox coefficients are finite
rooted-tree expressions in the contraction data, while the word coefficients
are finite by Theorem~\ref{thm:word-jet-lyapunov-schmidt}.  This proves the
finite-effectivity and base-change statements.
\end{proof}

\begin{corollary}[Effective finite Kuranishi coordinate comparison]\label{cor:kuranishi-coordinate-gap-finite}
On every constant-rank Fox/lci sector, the passage from the explicit classical
Kuranishi representative to any chosen intrinsic minimal Fox gauge can be
computed to arbitrary finite order by the linear recursion
\eqref{eq:kuranishi-coordinate-linear-recursion}.  Thus neither existence of a formal comparison nor extraction of its finite
jets remains problematic.  Theorem~\ref{thm:native-twisted-linfinity} below
removes the direct formula problem for the two-term Maurer--Cartan/$L_\infty[1]$
formal-slice branch by polarizing the explicit reduced equations themselves.
The later all-Sylow/common-source and raw-Herr comparison theorems supply the
cochain-level $A_\infty/L_\infty$ and tower comparison at the finite native
level; literal equality of a chosen unreduced raw-period representative is
only the optional strict gauge isolated near the end of the paper.
\end{corollary}

\begin{theorem}[Strictly transitive finite Kuranishi jet atlas]\label{thm:kuranishi-common-source-atlas}
Let $S$ be a constant-rank residual stratum and let
$\{\kappa_\alpha:T\to O\}_{\alpha\in\mathcal A}$ be a finite family of
word/lci Kuranishi gauges obtained from different finite presentations,
relator systems, splittings, or determinantal minor charts on $S$.  Fix one
intrinsic minimal Fox Kuranishi gauge
\[
 \kappa_*:T\longrightarrow O
\]
after the chosen linear identifications of the common tangent and obstruction
modules.  For every $\alpha$, choose the finite-recursive comparison data of
Theorem~\ref{thm:kuranishi-jet-coordinate-recursion}, normalized by
\begin{equation}\label{eq:kuranishi-common-source-comparison}
 \kappa_*\bigl(\Phi_\alpha(t)\bigr)
 =B_\alpha(t)\,\kappa_\alpha(t),
 \qquad
 d\Phi_{\alpha,0}=\operatorname{id}_T,
 \qquad
 B_\alpha(0)=\operatorname{id}_O.
\end{equation}
On an unobstructed component all Kuranishi maps vanish and the obstruction
frame may be omitted.  For obstructed components define
\begin{equation}\label{eq:kuranishi-direct-source-transition}
 \Psi_{\beta\alpha}
 :=\Phi_\beta^{-1}\circ\Phi_\alpha
\end{equation}
and
\begin{equation}\label{eq:kuranishi-direct-obstruction-transition}
 C_{\beta\alpha}(t)
 :=B_\beta\!\left(\Psi_{\beta\alpha}(t)\right)^{-1}B_\alpha(t).
\end{equation}
Then:
\begin{enumerate}[label=\textnormal{(\roman*)}]
\item The direct transition identifies the two reduced obstruction equations:
\begin{equation}\label{eq:kuranishi-direct-gauge-equivalence}
 \kappa_\beta\!\left(\Psi_{\beta\alpha}(t)\right)
 =C_{\beta\alpha}(t)\,\kappa_\alpha(t).
\end{equation}
Thus presentation and minor changes are explicit formal gauge changes once a
common intrinsic source has been fixed.

\item The transitions are strictly transitive:
\begin{equation}\label{eq:kuranishi-strict-source-cocycle}
 \Psi_{\alpha\alpha}=\operatorname{id},
 \qquad
 \Psi_{\gamma\beta}\circ\Psi_{\beta\alpha}
 =\Psi_{\gamma\alpha},
\end{equation}
and
\begin{equation}\label{eq:kuranishi-strict-obstruction-cocycle}
 C_{\alpha\alpha}=1,
 \qquad
 C_{\gamma\beta}\!\left(\Psi_{\beta\alpha}(t)\right)
 C_{\beta\alpha}(t)
 =C_{\gamma\alpha}(t).
\end{equation}
No additional triple-overlap obstruction occurs at the level of these chosen
formal Kuranishi gauges.

\item For every prescribed jet order $N$, the $N$-jets of
$\Psi_{\beta\alpha}$ and $C_{\beta\alpha}$ are finite matrix expressions.
They are obtained from the finite-recursive jets of $\Phi_\alpha,B_\alpha$
by formal inversion, composition, multiplication, and matrix inversion.  Each
of these operations is recursive in total degree, so only coefficients of
degree at most $N$ enter the $N$-jet.  After a finite Fitting/minor refinement,
all displayed coefficients are rational expressions in the word/lci jets,
the chosen Fox contraction matrices, and finitely many inverse minors.

\item On a family chart on which the chosen minors remain invertible, the
finite jet atlas commutes with flat base change.  Across a rank jump the common
perfect deformation object still specializes; only the selected minimal gauges
and their transition formulas may require passage to a different Fitting
stratum.

\item Replacing the common intrinsic gauge $\kappa_*$, the contractions, or
the recursive solutions replaces the displayed strict atlas by a chartwise
formal-gauge-equivalent atlas.  Hence strict transitivity is a property of the
chosen common-source realization, while the underlying formal deformation
stack is independent of those choices.
\end{enumerate}
\end{theorem}

\begin{proof}
Equation~\eqref{eq:kuranishi-common-source-comparison} for $\alpha$ and
$\beta$, together with
$\Phi_\alpha=\Phi_\beta\Psi_{\beta\alpha}$, gives
\[
 B_\beta\!\left(\Psi_{\beta\alpha}(t)\right)
 \kappa_\beta\!\left(\Psi_{\beta\alpha}(t)\right)
 =B_\alpha(t)\kappa_\alpha(t),
\]
which is exactly~\eqref{eq:kuranishi-direct-gauge-equivalence}.
The source cocycle follows by cancellation through the common maps
$\Phi_\alpha$:
\[
 \Phi_\gamma^{-1}\Phi_\beta\,
 \Phi_\beta^{-1}\Phi_\alpha
 =\Phi_\gamma^{-1}\Phi_\alpha.
\]
Substitution into~\eqref{eq:kuranishi-direct-obstruction-transition} cancels
the intermediate $B_\beta$ and proves
\eqref{eq:kuranishi-strict-obstruction-cocycle}.

Theorem~\ref{thm:kuranishi-jet-coordinate-recursion} computes every fixed jet
of $\Phi_\alpha$ and $B_\alpha$ by finite linear algebra.  A formal map with
identity linear part has an inverse whose homogeneous degree-$q$ coefficient
is obtained recursively from the coefficients of degrees at most $q$; formal
composition and matrix inversion satisfy the same finite-degree property.
This proves~(iii).  The base-change statement follows on the same determinantal
charts.  Finally, changing the common source or any chosen solution changes
each chart-to-source comparison by a formal gauge, which conjugates the
resulting transition system and leaves the represented formal deformation
stack unchanged.
\end{proof}

\begin{corollary}[Transitive finite presentation atlas]\label{cor:kuranishi-presentation-atlas}
On every finite constant-rank residual cover for which the word/lci and Fox
comparisons are available, changes of generators, relators, lci reductions,
minor charts, and Kuranishi splittings admit strictly transitive transition
jets to every finite order after one common intrinsic gauge is chosen.  Thus existence, finite computability, and cocycle coherence of the Kuranishi
coordinate changes are already closed.  Direct intrinsic twisted
$A_\infty/L_\infty$ formulas that avoid the common minimal gauge are developed
later by the native ordered and common-source constructions; they are a choice
of strict formula rather than an additional presentation-independence
obstruction.
\end{corollary}

\begin{theorem}[Strict presentation--tower interchange on Kuranishi jets]\label{thm:kuranishi-presentation-tower-interchange}
Let $S$ be a constant-rank residual stratum carrying the presentation/minor
Kuranishi atlas of Theorem~\ref{thm:kuranishi-common-source-atlas}.  Let
$\mathcal T$ be a finite collection of cohomological charts on the same
stratum, and assume that the common-source identifications of
Theorem~\ref{thm:common-source-chart-change} restrict in degrees one and two to
base-change-compatible isomorphisms
\[
 \lambda_{\tau,1}:T\xrightarrow{\sim}T_\tau,
 \qquad
 \lambda_{\tau,2}:O\xrightarrow{\sim}O_\tau,
 \qquad \tau\in\mathcal T,
\]
where $T$ and $O$ are the common intrinsic tangent and obstruction modules.
For every presentation index $\alpha$ and tower index $\tau$, define
\begin{equation}\label{eq:presentation-tower-kuranishi-chart}
 \kappa_{\alpha,\tau}(t_\tau)
 :=\lambda_{\tau,2}\,
   \kappa_\alpha\!\left(\lambda_{\tau,1}^{-1}t_\tau\right).
\end{equation}
For tower indices $\tau,\sigma$ put
\begin{equation}\label{eq:kuranishi-tower-linear-transitions}
 L_{\sigma\tau}:=\lambda_{\sigma,1}\lambda_{\tau,1}^{-1},
 \qquad
 M_{\sigma\tau}:=\lambda_{\sigma,2}\lambda_{\tau,2}^{-1}.
\end{equation}
For presentation indices $\alpha,\beta$, transport the transitions of
Theorem~\ref{thm:kuranishi-common-source-atlas} to tower $\tau$ by
\begin{align}
 \Psi^{(\tau)}_{\beta\alpha}
 &:=\lambda_{\tau,1}\Psi_{\beta\alpha}\lambda_{\tau,1}^{-1},
 \label{eq:kuranishi-presentation-transition-in-tower}\\
 C^{(\tau)}_{\beta\alpha}(t_\tau)
 &:=\lambda_{\tau,2}
 C_{\beta\alpha}\!\left(\lambda_{\tau,1}^{-1}t_\tau\right)
 \lambda_{\tau,2}^{-1}.
 \label{eq:kuranishi-obstruction-transition-in-tower}
\end{align}
Then:
\begin{enumerate}[label=\textnormal{(\roman*)}]
\item Tower change and presentation change separately preserve the Kuranishi
zero loci:
\begin{align}
 \kappa_{\alpha,\sigma}(L_{\sigma\tau}t_\tau)
 &=M_{\sigma\tau}\kappa_{\alpha,\tau}(t_\tau),
 \label{eq:kuranishi-tower-gauge-equivalence}\\
 \kappa_{\beta,\tau}(\Psi^{(\tau)}_{\beta\alpha}(t_\tau))
 &=C^{(\tau)}_{\beta\alpha}(t_\tau)
   \kappa_{\alpha,\tau}(t_\tau).
 \label{eq:kuranishi-presentation-gauge-equivalence-in-tower}
\end{align}

\item The two directions satisfy a strict interchange law:
\begin{equation}\label{eq:kuranishi-source-interchange}
 L_{\sigma\tau}\Psi^{(\tau)}_{\beta\alpha}
 =\Psi^{(\sigma)}_{\beta\alpha}L_{\sigma\tau},
\end{equation}
and
\begin{equation}\label{eq:kuranishi-obstruction-interchange}
 M_{\sigma\tau}C^{(\tau)}_{\beta\alpha}(t_\tau)
 =C^{(\sigma)}_{\beta\alpha}(L_{\sigma\tau}t_\tau)
  M_{\sigma\tau}.
\end{equation}
Hence no mixed presentation--tower obstruction appears at the level of the
chosen formal Kuranishi gauges.

\item Define the combined transition from $(\alpha,\tau)$ to
$(\beta,\sigma)$ by
\begin{align}
 \widehat\Psi_{(\beta,\sigma),(\alpha,\tau)}
 &:=\Psi^{(\sigma)}_{\beta\alpha}L_{\sigma\tau}
   =L_{\sigma\tau}\Psi^{(\tau)}_{\beta\alpha},
 \label{eq:kuranishi-product-source-transition}\\
 \widehat C_{(\beta,\sigma),(\alpha,\tau)}(t_\tau)
 &:=C^{(\sigma)}_{\beta\alpha}(L_{\sigma\tau}t_\tau)M_{\sigma\tau}
   =M_{\sigma\tau}C^{(\tau)}_{\beta\alpha}(t_\tau).
 \label{eq:kuranishi-product-obstruction-transition}
\end{align}
Then the combined transitions are strictly transitive on the product index
$\mathcal A\times\mathcal T$.  In particular, for
$(\alpha,\tau),(\beta,\sigma),(\gamma,\upsilon)$,
\begin{equation}\label{eq:kuranishi-product-source-cocycle}
 \widehat\Psi_{(\gamma,\upsilon),(\beta,\sigma)}
 \widehat\Psi_{(\beta,\sigma),(\alpha,\tau)}
 =\widehat\Psi_{(\gamma,\upsilon),(\alpha,\tau)},
\end{equation}
and
\begin{align}\label{eq:kuranishi-product-obstruction-cocycle}
 &\widehat C_{(\gamma,\upsilon),(\beta,\sigma)}
 \!\left(\widehat\Psi_{(\beta,\sigma),(\alpha,\tau)}(t_\tau)\right)
 \widehat C_{(\beta,\sigma),(\alpha,\tau)}(t_\tau)\\
 &\hspace{5.5em}=
 \widehat C_{(\gamma,\upsilon),(\alpha,\tau)}(t_\tau).
\end{align}

\item For every prescribed jet order $N$, all $N$-jets in the presentation
direction are finite by Theorem~\ref{thm:kuranishi-common-source-atlas}, while
the chart maps $L_{\sigma\tau},M_{\sigma\tau}$ are linear.  Thus every combined
transition jet is a finite matrix expression.  On determinantal tower charts
it is rational in the same inverse minors and in the presentation-atlas
minors, and it commutes with every flat base change preserving those charts.

\item This strict product atlas is a statement at the minimal
Kuranishi/formal-deformation level.  Realizing the chart direction on finite
Fox, Herr, Koszul, or Cayley--Hamilton chain models still
uses the homotopy-coherent $A_\infty/L_\infty$ transitions of
Theorems~\ref{thm:common-source-multiplicative-tower} and
\ref{thm:common-source-lax-monoidal-tower}.  The theorem therefore removes a
mixed coherence existence problem, but it does not manufacture direct
period-parametrix-free or presentation-native chain formulas.
\end{enumerate}
\end{theorem}

\begin{proof}
Equation~\eqref{eq:kuranishi-tower-gauge-equivalence} is immediate from
\eqref{eq:presentation-tower-kuranishi-chart}.  Conjugating
\eqref{eq:kuranishi-direct-gauge-equivalence} by the two $\lambda_{\tau,i}$
gives~\eqref{eq:kuranishi-presentation-gauge-equivalence-in-tower}.
The source interchange is ordinary cancellation:
\[
 (\lambda_{\sigma,1}\lambda_{\tau,1}^{-1})
 (\lambda_{\tau,1}\Psi_{\beta\alpha}\lambda_{\tau,1}^{-1})
 =\lambda_{\sigma,1}\Psi_{\beta\alpha}\lambda_{\tau,1}^{-1}
 =
 (\lambda_{\sigma,1}\Psi_{\beta\alpha}\lambda_{\sigma,1}^{-1})
 (\lambda_{\sigma,1}\lambda_{\tau,1}^{-1}).
\]
The same calculation in the obstruction module, with the argument of
$C_{\beta\alpha}$ transported by $\lambda_{\tau,1}^{-1}$, gives
\eqref{eq:kuranishi-obstruction-interchange}.  Thus the two definitions in
\eqref{eq:kuranishi-product-source-transition} and
\eqref{eq:kuranishi-product-obstruction-transition} agree.

For three combined indices, the source cocycle reduces to
\[
 \lambda_{\upsilon,1}
 \Psi_{\gamma\beta}\Psi_{\beta\alpha}
 \lambda_{\tau,1}^{-1}
 =\lambda_{\upsilon,1}\Psi_{\gamma\alpha}\lambda_{\tau,1}^{-1}
\]
by~\eqref{eq:kuranishi-strict-source-cocycle}; the obstruction cocycle reduces
in the same way to~\eqref{eq:kuranishi-strict-obstruction-cocycle}.  The
finite-effectivity statement follows because tower transport is linear and
the presentation transitions have finite jets by
Theorem~\ref{thm:kuranishi-common-source-atlas}.  The final assertion is
precisely the distinction between the strict common minimal model and its
homotopy-coherent finite-chain realizations in
Theorems~\ref{thm:common-source-multiplicative-tower} and
\ref{thm:common-source-lax-monoidal-tower}.
\end{proof}

\begin{corollary}[Presentation and tower independence commute at the formal deformation level]\label{cor:kuranishi-presentation-tower-product}
After choosing one common intrinsic minimal Kuranishi gauge and the
common-source cohomology coordinates of the chart atlas, finite presentation,
minor, and chart changes form a strictly transitive product atlas on formal
Kuranishi germs.  Thus presentation independence and tower independence do
not create an additional mixed obstruction.  Theorem~\ref{thm:native-twisted-linfinity}
upgrades this formal atlas to strict completed-Koszul transitions and direct
two-term $L_\infty[1]$ tensors.  The stronger native explicitization is treated later in separate layers:
direct ordered chain formulas, third-stage syzygy coherence, intrinsic
common-source cochain gauges, and raw-Herr Tate-dual comparison.  Consequently
this product-atlas result creates no residual mixed obstruction; only the
optional choice of especially strict raw representatives remains
noncanonical.
\end{corollary}


\begin{theorem}[Ordered all-arity twisted Fox--Magnus transport tensors]\label{thm:ordered-twisted-fox-magnus}
Retain a finite-word residual chart as in
Theorem~\ref{thm:arbitrary-residual-word-jets}.  Thus
\[
 \bar\rho:F\longrightarrow \mathrm{GL}_s(k),
 \qquad
 A_i:=\bar\rho(x_i),
\]
and let
\[
 w=a_1\cdots a_L,
 \qquad
 a_\ell=x_{j_\ell}^{\epsilon_\ell},
 \qquad
 \epsilon_\ell\in\{\pm1\}.
\]
For an arity $n\ge1$, choose independent tangent-direction matrices
\[
 X_i^{(1)},\ldots,X_i^{(n)}\in M_s(k)
\]
and noncommuting color variables $z_1,\ldots,z_n$.  In the completed free
associative color algebra put
\[
 Z_i:=\sum_{a=1}^n z_aX_i^{(a)},
 \qquad
 \rho_{\mathbf z}(x_i):=(1+Z_i)A_i.
\]
All statements below may equivalently be read modulo color degree $n+1$, so
only a finite nilpotent algebra is required.

Define the ordered arity-$n$ transport tensor by coefficient extraction
\begin{equation}\label{eq:ordered-twisted-word-jet-definition}
 \mathfrak J_n^{\bar\rho}
 \bigl(w;X^{(1)},\ldots,X^{(n)}\bigr)
 :=[z_1\cdots z_n]\,
 \rho_{\mathbf z}(w)\bar\rho(w)^{-1}.
\end{equation}
Set $\mathfrak J_0^{\bar\rho}(w):=1$.  Then:
\begin{enumerate}[label=\textnormal{(\roman*)}]
\item The tensor in
\eqref{eq:ordered-twisted-word-jet-definition} is multilinear in the $n$
ordered tangent directions and has the closed finite word formula
\begin{equation}\label{eq:ordered-twisted-word-jet}
 \boxed{
 \mathfrak J_n^{\bar\rho}
 \bigl(w;X^{(1)},\ldots,X^{(n)}\bigr)
 =
 \left(
 \sum_{0=b_0\le b_1\le\cdots\le b_L=n}
 \prod_{\ell=1}^L
 C_{\ell; b_{\ell-1},b_\ell}
 \right)\bar\rho(w)^{-1}.
 }
\end{equation}
Writing $q=b-a$, the local ordered block is
\begin{equation}\label{eq:ordered-letter-block}
 C_{\ell;a,b}:=
 \begin{cases}
 A_{j_\ell},
   &\epsilon_\ell=+1,\ q=0,\\
 X_{j_\ell}^{(b)}A_{j_\ell},
   &\epsilon_\ell=+1,\ q=1,\\
 0,
   &\epsilon_\ell=+1,\ q\ge2,\\
 A_{j_\ell}^{-1},
   &\epsilon_\ell=-1,\ q=0,\\
 (-1)^qA_{j_\ell}^{-1}
 X_{j_\ell}^{(a+1)}\cdots X_{j_\ell}^{(b)},
   &\epsilon_\ell=-1,\ q\ge1.
 \end{cases}
\end{equation}
Thus a fixed arity uses only weak compositions of the ordered color word
$1\cdots n$ among the finitely many letters of $w$.  In particular, no bar
diagonal, contraction, rooted tree, or division by a factorial occurs.

\item The ordered tensors satisfy the exact twisted concatenation rule
\begin{align}
 &\mathfrak J_n^{\bar\rho}
 \bigl(uv;X^{(1)},\ldots,X^{(n)}\bigr)
 \notag\\
 &\qquad=
 \sum_{r=0}^n
 \mathfrak J_r^{\bar\rho}
 \bigl(u;X^{(1)},\ldots,X^{(r)}\bigr)
 \operatorname{Ad}_{\bar\rho(u)}\!\left(
 \mathfrak J_{n-r}^{\bar\rho}
 \bigl(v;X^{(r+1)},\ldots,X^{(n)}\bigr)
 \right).
 \label{eq:ordered-twisted-concatenation}
\end{align}
This is an all-arity noncommutative Fox rule: the color word is split only
into a prefix and a suffix, not shuffled.

\item On the diagonal one recovers exactly the ordinary homogeneous word jet
of Theorem~\ref{thm:arbitrary-residual-word-jets}:
\begin{equation}\label{eq:ordered-diagonal-recovers-word-jet}
 \mathfrak J_n^{\bar\rho}(w;X,\ldots,X)
 =J_n^{\bar\rho}(w;X).
\end{equation}
For $n=1$ this is the evaluated twisted Fox Jacobian.  If the residual
representation is scalar and the relators have exact Zassenhaus depth $n$,
then after the usual relation-trace projection the first nonzero ordered
tensor is the matrix-colored Magnus--Fox tensor of
Theorem~\ref{thm:scalar-magnus-deformation}.

\item Formula~\eqref{eq:ordered-twisted-word-jet} is valid over every
characteristic in which the residual word chart is defined.  Unlike symmetric
polarization, it requires no inverse of $n!$.  On a family chart, the same
formula commutes with every base change carrying the residual transport
matrices and the chosen finite word data.
\end{enumerate}
\end{theorem}

\begin{proof}
For a positive letter one has
\[
 \rho_{\mathbf z}(x_j)=(1+Z_j)A_j,
\]
so its color degree is at most one.  For a negative letter,
\[
 \rho_{\mathbf z}(x_j)^{-1}
 =A_j^{-1}(1+Z_j)^{-1}
 =A_j^{-1}\sum_{q\ge0}(-Z_j)^q.
\]
The coefficient of the consecutive color block
$z_{a+1}\cdots z_b$ in the degree-$q=b-a$ term is exactly the final line of
\eqref{eq:ordered-letter-block}.  To obtain the coefficient of the full word
$z_1\cdots z_n$ in the product of the $L$ letter series, one chooses the cut
positions
\[
 0=b_0\le b_1\le\cdots\le b_L=n.
\]
This gives~\eqref{eq:ordered-twisted-word-jet} and multilinearity.

Put
\[
 R_w(\mathbf z):=\rho_{\mathbf z}(w)\bar\rho(w)^{-1}.
\]
For two words,
\[
 R_{uv}(\mathbf z)
 =R_u(\mathbf z)\,
 \operatorname{Ad}_{\bar\rho(u)}R_v(\mathbf z).
\]
The coefficient of a fixed noncommutative color word in a product splits
uniquely at one cut position, proving
\eqref{eq:ordered-twisted-concatenation}.  If all colored tangent matrices are
set equal, every weak composition contributes exactly the corresponding
homogeneous summand in
Equation~\eqref{eq:arbitrary-residual-word-jet}; this proves
\eqref{eq:ordered-diagonal-recovers-word-jet}.  The scalar/Zassenhaus
specialization is then Equation~\eqref{eq:matrix-colored-higher-fox}.  No step
uses division, which proves the final assertion.
\end{proof}

\begin{theorem}[Native noncommutative $A_\infty$ formal slice]\label{thm:native-ordered-ainfinity}
Let the ordered word tensors of
Theorem~\ref{thm:ordered-twisted-fox-magnus} be assembled into the homogeneous
coefficients
\[
 F_r:V^{\otimes r}\longrightarrow W,
 \qquad r\ge1,
\]
of a finite-word relation chart.  The same statement applies after any lci
row reduction whose chosen row operations have been lifted over the central
coefficient base.  Let
\[
 V=T\oplus N,
 \qquad
 W=I\oplus O,
 \qquad
 F_1(t,n)=(An,0),
 \qquad
 A:N\xrightarrow{\sim}I
\]
be the constant-rank splitting used in
Theorem~\ref{thm:word-jet-lyapunov-schmidt}.

Define ordered maps $h_q:T^{\otimes q}\to N$ recursively.  Put
\[
 U_1:=\iota_T:T\longrightarrow V,
 \qquad
 U_q:=\iota_N h_q\quad(q\ge2).
\]
For $q\ge2$ set
\begin{equation}\label{eq:native-ordered-ls-recursion}
 \boxed{
 h_q
 =-A^{-1}
 \sum_{r=2}^q
 \ \sum_{\substack{q_1+\cdots+q_r=q\\q_j\ge1}}
 (F_I)_r
 \bigl(U_{q_1}\otimes\cdots\otimes U_{q_r}\bigr).
 }
\end{equation}
Define the ordered reduced obstruction tensors
\begin{equation}\label{eq:native-ordered-kuranishi-tensor}
 \boxed{
 a_{\alpha,q}
 :=
 \sum_{r=2}^q
 \ \sum_{\substack{q_1+\cdots+q_r=q\\q_j\ge1}}
 (F_O)_r
 \bigl(U_{q_1}\otimes\cdots\otimes U_{q_r}\bigr)
 :T^{\otimes q}\longrightarrow O.
 }
\end{equation}
Then:
\begin{enumerate}[label=\textnormal{(\roman*)}]
\item Equations~\eqref{eq:native-ordered-ls-recursion} and
\eqref{eq:native-ordered-kuranishi-tensor} are the unique degree-by-degree
solution of the noncommutative formal equations
\[
 F_I(t,h(t))=0,
 \qquad
 \kappa_\alpha^{\mathrm{ord}}(t)
 :=F_O(t,h(t))
 =\sum_{q\ge2}a_{\alpha,q}(t^{\otimes q}).
\]
At every fixed degree they are finite planar-composition formulas in the
ordered twisted word tensors and the single inverse minor $A^{-1}$.

\item Use the suspended/bar $A_\infty$ convention: put the tangent term in
suspended degree zero and the obstruction term in suspended degree one.  Define
\begin{equation}\label{eq:native-formal-slice-ainfinity}
 b_q^{\alpha}\big|_{T^{\otimes q}}:=a_{\alpha,q},
 \qquad q\ge2,
\end{equation}
put $b_1=0$, and declare every $b_q$ with at least one obstruction input to be
zero.  Then the bar coderivation $b=\sum b_q$ squares to zero.  Its
Maurer--Cartan equation is exactly
\[
 \sum_{q\ge2}a_{\alpha,q}(t^{\otimes q})=0.
\]
Desuspension gives an ordinary nonunital minimal $A_\infty$ algebra with $T$
in cohomological degree one and $O$ in degree two; the usual universal arity
signs are inserted by the chosen suspension convention.  All Stasheff
identities hold strictly because every nonzero inner multiplication lands in
the obstruction term, while every possible outer multiplication with an
obstruction input is zero.

\item Let
\[
 a_\alpha^\vee:O^\vee
 \longrightarrow
 \widehat T_{\ge2}(T^\vee),
 \qquad
 a_\alpha^\vee:=\sum_{q\ge2}a_{\alpha,q}^\vee.
\]
The completed bar-dual relation algebra of the two-term chart is
\begin{equation}\label{eq:native-nc-relation-algebra}
 \boxed{
 \mathfrak R_\alpha^{\mathrm{nc}}
 :=
 \widehat T(T^\vee)\Big/
 \overline{\left\langle a_\alpha^\vee(O^\vee)\right\rangle_{\mathrm{2s}}}.
 }
\end{equation}
It is exactly the completed noncommutative relation algebra obtained by
eliminating the transverse variables from the ordered word equations.  Thus
\eqref{eq:native-formal-slice-ainfinity} is a direct associative
formal-slice higher-operation formula, rather than a transferred rooted-tree
formula.  This is the word-relation-side analogue of the standard
$A_\infty$ presentations of noncommutative deformation hulls
\cite{WangErickson}; identifying it with a chosen intrinsic cochain minimal
model still requires the corresponding gauge comparison.

\item No factorial occurs.  Hence the ordered formal-slice $A_\infty$ chart is
available in residual characteristic as soon as the linear minor $A$ is
invertible.  In characteristic zero, its commutative abelianization is the
classical reduced Kuranishi series of
Theorem~\ref{thm:word-jet-lyapunov-schmidt}.

\item For the characteristic-zero shifted $L_\infty$ tensors of
Theorem~\ref{thm:native-twisted-linfinity} one has the exact symmetrization
identity
\begin{equation}\label{eq:native-ainfty-linfinity-symmetrization}
 \boxed{
 q_{\alpha,n}(x_1,\ldots,x_n)
 =\sum_{\sigma\in S_n}
 a_{\alpha,n}
 (x_{\sigma(1)},\ldots,x_{\sigma(n)}).
 }
\end{equation}
For degree-one inputs this is the unsigned tensor underlying the usual
$A_\infty$-to-shifted-$L_\infty$ antisymmetrization; depending on suspension
convention there is at most the universal arity sign.  Thus the theorem
retains the ordered information that is lost by the commutative Kuranishi
diagonal.
\end{enumerate}
\end{theorem}

\begin{proof}
In the first component of the relation map the only linear contribution is
$Ah_q$ in degree $q$.  Every nonlinear degree-$q$ contribution is obtained by
choosing an ordered composition $q=q_1+\cdots+q_r$ and substituting the
already known maps $U_{q_j}$ into $(F_I)_r$.  Isolating $Ah_q$ gives
\eqref{eq:native-ordered-ls-recursion}.  Since the linear projection to $O$
is zero, the same ordered-composition sum with $(F_O)_r$ gives
\eqref{eq:native-ordered-kuranishi-tensor}.  This proves existence and
uniqueness recursively.

For~(ii), in the suspended convention every $b_q$ has degree one.  In a
Stasheff/bar-coderivation composite, a nonzero inner operation has only
tangent inputs and therefore lands in the obstruction term.  The outer
operation then contains an obstruction input and vanishes by definition.
Hence $b^2=0$ term by term.  Desuspension is the equivalent ordinary
$A_\infty$ statement.  Dualizing the ordered relation series gives exactly
the closed ideal in \eqref{eq:native-nc-relation-algebra}; this is the
standard bar-dual presentation of such a two-term minimal relation algebra.

The absence of division is visible in
\eqref{eq:native-ordered-ls-recursion}.  Abelianization identifies every
colored diagonal with the ordinary word series by
\eqref{eq:ordered-diagonal-recovers-word-jet}, so the commutative reduced
series is the one from Theorem~\ref{thm:word-jet-lyapunov-schmidt}.
Finally, if
\[
 \kappa_{\alpha,n}(x)=a_{\alpha,n}(x,\ldots,x),
\]
then the characteristic-zero finite-difference polarization of the diagonal
is the sum of all ordered multilinear coefficients, giving
\eqref{eq:native-ainfty-linfinity-symmetrization}.
\end{proof}

\begin{proposition}[Loss of ordered $A_\infty$ data under Kuranishi diagonalization]\label{prop:kuranishi-diagonal-ainfty-no-go}
Let $E$ have characteristic zero, let $\dim_E T\ge2$, let $O\ne0$, and fix
$n\ge2$.  The diagonal map
\[
 \operatorname{Diag}_n:
 \operatorname{Hom}_E(T^{\otimes n},O)
 \longrightarrow
 \operatorname{Poly}_n(T,O),
 \qquad
 a\longmapsto\bigl[x\mapsto a(x,\ldots,x)\bigr]
\]
factors through total symmetrization.  Its kernel is nonzero.  Consequently
no construction whose input is only the commutative homogeneous Kuranishi
coefficient $\kappa_n(x)$, even together with all of its ordinary
polarizations, can reconstruct a genuinely nonsymmetric $n$-ary
$A_\infty$ tensor.
\end{proposition}

\begin{proof}
For every multilinear $a$, define
\[
 \operatorname{Sym}(a)(v_1,\ldots,v_n)
 :=\frac1{n!}\sum_{\sigma\in S_n}
 a(v_{\sigma(1)},\ldots,v_{\sigma(n)}).
\]
Then
\[
 \operatorname{Sym}(a)(x,\ldots,x)=a(x,\ldots,x),
\]
so the diagonal depends only on the total symmetrization of $a$.  Choose
linearly independent $\lambda,\mu\in T^\vee$, a nonzero $o\in O$, and, for
$n>2$, any $\nu_3,\ldots,\nu_n\in T^\vee$.  Then
\[
 a(v_1,\ldots,v_n)
 :=o\,
 \bigl(\lambda(v_1)\mu(v_2)-\mu(v_1)\lambda(v_2)\bigr)
 \prod_{j=3}^n\nu_j(v_j)
\]
is nonzero but satisfies $a(x,\ldots,x)=0$.  Hence
$\ker(\operatorname{Diag}_n)\ne0$.
\end{proof}

\begin{corollary}[Ordered presentation comparison for native $A_\infty$ slices]\label{cor:native-ordered-ainfinity-boundary}
On every finite-word constant-rank residual chart, genuinely twisted ordered
higher deformation tensors are now available directly in all arities from
Equations~\eqref{eq:ordered-twisted-word-jet} and
\eqref{eq:native-ordered-ls-recursion}--
\eqref{eq:native-ordered-kuranishi-tensor}.  They require neither homotopy
transfer nor a chosen bar/Fox contraction, and they remain meaningful in
residual characteristic.  Proposition~\ref{prop:kuranishi-diagonal-ainfty-no-go}
shows why the commutative presentation--tower Kuranishi atlas alone could
never supply this ordered information.

The next theorem constructs the ordered transition formulas for finite marked presentation certificates, including the induced syzygies.  The subsequent path-homotopy, third-coherence, and filtered-decomposable contraction theorems then close the presentation side at every arity required by the native atlas.  The individual intrinsic minimal cochain $A_\infty$ tensors remain gauge dependent; what is presentation-independent is their $A_\infty$ equivalence class, or equivalently the corresponding noncommutative deformation hull with its coherent coordinate atlas.
\end{corollary}


\begin{theorem}[Certified native presentation transition for ordered word relations]\label{thm:certified-native-ordered-transition}
Let
\[
 \mathcal P_\alpha=\langle x_1,\ldots,x_d\mid r_1,\ldots,r_m\rangle,
 \qquad
 \mathcal P_\beta=\langle y_1,\ldots,y_e\mid s_1,\ldots,s_n\rangle
\]
be finite-word transport presentations on residual charts carrying compatible
representations
\[
 \bar\rho_\alpha:F_x\to \mathrm{GL}_s(k),
 \qquad
 \bar\rho_\beta:F_y\to \mathrm{GL}_s(k).
\]
A \emph{finite marked presentation morphism certificate}
$\mathfrak c:\alpha\leftarrow\beta$ consists of:
\begin{enumerate}[label=\textnormal{(\alph*)}]
\item generator words $u_i(y)\in F_y$ defining a homomorphism
$\phi:F_x\to F_y$ with
\[
 \bar\rho_\alpha(x_i)=\bar\rho_\beta(u_i);
\]
\item for every source relator $r_a$, a finite normal-closure identity
\begin{equation}\label{eq:certified-relator-identity}
 \phi(r_a)
 =\prod_{\nu=1}^{L_a}
 c_{a\nu}s_{b_{a\nu}}^{\varepsilon_{a\nu}}c_{a\nu}^{-1},
 \qquad
 \varepsilon_{a\nu}\in\{\pm1\},
\end{equation}
inside $F_y$.
\end{enumerate}
Let $Y=(Y_1,\ldots,Y_e)$ be the left-trivialized matrix coordinates on the
$\beta$ chart,
\[
 \rho_Y(y_j)=(1+Y_j)\bar\rho_\beta(y_j),
\]
and put
\[
 f_b^\beta(Y):=\rho_Y(s_b)-1.
\]
The completed inverse $(1+f_b^\beta)^{-1}$ exists because
$f_b^\beta$ has zero constant term.

Then the certificate has the following direct ordered transition formulas.
\begin{enumerate}[label=\textnormal{(\roman*)}]
\item The source substitution is the noncommutative formal map
\begin{equation}\label{eq:certified-native-source-substitution}
 \boxed{
 \Phi_{\alpha\beta,i}(Y)
 :=\rho_Y(u_i)\bar\rho_\beta(u_i)^{-1}-1.
 }
\end{equation}
Its ordered arity-$q$ Taylor tensor is exactly
\begin{equation}\label{eq:certified-source-ordered-tensor}
 [z_1\cdots z_q]\,\Phi_{\alpha\beta,i}(Z)
 =\mathfrak J_q^{\bar\rho_\beta}
   (u_i;Y^{(1)},\ldots,Y^{(q)}),
\end{equation}
so every fixed arity is given by the finite word formula
\eqref{eq:ordered-twisted-word-jet}.

\item For a fixed identity \eqref{eq:certified-relator-identity}, abbreviate
\begin{align*}
 R_b(Y)&:=1+f_b^\beta(Y),\\
 C_{a\nu}(Y)&:=\rho_Y(c_{a\nu}),\\
 P_{a\nu}(Y)&:=
 C_{a\nu}(Y)R_{b_{a\nu}}(Y)^{\varepsilon_{a\nu}}
 C_{a\nu}(Y)^{-1},\\
 P_{a,<\nu}(Y)&:=P_{a1}(Y)\cdots P_{a,\nu-1}(Y),
 \qquad P_{a,<1}:=1.
\end{align*}
Define
\begin{equation}\label{eq:certified-relator-U-V}
 (U_{a\nu},V_{a\nu})
 :=
 \begin{cases}
 \bigl(P_{a,<\nu}C_{a\nu},\ C_{a\nu}^{-1}\bigr),
 &\varepsilon_{a\nu}=+1,\\[3pt]
 \bigl(-P_{a,<\nu}C_{a\nu}R_{b_{a\nu}}^{-1},\ C_{a\nu}^{-1}\bigr),
 &\varepsilon_{a\nu}=-1.
 \end{cases}
\end{equation}
Then one has the exact two-sided syzygy identity
\begin{equation}\label{eq:certified-relator-telescoping}
 \boxed{
 f_a^\alpha\bigl(\Phi_{\alpha\beta}(Y)\bigr)
 =\sum_{\nu=1}^{L_a}
 U_{a\nu}(Y)\,
 f_{b_{a\nu}}^\beta(Y)\,
 V_{a\nu}(Y).
 }
\end{equation}
Thus presentation change acts directly on the full noncommutative relation
ideal, not only on its commutative zero locus.

\item Let $e_a^\alpha$ and $e_b^\beta$ denote matrix-valued degree-one
semi-free relation symbols, interpreted entrywise if scalar generators are
preferred, with homological differential
\[
 \partial e_a^\alpha=f_a^\alpha,
 \qquad
 \partial e_b^\beta=f_b^\beta,
 \qquad
 \partial Y_j=0.
\]
Equations \eqref{eq:certified-native-source-substitution} and
\eqref{eq:certified-relator-telescoping} define a strict continuous map of
completed semi-free relation dg algebras by
\begin{equation}\label{eq:certified-relation-dga-map}
 e_a^\alpha
 \longmapsto
 \sum_{\nu=1}^{L_a}
 U_{a\nu}(Y)e_{b_{a\nu}}^\beta V_{a\nu}(Y).
\end{equation}
Indeed, applying $\partial$ to the right-hand side gives exactly
\eqref{eq:certified-relator-telescoping}.  If two marked certificates are
composed by literal word substitution and substitution of their chosen
normal-closure identities, the associated dg maps compose strictly.  Every
homogeneous coefficient is a finite planar word expression; no contraction,
rooted-tree transfer, or factorial division occurs.

\item The two elementary stabilization moves have explicit syzygy-complete
forms.

For a generator stabilization adjoining $y$ with relation
$t=yu^{-1}$, write
\[
 \Phi_u(X)=\rho_X(u)\bar\rho(u)^{-1}-1
\]
and introduce the new coordinate
\begin{equation}\label{eq:tietze-generator-contractible-coordinate}
 Y':=(1+Y)(1+\Phi_u(X))^{-1}-1.
\end{equation}
Then the new relation is exactly $f_t=Y'$.  Hence the added generator/relation
pair is a semi-free contractible Koszul extension and is eliminated by the
identity linear minor in the native Lyapunov--Schmidt recursion.

For a redundant-relator stabilization adjoining
\[
 q=\prod_{\nu=1}^{L}c_\nu s_{b_\nu}^{\varepsilon_\nu}c_\nu^{-1},
\]
let
\[
 S_q:=\sum_{\nu=1}^{L}U_\nu e_{b_\nu}V_\nu
\]
be the degree-one expression supplied by
\eqref{eq:certified-relator-U-V}.  Adjoin one degree-two syzygy symbol
$\eta_q$ and set
\begin{equation}\label{eq:tietze-redundant-relator-syzygy}
 \boxed{
 \partial\eta_q=e_q-S_q.
 }
\end{equation}
Equation \eqref{eq:certified-relator-telescoping} gives
$\partial^2\eta_q=0$.  After the triangular change
\[
 e_q':=e_q-S_q,
\]
the new pair is
\[
 \partial\eta_q=e_q',
 \qquad
 \partial e_q'=0.
\]
It is a relative contractible semi-free pair.  Explicitly, on the ideal
containing at least one of the new letters, the standard leftmost-letter
homotopy replaces the first $e_q'$ by $\eta_q$ and kills a first $\eta_q$,
with the Koszul sign from the preceding old word; it satisfies
$\partial H+H\partial=\mathrm{id}$.  Thus adding or removing a redundant
relator is a strict deformation retract once its induced degree-two syzygy is
retained.

\item Consequently every finite zigzag of such certified generator and
relator moves produces a finite, parametrix-free quasi-isomorphism between the
syzygy-complete ordered relation dg algebras.  On constant-rank strata, apply
the explicit noncommutative Lyapunov--Schmidt elimination of
Theorem~\ref{thm:native-ordered-ainfinity} at the two ends.  Conjugating the
preceding dg map by those triangular eliminations gives an ordered
$A_\infty$ quasi-isomorphism between the reduced two-term formal-slice charts.
Every arity of this $A_\infty$ transition is a finite planar composition of:
\[
 \text{the word tensors }\mathfrak J_q,
 \qquad
 A_\alpha^{-1},A_\beta^{-1},
 \qquad
 U_{a\nu},V_{a\nu},
\]
and lower-degree transition tensors.  Thus no intrinsic minimal-model
parametrix is required to write the ordered presentation transition itself.
\end{enumerate}
\end{theorem}

\begin{proof}
Part~(i) is coefficient extraction from the same normalized word series used
in Theorem~\ref{thm:ordered-twisted-fox-magnus}.  For~(ii), the exact group
identity \eqref{eq:certified-relator-identity} gives
\[
 \rho_Y(\phi(r_a))
 =P_{a1}\cdots P_{aL_a}.
\]
Use the noncommutative telescoping identity
\[
 P_1\cdots P_L-1
 =\sum_{\nu=1}^{L}P_1\cdots P_{\nu-1}(P_\nu-1).
\]
For $\varepsilon=+1$,
\[
 C R_b C^{-1}-1=Cf_bC^{-1},
\]
while for $\varepsilon=-1$,
\[
 C R_b^{-1}C^{-1}-1
 =-CR_b^{-1}f_bC^{-1}.
\]
Substitution gives exactly
\eqref{eq:certified-relator-U-V}--\eqref{eq:certified-relator-telescoping}.
Part~(iii) follows by applying the differential to
\eqref{eq:certified-relation-dga-map}; strict compatibility with composition
is built into literal substitution of the marked certificates.

For the generator move, Equation
\eqref{eq:tietze-generator-contractible-coordinate} is an invertible
triangular change because
\[
 1+Y=(1+Y')(1+\Phi_u(X)).
\]
The defining relation $yu^{-1}$ therefore becomes the coordinate $Y'$ and is
removed by an identity linear block.  For the redundant-relator move,
Equation \eqref{eq:tietze-redundant-relator-syzygy} is a differential because
its square is the telescoping relation identity.  The triangular variable
$e_q'$ makes the added semi-free piece the tensor extension by the acyclic
two-term complex $\eta_q\mapsto e_q'$.  The standard tensor trick, using the
leftmost new letter, contracts the relative augmentation ideal and needs no
division.  This proves~(iv).

Finally, the noncommutative Lyapunov--Schmidt changes are triangular formal
coordinate changes whose coefficients are recursively given by
\eqref{eq:native-ordered-ls-recursion}.  Conjugating a finite marked dg
transition by two such changes therefore gives, degree by degree, only finite
planar substitution sums.  Dual bar coordinates are precisely the Taylor
components of the resulting $A_\infty$ morphism.  Since every elementary
stabilization is a quasi-isomorphism, so is every finite certified zigzag.
\end{proof}

\begin{corollary}[Ordered presentation invariance with finite syzygy certificates]\label{cor:ordered-presentation-certified-invariance}
On finite-word constant-rank residual charts, the ordered native deformation
hull is no longer tied to one presentation in the following effective sense.
Every finite marked presentation equivalence for which the generator
substitutions, normal-closure identities, and elementary stabilization
syzygies are supplied determines a parametrix-free ordered $A_\infty$
quasi-isomorphism in every arity.  The formulas are finite and remain valid in
residual characteristic.

What is not made canonical by Theorem~\ref{thm:certified-native-ordered-transition}
is the choice of a path between two arbitrary profinite presentations.  A
direct path-to-path homotopy requires the next layer of identities among
relations.  Theorem~\ref{thm:native-ordered-path-homotopy} below shows that
this next layer is itself native once a finite identity-among-relations
certificate is supplied.  The separate tower-side gap remains the
period-parametrix-free Fox--Herr comparison.
\end{corollary}


\begin{theorem}[Native path homotopy from a finite identity-among-relations certificate]\label{thm:native-ordered-path-homotopy}
Let $\mathfrak c_0,\mathfrak c_1:\alpha\leftarrow\beta$ be two finite marked
presentation morphism certificates as in
Theorem~\ref{thm:certified-native-ordered-transition}, and let
\[
 F_0,F_1:\mathcal R_\alpha^{\mathrm{ord}}\longrightarrow
 \mathcal R_\beta^{\mathrm{syz}}
\]
be their strict maps on completed ordered relation dg algebras.  Here
$\mathcal R_\beta^{\mathrm{syz}}$ is a chosen syzygy-complete semi-free target
through homological degree two, so every degree-one identity-among-relations
cycle occurring below is represented by a specified degree-two filling.  Write
\[
 F_k(X_i)=\Phi_i^{(k)},\qquad k=0,1,
\]
and let $u_i^{(k)}(y)$ be the corresponding generator word.

A \emph{finite marked two-certificate}
$\mathfrak h:\mathfrak c_0\Rightarrow\mathfrak c_1$ consists of the following
finite data.
\begin{enumerate}[label=\textnormal{(\alph*)}]
\item For every source generator $x_i$, a normal-closure identity
\begin{equation}\label{eq:two-certificate-generator-identity}
 w_i:=u_i^{(1)}(u_i^{(0)})^{-1}
 =\prod_{\lambda=1}^{M_i}
 d_{i\lambda}s_{c_{i\lambda}}^{\delta_{i\lambda}}d_{i\lambda}^{-1},
 \qquad \delta_{i\lambda}\in\{\pm1\}.
\end{equation}
Apply the two-sided telescoping formula
\eqref{eq:certified-relator-U-V} to this identity and put
\begin{equation}\label{eq:two-certificate-generator-syzygy}
 D_i:=\sum_{\lambda=1}^{M_i}
 U_{i\lambda}e_{c_{i\lambda}}^\beta V_{i\lambda},
 \qquad
 \partial D_i=\rho_Y(w_i)-1.
\end{equation}
\item For every source relator symbol $e_a^\alpha$, a specified
$\Xi_a\in(\mathcal R_\beta^{\mathrm{syz}})_2$ satisfying
\begin{equation}\label{eq:two-certificate-relation-filling}
 \partial\Xi_a
 =F_1(e_a^\alpha)-F_0(e_a^\alpha)
   -\mathsf H\!\left(f_a^\alpha\right),
\end{equation}
where $\mathsf H$ is the twisted derivation determined below by the generator
part of the certificate.  Equation
\eqref{eq:two-certificate-relation-filling} is the chain-level form of an
identity among relations.
\end{enumerate}

Then the two-certificate determines a direct, parametrix-free homotopy from
$F_0$ to $F_1$ as follows.
\begin{enumerate}[label=\textnormal{(\roman*)}]
\item Define on degree-zero generators
\begin{equation}\label{eq:native-generator-homotopy}
 \boxed{
 \mathsf H(X_i):=D_i\bigl(1+\Phi_i^{(0)}\bigr).
 }
\end{equation}
Since
\[
 1+\Phi_i^{(1)}
 =\rho_Y(w_i)\bigl(1+\Phi_i^{(0)}\bigr),
\]
one has the exact identity
\begin{equation}\label{eq:native-generator-homotopy-boundary}
 \partial\mathsf H(X_i)=\Phi_i^{(1)}-\Phi_i^{(0)}.
\end{equation}
Extend $\mathsf H$ continuously as the degree-$+1$
$(F_1,F_0)$-derivation
\begin{equation}\label{eq:twisted-derivation-leibniz}
 \mathsf H(ab)
 =\mathsf H(a)F_0(b)+(-1)^{|a|}F_1(a)\mathsf H(b).
\end{equation}
In particular, for every invertible degree-zero matrix word $A$,
\begin{equation}\label{eq:twisted-derivation-inverse}
 \mathsf H(A^{-1})
 =-F_1(A)^{-1}\mathsf H(A)F_0(A)^{-1}.
\end{equation}

\item If $W=A_1\cdots A_L$ is any degree-zero source word written as a
product of letter blocks, then
\begin{equation}\label{eq:mixed-fox-homotopy-word}
 \boxed{
 \mathsf H(W)
 =\sum_{\ell=1}^{L}
 F_1(A_1\cdots A_{\ell-1})\,
 \mathsf H(A_\ell)\,
 F_0(A_{\ell+1}\cdots A_L).
 }
\end{equation}
Thus $\mathsf H(f_a^\alpha)$ in
\eqref{eq:two-certificate-relation-filling} is a finite mixed
prefix--suffix Fox formula.  Every ordered Taylor coefficient is obtained by
coefficient extraction from the same finite words and is therefore a finite
planar expression in the tensors $\mathfrak J_q$, the two certificate word
lists, and the lower ordered coefficients.

\item Put
\begin{equation}\label{eq:native-relation-homotopy}
 \mathsf H(e_a^\alpha):=\Xi_a.
\end{equation}
Then
\begin{equation}\label{eq:native-dga-homotopy}
 \boxed{
 \partial\mathsf H+\mathsf H\partial=F_1-F_0
 }
\end{equation}
on the entire completed semi-free relation dg algebra.
Indeed, the right-hand side of
\eqref{eq:two-certificate-relation-filling} is automatically a cycle, so the
only additional datum needed beyond the generator normal-closure identities
is precisely its degree-two syzygy filling.

\item The construction is stable under strict whiskering.  If
$G:\mathcal R_\beta^{\mathrm{syz}}\to\mathcal R_\gamma^{\mathrm{syz}}$
and $L:\mathcal R_\delta^{\mathrm{ord}}\to\mathcal R_\alpha^{\mathrm{ord}}$
are strict certified transition maps, then $G\mathsf H$ and $\mathsf HL$ are
the corresponding native homotopies between $GF_0,GF_1$ and
$F_0L,F_1L$.  Vertical composition may be represented by finite Moore
strings of such elementary two-certificates, with concatenation as a strict
operation.  Hence a finite certificate graph carries explicit two-morphisms
without choosing a common intrinsic contraction.

\item On constant-rank strata, conjugate the endpoint maps and the homotopy by
the triangular noncommutative Lyapunov--Schmidt eliminations of
Theorem~\ref{thm:native-ordered-ainfinity}.  Bar duality sends
$\mathsf H$ to an $A_\infty$ homotopy between the two ordered reduced
$A_\infty$ transition morphisms.  Every fixed arity is again a finite planar
formula and contains no period parametrix, rooted-tree transfer, or factorial
division.
\end{enumerate}
\end{theorem}

\begin{proof}
For~(i), Equation~\eqref{eq:two-certificate-generator-syzygy} and the factorization
\[
 \rho_Y(u_i^{(1)})
 =\rho_Y(w_i)\rho_Y(u_i^{(0)})
\]
give
\[
 \partial\mathsf H(X_i)
 =(\rho_Y(w_i)-1)(1+\Phi_i^{(0)})
 =\Phi_i^{(1)}-\Phi_i^{(0)}.
\]
The twisted Leibniz rule uniquely extends $\mathsf H$ over the completed free
algebra.  Applying it to $AA^{-1}=1$ gives
\eqref{eq:twisted-derivation-inverse}, and induction on word length gives
\eqref{eq:mixed-fox-homotopy-word}.

The graded commutator $\partial\mathsf H+\mathsf H\partial$ is itself an
$(F_1,F_0)$-derivation of degree zero.  By
\eqref{eq:native-generator-homotopy-boundary} it agrees with $F_1-F_0$ on all
degree-zero generators, hence on every degree-zero word.  Therefore
\begin{align*}
 \partial\Bigl(
 F_1(e_a^\alpha)-F_0(e_a^\alpha)
 -\mathsf H(f_a^\alpha)
 \Bigr)
 &=F_1(f_a^\alpha)-F_0(f_a^\alpha)
   -(F_1-F_0)(f_a^\alpha)\\
 &=0.
\end{align*}
The supplied identity-among-relations filling
\eqref{eq:two-certificate-relation-filling} and
\eqref{eq:native-relation-homotopy} therefore make
\eqref{eq:native-dga-homotopy} hold on each $e_a^\alpha$, and hence on the
whole semi-free algebra.  This proves~(ii)--(iii).

Whiskering preserves the twisted derivation rule and the homotopy equation;
finite Moore strings make vertical concatenation literal.  Finally, the
Lyapunov--Schmidt changes are triangular formal coordinate changes with
finite ordered coefficients.  Conjugating the endpoint maps and
$\mathsf H$ therefore preserves finite planarity degree by degree, and the
standard bar description identifies the resulting degree-$+1$ coderivation
with an $A_\infty$ homotopy.  This proves~(iv)--(v).
\end{proof}


\begin{proposition}[Determinantal native filling of the two-certificate]\label{prop:native-two-certificate-minor-filling}
Retain Theorem~\ref{thm:native-ordered-path-homotopy}.  Suppose the relevant
part of the degree-two differential of the chosen finite syzygy-complete
target is represented, in finite free frames over the commutative coefficient
chart, by a formal matrix
\begin{equation}\label{eq:syzygy-differential-matrix}
 S(Y)=S_0+S_{\ge1}(Y):E_2\longrightarrow E_1,
\end{equation}
where $S_{\ge1}(0)=0$.  Assume the residual matrix $S_0$ has constant rank
$r$ on the selected coefficient stratum.  Choose row and column sets $I,J$ of
cardinality $r$ such that
\[
 A_{I,J}:=(S_0)_{I,J}
\]
is invertible, and define the residual syzygy section
\begin{equation}\label{eq:residual-syzygy-section}
 \sigma_0:=\iota_J A_{I,J}^{-1}\pi_I
 \quad\text{on }\operatorname{im}S_0.
\end{equation}
For the discrepancy cycle
\begin{equation}\label{eq:two-certificate-discrepancy}
 Z_a:=F_1(e_a^\alpha)-F_0(e_a^\alpha)
      -\mathsf H(f_a^\alpha),
\end{equation}
write homogeneous ordered expansions
\[
 Z_a=\sum_{q\ge0}Z_{a,q},
 \qquad
 S_{\ge1}=\sum_{q\ge1}S_q,
 \qquad
 \Xi_a=\sum_{q\ge0}\Xi_{a,q}.
\]
Then the filling equation $S(Y)\Xi_a=Z_a$ is solved recursively by
\begin{equation}\label{eq:two-certificate-minor-recursion}
 \boxed{
 \Xi_{a,q}
 =\sigma_0\!\left(
 Z_{a,q}-
 \sum_{j=1}^{q}S_j\Xi_{a,q-j}
 \right),
 \qquad q\ge0.
 }
\end{equation}
At each stage the parenthesized term lies in $\operatorname{im}S_0$ because
the lower degrees have already been solved and the target is syzygy-complete.
Consequently every prescribed ordered Taylor coefficient of every $\Xi_a$ is
a finite planar expression in the certificate word tensors, the coefficients
of $S$, and the single inverse residual minor $A_{I,J}^{-1}$.

If one forgets the noncommuting color probes and works in the underlying
commutative formal coefficient chart on which the corresponding rank-$r$
minor $S_{I,J}(Y)$ remains invertible and the determinantal rank is exactly
$r$, the recursion resums to the ordinary right inverse
\begin{equation}\label{eq:syzygy-minor-section}
 \sigma_{I,J}(Y)
 =\iota_JS_{I,J}(Y)^{-1}\pi_I
 \quad\text{on }\operatorname{im}S(Y),
\end{equation}
so that $\Xi_a=\sigma_{I,J}(Y)Z_a$.  Both forms commute with flat base change
preserving the chosen minor.
\end{proposition}

\begin{proof}
The chosen residual minor makes the $J$-columns of $S_0$ a basis of
$\operatorname{im}S_0$ on the rank-$r$ stratum, so
$S_0\sigma_0=\mathrm{id}$ on that image.  In homogeneous degree $q$, the
equation $S\Xi_a=Z_a$ reads
\[
 S_0\Xi_{a,q}
 =Z_{a,q}-\sum_{j=1}^{q}S_j\Xi_{a,q-j}.
\]
The right-hand side is in $\operatorname{im}S_0$: this is the degree-$q$
compatibility condition obtained after the lower-degree equations have been
solved, and it holds because the syzygy-complete target contains a filling of
$Z_a$.  Applying $\sigma_0$ gives
\eqref{eq:two-certificate-minor-recursion}.  Induction proves the ordered
formal solution degree by degree and shows finiteness of every coefficient.

On the commutative determinantal chart, exact rank $r$ means that all
$(r+1)$-minors vanish after localizing at the chosen invertible $r$-minor.
The Schur-complement identities then give
\[
 S
 =S\iota_JS_{I,J}^{-1}\pi_I S.
\]
Hence Equation~\eqref{eq:syzygy-minor-section} is a right inverse on
$\operatorname{im}S$, and its formal expansion is exactly the preceding
recursion.  Flat base change preserves these identities while the minor
remains invertible.
\end{proof}


\begin{proposition}[Explicit nonlinear kernel frame from one syzygy minor]\label{prop:native-syzygy-kernel-frame}
Retain Proposition~\ref{prop:native-two-certificate-minor-filling}, and work on
the commutative determinantal chart on which the matrix
\[
 S(Y):E_2\longrightarrow E_1
\]
has exact rank $r$ and the $r\times r$ block $S_{I,J}(Y)$ is invertible.  Put
$K:=J^c$.  Define
\begin{equation}\label{eq:native-syzygy-kernel-frame}
 \boxed{
 B_{I,J}(Y)
 :=\iota_K
   -\iota_J S_{I,J}(Y)^{-1}S_{I,K}(Y)
 :A^{K}\longrightarrow E_2.
 }
\end{equation}
Then
\begin{equation}\label{eq:native-syzygy-kernel-exact}
 S(Y)B_{I,J}(Y)=0,
 \qquad
 \pi_KB_{I,J}(Y)=\mathrm{id}_{A^K},
 \qquad
 \operatorname{im}B_{I,J}(Y)=\ker S(Y).
\end{equation}
Consequently
\begin{equation}\label{eq:native-degree-three-carrier}
 0\longrightarrow A^K
 \xrightarrow{\ B_{I,J}(Y)\ }
 E_2\xrightarrow{\ S(Y)\ }\operatorname{im}S(Y)
 \longrightarrow0
\end{equation}
is an explicit finite free degree-three syzygy carrier on the chosen
rank stratum.  Its inverse on the kernel is simply
\begin{equation}\label{eq:native-kernel-coordinate-section}
 \tau_{I,J}:=\pi_K\big|_{\ker S(Y)}.
\end{equation}
No new matrix inverse beyond the same minor $S_{I,J}^{-1}$ is required.

If $(I',J')$ is a second invertible rank-$r$ minor and $K'=(J')^c$, then on
the overlap the two kernel frames have the strict transition
\begin{equation}\label{eq:native-syzygy-frame-transition}
 C_{(I',J'),(I,J)}
 :=\pi_{K'}B_{I,J}:A^K\longrightarrow A^{K'},
\end{equation}
with
\begin{equation}\label{eq:native-syzygy-frame-transition-identity}
 B_{I',J'}C_{(I',J'),(I,J)}=B_{I,J}.
\end{equation}
Hence the transitions are invertible and satisfy the strict cocycle law on
triple overlaps.

Every ordered noncommutative Taylor coefficient of $B_{I,J}$ and of the
transition matrices is finite.  Indeed, if
\[
 S_{I,J}=A+C,
 \qquad A=(S_0)_{I,J},
 \qquad C(0)=0,
\]
then the ordered coefficients $Q_q$ of $S_{I,J}^{-1}$ are given by
\begin{equation}\label{eq:native-syzygy-inverse-ordered-recursion}
 Q_0=A^{-1},
 \qquad
 Q_q=-A^{-1}\sum_{j=1}^{q}C_jQ_{q-j}
 \quad(q\ge1),
\end{equation}
and substitution in \eqref{eq:native-syzygy-kernel-frame} gives the ordered
kernel-frame coefficients without factorial division.
\end{proposition}

\begin{proof}
Write vectors of $E_2$ in the column decomposition $J\sqcup K$.  The $I$-rows
of $S B_{I,J}$ vanish by construction.  Since $S$ has exact rank $r$ after
localizing at the invertible block $S_{I,J}$, every $(r+1)$-minor vanishes.
Equivalently, the Schur complement vanishes:
\[
 S_{I^c,K}-S_{I^c,J}S_{I,J}^{-1}S_{I,K}=0.
\]
Thus all rows of $S B_{I,J}$ vanish.  The identity
$\pi_KB_{I,J}=\mathrm{id}$ proves injectivity.  Conversely, if $v\in\ker S$,
the $I$-row equation gives
\[
 v_J=-S_{I,J}^{-1}S_{I,K}v_K,
\]
so $v=B_{I,J}v_K$.  This proves
\eqref{eq:native-syzygy-kernel-exact}--\eqref{eq:native-kernel-coordinate-section}.

For two minors, both $B_{I,J}$ and $B_{I',J'}$ identify their source with the
same module $\ker S$.  Therefore
\[
 B_{I',J'}\pi_{K'}B_{I,J}=B_{I,J}.
\]
Swapping the two frames gives the inverse transition, and inserting a third
frame proves the cocycle law literally.  Finally,
\eqref{eq:native-syzygy-inverse-ordered-recursion} is coefficient comparison
in $(A+C)\sum_qQ_q=1$ in the completed ordered power-series algebra.
\end{proof}

\begin{theorem}[Determinantal native third coherence between competing two-certificates]\label{thm:native-ordered-third-coherence}
Retain Theorem~\ref{thm:native-ordered-path-homotopy} and suppose the target
syzygy carrier satisfies the exact-rank hypotheses of
Propositions~\ref{prop:native-two-certificate-minor-filling} and
\ref{prop:native-syzygy-kernel-frame}.  Let
\[
 \mathsf H_0,\mathsf H_1:F_0\Rightarrow F_1
\]
be the degree-$+1$ twisted derivation homotopies produced by two competing
finite marked two-certificates with the same endpoint transition maps.  Put
\[
 \Delta\mathsf H:=\mathsf H_1-\mathsf H_0.
\]
Then the difference admits a direct degree-$+2$ native homotopy using only the
same syzygy minor.

\begin{enumerate}[label=\textnormal{(\roman*)}]
\item For every degree-zero source generator define
\begin{equation}\label{eq:native-third-generator-filling}
 \Theta_i
 :=\sigma_{I,J}(Y)\,\Delta\mathsf H(X_i)\in E_2,
\end{equation}
where $\sigma_{I,J}$ is the right inverse on $\operatorname{im}S$ from
\eqref{eq:syzygy-minor-section}.  The difference
$\Delta\mathsf H(X_i)$ is a degree-one cycle because the two homotopies have
the same endpoints; syzygy completeness places it in $\operatorname{im}S$.
Hence
\begin{equation}\label{eq:native-third-generator-boundary}
 \partial\Theta_i=\Delta\mathsf H(X_i).
\end{equation}

Define $\mathsf K(X_i):=\Theta_i$ and extend $\mathsf K$ continuously as the
degree-$+2$ $(F_1,F_0)$-derivation
\begin{equation}\label{eq:degree-two-twisted-derivation}
 \mathsf K(ab)
 =\mathsf K(a)F_0(b)+F_1(a)\mathsf K(b).
\end{equation}
For a degree-zero word $W=A_1\cdots A_L$ this gives the finite mixed Fox
formula
\begin{equation}\label{eq:degree-two-mixed-fox-word}
 \mathsf K(W)
 =\sum_{\ell=1}^{L}
 F_1(A_1\cdots A_{\ell-1})\,
 \mathsf K(A_\ell)\,
 F_0(A_{\ell+1}\cdots A_L),
\end{equation}
with
\[
 \mathsf K(A^{-1})
 =-F_1(A)^{-1}\mathsf K(A)F_0(A)^{-1}.
\]

\item For each source relator symbol put
\begin{equation}\label{eq:native-third-relation-cycle}
 W_a
 :=\Delta\mathsf H(e_a^\alpha)
   +\mathsf K(f_a^\alpha)\in E_2.
\end{equation}
Then $W_a\in\ker S(Y)$.  Define its degree-three kernel coordinate by
\begin{equation}\label{eq:native-third-relation-filling}
 \Omega_a:=\tau_{I,J}(W_a)=\pi_KW_a\in A^K,
 \qquad
 \mathsf K(e_a^\alpha):=\Omega_a,
\end{equation}
where the differential on $A^K$ is the explicit kernel frame
$B_{I,J}$ of \eqref{eq:native-degree-three-carrier}.  Then
\begin{equation}\label{eq:native-third-homotopy-equation}
 \boxed{
 \partial\mathsf K-\mathsf K\partial
 =\mathsf H_1-\mathsf H_0.
 }
\end{equation}
Thus no separately supplied identity-among-identities filling is needed on a
constant-rank determinantal syzygy chart.

\item Every fixed ordered arity of $\mathsf K$ is a finite planar expression
in the two marked two-certificates, the word tensors $\mathfrak J_q$, the
coefficients of $S$, and the single residual inverse minor
$A_{I,J}^{-1}$.  In ordered formal coordinates one may use the recursions
\eqref{eq:two-certificate-minor-recursion} and
\eqref{eq:native-syzygy-inverse-ordered-recursion}; no factorial, bar
contraction, period parametrix, or rooted-tree transfer occurs.

\item Changing the syzygy minor does not create a new existential coherence
problem.  On an overlap, transport $\Omega_a$ by the strict kernel-frame
matrix \eqref{eq:native-syzygy-frame-transition}.  The two resulting
representatives have the same image in $\ker S$, and the frame transitions
satisfy the strict cocycle law.  Hence the determinantal syzygy atlas carries
native third-coherence data compatible with change of minor.

\item Conjugating by the endpoint noncommutative Lyapunov--Schmidt
eliminations and passing to bar dual coordinates sends $\mathsf K$ to an
explicit homotopy between the two ordered $A_\infty$ homotopies induced by
$\mathsf H_0$ and $\mathsf H_1$.  This statement concerns the finite
presentation-certificate branch; it does not assert that the full intrinsic
cochain $A_\infty$ mapping space is truncated at this stage.
\end{enumerate}
\end{theorem}

\begin{proof}
Because both $\mathsf H_0$ and $\mathsf H_1$ satisfy
\[
 \partial\mathsf H_j+\mathsf H_j\partial=F_1-F_0,
\]
their difference obeys
\begin{equation}\label{eq:delta-H-cycle-equation}
 \partial\Delta\mathsf H+\Delta\mathsf H\partial=0.
\end{equation}
On $X_i$ this says that $\Delta\mathsf H(X_i)$ is a degree-one cycle.
Equation~\eqref{eq:native-third-generator-boundary} follows from the right
inverse property of $\sigma_{I,J}$.  The graded commutator
\[
 \partial\mathsf K-\mathsf K\partial
\]
is a degree-$+1$ $(F_1,F_0)$-derivation, so by
\eqref{eq:native-third-generator-boundary} it agrees with
$\Delta\mathsf H$ on every degree-zero word.  In particular,
\begin{equation}\label{eq:native-third-word-boundary}
 \partial\mathsf K(f_a^\alpha)
 =\Delta\mathsf H(f_a^\alpha).
\end{equation}
Now \eqref{eq:delta-H-cycle-equation} on $e_a^\alpha$ gives
\[
 \partial\Delta\mathsf H(e_a^\alpha)
 =-\Delta\mathsf H(f_a^\alpha).
\]
Together with \eqref{eq:native-third-word-boundary}, this yields
\[
 \partial W_a=0,
\]
so $W_a\in\ker S$.  Proposition~\ref{prop:native-syzygy-kernel-frame} gives
\[
 \partial\Omega_a
 =B_{I,J}\pi_KW_a
 =W_a.
\]
Therefore, on a relation generator,
\[
 (\partial\mathsf K-\mathsf K\partial)(e_a^\alpha)
 =W_a-\mathsf K(f_a^\alpha)
 =\Delta\mathsf H(e_a^\alpha),
\]
which proves \eqref{eq:native-third-homotopy-equation}.  The ordered
finiteness follows from the displayed inverse-minor recursions and mixed Fox
formula.  Strict frame compatibility is
\eqref{eq:native-syzygy-frame-transition-identity}.  The final assertion is
the same triangular conjugation/bar-duality argument used in
Theorem~\ref{thm:native-ordered-path-homotopy}.
\end{proof}

\begin{corollary}[Finite native two- and three-coherence on determinantal syzygy strata]\label{cor:native-two-coherence-determinantal}
On a finite-word constant-rank relation chart equipped with a finite
constant-rank syzygy carrier, a marked generator-level two-certificate plus
one invertible residual syzygy minor determines the entire path homotopy of
Theorem~\ref{thm:native-ordered-path-homotopy}.  No separate list of
$\Xi_a$ is needed.  The same minor also constructs the nonlinear kernel frame
$B_{I,J}$ and hence the degree-three carrier.  Consequently any two competing
marked two-certificates with the same endpoint maps admit the explicit third
coherence of Theorem~\ref{thm:native-ordered-third-coherence}, again with no
new inverse matrix.

Thus the local determinantal presentation atlas now contains explicit
one-morphisms, two-morphisms, and the first homotopies between two-morphisms.
Moreover, changes of syzygy minor are strictly cocyclic on the kernel frame.
The next issue is to distinguish decomposable corrections from genuinely new
crossed-resolution data.  Proposition~\ref{prop:decomposable-native-coherence-contraction}
below contracts the relative decomposable ideal generated by this finite carrier,
while Theorem~\ref{thm:relative-decomposable-ainfty-lifting} explains the extra
quotient-level hypothesis needed before one may recursively promote that
contraction to direct all-arity word formulas.  On the tower side, the linear raw Herr--Fox comparison and representative-level multiplicative comparison are already closed.  Theorem~\ref{thm:raw-herr-phantom-reduced-normalization} also fixes their finite-native linear class canonically; literal equality on one unreduced Robba representative is only the optional topological gauge of Remark~\ref{rem:raw-herr-exact-gauge-boundary}.
\end{corollary}

\begin{remark}[Crossed-resolution meaning of the two-certificate]\label{rem:crossed-resolution-two-certificate}
For an ordinary group presentation, identities among relations are organized
by the kernel of the free crossed module on the relators, and free crossed
resolutions make the subsequent syzygy stages explicit
\cite{BrownRazakCrossed,HuebschmannCrossed}.  The role of
$\Xi_a$ above is exactly the evaluated completed-transport shadow of such a
finite degree-two crossed-resolution certificate.  The theorem does not
assert that a finite generating set of these certificates exists uniformly
for every profinite presentation; it says that once a finite certificate is
present, its ordered deformation homotopy has a closed word formula.
\end{remark}

\begin{corollary}[Scope beyond linear third coherence]\label{cor:native-ordered-third-coherence-boundary}
Theorems~\ref{thm:certified-native-ordered-transition},
\ref{thm:native-ordered-path-homotopy}, and
\ref{thm:native-ordered-third-coherence} give explicit ordered
one-morphisms, two-morphisms, and homotopies between competing two-morphisms
on finite determinantal syzygy charts.  Proposition~\ref{prop:native-syzygy-kernel-frame}
constructs the needed degree-three carrier and its overlap maps from the same
invertible syzygy minor, so the previously isolated ``next crossed-syzygy''
obstruction is not an additional existential datum on these strata.

The next theorem shows that, at the level of the length-two linear
full-transport resolutions themselves, these three stages already exhaust the
independent homotopy data.  What may remain beyond them belongs to the
multiplicative/decomposable ordered relation algebra or to a demand for direct
word formulas, not to the linear presentation atlas.
\end{corollary}


\begin{theorem}[Contractibility of the identity linear comparison space]\label{thm:presentation-linear-comparison-contractible}
Let $\Lambda$ be a complete coefficient/transport algebra and let
\[
 0\to P_2\to P_1\to P_0\to A\to0,
 \qquad
 0\to Q_2\to Q_1\to Q_0\to A\to0
\]
be augmented finite-projective resolutions in the pseudocompact derived category.
Then the derived mapping-space component consisting of augmentation-preserving
maps $P_\bullet\to Q_\bullet$ that induce $\operatorname{id}_A$ is contractible.

In a strict length-two chain model, a comparison is represented by a chain map,
a path between two comparisons by a degree-$+1$ chain homotopy, and a homotopy
between such paths by a degree-$+2$ map.  There are no degree-raising strict maps
above $+2$.  Thus the explicit native path and third-coherence formulas of
Theorems~\ref{thm:native-ordered-path-homotopy} and
\ref{thm:native-ordered-third-coherence} provide finite strict witnesses for the
only nonzero chain-homotopy degrees, while all higher linear coherence is unique
up to contractible choice in the derived mapping space.
\end{theorem}

\begin{proof}
Both $P_\bullet$ and $Q_\bullet$ represent the trivial $\Lambda$-module $A$.
Hence in the derived $\infty$-category
\[
 \operatorname{Map}_{\Lambda}(P_\bullet,Q_\bullet)
 \simeq
 \operatorname{Map}_{\Lambda}(A,A).
\]
For $j>0$ one has
\[
 \pi_j\operatorname{Map}_{\Lambda}(A,A)
 \cong
 \operatorname{Ext}^{-j}_{\Lambda}(A,A)=0,
\]
because $A$ lies in the heart of the standard derived $t$-structure.  On
$\pi_0$, the requirement that the augmentation-induced endomorphism be
$\operatorname{id}_A$ selects the single identity class; equivalently, the
comparison theorem for projective resolutions says that any two such strict
lifts are chain homotopic.  The selected component is therefore contractible.
Finally, because both strict resolutions are concentrated in homological degrees
$0,1,2$, their continuous mapping complex has no degree-raising cochains above
$+2$, which gives the stated finite strict realization of the first two homotopy
levels.
\end{proof}

\begin{corollary}[No higher linear presentation obstruction]\label{cor:no-higher-linear-presentation-obstruction}
The augmentation-preserving identity component of the linear Fox/full-transport comparison space is contractible.  In strict length-two chain coordinates, the only nonzero degree-raising homotopy operators occur in degrees $+1$ and $+2$.  Direct word-native formulas in arbitrarily high coherence arity are therefore not required to globalize the finite atlas.  Multiplicative and simplicial higher coherence is supplied through the common intrinsic cochain object by the common-source and Cech--Stasheff constructions.  Proposition~\ref{prop:decomposable-native-coherence-contraction} and Theorem~\ref{thm:relative-decomposable-ainfty-lifting} below isolate what can additionally be made direct inside the ordered relation algebra without turning that optional explicitization into a mainline hypothesis.
\end{corollary}


\begin{proposition}[Filtered contraction of the relative decomposable ideal]\label{prop:decomposable-native-coherence-contraction}
Work on a finite-word formal determinantal syzygy chart satisfying
Propositions~\ref{prop:native-two-certificate-minor-filling} and
\ref{prop:native-syzygy-kernel-frame}.  Put
\[
 \sigma:=\sigma_{I,J}(Y),\qquad
 B:=B_{I,J}(Y),\qquad
 \tau:=\pi_K|_{\ker S(Y)}.
\]
The exact syzygy carrier
\[
 0\longrightarrow A^K\xrightarrow{B}E_2\xrightarrow{S}\operatorname{im}S\longrightarrow0
\]
admits the explicit contraction
\begin{equation}\label{eq:syzygy-three-term-contraction}
 h_1:=\sigma:\operatorname{im}S\to E_2,
 \qquad
 h_2:=\tau(1-\sigma S):E_2\to A^K,
\end{equation}
for which
\[
 Sh_1=1,\qquad Bh_2+h_1S=1,\qquad h_2B=1.
\]

Let $\mathfrak B_{\mathrm{dec}}$ be the closed relative tensor/operadic ideal in
the completed semi-free ordered relation algebra generated by this contractible
syzygy carrier, and filter it by positive residual-augmentation together with
positive syzygy-letter degree.  On the associated graded, the standard tensor
trick, contracting the leftmost syzygy letter, extends
\eqref{eq:syzygy-three-term-contraction} to a contraction $h_0$ of
$\operatorname{gr}\mathfrak B_{\mathrm{dec}}$.  Suppose the full differential is
written
\[
 \partial=\partial_0+\delta
\]
with $\delta$ strictly raising this complete filtration, as it does for the
positive-order transport corrections on the chosen formal chart.  Then the
complete filtered perturbation lemma gives a contraction
\begin{equation}\label{eq:decomposable-coherence-perturbation-contraction}
 \boxed{
 h_{\mathrm{dec}}
 :=h_0(1+\delta h_0)^{-1}
 =\sum_{m\ge0}(-1)^m h_0(\delta h_0)^m .}
\end{equation}
In particular every closed element of $\mathfrak B_{\mathrm{dec}}$ is filled
inside $\mathfrak B_{\mathrm{dec}}$.  Modulo filtration degree $>N$, only
finitely many terms of \eqref{eq:decomposable-coherence-perturbation-contraction}
contribute, so each prescribed $N$-jet is a finite planar expression in the
certificate word tensors, the coefficients of $S(Y)$, and the same inverse
minor $S_{I,J}(Y)^{-1}$.
\end{proposition}

\begin{proof}
The three contraction identities follow from $S\sigma=1$ on
$\operatorname{im}S$, from $1-\sigma S$ having image in
$\ker S=\operatorname{im}B$, and from $\tau B=1$.  Because the ordered relation
algebra is a relative semi-free completed tensor algebra in the certificate and
syzygy letters, the ideal generated by the contractible syzygy carrier has the
standard leftmost-letter contraction on the associated graded.  The positive
transport correction vanishes at the residual point and therefore raises the
chosen filtration.  The complete filtered perturbation lemma then applies and
gives \eqref{eq:decomposable-coherence-perturbation-contraction}; filtration
growth gives the finite-jet assertion \cite{ManettiPerturbation,VokrinekHPT}.
\end{proof}

\begin{theorem}[Relative filtered lifting criterion for direct presentation formulas]\label{thm:relative-decomposable-ainfty-lifting}
Retain Proposition~\ref{prop:decomposable-native-coherence-contraction}.  Let
$\mathfrak B$ be the completed bar/operadic complex in which one is recursively
constructing a direct ordered $A_\infty$ transition or a simplicial coherence
cell, and let
\[
 \overline{\mathfrak B}:=\mathfrak B/\mathfrak B_{\mathrm{dec}}.
\]
Assume all components below total arity $N$ have been constructed and let
$\mathcal O_N\in\mathfrak B$ be the usual arity-$N$ obstruction.  If its image
in $\overline{\mathfrak B}$ vanishes, then $\mathcal O_N$ lies in
$\mathfrak B_{\mathrm{dec}}$, is closed by the lower Stasheff identities, and
\begin{equation}\label{eq:relative-decomposable-ainfty-filler}
 \mathcal F_N:=-h_{\mathrm{dec}}(\mathcal O_N)
\end{equation}
fills the arity-$N$ equation, with the sign adjusted to the chosen bar
convention.  The same statement holds for $A_\infty$ homotopies and higher
simplicial fillers.

Consequently, if the quotient-level $A_\infty$ and coherence identities are
already satisfied in every arity, the entire remaining decomposable correction
lifts recursively by native finite formulas.  This theorem does \emph{not}
assert that the quotient-level identities hold automatically for an arbitrary
ordered presentation model; verifying that stronger direct-word statement is a
separate optional explicitization problem.
\end{theorem}

\begin{proof}
The arity-$N$ $A_\infty$ or coherence equation is linear in its new unknown
component after all lower components are fixed.  Its known part
$\mathcal O_N$ is closed by the lower identities.  The quotient hypothesis says
exactly that this cycle belongs to the contracted relative ideal
$\mathfrak B_{\mathrm{dec}}$.  Proposition~\ref{prop:decomposable-native-coherence-contraction}
therefore supplies the filler \eqref{eq:relative-decomposable-ainfty-filler}.
Induction gives the final assertion whenever the quotient condition is known in
all arities.
\end{proof}

\begin{remark}[Boundary of direct word-native higher coherence]\label{rem:direct-word-higher-coherence-boundary}
The finite native atlas does not require the quotient hypothesis of
Theorem~\ref{thm:relative-decomposable-ainfty-lifting}: its multiplicative and
simplicial higher coherence is already supplied through the common intrinsic
cochain object by Theorems~\ref{thm:common-source-multiplicative-tower} and
\ref{thm:cech-stasheff-gauges}.  The relative lifting theorem only says that,
whenever a direct ordered presentation realizes the required quotient-level
identity, all decomposable corrections are then forced by the existing finite
syzygy carrier.  Thus a possible nondecomposable crossed-resolution class is a
question about a stronger purely word-native presentation, not an obstruction to
the atlas constructed in this paper.
\end{remark}

\begin{theorem}[Native explicit formula for the twisted shifted $L_\infty$ deformation chart]\label{thm:native-twisted-linfinity}
Work over a characteristic-zero coefficient field $E$ and on a constant-rank
finite-word/lci stratum on which
Theorems~\ref{thm:arbitrary-residual-word-jets},
\ref{thm:word-jet-lyapunov-schmidt}, and
\ref{thm:kuranishi-presentation-tower-interchange} apply.  Let
\[
 F=(F_I,F_O):T\oplus N\longrightarrow I\oplus O
\]
be one of the finite lci relation maps after the linear splitting used in
Theorem~\ref{thm:word-jet-lyapunov-schmidt}, and let
\[
 h:T\longrightarrow N,
 \qquad
 \kappa_\alpha(t)=F_O(t,h(t))
 =\sum_{n\ge2}\kappa_{\alpha,n}(t)
\]
be the resulting reduced obstruction map on a presentation/minor chart
$\alpha$.  Here $\kappa_{\alpha,n}$ is homogeneous of degree $n$.

For every $n\ge2$, define the symmetric polarization
\begin{equation}\label{eq:native-mc-polarization}
 q_{\alpha,n}(x_1,\ldots,x_n)
 :=
 \sum_{S\subseteq\{1,\ldots,n\}}
 (-1)^{n-|S|}
 \kappa_{\alpha,n}\!\left(\sum_{i\in S}x_i\right),
 \qquad x_i\in T,
\end{equation}
with the empty-set summand equal to $\kappa_{\alpha,n}(0)=0$.  Then:
\begin{enumerate}[label=\textnormal{(\roman*)}]
\item The map $q_{\alpha,n}:\operatorname{Sym}^n_E T\to O$ is $E$-multilinear and
\begin{equation}\label{eq:native-mc-diagonal}
 \kappa_{\alpha,n}(x)
 =\frac1{n!}q_{\alpha,n}(x,\ldots,x).
\end{equation}
Declare every higher bracket with at least one input in $O$ to be zero, put
$q_1=0$ on the reduced slice, and use the symmetric $L_\infty[1]$ convention.
Then the maps $q_{\alpha,n}$ define a minimal two-term $L_\infty[1]$ algebra
on $T\oplus O[-1]$ whose Maurer--Cartan equation is exactly
\[
 \kappa_\alpha(t)=0.
\]
Equivalently, its completed Chevalley--Eilenberg algebra is the completed
Koszul dg algebra of the reduced equations $\kappa_\alpha$.

\item The formula is native to the transport word chart.  For every fixed
arity $n$, its coefficients are finite expressions obtained by the following
finite chain:
\begin{equation}\label{eq:native-mc-explicit-chain}
 \text{finite transport words}
 \xrightarrow{\ \eqref{eq:arbitrary-residual-word-jet}\ }
 J_q^{\bar\rho}
 \xrightarrow{\ \eqref{eq:word-ls-recursion}\ }
 h_q
 \xrightarrow{\ \eqref{eq:word-ls-kuranishi-jet}\ }
 \kappa_{\alpha,q}
 \xrightarrow{\ \eqref{eq:native-mc-polarization}\ }
 q_{\alpha,q}.
\end{equation}
Thus $q_{\alpha,n}$ uses only the residual transport matrices
$\bar\rho(x_i)^{\pm1}$, the finite relator syntax, the lci row-reduction
matrices when present, one inverse Jacobian minor from the
Lyapunov--Schmidt splitting, and finite additions and multiplications.  It
uses no bar diagonal, homotopy-transfer tree, Fox contraction, period
parametrix, or additional higher component.

\item The construction already exists before the minimal reduction.  If
$g=\sum_{n\ge1}g_n:V\to W$ is a finite lci regular-sequence map from
Theorem~\ref{thm:transport-explicit-lci-row-reduction}, then
\begin{equation}\label{eq:native-lci-polarization}
 Q_n(v_1,\ldots,v_n)
 :=\sum_{S\subseteq\{1,\ldots,n\}}
 (-1)^{n-|S|}g_n\!\left(\sum_{i\in S}v_i\right)
\end{equation}
with $Q_1=g_1=dg|_0$ gives a two-term $L_\infty[1]$ presentation whose completed
Chevalley--Eilenberg algebra is precisely the lci Koszul carrier.  The
Lyapunov--Schmidt recursion removes the contractible linear pair
$N\xrightarrow{\sim}I$ and produces the minimal slice
\eqref{eq:native-mc-polarization}.  Hence the native higher tensors are not
imported from homotopy transfer; they are the Taylor coefficients of the
already explicit derived lci equations.

\item Presentation/minor changes lift strictly to the corresponding completed
Koszul dg algebras.  Indeed, for the transitions of
Theorem~\ref{thm:kuranishi-common-source-atlas},
\[
 \kappa_\beta\!\left(\Psi_{\beta\alpha}(t)\right)
 =C_{\beta\alpha}(t)\kappa_\alpha(t),
\]
define the dg-algebra map by pullback on degree-zero coordinates
$t_\beta\mapsto\Psi_{\beta\alpha}(t_\alpha)$ and on the column $e_\beta$ of
Koszul generators by
\begin{equation}\label{eq:native-koszul-transition}
 e_\beta\longmapsto C_{\beta\alpha}(t_\alpha)e_\alpha.
\end{equation}
Because the base coordinates have zero differential,
\[
 d\bigl(C_{\beta\alpha}(t)e_\alpha\bigr)
 =C_{\beta\alpha}(t)\kappa_\alpha(t)
 =\kappa_\beta\!\left(\Psi_{\beta\alpha}(t)\right).
\]
The strict source and obstruction cocycles
\eqref{eq:kuranishi-strict-source-cocycle}--\eqref{eq:kuranishi-strict-obstruction-cocycle}
therefore give strict cocycles of completed Koszul dg algebras.  Dually, the
native two-term $L_\infty[1]$ charts form a strictly transitive finite-jet
atlas; the individual tensors $q_{\alpha,n}$ transform triangularly because a
nonlinear source change mixes lower arities, but no additional coherence datum
is required.

\item Tower transport is linear and hence acts directly on every arity.  With
notation from Theorem~\ref{thm:kuranishi-presentation-tower-interchange}, set
\[
 q_{\alpha,\tau,n}(x_1,\ldots,x_n)
 :=\lambda_{\tau,2}\,
 q_{\alpha,n}(\lambda_{\tau,1}^{-1}x_1,\ldots,
              \lambda_{\tau,1}^{-1}x_n).
\]
Then
\begin{equation}\label{eq:native-linfinity-tower-covariance}
 q_{\alpha,\sigma,n}
 (L_{\sigma\tau}x_1,\ldots,L_{\sigma\tau}x_n)
 =M_{\sigma\tau}
 q_{\alpha,\tau,n}(x_1,\ldots,x_n),
\end{equation}
and the presentation--tower interchange is strict at the completed Koszul
level as well.

\item In arity two, after the chosen tangent/obstruction identification, the
native minimal tensor represents the ordinary cup--commutator obstruction;
on a Fox chart it agrees with the antisymmetrization of the twisted tensor-Fox
cup matrix of Equation~\eqref{eq:twisted-tensor-fox-cup}.  If the residual
representation is scalar and the defining relations have exact Zassenhaus
depth $n$, all lower homogeneous terms vanish, so identity-linear source and
obstruction-frame changes cannot alter the first nonzero homogeneous tensor.
Equation~\eqref{eq:native-mc-polarization} therefore recovers the
matrix-colored Magnus--Fox obstruction of
Theorem~\ref{thm:scalar-magnus-deformation} as the first nonzero native
$L_\infty[1]$ deformation tensor.
\end{enumerate}
\end{theorem}

\begin{proof}
For a homogeneous degree-$n$ polynomial $P:T\to O$ over a characteristic-zero
field, the iterated finite difference at the origin gives its symmetric
polarization
\[
 \operatorname{Pol}_n(P)(x_1,\ldots,x_n)
 =\sum_{S\subseteq\{1,\ldots,n\}}
 (-1)^{n-|S|}P\!\left(\sum_{i\in S}x_i\right).
\]
This is multilinear and satisfies
$P(x)=\operatorname{Pol}_n(P)(x,\ldots,x)/n!$, proving
\eqref{eq:native-mc-diagonal}.  Apply this to every homogeneous part of
$\kappa_\alpha$.

In the symmetric $L_\infty[1]$ convention, every $q_n$ has degree one.  The
only nonzero brackets just defined have all inputs in the tangent term and
output in the obstruction term, while every bracket with an obstruction input
is zero.  Hence every composite appearing in an $L_\infty[1]$ Jacobi identity
vanishes: an inner nonzero bracket lands in the obstruction term and cannot be
inserted into another nonzero bracket.  Thus the identities hold.  The dual
completed Chevalley--Eilenberg differential sends the obstruction generators
to the components of $\kappa_\alpha$ and annihilates the tangent coordinates,
which is exactly the completed Koszul differential.  The same argument with
the linear term included proves~(iii).

The explicitness in~(ii) is the direct composition of
Theorems~\ref{thm:arbitrary-residual-word-jets} and
\ref{thm:word-jet-lyapunov-schmidt}.  No transfer data enter those recursions.
The linear pair $N\xrightarrow{\sim}I$ is contractible, and solving it by
\eqref{eq:word-ls-recursion} gives the equivalent reduced Koszul model, so the
minimal slice has the displayed brackets.

For~(iv), the differential calculation following
\eqref{eq:native-koszul-transition} proves that the coordinate/frame change is
a dg-algebra map.  Its inverse is obtained from
$\Psi_{\alpha\beta}$ and $C_{\alpha\beta}$, and the strict cocycle formulas of
Theorem~\ref{thm:kuranishi-common-source-atlas} give strict composition.
Equation~\eqref{eq:native-linfinity-tower-covariance} follows immediately by
polarizing Equation~\eqref{eq:presentation-tower-kuranishi-chart}; the strict
mixed interchange is inherited from
Equations~\eqref{eq:kuranishi-source-interchange} and
\eqref{eq:kuranishi-obstruction-interchange}.

Finally, the quadratic Maurer--Cartan obstruction of the adjoint dg-Lie
algebra is the cup--commutator bracket, and
Theorem~\ref{thm:twisted-tensor-fox-cup} computes its Fox coefficient product.
If all homogeneous relation terms below degree $n$ vanish, a formal source
change with identity linear part and an obstruction-frame change equal to the
identity at the origin preserve the first nonzero homogeneous term.  The
scalar/Zassenhaus specialization is therefore exactly the tensor of
Theorem~\ref{thm:scalar-magnus-deformation} after the standard
$A_\infty$-to-$L_\infty$ antisymmetrization convention.
\end{proof}

\begin{corollary}[Native deformation-level higher explicitization]\label{cor:native-deformation-explicitization}
At every characteristic-zero point lying on a finite-word/lci constant-rank
chart, the full formal-slice deformation equation admits a direct finite
higher-operation formula in every arity: its two-term minimal
$L_\infty[1]$ tensors are the polarizations
\eqref{eq:native-mc-polarization} of the explicitly reduced transport-word
jets.  Presentation/minor changes are strict completed-Koszul dg-algebra
isomorphisms, and chart changes act by the linear covariance
\eqref{eq:native-linfinity-tower-covariance}.  Thus the native explicit-formula
problem is closed for the Maurer--Cartan/$L_\infty$ formal-slice branch.

The stronger branches are handled later without reopening this formal-slice
statement.  The native presentation branch gives explicit one-, two-, and
three-stage data, while Theorem~\ref{thm:presentation-linear-comparison-contractible}
shows that the identity component of the derived linear comparison space is contractible; these explicit low-degree witnesses therefore carry no further linear invariant.
All multiplicative and simplicial higher coherence required by the finite atlas
is supplied through the common intrinsic cochain object; the stronger demand for
direct ordered-word formulas is governed only by the relative criterion of
Theorem~\ref{thm:relative-decomposable-ainfty-lifting}.  Intrinsic cochain and
Fox--Herr comparisons are supplied by the all-Sylow/common-source and raw-Herr
Tate-dual constructions, and Cayley--Hamilton evaluation makes every prescribed
finite higher component finite.
\end{corollary}

\begin{warning}[Native slice brackets are gauges, not canonical cochain brackets]\label{warn:word-jets-not-higher-brackets}
The matrices $J_q^{\bar\rho}(r;X)$ and the reduced series
$\kappa_{\mathrm{word}}$ depend on the chosen generators, relator words,
left trivialization, splittings, inverse minor, and framed coordinates.  For
nontrivial adjoint action they should not be identified term by term with a
canonical transferred cochain $A_\infty/L_\infty$ model.  Theorem~\ref{thm:native-twisted-linfinity}
does, however, turn each reduced word/lci chart into a genuine two-term
$L_\infty[1]$ formal-slice model, and the strict Kuranishi transitions turn
these models into a coherent atlas.  The invariant object is still the
formal deformation problem or its intrinsic cochain dg-Lie algebra; different
native slice tensors are related by the explicit formal gauge transformations
of Theorems~\ref{thm:kuranishi-common-source-atlas} and
\ref{thm:kuranishi-presentation-tower-interchange}.  The scalar/Zassenhaus
case of Theorem~\ref{thm:scalar-magnus-deformation} is stronger because all
lower homogeneous terms vanish and the first nonzero tensor itself is
unchanged by identity-linear gauge transformations after the chosen
relation-trace normalization.
\end{warning}

\begin{corollary}[Explicit symplectic Demu\v{s}kin chart]\label{cor:demushkin-symplectic-chart}
Assume that $p$ is odd and that $P$ has the standard Demu\v{s}kin presentation
\begin{equation}\label{eq:standard-demushkin-relator}
r
=
x_1^q[x_1,x_2][x_3,x_4]\cdots[x_{d-1},x_d],
\qquad
q=0\ \text{or}\ q=p^f.
\end{equation}
Then, modulo terms of Magnus degree at least three,
\begin{equation}\label{eq:demushkin-quadratic-symbol}
\mathcal M(r)
\equiv
1+
\sum_{k=1}^{d/2}
\left(
X_{2k-1}X_{2k}-X_{2k}X_{2k-1}
\right).
\end{equation}
Consequently one may normalize $\omega$ so that
\[
\chi_{2k-1}\smile\chi_{2k}=\omega,
\qquad
\chi_{2k}\smile\chi_{2k-1}=-\omega,
\]
and all other products of basis elements vanish.  The binary Cartan symbol is therefore the standard nondegenerate symplectic matrix.
\end{corollary}

\begin{proof}
For the commutator convention $[x,y]=x^{-1}y^{-1}xy$, its Magnus expansion begins
\[
\mathcal M([x,y])
=1+XY-YX+\text{terms of degree at least three}.
\]
Products of the commutator factors add their degree-two symbols.  In characteristic $p$, the first nonconstant term of $\mathcal M(x_1^{p^f})=(1+X_1)^{p^f}$ has degree $p^f\ge p\ge3$, while the term is absent when $q=0$.  Formula \eqref{eq:demushkin-quadratic-symbol} follows, and Theorem~\ref{thm:quadratic-fox-cup} gives the asserted products.  Nondegeneracy agrees with the Poincar\'e-duality characterization of Demu\v{s}kin groups \cite{LabuteClassification}.
\end{proof}


\subsection{Integral Fox torsion and Bockstein lifting depth}

The mod-$p$ quadratic symbol forgets the power term, but the integral Fox complex does not.  Let
\[
\varepsilon_{\Zp}:\Zp[[F]]\longrightarrow\Zp
\]
be augmentation and define the integral linear Fox vector
\begin{equation}\label{eq:integral-linear-fox-vector}
\lambda(r)
:=
\left(
\varepsilon_{\Zp}(\partial_1r),\ldots,
\varepsilon_{\Zp}(\partial_dr)
\right)
\in\Zp^d.
\end{equation}
Its entries are the $p$-adic exponent sums of the chosen generators in the relation.  Under a change of minimal generators and a change of the normal generator of the relation module, $\lambda(r)$ is multiplied on the right by an element of $\mathrm{GL}_d(\Zp)$ and by an overall unit.  Hence the ideal
\begin{equation}\label{eq:integral-relation-ideal}
\mathfrak q(P)
:=
\bigl(
\varepsilon_{\Zp}(\partial_1r),\ldots,
\varepsilon_{\Zp}(\partial_dr)
\bigr)
\subseteq\Zp
\end{equation}
is intrinsic.

\begin{theorem}[The integral linear Fox symbol recovers the Demu\v{s}kin invariant]\label{thm:integral-fox-recovers-q}
Let $p$ be odd and let $P$ have the standard Demu\v{s}kin presentation \eqref{eq:standard-demushkin-relator}.  Then
\begin{equation}\label{eq:standard-integral-fox-vector}
\lambda(r)=(q,0,\ldots,0),
\qquad
\mathfrak q(P)=(q).
\end{equation}
Moreover,
\begin{equation}\label{eq:demushkin-abelianization-from-fox}
P^{\mathrm{ab}}
\cong
\Zp^{d-1}\oplus\Zp/q\Zp,
\end{equation}
with the convention $\Zp/0\Zp=\Zp$.  Thus the integral first Fox derivatives recover $q$, including the case $q=0$.
\end{theorem}

\begin{proof}
Augmentation of a Fox derivative is the exponent sum of the corresponding generator.  Every commutator has zero exponent sum in every generator, while $x_1^q$ has exponent sum $q$ in $x_1$ and zero exponent sum in the other generators.  This proves \eqref{eq:standard-integral-fox-vector}.

Abelianizing the one-relator presentation gives the cokernel of the map
\[
\Zp\longrightarrow\Zp^d,
\qquad
1\longmapsto\lambda(r).
\]
For the vector in \eqref{eq:standard-integral-fox-vector}, its Smith normal form is the diagonal entry $q$ followed by zeros, which gives \eqref{eq:demushkin-abelianization-from-fox}.  The ideal generated by the entries is invariant under the allowed basis and relation changes, so it recovers the intrinsic invariant $q(P)$ of Labute's classification \cite{LabuteClassification}.
\end{proof}

Cohomologically, the same invariant appears by applying the integral Fox resolution to the trivial module $\Zp$.  In the chosen basis the cochain complex is
\begin{equation}\label{eq:integral-demushkin-trivial-complex}
\Zp
\xrightarrow{0}
\Zp^d
\xrightarrow{(a_1,\ldots,a_d)\mapsto qa_1}
\Zp.
\end{equation}
For $m\ge1$, let
\begin{equation}\label{eq:bockstein-lifting-filtration}
L_m(P)
:=
\operatorname{Im}\!\left(
H^1(P,\mathbf Z/p^m)
\longrightarrow
H^1(P,\mathbf F_p)
\right).
\end{equation}
This filtration measures actual integral liftability rather than merely the first connecting homomorphism.

\begin{proposition}[Bockstein lifting depth]\label{prop:bockstein-lifting-depth}
Let $\chi_1,\ldots,\chi_d$ be the mod-$p$ basis used in Theorem~\ref{thm:quadratic-fox-cup}.
If $q=p^f\ne0$, then
\begin{equation}\label{eq:bockstein-lifting-depth-formula}
L_m(P)
=
\begin{cases}
H^1(P,\mathbf F_p),&1\le m\le f,\\[2pt]
\operatorname{span}_{\mathbf F_p}\{\chi_2,\ldots,\chi_d\},&m\ge f+1.
\end{cases}
\end{equation}
If $q=0$, then $L_m(P)=H^1(P,\mathbf F_p)$ for every $m$.  Consequently
\begin{equation}\label{eq:q-from-lifting-depth}
f
=
\max\{m:\ L_m(P)=H^1(P,\mathbf F_p)\}
\end{equation}
when $q=p^f$, while the right-hand side is infinite when $q=0$.

For the ordinary Bockstein
\[
\beta_1:H^1(P,\mathbf F_p)\longrightarrow H^2(P,\mathbf F_p)
\]
associated with $0\to\mathbf F_p\to\mathbf Z/p^2\to\mathbf F_p\to0$, and with the Fox orientation of the relation $2$-cell,
\begin{equation}\label{eq:first-bockstein-demushkin}
\beta_1(\chi_1)
=
\left(\frac{q}{p}\bmod p\right)\omega,
\qquad
\beta_1(\chi_i)=0\quad(i>1).
\end{equation}
Thus the first Bockstein detects precisely the case $v_p(q)=1$, whereas the full lifting filtration \eqref{eq:bockstein-lifting-depth-formula} recovers every finite value of $v_p(q)$.
\end{proposition}

\begin{proof}
Reducing \eqref{eq:integral-demushkin-trivial-complex} modulo $p^m$, a vector $(a_1,\ldots,a_d)$ is a $1$-cocycle exactly when
\[
q a_1\equiv0\pmod{p^m}.
\]
If $q=p^f$ and $m\le f$, this imposes no condition.  If $m\ge f+1$, it forces $a_1$ to be divisible by $p^{m-f}$, so its reduction modulo $p$ has zero $\chi_1$-component.  The other coordinates remain arbitrary.  This proves \eqref{eq:bockstein-lifting-depth-formula} and \eqref{eq:q-from-lifting-depth}; the case $q=0$ is immediate.

For \eqref{eq:first-bockstein-demushkin}, lift the mod-$p$ cocycle $\chi_1$ to the first basis vector over $\mathbf Z/p^2$.  Its differential is $q$ times the relation $2$-cochain.  Under the identification $p(\mathbf Z/p^2)\cong\mathbf F_p$, division by $p$ gives the displayed coefficient.  The remaining basis vectors have zero differential already over $\Zp$.
\end{proof}

\begin{corollary}[A finite faithful symbol in the odd-prime Demu\v{s}kin sector]\label{cor:faithful-demushkin-symbol}
Among odd-prime Demu\v{s}kin groups, the finite datum
\begin{equation}\label{eq:faithful-demushkin-symbol}
\left(
\dim_{\mathbf F_p}H^1(P,\mathbf F_p),
\sigma_2(r),
\mathfrak q(P)
\right)
\end{equation}
determines the isomorphism class of $P$.  Equivalently, the symplectic quadratic Fox symbol records the duality pairing, while the integral linear Fox symbol---or the Bockstein lifting filtration---records $q$.  No higher $A_\infty$ operation is required for group-level classification in this standard odd-prime one-relator sector.
\end{corollary}

\begin{proof}
Theorem~\ref{thm:quadratic-fox-cup} and Corollary~\ref{cor:demushkin-symplectic-chart} recover the nondegenerate cup form and the rank.  Theorem~\ref{thm:integral-fox-recovers-q} recovers $q$.  Labute's odd-prime classification states that the rank and $q$ determine the Demu\v{s}kin group up to isomorphism \cite{LabuteClassification}.
\end{proof}

\begin{warning}[Quadratic-only no-go, repaired integrally]\label{warn:quadratic-symbol-q-nogo}
For odd $p$, every standard Demu\v{s}kin relation \eqref{eq:standard-demushkin-relator} has the same degree-two mod-$p$ Fox symbol, independently of $q$.  Hence the binary cup-product pairing alone does not classify the transport group.  This does \emph{not} imply that higher operations are necessary in this sector: Theorem~\ref{thm:integral-fox-recovers-q} and Proposition~\ref{prop:bockstein-lifting-depth} show that the integral linear Fox symbol, equivalently the full Bockstein lifting depth, restores the missing power invariant.  Higher operations remain relevant for representation-level reconstruction, for non-one-relator charts, and for comparisons that discard the integral relation vector.
\end{warning}

Thus the odd-prime Demu\v{s}kin chart is finite and group-faithful: integral degree one records the power relation and mod-$p$ degree two the duality pairing.

% Versioned self-contained section source; integrated verbatim into fixed124.
\subsection{Dyadic Fox--Bockstein--orientation completion}
\label{subsec:dyadic-fox-orientation}

The odd-prime symplectic normal form does not extend verbatim to
$p=2$: a square term already contributes in Magnus degree two and
need not vanish in the cup form.  More importantly, for the
exceptional Demu\v{s}kin invariant $q=2$, the quadratic Fox form and
the ordinary integral exponent-sum/Bockstein information do not
record the canonical dualizing-orientation image.  We now close this
dyadic classification gap without changing the finite Fox
resolution or the all-Sylow descent \cite{LabuteClassification,NSW}.

\begin{theorem}[Four dyadic Demu\v{s}kin forms and their quadratic symbols]
\label{thm:dyadic-relator-cases}
Let $P$ be an infinite finitely generated Demu\v{s}kin pro-$2$
group of rank $d$, with invariant $q=q(P)$, and write
$[x,y]=x^{-1}y^{-1}xy$.  In a suitable minimal presentation
$P=\langle x_1,\ldots,x_d\mid r\rangle_{\mathrm{pro}-2}$,
the relation $r$ is in exactly one of the following normal-form
families:
\begin{enumerate}[label=\textnormal{(\roman*)}]
\item If $q\ne2$, then $d$ is even and
\[
r=x_1^q[x_1,x_2][x_3,x_4]\cdots[x_{d-1},x_d],
\qquad q=0\ \text{or}\ q=2^f,\ f\ge2.
\]
\item If $q=2$, $d$ is odd and $d\ge3$, then
\[
r=x_1^2x_2^{2^f}[x_2,x_3][x_4,x_5]\cdots[x_{d-1},x_d],
\quad f\in\{2,3,\ldots,\infty\}.
\]
\item If $q=2$, $d$ is even and the canonical-orientation
image has square-index $2$, then
\[
r=x_1^{2+2^f}[x_1,x_2][x_3,x_4]\cdots[x_{d-1},x_d],
\quad f\in\{2,3,\ldots,\infty\}.
\]
\item If $q=2$, $d\ge4$ is even and the canonical-orientation
image has square-index $4$, then
\[
r=x_1^2[x_1,x_2]x_3^{2^f}[x_3,x_4]\cdots[x_{d-1},x_d],
\quad f\in\{2,3,\ldots\}.
\]
\end{enumerate}
Here $2^\infty=0$ means that the corresponding factor is absent.
With $C_{ij}=X_iX_j+X_jX_i\in\mathbf F_2\langle X_1,\ldots,X_d\rangle$,
their degree-two Magnus symbols are
\begin{equation}\label{eq:dyadic-quadratic-symbols}
\sigma_2(r)=
\begin{cases}
\displaystyle\sum_{j=1}^{d/2}C_{2j-1,2j}, &(i),\\[2pt]
\displaystyle X_1^2+\sum_{j=1}^{(d-1)/2}C_{2j,2j+1},
 &(ii),\\[2pt]
\displaystyle X_1^2+\sum_{j=1}^{d/2}C_{2j-1,2j},
 &(iii),(iv).
\end{cases}
\end{equation}
The induced Fox cup form is nondegenerate, but in cases
(ii)--(iv) it is \emph{nonalternating}: $\chi_1^2=\omega$.

For $P=G_K(2)$ with $K/\mathbf Q_2$ finite,
one always has $q\ge2$, $d=[K:\mathbf Q_2]+2$,
and an open orientation image, so any occurring $f$ is finite.
\end{theorem}

\begin{proof}
The four presentations and the square-index distinction are
the Demu\v{s}kin--Serre--Labute classification
\cite{LabuteClassification,NSW}.
Over $\mathbf F_2$ one has
$\mathcal M([x_i,x_j])=1+C_{ij}+O(\deg3)$ and
$\mathcal M(x_i^{2^f})=1+X_i^{2^f}$.
For $f\ge2$ the latter does not enter quadratic degree.
The factors $x_1^2$ and $x_1^{2+2^f}$, however,
both contribute $X_1^2$.
This proves \eqref{eq:dyadic-quadratic-symbols}.
The quadratic Fox cup theorem
(Theorem~\ref{thm:quadratic-fox-cup})
identifies its coefficients with the cup matrix.
Its blocks are symplectic $2$-blocks together,
in the exceptional cases, with either the isolated
diagonal $1$ or the nonalternating block
$\bigl(\begin{smallmatrix}1&1\\1&0\end{smallmatrix}\bigr)$;
in all cases the form is nondegenerate.
For the local assertion, $\mu_2\subset K$ implies that
$G_K(2)$ is Demu\v{s}kin, and its dualizing character
is cyclotomic.  Its image in $\mathbb Z_2^\times$ is open
because $K$ is finite.
\end{proof}

Let $\partial_i$ be the completed \emph{left} Fox derivative,
and for a continuous character $\theta:P\to\mathbb Z_2^\times$
write
\[
\varepsilon_\theta:\mathbb Z_2[[F_d]]\longrightarrow\mathbb Z_2,
\quad x_i\longmapsto\theta(x_i)
\]
for the associated twisted augmentation.  The system
\begin{equation}\label{eq:dyadic-twisted-fox-orientation}
\boxed{\quad\varepsilon_\theta(\partial_i r)=0
\quad(1\le i\le d)\quad}
\end{equation}
says that every continuous $\theta$-crossed derivation of
the free pro-$2$ group vanishes on the relator.  It is
presentation-covariant, although its coordinates are not
generator-independent.

\begin{theorem}[Canonical dyadic orientation from twisted Fox equations]
\label{thm:dyadic-twisted-fox-orientation}
On each normal form of Theorem~\ref{thm:dyadic-relator-cases},
Equation~\eqref{eq:dyadic-twisted-fox-orientation} admits
a unique solution in generator values
$\theta(x_i)\in\mathbb Z_2^\times$.
It is the canonical dualizing orientation $\theta_P$.
All unmentioned generator values are $1$, and the nontrivial
values are
\begin{align}
\textnormal{(i)}\quad&
 \theta(x_1)=1,\quad
 \theta(x_2)=(1-q)^{-1};\label{eq:dyadic-orientation-i}\\
\textnormal{(ii)}\quad&
 \theta(x_1)=-1,\quad\theta(x_2)=1,\quad
 \theta(x_3)=(1-2^f)^{-1};\label{eq:dyadic-orientation-ii}\\
\textnormal{(iii)}\quad&
 \theta(x_1)=1,\quad
 \theta(x_2)=(-1-2^f)^{-1};
 \label{eq:dyadic-orientation-iii}\\
\textnormal{(iv)}\quad&
 \theta(x_1)=1,\quad\theta(x_2)=-1,\quad
 \theta(x_3)=1,\quad\theta(x_4)=(1-2^f)^{-1}.
 \label{eq:dyadic-orientation-iv}
\end{align}
The corresponding closed orientation images are
\begin{equation}\label{eq:dyadic-orientation-images}
\begin{aligned}
\textnormal{(i)}\quad&
 1+q\mathbb Z_2\ (q\ge4),\quad\{1\}\ (q=0),\\
\textnormal{(ii)}\quad&
 \{\pm1\}(1+2^f\mathbb Z_2),\\
\textnormal{(iii)}\quad&
 \overline{\langle(-1-2^f)^{-1}\rangle},\\
\textnormal{(iv)}\quad&
 \{\pm1\}(1+2^f\mathbb Z_2).
\end{aligned}
\end{equation}
For $f=\infty$ use $1+2^f\mathbb Z_2=\{1\}$ and
$(-1-2^f)^{-1}=-1$.
\end{theorem}

\begin{proof}
With the commutator convention above,
\[
\varepsilon_\theta(\partial_x[x,y])
=\theta(x)^{-1}(\theta(y)^{-1}-1),
\quad
\varepsilon_\theta(\partial_y[x,y])
=\theta(x)^{-1}\theta(y)^{-1}(\theta(x)-1).
\]
Multiplication by a preceding word only adds a unit factor.
A pure commutator block therefore forces both orientation
values to $1$.  In a block $x^m[x,y]$, the $y$-derivative
forces $\theta(x)=1$, after which the $x$-derivative
becomes $m+\theta(y)^{-1}-1$; hence
$\theta(y)=(1-m)^{-1}$.
For the odd-rank family, the first derivative of $x_1^2$
forces $\theta(x_1)=-1$, then the $(x_2,x_3)$ block
forces $\theta(x_2)=1$ and
$\theta(x_3)=(1-2^f)^{-1}$.
For family (iv), the first squared commutator block gives
$(\theta(x_1),\theta(x_2))=(1,-1)$,
followed by the ordinary $x_3^{2^f}[x_3,x_4]$ block.
This gives the four displayed lists and proves uniqueness.
The Fox identity $r-1=\sum_i\partial_i r\,(x_i-1)$
ensures that the resulting map kills $r$, and
the known canonical orientation in Labute's classification
has precisely these normal-form values
\cite{LabuteClassification,NSW}.
For finite $f$, $(1-2^f)^{-1}$ generates
$1+2^f\mathbb Z_2$ topologically, as does
$(1-q)^{-1}$ when $q\ge4$.
The remaining image claims are immediate.
\end{proof}

\begin{corollary}[Finite classification-faithful dyadic symbol]
\label{thm:dyadic-classification-certificate}
The presentation-covariant certificate
\begin{equation}\label{eq:dyadic-classification-certificate}
\boxed{\quad
\mathfrak C_2(P)=
\bigl(d,\,[\sigma_2(r)],\,
\mathfrak q(P),\,\operatorname{im}\theta_P\bigr)
\quad}
\end{equation}
determines the isomorphism class of every
Demu\v{s}kin pro-$2$ group in
Theorem~\ref{thm:dyadic-relator-cases}.
Here $\mathfrak q(P)$ is the ideal of the augmented
integral first Fox derivatives and $\theta_P$
is determined by the twisted Fox system.
On all $q=2$ normal forms one has
$\mathfrak q(P)=(2)$ and the same diagonal-square
information; the orientation image records the extra
Labute index and parameter $f$.
\end{corollary}

\begin{proof}
The quadratic Fox theorem gives the mod-$2$ cup form,
and the augmented first derivatives give the
abelianization's torsion exponent ideal.
Theorem~\ref{thm:dyadic-twisted-fox-orientation}
recovers the canonical orientation.
Labute's classification makes $(d,\operatorname{im}\theta_P)$
a complete isomorphism invariant in the finite-rank
Demu\v{s}kin sector \cite{LabuteClassification};
the displayed richer certificate is therefore faithful.
The individual matrix coordinates need not be
presentation-invariant.
\end{proof}

\begin{proposition}[The explicit $G_{\mathbf Q_2}(2)$ anchor]
\label{prop:dyadic-q2-explicit-anchor}
Let
\[
P_{\mathbf Q_2}=G_{\mathbf Q_2}(2)
=\langle a,s,y\mid a^2s^4[s,y]\rangle_{\mathrm{pro}-2}.
\]
For the Magnus variables $A,S,Y$, the mod-$2$
Fox--Cartan and integral Fox--Bockstein data are
\begin{equation}\label{eq:dyadic-q2-anchor-data}
\begin{gathered}
\sigma_2(r)=A^2+SY+YS,\qquad
B_{\rm cup}=
 \begin{pmatrix}1&0&0\\0&0&1\\0&1&0\end{pmatrix},\\
\lambda(r)=(2,4,0),\quad \mathfrak q(P_{\mathbf Q_2})=(2),\\
\beta_1(\chi_a)=\omega,\qquad
\beta_1(\chi_s)=\beta_1(\chi_y)=0,\\
(\theta(a),\theta(s),\theta(y))=
(-1,1,-1/3),\quad
\operatorname{im}\theta=\mathbb Z_2^\times.
\end{gathered}
\end{equation}
For trivial integral coefficients, its finite
Fox cochain complex is
\[
\mathbb Z_2\xrightarrow{0}\mathbb Z_2^3
 \xrightarrow{(2,4,0)}\mathbb Z_2.
\]
Its lifting filtration satisfies
$L_1=H^1(P_{\mathbf Q_2},\mathbf F_2)$ and
$L_m=\operatorname{span}_{\mathbf F_2}
\{\chi_s,\chi_y\}$ for every $m\ge2$.
\end{proposition}

\begin{proof}
The displayed relator is the rank-$3$, $q=2$,
$f=2$ normal form for $G_{\mathbf Q_2}(2)$
\cite{LabuteClassification,NSW,RoeTurtureanGQ2}.
Modulo $2$, $a^2$ contributes $A^2$,
$s^4$ has no quadratic term, and $[s,y]$
contributes $SY+YS$.
The quadratic Fox cup theorem gives the matrix.
The exponent sums are $(2,4,0)$.
In the integral trivial-coefficient Fox complex,
a $1$-cochain $(u,v,w)$ is closed modulo $2^m$
exactly when $2u+4v\equiv0\pmod{2^m}$.
For $m=2$, dividing its coboundary by $2$
gives the ordinary Bockstein $(1,0,0)$ mod $2$;
for all $m\ge2$, precisely the reductions with
$u=0$ lift, proving the filtration.
Finally the twisted Fox equations imply
$\theta(a)=-1$, $\theta(s)=1$ and
$4+\theta(y)^{-1}-1=0$, so
$\theta(y)=-1/3$.
Its image together with $-1$ is
$\{\pm1\}(1+4\mathbb Z_2)=\mathbb Z_2^\times$.
\end{proof}

\begin{warning}[The ordinary dyadic Bockstein depth is insufficient]
\label{warn:dyadic-bockstein-depth-insufficient}
For each finite $f\ge2$, the rank-three relations
\[
r_f=x_1^2x_2^{2^f}[x_2,x_3]
\]
have the same $\sigma_2(r_f)$, integral
Fox ideal $(2)$, first Bockstein, and full
untwisted lifting filtration
\[
L_1=H^1(P_f,\mathbf F_2),\qquad
L_m=\operatorname{span}\{\chi_2,\chi_3\}
\quad(m\ge2).
\]
Nevertheless their orientation images
$\{\pm1\}(1+2^f\mathbb Z_2)$ vary with $f$.
Thus the original odd-prime symbol
$(d,\sigma_2,\mathfrak q)$ does not classify
the exceptional dyadic groups.
The missing input is the \emph{twisted}
Fox orientation, not another untwisted
Bockstein depth.
\end{warning}

\begin{remark}[Interface with all-Sylow finite closure]
\label{rem:dyadic-all-sylow-interface}
For any residual sector at $p=2$, the preimage
$H$ of a residual Sylow $2$-subgroup has odd
index in $G_K$ and is the absolute Galois
group of a finite dyadic extension $K_H$.
Because $\mu_2\subset K_H$, the maximal
pro-$2$ quotient $H(2)$ is always
Demu\v{s}kin, not the free branch.
Choosing the normal form determined by its
cyclotomic orientation and evaluating its
completed Fox derivatives gives exactly
the bounded finite-free coefficient chart
of Theorem~\ref{thm:family-sylow-fox}.
The odd index supplies the integral
transfer idempotent from
Theorem~\ref{thm:sylow-fox-retract}.
Hence the additional dyadic classification
data alter only the explicit finite
Fox matrices, not perfectness, derived base
change, Cayley--Hamilton transport, or
the recursive globalization theorem.
The maximal pro-$2$ presentation is not
misidentified with a marked presentation
of the entire absolute Galois group.
\end{remark}

% Dyadic Fox--Wu--Tate extension for the fixed125 integrated manuscript.
\subsection{Mod-four Wu classes and exact twisted Fox--Bockstein depth}
\label{subsec:dyadic-fox-wu-tate}

The canonical orientation of the preceding subsection admits three further
finite interpretations: its reduction modulo $4$ is the characteristic
class of the mod-$2$ cup form; its full integral Fox row is the exact
obstruction to lifting $H^1$ with arbitrary rank-one twists; and its
vanishing row produces an integral top-cell trace for local Tate
duality.  These observations separate the \emph{ordinary} Bockstein
loss of information from the \emph{twisted} Kummer lifting property.

\begin{remark}[Disjointness at the even $q=2$ boundary]
\label{rem:dyadic-even-infinity-boundary}
One may formally write $f=\infty$ in the even-rank
two-generator-image normal form (iv) of
Theorem~\ref{thm:dyadic-relator-cases}.
But then $x_3^{2^\infty}=1$ and the resulting
relator is exactly the $f=\infty$ relation of
the even-rank cyclic-image family (iii),
namely $x_1^2[x_1,x_2]\cdots[x_{d-1},x_d]$.
Its orientation image is $\{\pm1\}$ and
has square-index $2$, not $4$.
The restriction to finite $f$ in family (iv)
is therefore a disjoint normalization and
excludes no Demu\v{s}kin isomorphism type.
\end{remark}

\begin{theorem}[The dyadic Fox--Wu/Bockstein identity]
\label{thm:dyadic-fox-wu-bockstein}
Let $P$ be a finite-rank infinite Demu\v{s}kin pro-$2$ group
and let $\theta=\theta_P:P\to\mathbb Z_2^\times$
be its canonical dualizing orientation.  Define
\begin{equation}\label{eq:dyadic-fox-wu-class}
w_P:P\longrightarrow\mathbf F_2,\qquad
w_P(g)=\frac{\theta(g)-1}{2}\bmod2.
\end{equation}
This is a continuous group homomorphism, hence a class in
$H^1(P,\mathbf F_2)$.  For every
$\chi\in H^1(P,\mathbf F_2)$ one has
\begin{equation}\label{eq:dyadic-fox-wu-identity}
\boxed{\qquad
\beta_1(\chi)=\chi\smile\chi=\chi\smile w_P
\quad\text{in }H^2(P,\mathbf F_2).\qquad}
\end{equation}
In the normal-form basis of
Theorem~\ref{thm:dyadic-relator-cases},
\begin{equation}\label{eq:dyadic-wu-basis-values}
w_P=
\begin{cases}
0,&q\ne2,\\
\chi_1,&q=2,\ d\text{ odd (case (ii))},\\
\chi_2,&q=2,\ d\text{ even (cases (iii),(iv))}.
\end{cases}
\end{equation}
For a finite extension $K/\mathbf Q_2$ and
$P=G_K(2)$, the class $w_P$ equals the
Kummer character $\kappa_{-1}$ of $-1\in K^\times$:
\begin{equation}\label{eq:dyadic-kummer-minus-one}
\chi\smile\chi=\chi\smile\kappa_{-1}.
\end{equation}
In particular, $w_P=0$ precisely when $\mu_4\subset K$
in the local Galois case.
\end{theorem}

\begin{proof}
Every $2$-adic unit is $1$ modulo $2$, so the formula
\eqref{eq:dyadic-fox-wu-class} is defined.  Multiplication
of two elements $1+2a$ and $1+2b$ gives
$1+2(a+b)+4ab$, proving additivity modulo $2$.
The ordinary connecting Bockstein for
$0\to\mathbf F_2\to\mathbf Z/4\to\mathbf F_2\to0$
equals the degree-one cup square:
$\beta_1(\chi)=\chi\smile\chi$.
Theorem~\ref{thm:dyadic-twisted-fox-orientation}
gives the reduction of $\theta$ modulo $4$:
it is trivial in case (i), is $-1$ on $x_1$ in
case (ii), and is $-1$ on $x_2$ in
cases (iii),(iv), with all other generator values
equal to $1$ modulo $4$.
This proves \eqref{eq:dyadic-wu-basis-values}.

The quadratic Fox cup matrices
\eqref{eq:dyadic-quadratic-symbols}
have zero diagonal in case (i), and otherwise
their only nonzero diagonal entry is $B_{11}=1$.
In case (ii), the first generator is orthogonal
to every other basis vector, so
$B(\chi,\chi)=B(\chi,\chi_1)$ for every $\chi$.
In cases (iii),(iv), the first block is
$\bigl(\begin{smallmatrix}1&1\\1&0\end{smallmatrix}\bigr)$,
so $B(\chi,\chi)=B(\chi,\chi_2)$.
This proves \eqref{eq:dyadic-fox-wu-identity}
basis-independently by bilinearity.

For $P=G_K(2)$ the orientation is the restricted
$2$-adic cyclotomic character.  Reducing it modulo
$4$ records the action of $G_K$ on a square root
$i$ of $-1$, namely the quadratic extension
$K(i)/K$.  Its mod-$2$ character is precisely
the Kummer class $\kappa_{-1}$.  Hence
$w_P=\kappa_{-1}$ and
\eqref{eq:dyadic-kummer-minus-one} follows.
The character is zero exactly when $i\in K$.
\end{proof}

For an arbitrary continuous unit-valued character
$\psi:P\to\mathbb Z_2^\times$, retain a fixed
minimal one-relator presentation
$P=\langle x_1,\ldots,x_d\mid r\rangle_{\mathrm{pro}-2}$
and define the \emph{twisted top-cell Fox row}
\begin{equation}\label{eq:dyadic-psi-fox-row}
J_\psi:=
\bigl(\varepsilon_\psi(\partial_1r),\ldots,
\varepsilon_\psi(\partial_dr)\bigr)\in\mathbb Z_2^d,
\qquad
a(\psi):=\min_i v_2((J_\psi)_i),
\end{equation}
where $v_2(0)=\infty$.  When $a(\psi)<\infty$,
write
\begin{equation}\label{eq:dyadic-psi-leading-row}
j_\psi=
2^{-a(\psi)}J_\psi\bmod2
\in(\mathbf F_2)^d\setminus\{0\}.
\end{equation}
The ideal generated by $J_\psi$ is invariant under
Fox presentation change; individual row coordinates
and the hyperplane below transform covariantly.

\begin{theorem}[Exact twisted Fox lifting-depth obstruction]
\label{thm:dyadic-twisted-fox-lifting-depth}
Let $P$ and $\psi$ be as above.  Let
$\mathbb Z_2(\psi)$ be the rank-one integral module
on which $P$ acts through $\psi$ and put
\begin{equation}\label{eq:dyadic-psi-lifting-def}
L_m^\psi(P)=
\operatorname{im}\left(
H^1(P,\mathbb Z_2(\psi)/2^m)
\longrightarrow H^1(P,\mathbf F_2)\right),
\quad m\ge1.
\end{equation}
Then the rank-one twisted Fox cochain complex is
\begin{equation}\label{eq:dyadic-psi-three-term-fox}
\mathbb Z_2
\xrightarrow{\ \bigl(\psi(x_i)-1\bigr)_i\ }
\mathbb Z_2^d
\xrightarrow{\ J_\psi\ }
\mathbb Z_2.
\end{equation}
Consequently
\begin{equation}\label{eq:dyadic-psi-h2}
H^2(P,\mathbb Z_2(\psi))
\cong
\begin{cases}
\mathbb Z_2,&a(\psi)=\infty,\\
\mathbb Z_2/2^{a(\psi)}\mathbb Z_2,
  &a(\psi)<\infty,
\end{cases}
\end{equation}
and the entire lifting filtration is
\begin{equation}\label{eq:dyadic-psi-lifting-formula}
L_m^\psi(P)=
\begin{cases}
H^1(P,\mathbf F_2),
 &m\le a(\psi),\\
\ker\bigl(j_\psi:(\mathbf F_2)^d\to\mathbf F_2\bigr),
 &m>a(\psi),\quad a(\psi)<\infty.
\end{cases}
\end{equation}
Here the identification $H^1(P,\mathbf F_2)
\cong(\mathbf F_2)^d$ uses the minimal generators.
For all $\psi$, either $a(\psi)=\infty$ or
$a(\psi)\ge1$.  Moreover, the following conditions
are equivalent:
\begin{equation}\label{eq:dyadic-psi-orientation-criterion}
\boxed{\quad
\psi=\theta_P
\quad\Longleftrightarrow\quad
J_\psi=0
\quad\Longleftrightarrow\quad
H^2(P,\mathbb Z_2(\psi))\simeq\mathbb Z_2
\quad\Longleftrightarrow\quad
L_m^\psi(P)=H^1(P,\mathbf F_2)\ \forall m.
\quad}
\end{equation}
Thus the canonical orientation is also the unique
integral unit twist with \emph{unrestricted all-level
mod-$2$ degree-one liftability}.
\end{theorem}

\begin{proof}
Apply the completed finite free Fox resolution
\eqref{eq:fox-chain-complex} to the rank-one module
$\mathbb Z_2(\psi)$.  The displayed differentials
follow from the Fox cochain formulas
\eqref{eq:fox-d0}--\eqref{eq:fox-d1}.
Since the top cohomology is the cokernel of
the single row $J_\psi$, its Smith invariant
over the discrete valuation ring $\mathbb Z_2$
gives \eqref{eq:dyadic-psi-h2}.

The presentation is minimal and every unit
$\psi(x_i)$ is $1$ modulo $2$, so both Fox
differentials vanish modulo $2$.  In particular,
$H^1(P,\mathbf F_2)\cong(\mathbf F_2)^d$ and
$J_\psi\equiv0\bmod2$.
If $m\le a(\psi)$, the twisted top differential
vanishes modulo $2^m$, so every mod-$2$ degree-one
cocycle lifts.  If $m>a(\psi)$, reducing the
equation $J_\psi v=0\bmod2^m$ after dividing by
$2^{a(\psi)}$ gives $j_\psi\bar v=0$.
Conversely, if $j_\psi\bar v=0$, choose an index
where $2^{-a(\psi)}(J_\psi)_i$ is a unit.  Keeping
all other entries fixed and adjusting that entry
modulo $2^{m-a(\psi)}$ solves the lifted equation
without changing $\bar v$.  Hence the image is
exactly the stated hyperplane.

Theorem~\ref{thm:dyadic-twisted-fox-orientation}
identifies $\theta_P$ as the \emph{unique} unit
character annihilating the entire twisted Fox row.
This proves the first two equivalences.
Equation~\eqref{eq:dyadic-psi-h2} proves the third,
and \eqref{eq:dyadic-psi-lifting-formula} the fourth:
if $a(\psi)<\infty$, its leading row is nonzero
and full liftability fails first at $m=a(\psi)+1$.
\end{proof}

\begin{corollary}[The orientation repairs every untwisted Bockstein loss]
\label{cor:dyadic-orientation-kummer-lifts}
For every such $P$ and every $m\ge1$,
\[
H^1(P,\mathbb Z_2(\theta_P)/2^m)
\twoheadrightarrow H^1(P,\mathbf F_2).
\]
If the trivial twist has invariant
$\mathfrak q(P)=(2^f)$ with finite $f$,
then its lifting filtration loses exactly
one mod-$2$ direction at level $f+1$,
whereas the canonical twist loses none.
For the local Galois groups $G_K(2)$,
the twisted surjectivity is the finite-Fox
realization of the Kummer lifting property.
\end{corollary}

\begin{proof}
For the canonical orientation the top Fox row
vanishes integrally by
Theorem~\ref{thm:dyadic-twisted-fox-orientation};
apply \eqref{eq:dyadic-psi-lifting-formula}.
For the trivial twist $J_1$ is the integral
exponent-sum row, whose ideal is $(2^f)$;
its leading reduction is nonzero, giving
the claimed codimension-one loss.
\end{proof}

\begin{theorem}[Integral dyadic Kummer--Tate groups from one Fox row]
\label{thm:dyadic-fox-kummer-tate}
Let $K/\mathbf Q_2$ be finite of degree $n$,
$P_K=G_K(2)$, and put
\[
q_K=|\mu_{2^\infty}(K)|=2^{f_K}\ge2,
\qquad d=n+2.
\]
The cyclotomic orientation $\theta_K$ of $P_K$
has both Fox rows
\begin{equation}\label{eq:dyadic-kummer-fox-rows}
J_{\theta_K}=0,\qquad
\bigl(\theta_K(x_i)-1\bigr)_i
\text{ generating the ideal }(q_K)\subset\mathbb Z_2.
\end{equation}
Thus its Fox cochain complex with cyclotomic
coefficients computes
\begin{equation}\label{eq:dyadic-kummer-tate-groups}
\boxed{
\begin{aligned}
H^0(G_K,\mathbb Z_2(1))&=0,\\
H^1(G_K,\mathbb Z_2(1))
  &\cong\mathbb Z_2^{\,n+1}
     \oplus\mathbb Z_2/q_K\mathbb Z_2,\\
H^2(G_K,\mathbb Z_2(1))&\cong\mathbb Z_2.
\end{aligned}}
\end{equation}
The untwisted Fox complex likewise gives
\begin{equation}\label{eq:dyadic-untwisted-tate-groups}
H^0(G_K,\mathbb Z_2)=\mathbb Z_2,\qquad
H^1(G_K,\mathbb Z_2)\cong\mathbb Z_2^{\,n+1},\qquad
H^2(G_K,\mathbb Z_2)\cong
   \mathbb Z_2/q_K\mathbb Z_2.
\end{equation}
In the $K=\mathbf Q_2$ basis of
Proposition~\ref{prop:dyadic-q2-explicit-anchor},
the entire twisted Fox differential is
\begin{equation}\label{eq:dyadic-q2-twisted-complex}
\mathbb Z_2
\xrightarrow{\;(-2,\ 0,\ -4/3)^{\rm t}\;}
\mathbb Z_2^3
\xrightarrow{\;(0,\ 0,\ 0)\;}
\mathbb Z_2.
\end{equation}
In particular,
$H^1(G_{\mathbf Q_2},\mathbb Z_2(1))
\cong\mathbb Z_2^2\oplus\mathbb Z/2$
and the top twisted cohomology is torsion-free,
whereas
$H^2(G_{\mathbf Q_2},\mathbb Z_2)
\cong\mathbb Z/2$.
\end{theorem}

\begin{proof}
The maximal pro-$2$ local Galois quotient
is Demu\v{s}kin, and its canonical orientation
is the cyclotomic character
\cite{LabuteClassification,NSW}.
Theorem~\ref{thm:dyadic-twisted-fox-orientation}
gives $J_{\theta_K}=0$.
By definition of the cyclotomic action,
$\theta_K(G_K)$ acts trivially on all $q_K$-th
roots of unity but not on all $2q_K$-th
roots, since $q_K$ is the exact order
of the $2$-primary roots of unity in $K$.
The ideal generated by the differences
$\theta_K(g)-1$ is therefore exactly $(q_K)$.
As the chosen $x_i$ topologically generate
$P_K$, the same ideal is generated by
$\theta_K(x_i)-1$.
Consequently the twisted Fox complex has
an injective first differential whose image
has one Smith invariant $q_K$, and a zero
second differential.  Its cohomology is
\eqref{eq:dyadic-kummer-tate-groups}.
This agrees with local Kummer theory,
$H^1(G_K,\mathbb Z_2(1))
\simeq\widehat{K^\times}^{\,2}$.

For the trivial twist, the first differential
is zero, and the second differential is
the exponent-sum row with ideal $(q_K)$,
by local class field theory and the
Demu\v{s}kin abelianization invariant.
Its kernel is free of rank $d-1$ and its
cokernel is $\mathbb Z_2/q_K$, proving
\eqref{eq:dyadic-untwisted-tate-groups}.
Finally the known $G_{\mathbf Q_2}(2)$
orientation takes values $(-1,1,-1/3)$
on $(a,s,y)$; subtracting $1$ gives
$(-2,0,-4/3)$, and its twisted Fox row
vanishes.  Formula
\eqref{eq:dyadic-q2-twisted-complex}
and the claimed groups follow by Smith
normal form.
\end{proof}

\begin{corollary}[Exact cyclotomic-power depth at $\mathbf Q_2$]
\label{cor:dyadic-q2-cyclotomic-power-depth}
Retain the marked generators $(a,s,y)$ of
Proposition~\ref{prop:dyadic-q2-explicit-anchor}.
For any odd integer $k\ge1$, let
$\psi_k=\theta_{\mathbf Q_2}^{\,k}$.  Its entire
twisted top Fox row is
\begin{equation}\label{eq:dyadic-q2-cyclotomic-power-row}
J_{\psi_k}=(0,\ 3-3^k,\ 0).
\end{equation}
For $k=1$ this row is zero.  For every odd
$k\ge3$ its first nonvanishing depth is
\[
a(\psi_k)=2+v_2(k-1),
\]
and therefore
\begin{equation}\label{eq:dyadic-q2-cyclotomic-power-lifting}
L_m^{\psi_k}(P_{\mathbf Q_2})=
\begin{cases}
H^1(P_{\mathbf Q_2},\mathbf F_2),
 &m\le 2+v_2(k-1),\\
\operatorname{span}_{\mathbf F_2}
   \{\chi_a,\chi_y\},
 &m>2+v_2(k-1).
\end{cases}
\end{equation}
In particular
\[
H^2(G_{\mathbf Q_2},\mathbb Z_2(k))
\cong \mathbb Z_2/2^{\,2+v_2(k-1)}\mathbb Z_2
\qquad(k\ge3\text{ odd}),
\]
while $H^2(G_{\mathbf Q_2},\mathbb Z_2(1))
\cong\mathbb Z_2$.
The exponent $2+v_2(k-1)$ is an exact
orientation-sensitive $2$-adic lifting depth,
not a new invariant of the untwisted group.
\end{corollary}

\begin{proof}
For odd $k$ one has
$\psi_k(a)=-1$, $\psi_k(s)=1$,
and $\psi_k(y)=-3^{-k}$.
The first Fox derivative of $a^2$
evaluates to $1+\psi_k(a)=0$.
The $y$-derivative of $s^4[s,y]$
vanishes because $\psi_k(s)=1$.
Finally the $s$-derivative evaluates to
$4+\psi_k(y)^{-1}-1=3-3^k$.
This proves \eqref{eq:dyadic-q2-cyclotomic-power-row}.
For even $n=k-1\ge2$, the $2$-adic
lifting-the-exponent formula gives
$v_2(3^n-1)=2+v_2(n)$.
The leading Fox row therefore reduces to
the $\chi_s$ coordinate, and
Theorem~\ref{thm:dyadic-twisted-fox-lifting-depth}
gives the precise filtration and top
cohomology claimed.
\end{proof}
\begin{proposition}[Transfer-normalized integral Fox--Tate pairing]
\label{prop:dyadic-sylow-fox-tate-pairing}
Let $E/\mathbf Q_2$ be finite, let
$T$ be a finite free $\mathcal O_E$-lattice
with continuous $G_K$-action, and put
$T^\vee(1)=\operatorname{Hom}_{\mathcal O_E}
(T,\mathcal O_E)(1)$.
Choose a residual Sylow-$2$ preimage
$H\subseteq G_K$ of odd index $m$, and let
$K_H$ be its fixed field.
The orientation-twisted top Fox complex of
$P_H=H(2)$ has zero final differential and
a top-degree trace
\[
\operatorname{tr}_{H}^{\rm Fox}:
C^2_{\rm Fox}(P_H,\mathcal O_E(1))
\longrightarrow\mathcal O_E
\]
normalized to the local invariant
$H^2(G_{K_H},\mathcal O_E(1))
\simeq\mathcal O_E$.
With any Fox diagonal approximation and the
coefficient evaluation
$T\otimes T^\vee(1)\to\mathcal O_E(1)$,
the formula
\begin{equation}\label{eq:dyadic-transfer-tate-pairing}
\boxed{\quad
\langle a,b\rangle_K^{\rm Fox}
=
\frac1m\operatorname{tr}_{H}^{\rm Fox}
\left(
\operatorname{res}_{H}(a)
  \smile_{\rm Fox}
\operatorname{res}_{H}(b)
\right)\quad}
\end{equation}
represents the intrinsic integral local
Tate pairing of perfect complexes
\[
R\Gamma(G_K,T)\otimes_{\mathcal O_E}^{\mathbf L}
R\Gamma(G_K,T^\vee(1))
\longrightarrow\mathcal O_E[-2].
\]
In particular it is perfect in the derived
sense, independent of the Sylow subgroup and
presentation, and requires no integral
termwise $\Delta_K$-invariants.
Its finite Fox matrix representative depends
on the chosen diagonal approximation and
normalization; the derived pairing does not.
\end{proposition}

\begin{proof}
Theorem~\ref{thm:dyadic-fox-kummer-tate},
applied to $K_H$ and extended to $\mathcal O_E$,
gives $H^2(P_H,\mathcal O_E(1))\cong\mathcal O_E$
and shows that the twisted top Fox differential
is zero.  Choose its generator so that the
resulting trace coincides with the usual
local invariant under the Fox--Galois comparison.
The finite Fox diagonal and tensor Fox
formula of
Theorem~\ref{thm:twisted-tensor-fox-cup}
represent the intrinsic cup product.
For a finite extension $K_H/K$ of degree
$m$, compatibility of local invariants
with restriction gives
\[
\operatorname{inv}_{K_H}(\operatorname{res}_{H}\eta)
=m\,\operatorname{inv}_{K}(\eta)
\quad
(\eta\in H^2(G_K,\mathcal O_E(1))).
\]
Since $m$ is odd, division by $m$ is integral.
Hence \eqref{eq:dyadic-transfer-tate-pairing}
is the original cup/evaluation/invariant
pairing on $G_K$.
Integral local Tate duality identifies its
adjoint with a derived equivalence
\[
R\Gamma(G_K,T^\vee(1))
\simeq
R\operatorname{Hom}_{\mathcal O_E}
(R\Gamma(G_K,T),\mathcal O_E)[-2]
\]
\cite{NSW}.
Different Sylow preimages and Fox diagonals
give representatives of this same intrinsic
pairing.  In particular, no assertion that
normalized corestriction is itself a
multiplicative cochain map is used.
\end{proof}

% Dyadic Hilbert-marked and pro-2 idempotent-power effectivity extension.
% Standalone section; integrated into fixed126.

\subsection{Hilbert-marked Fox coordinates and residual profinite-power effectivity}
\label{subsec:dyadic-hilbert-marked-power}

The dyadic Fox--Wu--Tate completion has two complementary native
interfaces.  The first identifies the mod-$2$ Fox quadratic form in
marked $G_{\mathbf Q_2}(2)$ coordinates with the classical explicit
quadratic Hilbert pairing.  The second evaluates the profinite
idempotent exponents appearing in the marked full-$G_{\mathbf Q_2}$
presentation by convergent, residue-sector-specific \emph{integral}
matrix power series.  Only finite jets are claimed to be polynomial;
a global polynomial formula for the whole profinite-power operation
is neither assumed nor needed.

\begin{theorem}[Hilbert-marked Fox comparison at $\mathbf Q_2$]
\label{thm:dyadic-fox-hilbert-marked}
Choose the Roe--Turturean marked Demu\v{s}kin normalization
\cite{RoeTurtureanGQ2}
\[
P=G_{\mathbf Q_2}(2)
 =\langle a,s,y\mid a^2s^4[s,y]\rangle_{\mathrm{pro}-2},
\]
in which the unramified character
$\nu_{\rm ur}:P\twoheadrightarrow\mathbb Z_2$
takes values
\begin{equation}\label{eq:q2-unramified-marking}
 \nu_{\rm ur}(a)=-2,\qquad
 \nu_{\rm ur}(s)=1,\qquad
 \nu_{\rm ur}(y)=0.
\end{equation}
Let $\chi_a,\chi_s,\chi_y$ be the mod-$2$ characters
dual to these generators, and let
$\kappa_u\in H^1(G_{\mathbf Q_2},\mathbf F_2)
=H^1(P,\mathbf F_2)$ denote the Kummer class
of $u\in\mathbf Q_2^\times$.
Then
\begin{equation}\label{eq:q2-kummer-fox-markings}
 \boxed{\quad
 \kappa_{-1}=\chi_a,\qquad
 \kappa_{5}=\chi_s,\qquad
 \kappa_2=\chi_y.
 \quad}
\end{equation}
Thus the marked Fox basis and the Kummer basis
$(-1,5,2)$ have the \emph{same} nonalternating cup matrix
\begin{equation}\label{eq:q2-hilbert-fox-matrix}
 B_{\mathbf Q_2}=
 \begin{pmatrix}
 1&0&0\\
 0&0&1\\
 0&1&0
 \end{pmatrix},
\end{equation}
where entry $1$ means Hilbert symbol $-1$.
In particular, every quadratic Hilbert-symbol pairing on
$\mathbf Q_2^\times/\mathbf Q_2^{\times2}$ is reconstructed
from the completed Fox quadratic symbol and the two canonical
markings $\theta_P\bmod4$ and $\nu_{\rm ur}\bmod2$.
The equality $\chi_y=\kappa_2$ uses the \emph{full} Roe--Turturean
local-reciprocity marking of the generators, not just the
abstract quadratic cup matrix; see
Theorem~\ref{thm:dyadic-exact-hilbert-reciprocity-marking}.
\end{theorem}

\begin{proof}
The cyclotomic orientation values in
Proposition~\ref{prop:dyadic-q2-explicit-anchor}
are $(-1,1,-1/3)$ on $(a,s,y)$.
Since $-1/3\equiv1\bmod4$, the Fox--Wu identity
(Theorem~\ref{thm:dyadic-fox-wu-bockstein})
gives
\[
 \kappa_{-1}=w_P=(\theta_P-1)/2\bmod2=\chi_a.
\]
The extension $\mathbf Q_2(\sqrt5)/\mathbf Q_2$
is the unramified quadratic extension: for example,
$t^2-t-1$ is irreducible modulo $2$ and has
discriminant $5$.  Its Kummer character is
therefore the unique nontrivial unramified quadratic
character, namely $\nu_{\rm ur}\bmod2=\chi_s$
by \eqref{eq:q2-unramified-marking}.

Write $b=2^\beta v$ and $c=2^\gamma w$
with $v,w$ odd.  The explicit dyadic
Hilbert-symbol formula is
\begin{equation}\label{eq:q2-hilbert-explicit}
 (b,c)_2=(-1)^{
 \varepsilon(v)\varepsilon(w)
 +\beta\lambda(w)+\gamma\lambda(v)},
 \quad
 \varepsilon(v)=\frac{v-1}{2}\bmod2,\qquad
 \lambda(v)=\frac{v^2-1}{8}\bmod2,
\end{equation}
as in the standard local reciprocity computation
\cite{RoeTurtureanGQ2,NSW}.
For $v=-1,5,1$ respectively, the pairs
$(\varepsilon(v),\lambda(v))$ are
$(1,0),(0,1),(0,0)$.
Consequently the only nontrivial Hilbert
pairings among $(-1,5,2)$ are
\[
 (-1,-1)_2=-1,\qquad
 (5,2)_2=(2,5)_2=-1,
\]
giving \eqref{eq:q2-hilbert-fox-matrix}.
This equals the independent Fox cup matrix
of Proposition~\ref{prop:dyadic-q2-explicit-anchor}.

Write $\kappa_2=u\chi_a+v\chi_s+w\chi_y$.
Its pairing with $\kappa_{-1}=\chi_a$
vanishes, hence $u=0$ by the Fox cup matrix.
Its pairing with $\kappa_5=\chi_s$ is nonzero,
hence $w=1$.
Thus $\kappa_2=\chi_y+v\chi_s$ using only the cup
matrix and the coarse unramified mark.
Roe--Turturean's full reciprocity marking identifies
$\bar s=\operatorname{rec}(1/2)$, and the dyadic Hilbert
formula gives $(2,1/2)_2=(2,2)_2=+1$.
Therefore $v=\kappa_2(s)=0$, proving the exact
identity in \eqref{eq:q2-kummer-fox-markings}.
\end{proof}

\begin{remark}[Coarse versus full marking of the dyadic Hilbert basis]
\label{rem:q2-hilbert-shear-boundary}
The abstract Fox quadratic form and the coarse pair
$(\nu_{\mathrm{ur}}\bmod2,\theta\bmod4)$
admit a formal shear
$\chi_y\mapsto\chi_y+\epsilon\chi_s$.
This is a \emph{coordinate} ambiguity if the
local-reciprocity images of the generators are
not supplied; it is not intrinsic to the
fully marked Roe--Turturean chart.
Their Lemma~3.5 fixes
$\bar s=\operatorname{rec}(1/2)$ and
$\bar y=\operatorname{rec}(-3)$,
and thus fixes the shear to $\epsilon=0$.
The complete evaluation table and full-marked
Nielsen coordinates are
Theorem~\ref{thm:dyadic-exact-hilbert-reciprocity-marking}
and Corollary~\ref{cor:dyadic-full-marked-kummer-coordinates}.
\end{remark}

\begin{theorem}[Exact integral rank-one dyadic residual deformation ring]
\label{thm:dyadic-q2-rankone-integral-ring}
Let $\bar\rho=\mathbf1$ be the trivial
rank-one residual representation
of $G_{\mathbf Q_2}$ over $k_E$.
The framed universal deformation ring
over $\mathcal O_E$ is
\begin{equation}\label{eq:dyadic-q2-rankone-integral-ring}
\boxed{\quad
R_{\mathbf1}^{\square}
\simeq
\mathcal O_E[[u,v,w]]/
\bigl(u(u+2)\bigr).
\quad}
\end{equation}
Under the canonical normal form of
Proposition~\ref{prop:dyadic-q2-explicit-anchor},
the universal character is explicitly
\begin{equation}\label{eq:dyadic-q2-rankone-transport}
\rho(a)=(1+u)(1+v)^{-2},\qquad
\rho(s)=1+v,\qquad
\rho(y)=1+w.
\end{equation}
The displayed ring is a hypersurface
complete intersection and is finite
free of rank two over
$\mathcal O_E[[v,w]]$, with basis
$\{1,u\}$.  Its closed residual
fiber is
\begin{equation}\label{eq:dyadic-q2-rankone-special-fiber}
R_{\mathbf1}^{\square}/\varpi_E
\simeq
k_E[[u,v,w]]/(u^2),
\end{equation}
which is nonreduced.  After
inverting $2$ in the integral
ring, the two branches $u=0$
and $u=-2$ split by the explicit
idempotent $e=-u/2$.

The quadratic obstruction $u^2$
on the residual fiber is exactly
the nonzero Fox cup square
$\chi_a\smile\chi_a=\omega$:
a first-order residual deformation
with $u=\varepsilon$, $v=w=0$
over $k_E[\varepsilon]/(\varepsilon^2)$
does not lift with that first-order
part to $k_E[\varepsilon]/(\varepsilon^3)$.
Thus dyadic $H^2$ has a literal
integral hypersurface manifestation,
even though its characteristic-zero
rank-one adjoint obstruction vanishes.
\end{theorem}

\begin{proof}
Every rank-one lift of the trivial
residual representation has image
in $1+\mathfrak m_A$, a pro-$2$
group, and hence factors through
the maximal pro-$2$ quotient
$P_{\mathbf Q_2}$.
Abelianizing its defining relation
$a^2s^4[s,y]=1$ gives
$\rho(a)^2\rho(s)^4=1$.
Put
\[
1+u=\rho(a)\rho(s)^2,\qquad
1+v=\rho(s),\qquad
1+w=\rho(y).
\]
This change of complete local
variables has inverse
\eqref{eq:dyadic-q2-rankone-transport}.
The defining relation becomes
$(1+u)^2=1$, equivalently
$u(u+2)=0$, proving the exact
representing-ring identification.

Because $u^2+2u$ is a monic
distinguished polynomial of degree
two in $u$, Weierstrass division
gives the claimed free basis over
$\mathcal O_E[[v,w]]$.
The hypersurface equation is a
nonzero divisor in its regular
power-series ambient ring, so it
is a complete intersection.
Reduction modulo $\varpi_E$ gives
$u^2=0$, and in the localization
$R_{\mathbf1}^{\square}[1/2]$
the element $-u/2$ is idempotent
since $u^2=-2u$.
It separates $u=0$ and $u=-2$.

Finally, in characteristic two
$u=\varepsilon$ solves $u^2=0$
modulo $\varepsilon^2$, but a
lift of the same first-order term
has square $\varepsilon^2\ne0$
modulo $\varepsilon^3$, independent
of its second-order correction.
The Fox quadratic symbol
$A^2+SY+YS$ already identifies
the coefficient of this first
obstruction with
$\chi_a\smile\chi_a=\omega$,
as asserted.
\end{proof}

\begin{theorem}[Residue-sector formula for the profinite $2$-power projection]
\label{thm:dyadic-omega2-power-formula}
Let $(A,\mathfrak m,k)$ be a complete Noetherian
local $\mathcal O_E$-algebra with finite residue
field $k$ of characteristic $2$, and let
$X\in\mathrm{GL}_r(A)$ reduce to
$\bar X\in\mathrm{GL}_r(k)$ of order
\[
 N=2^a m,\qquad m\text{ odd}.
\]
Let
$\omega_2\in\widehat{\mathbf Z}
=\prod_\ell\mathbf Z_\ell$ be the idempotent
whose $2$-adic component is $1$ and whose
other components are $0$.
Choose any nonnegative integer $e<N$ with
\begin{equation}\label{eq:dyadic-omega2-CRT}
 e\equiv1\pmod{2^a},\qquad
 e\equiv0\pmod m,
\end{equation}
where the first congruence is vacuous when $a=0$,
and set
\[
 c=\frac{1-e}{N}\in\mathbb Z_2,
 \qquad S=X^N-I_r\in M_r(\mathfrak m).
\]
Then the intrinsic profinite power is the
following explicit convergent integral series:
\begin{equation}\label{eq:dyadic-omega2-binomial-power}
 \boxed{\quad
 X^{\omega_2}
 =X^e\sum_{j\ge0}
   \binom cj (X^N-I_r)^j.
 \quad}
\end{equation}
For every $t\ge1$ it has the finite-stage formula
\begin{equation}\label{eq:dyadic-omega2-mod-t}
 X^{\omega_2}\equiv
 X^e\sum_{j=0}^{t-1}\binom cj S^j
 \pmod{\mathfrak m^t}.
\end{equation}
The binomial coefficients lie in $\mathbb Z_2$
and require no inversion of $2$.
The result is independent of the chosen
CRT representative $e$, commutes with
conjugation, and is natural under continuous
local coefficient base change preserving the
residual chart.

For a variation $H\in M_r(A)$, put
\begin{equation}\label{eq:dyadic-omega2-power-derivative}
 H_N=\sum_{\ell=0}^{N-1}
  X^\ell H X^{N-1-\ell},\qquad
 H_e=\sum_{\ell=0}^{e-1}
  X^\ell H X^{e-1-\ell}.
\end{equation}
Then the full matrix Jacobian of the profinite
power is represented by the convergent series
\begin{equation}\label{eq:dyadic-omega2-jacobian}
\begin{aligned}
 D(X^{\omega_2})[H]
 ={}&
 H_e\sum_{j\ge0}\binom cj S^j\\
 &+
 X^e\sum_{j\ge1}\binom cj
 \sum_{\ell=0}^{j-1}
  S^\ell H_N S^{j-1-\ell}.
\end{aligned}
\end{equation}
Modulo $\mathfrak m^t$, the first sum
stops at $j=t-1$ and the second at $j=t$.
Thus both the power operation and its
noncommutative first derivative have
terminating matrix formulas at every
specified finite residual precision.
\end{theorem}

\begin{proof}
The kernel of
$\mathrm{GL}_r(A)\to\mathrm{GL}_r(k)$
is pro-$2$ in the $\mathfrak m$-adic
topology.  Hence the closure of the
cyclic group generated by $X$ is a
procyclic profinite group, whose odd-order
factor has order $m$, while its remaining
factor is pro-$2$.  Write its canonical
commuting decomposition as
$X=X_{(2)}X_{\rm odd}$.
Then $X^N=X_{(2)}^N$ belongs to
$1+M_r(\mathfrak m)$.
The congruences
\eqref{eq:dyadic-omega2-CRT} give
$X_{\rm odd}^e=1$ and
$e+Nc=1$ in $\mathbf Z_2$.
Consequently
\[
 X^e(X^N)^c
 =X_{(2)}^{e+Nc}X_{\rm odd}^e
 =X_{(2)}=X^{\omega_2}.
\]
The usual binomial expression
\[
 (I_r+S)^c=\sum_{j\ge0}\binom cj S^j
\]
agrees with continuous $\mathbf Z_2$
exponentiation: it follows by
approximating $c$ with integers and
using $\mathfrak m$-adic continuity.
Since $S\in M_r(\mathfrak m)$,
the $j$th summand lies in
$M_r(\mathfrak m^j)$ and the
series converges, giving
\eqref{eq:dyadic-omega2-binomial-power}
and \eqref{eq:dyadic-omega2-mod-t}.
The equality with intrinsic
$X^{\omega_2}$ proves independence,
conjugation covariance, and naturality.

For the derivative, the standard
noncommutative power rule gives
$D(X^N)[H]=H_N$ and
$D(X^e)[H]=H_e$.
Differentiate the convergent binomial
series termwise, using
\[
 D(S^j)[H]=
 \sum_{\ell=0}^{j-1}
 S^\ell H_N S^{j-1-\ell},
\]
and apply the Leibniz rule for its
leading factor $X^e$.  This is
\eqref{eq:dyadic-omega2-jacobian}.
The derivative of the $j$th binomial
term has at least $j-1$ factors of
$S\in\mathfrak m$, which gives
the asserted exact truncation
order modulo $\mathfrak m^t$.
\end{proof}

\begin{example}[Removal of an odd residual eigenvalue]
\label{ex:dyadic-omega2-odd-teich}
Let $E/\mathbf Q_2$ be unramified quadratic,
so $\mathcal O_E$ contains a primitive
third root of unity $\zeta_3$.
For $A=\mathcal O_E[[z]]$ take the rank-one
transport $X=\zeta_3(1+z)$.
Its residual order is $N=3$.
Choose $e=0$ and $c=1/3\in\mathbb Z_2$.
The formula gives
\[
 X^{\omega_2}
 =\left(X^3\right)^{1/3}
 =\left((1+z)^3\right)^{1/3}
 =1+z.
\]
Thus the explicit integral binomial
calculus really removes the odd
Teichm\"uller component while retaining
the complete pro-$2$ deformation factor.
\end{example}

\begin{theorem}[Finite-jet full-$G_{\mathbf Q_2}$ marked word solver]
\label{thm:dyadic-q2-marked-word-jet-solver}
Assume the Roe--Turturean marked
profinite presentation of $G_{\mathbf Q_2}$
\cite{RoeTurtureanGQ2}, with generators
$\sigma,\tau,x_0,x_1$, tame relation
$\tau^\sigma=\tau^2$, one additional
marked relation $r_{\rm wild}=1$,
and the requirement that the normal
closure of $x_0,x_1$ be pro-$2$.
In the displayed marked word,
the idempotent powers are
\[
 \sigma_2=\sigma^{\omega_2},
 \qquad
 u_i=(x_i\tau)^{\omega_2}\quad(i=0,1),
\]
and all remaining operations are
finite products, inverses, conjugations,
and ordinary integral powers of
these marked words.

Fix a genuine residual representation
$\bar\rho:G_{\mathbf Q_2}\to\mathrm{GL}_r(k_E)$,
and choose integral lifts of its four
residual generator matrices.
Let
\[
 A_{\rm word}
 =\mathcal O_E[[z_{\sigma,ij},z_{\tau,ij},
                z_{0,ij},z_{1,ij}]]_{1\le i,j\le r}
\]
be the \emph{ambient} formally smooth
completed matrix-coordinate ring on
these four lifted generators.
Each occurrence of $\omega_2$ on
$A_{\rm word}$ has the
convergent expression
\eqref{eq:dyadic-omega2-binomial-power}.
For every requested integer $t\ge1$:
\begin{enumerate}[label=\textnormal{(\roman*)}]
\item the full tame and marked wild
matrix relation maps are evaluated
modulo $\mathfrak m^t$ by a finite
sequence of matrix products and
the truncated binomial formula
\eqref{eq:dyadic-omega2-mod-t};
\item the entire first matrix Jacobian
of both marked relations is computed
modulo $\mathfrak m^t$ by the
finite derivative formula
\eqref{eq:dyadic-omega2-jacobian},
product rules, and inverse-matrix
differentiation;
\item the pro-$2$ marked-normal-closure
condition imposes no additional
equation on any Artin lift with
this genuine residual image;
\item the quotient of the ambient
$A_{\rm word}$ by the scalar entries
of the two full convergent relation
matrices is the completed framed
deformation ring $R_{\bar\rho}^{\square}$;
\item at the first-order residual
thickening the kernel of the
\emph{ambient} marked matrix
Jacobian is canonically the framed
cocycle space $Z^1(G_{\mathbf Q_2},
\operatorname{ad}\bar\rho)$.
Quotienting by infinitesimal
conjugation gives
$H^1(G_{\mathbf Q_2},
\operatorname{ad}\bar\rho)$.
\end{enumerate}
This is a \emph{terminating
formal-native algorithm} at every
prescribed finite order.  It does
not identify the naive two-relation
Jacobian cokernel with $H^2$ or
the naive two-relator Koszul complex
with an integral group-resolution.
The independent lci minimal-sequence
and intrinsic derived-cochain
constructions of
Theorem~\ref{thm:word-lci-reduction}
and Theorem~\ref{thm:family-sylow-fox}
remain necessary for those statements.
\end{theorem}

\begin{proof}
For each of the residual matrices of
$\sigma,x_0\tau,x_1\tau$ its order in
$\mathrm{GL}_r(k_E)$ is finite.
Apply Theorem~\ref{thm:dyadic-omega2-power-formula}
to these three matrices with their
own CRT exponents.  Every occurrence
of a profinite idempotent power in
the Roe--Turturean relation becomes
an integral convergent analytic
matrix expression on the completed
residual chart.
Since the remaining word operations
are finite, truncating each
constituent to the requested
$\mathfrak m$-adic order evaluates
the relation matrices by finitely
many operations.  The explicit
Jacobian in
\eqref{eq:dyadic-omega2-jacobian}
and the ordinary matrix product,
inverse, and conjugation rules
prove (i)--(ii).

The marked wild normal subgroup of
$G_{\mathbf Q_2}$ is pro-$2$,
so its image under the genuine
residual representation is a
finite $2$-group.  On every Artin
local lift, the reduction kernel
is a successive extension of
additive groups in characteristic
$2$, hence a $2$-group.
The normal closure of the lifted
wild image is therefore an
extension of finite $2$-groups,
so is again a $2$-group.
This is precisely the residual
automaticity theorem
(Theorem~\ref{thm:marked-residual-strictification}),
giving (iii).  Because the Roe--Turturean
presentation is universal among the
marked finite quotients, the same
residual automaticity identifies the
functor of solutions to the completed
two-relation equations in $A_{\rm word}$
with the framed deformation functor.
This proves (iv).

On the dual-number thickening
$k_E[\varepsilon]/(\varepsilon^2)$,
every lift is
$\rho_\varepsilon(g)
=(1+\varepsilon c(g))\bar\rho(g)$.
The defining representation
equation is the continuous
crossed-cocycle equation
$c(gh)=c(g)+g\cdot c(h)$.
By the marked presentation and
(iii), it is equivalent to the
linearization of the two
displayed marked relations,
computed by (ii).  This identifies
the tangent kernel with $Z^1$.
Infinitesimal conjugation is
exactly the coboundary action,
giving (v).
The last assertion follows from
the distinction between formal
relation equations and a complete
projective group resolution
already made in
Warning~\ref{warn:q2-marked-not-resolution}.
\end{proof}

\begin{warning}[Finite residual jets are not a globally polynomial presentation]
\label{warn:dyadic-marked-word-not-algebraic-polynomial}
The matrices in
\eqref{eq:dyadic-omega2-binomial-power}
are generally infinite convergent
power series on one chosen completed
residual chart.  Their reductions
at each finite order are literal
finite matrix expressions, but
a single finite polynomial in
unrestricted generic matrices
does not follow.
The exact native output of
Theorem~\ref{thm:dyadic-q2-marked-word-jet-solver}
is a recursively effective
formal marked-word atlas, with
strict integral tangent readout.
Global finite-type algebraization,
a uniform marked presentation for
all dyadic $K$, and a preferred
raw Herr-to-marked-word chain
quasi-isomorphism remain stronger
separate questions.  None is
needed by the established
all-Sylow--Cayley--Hamilton
derived faithful atlas.
\end{warning}

% fixed127 substantive insertion: reciprocity-marked Hilbert--Fox dictionary
% and the strict full-group marked Jacobian / pro-2 Fox chain comparison.
% This section is independent of all prior draft-file names and is included
% verbatim in the fixed127 master manuscript.

\subsection{Exact reciprocity marking and a strict dyadic marked-to-Fox bridge}
\label{subsec:dyadic-reciprocity-nielsen-bridge}

The abstract quadratic Fox form does not by itself fix a Kummer basis,
but the full reciprocity marking does.  On residual sectors with pro-$2$
image the marked tame direction also becomes contractible, leaving a
literal integral equivalence between the completed full-group
two-relation Jacobian and the genuine one-relator Fox resolution.
The resulting bridge is stronger than a comparison through the
intrinsic derived category, but is specific to the pro-$2$ residual
sectors of $G_{\mathbf Q_2}$.

\begin{theorem}[Exact reciprocity completion of the dyadic Hilbert--Fox basis]
\label{thm:dyadic-exact-hilbert-reciprocity-marking}
Fix the arithmetic local reciprocity map
$\operatorname{rec}:\mathbf Q_2^\times\to
G_{\mathbf Q_2}^{\mathrm{ab}}$,
so that $\operatorname{rec}(2)$ is arithmetic
Frobenius.  In the Roe--Turturean fully marked
Demu\v{s}kin identification
\cite{RoeTurtureanGQ2}, the pro-$2$
abelianizations of the standard generators satisfy
\begin{equation}\label{eq:dyadic-reciprocity-generator-images}
\boxed{\quad
\bar a=\operatorname{rec}(-4),\qquad
\bar s=\operatorname{rec}(1/2),\qquad
\bar y=\operatorname{rec}(-3).
\quad}
\end{equation}
Consequently the Hilbert--Fox relation
\eqref{eq:q2-kummer-fox-markings}
has \emph{zero} residual shear:
\begin{equation}\label{eq:dyadic-exact-kummer-fox}
\boxed{\quad
\kappa_{-1}=\chi_a,\qquad
\kappa_5=\chi_s,\qquad
\kappa_2=\chi_y.
\quad}
\end{equation}
More explicitly, for the ordered generators $(a,s,y)$
and ordered square classes $(-1,5,2)$, the matrix
of quadratic-character evaluations is exactly $I_3$:
\begin{equation}\label{eq:dyadic-exact-hilbert-evaluation}
\bigl(\kappa_u(g)\bigr)_{
u\in(-1,5,2),\ g\in(a,s,y)}
=
\begin{pmatrix}1&0&0\\0&1&0\\0&0&1\end{pmatrix},
\end{equation}
where $1\in\mathbf F_2$ means Hilbert symbol $-1$.
Writing a square class as
$u=(-1)^b5^c2^h$ with $b,c,h\in\mathbf F_2$,
the Hilbert pairing between
$u=(-1)^b5^c2^h$ and
$v=(-1)^{b'}5^{c'}2^{h'}$ becomes
\begin{equation}\label{eq:dyadic-square-class-hilbert-polar}
\boxed{\quad
(u,v)_2=(-1)^{\,bb'+ch'+hc'}.
\quad}
\end{equation}
This is exactly the bilinear pairing supplied by the
Fox degree-two symbol $A^2+SY+YS$.
\end{theorem}

\begin{proof}
The reciprocity identities
\eqref{eq:dyadic-reciprocity-generator-images}
are the marked abelianization of
Roe--Turturean, Lemma~3.5
\cite{RoeTurtureanGQ2}, and supply precisely
the information not read from the abstract cup matrix.
For a Kummer character, the evaluation on
$\operatorname{rec}(v)$ is the sign of the
quadratic norm-residue Hilbert symbol $(u,v)_2$;
using the opposite reciprocity normalization
would invert $\operatorname{rec}(v)$, but has
no effect on a quadratic character.
Apply Equation~\eqref{eq:q2-hilbert-explicit}
to the pairs
\[
 (u,v),\qquad
 u\in\{-1,5,2\},\quad
 v\in\{-4,1/2,-3\}.
\]
For $u=-1$, the symbols are $(-1,+1,+1)$;
for $u=5$, they are $(+1,-1,+1)$;
and for $u=2$, they are $(+1,+1,-1)$.
In particular,
$(2,1/2)_2=(2,2)_2=+1$, which fixes the
formerly undetermined value
$\kappa_2(s)=0$.
This gives \eqref{eq:dyadic-exact-hilbert-evaluation}
and \eqref{eq:dyadic-exact-kummer-fox}.
Finally the Hilbert pairing in the square-class
basis $(-1,5,2)$ has matrix
\eqref{eq:q2-hilbert-fox-matrix}.
Multiplying by the two coordinate column vectors
gives exactly
\eqref{eq:dyadic-square-class-hilbert-polar}.
\end{proof}

\begin{corollary}[Exact full-marked pro-$2$ Kummer coordinates]
\label{cor:dyadic-full-marked-kummer-coordinates}
Under the Nielsen transform
\begin{equation}\label{eq:dyadic-nielsen-full-generator-map}
\sigma=s,\qquad x_1=y,\qquad
x_0=s^{-2}a^{-1},
\qquad
a=x_0^{-1}\sigma^{-2},
\end{equation}
the pro-$2$ abelianized generator markings are
\begin{equation}\label{eq:dyadic-full-pro2-reciprocity}
\bar x_0=\operatorname{rec}(-1),\qquad
\bar\sigma=\operatorname{rec}(1/2),\qquad
\bar x_1=\operatorname{rec}(-3).
\end{equation}
Hence the three quadratic Kummer classes
are the dual basis of $(x_0,\sigma,x_1)$
in the order $(-1,5,2)$:
\begin{equation}\label{eq:dyadic-full-generator-evaluation}
\bigl(\kappa_u(g)\bigr)_{u=(-1,5,2),\,g=(x_0,\sigma,x_1)}
=I_3.
\end{equation}
In this ordering the mod-$2$ quadratic Fox
symbol of the full-marked pro-$2$ relation
is
\[
X_0^2+X_1\Sigma+\Sigma X_1.
\]
The $\nu_{\rm ur}$ and $\theta$ markings,
together with the fixed reciprocity classes,
therefore remove the last Hilbert--Fox
basis ambiguity without changing the
underlying cup form.
\end{corollary}

\begin{proof}
On abelianizations the identity
$x_0=s^{-2}a^{-1}$ gives
\[
 \bar x_0=
 \operatorname{rec}\bigl((1/2)^{-2}(-4)^{-1}\bigr)
 =\operatorname{rec}(-1).
\]
The two other identities follow immediately from
\eqref{eq:dyadic-reciprocity-generator-images}.
The Hilbert computations in
Theorem~\ref{thm:dyadic-exact-hilbert-reciprocity-marking}
then prove \eqref{eq:dyadic-full-generator-evaluation}.
For the quadratic relation, in the notation
$P_0=x_0^{\sigma^2}x_0$ and
$Q_0=[x_1,\sigma]$, its degree-two Magnus terms
are $X_0^2$ and $X_1\Sigma+\Sigma X_1$:
conjugation by $\sigma^2$ changes the
degree-one Magnus component of $x_0$
only in degree at least three.
\end{proof}

\begin{theorem}[Pro-$2$ residual compression of the marked full group]
\label{thm:dyadic-full-marked-pro2-compression}
Let $E/\mathbf Q_2$ be finite,
$\bar\rho:G_{\mathbf Q_2}\to\mathrm{GL}_r(k_E)$
be a genuine residual representation
whose finite image is a $2$-group,
and let $A$ be a complete Noetherian
local $\mathcal O_E$-algebra with residue
field $k_E$.  Every framed continuous
lift $\rho_A$ of $\bar\rho$ has pro-$2$
image and factors through
$P=G_{\mathbf Q_2}(2)$.
In the Roe--Turturean full marked presentation
\cite{RoeTurtureanGQ2}, the identities
\begin{equation}\label{eq:dyadic-pro2-marked-specialization}
\rho_A(\tau)=I_r,\qquad
\rho_A(\sigma^{\omega_2})=\rho_A(\sigma),
\qquad
\rho_A((x_i\tau)^{\omega_2})=\rho_A(x_i)
\quad (i=0,1)
\end{equation}
hold.  The tame relator is thereby eliminated,
and the full marked wild relator becomes
\begin{equation}\label{eq:dyadic-full-pro2-relator}
\boxed{\quad
R=P_0Q_0,\quad
P_0=x_0^{\sigma^2}x_0,\quad
Q_0=[x_1,\sigma].
\quad}
\end{equation}
This gives an equality of framed deformation
functors and therefore an isomorphism of
complete local representing rings
\begin{equation}\label{eq:dyadic-pro2-deformation-ring-equality}
R^\square_{\bar\rho}(G_{\mathbf Q_2})
\xrightarrow{\;\sim\;}
R^\square_{\bar\rho}(P).
\end{equation}
The two full-group marked equations on
this residual sector are ordinary
finite-word matrix equations once the
canonical pro-$2$ specialization is made.
No infinite profinite binomial series is
needed \emph{on this sector}; outside
pro-$2$ residual sectors, the formula of
Theorem~\ref{thm:dyadic-omega2-power-formula}
is still necessary.
\end{theorem}

\begin{proof}
Reduction identifies the image of $\rho_A$
modulo $\mathfrak m_A$ with the finite
$2$-group $\operatorname{im}\bar\rho$.
The congruence kernel of
$\mathrm{GL}_r(A)\to\mathrm{GL}_r(k_E)$
is pro-$2$ for the maximal-ideal
topology.  Therefore the compact image
$\rho_A(G_{\mathbf Q_2})$ is an extension
of pro-$2$ groups and is pro-$2$.
The universal property of the maximal
pro-$2$ quotient proves factorization.

In a pro-$2$ group, a relation
$g^{-1}tg=t^2$ forces $t=1$:
in any finite $2$-group quotient, a
nontrivial element cannot be conjugate
to its square since the two have
different orders.  Apply this to the
tame relation
$\tau^\sigma=\tau^2$.
The profinite idempotent $\omega_2$
acts as the identity on all pro-$2$
elements, so
$\sigma_2=\sigma$ and
$u_i=(x_i\tau)^{\omega_2}=x_i$.
The remaining marked words then simplify
successively as
$d_0=c_0=d_g=h_c=1$,
$g_0=\sigma^2$, and
$h_0=x_0^{\sigma^2}x_0$.
The full wild relation
$h_0u_1^{-1}x_1^\sigma c_0=1$
is precisely
\eqref{eq:dyadic-full-pro2-relator}.
The marked condition on the normal closure
of $x_0,x_1$ is automatically satisfied
in every pro-$2$ quotient.
Conversely every representation of the
one-relator pro-$2$ group in this residual
chart defines an admissible full-group
marked representation by putting
$\tau=1$.  Hence the framed lift functors
are naturally isomorphic on every
Artinian coefficient algebra; continuity
and completeness give
\eqref{eq:dyadic-pro2-deformation-ring-equality}.
The last assertion follows since all
remaining operations in $P_0Q_0$
are ordinary finite words.
\end{proof}

\begin{proposition}[Closed Fox row of the marked pro-$2$ relation]
\label{prop:dyadic-closed-marked-fox-row}
For the marked pro-$2$ presentation with
generators $(x_0,\sigma,x_1)$ and
$R=P_0Q_0$, where
\[
P_0=\sigma^{-2}x_0\sigma^2x_0,\qquad
Q_0=x_1^{-1}\sigma^{-1}x_1\sigma,
\]
the three completed \emph{left} Fox derivatives
have finite closed word expressions
\begin{equation}\label{eq:dyadic-closed-marked-fox-row}
\boxed{
\begin{aligned}
\partial_{x_0}R
  &=\sigma^{-2}+\sigma^{-2}x_0\sigma^2,\\
\partial_{\sigma}R
  &=-\sigma^{-1}-\sigma^{-2}
     +\sigma^{-2}x_0(1+\sigma)\\
  &\hspace{1em}
     +P_0(-x_1^{-1}\sigma^{-1}
           +x_1^{-1}\sigma^{-1}x_1),\\
\partial_{x_1}R
  &=P_0(-x_1^{-1}+x_1^{-1}\sigma^{-1}).
\end{aligned}}
\end{equation}
Their images in $\mathbb Z_2[[P]]$,
evaluated in any finite projective
coefficient module, give the complete
marked differential
$d^1_{\rm mark}:M^3\to M$
appearing in
Theorem~\ref{thm:dyadic-strict-nielsen-fox-chain}.
It is a literal finite integral
noncommutative matrix row, not a
numerically fitted cohomology operator.
\end{proposition}

\begin{proof}
The completed left Fox rule is
$\partial_i(uv)=\partial_i(u)+u\partial_i(v)$
and
$\partial_i(g^{-1})=-g^{-1}\partial_i(g)$.
Apply it first to
$P_0=\sigma^{-2}x_0\sigma^2x_0$:
\[
\partial_{x_0}P_0
 =\sigma^{-2}+\sigma^{-2}x_0\sigma^2,
\quad
\partial_\sigma P_0
 =-\sigma^{-1}-\sigma^{-2}
    +\sigma^{-2}x_0(1+\sigma).
\]
Next
\[
\partial_{x_1}Q_0
 =-x_1^{-1}+x_1^{-1}\sigma^{-1},
\quad
\partial_{\sigma}Q_0
 =-x_1^{-1}\sigma^{-1}
     +x_1^{-1}\sigma^{-1}x_1.
\]
The formula $\partial_i(P_0Q_0)
=\partial_iP_0+P_0\partial_iQ_0$
gives
\eqref{eq:dyadic-closed-marked-fox-row}.
Its evaluation is the Fox differential
from \eqref{eq:fox-d1}.
\end{proof}

\begin{theorem}[Strict integral Nielsen--Fox chain equivalence]
\label{thm:dyadic-strict-nielsen-fox-chain}
Retain Theorem~\ref{thm:dyadic-full-marked-pro2-compression},
let $B$ be a complete local coefficient algebra,
and let $M$ be a finite projective $B$-module
with continuous $P$-action.
Set
\[
S=\rho_M(s),\quad A_0=\rho_M(a),\quad
X_0=\rho_M(x_0)=S^{-2}A_0^{-1},
\quad
\Sigma=\rho_M(\sigma)=S,
\]
and let $Q_0=[x_1,\sigma]$ as above.
For a degree-one standard Fox coordinate
$(v_a,v_s,v_y)\in M^3$ define
\begin{equation}\label{eq:dyadic-strict-nielsen-s1}
\mathsf N^1(v_a,v_s,v_y)=
\left(
 -S^{-2}A_0^{-1}v_a
 -(S^{-1}+S^{-2})v_s,\quad
 v_s,\quad v_y
\right),
\end{equation}
where the target is in $(x_0,\sigma,x_1)$
coordinate order.  Put
\begin{equation}\label{eq:dyadic-strict-nielsen-s2}
\mathsf N^0=\operatorname{id}_M,\qquad
\mathsf N^2=-\rho_M(Q_0)^{-1}.
\end{equation}
Then $\mathsf N^\bullet$ is an \emph{invertible
strict integral chain map} from the
standard one-relator Fox complex of
$a^2s^4[s,y]$ to the marked pro-$2$
Fox complex of $P_0Q_0$:
\begin{equation}\label{eq:dyadic-strict-nielsen-chain-square}
\mathsf N^1 d_{\rm std}^0
=d_{\rm mark}^0\mathsf N^0,
\qquad
d_{\rm mark}^1\mathsf N^1
=\mathsf N^2d_{\rm std}^1.
\end{equation}
Its degree-one inverse is explicit:
\begin{equation}\label{eq:dyadic-strict-nielsen-inverse}
\begin{aligned}
v_s&=v_\sigma,\qquad v_y=v_{x_1},\\
v_a&=-X_0^{-1}v_{x_0}
     -X_0^{-1}(\Sigma^{-1}+\Sigma^{-2})v_\sigma.
\end{aligned}
\end{equation}
The chain isomorphism commutes strictly with
integral coefficient base change and
preserves the intrinsic cup pairing
on cohomology.  The chosen Fox diagonals
need not be strictly identified as
multiplicative chain maps.
\end{theorem}

\begin{proof}
For the left crossed-derivation convention,
$c(gh)=c(g)+\rho_M(g)c(h)$ and
$c(g^{-1})=-\rho_M(g)^{-1}c(g)$.
By the Nielsen identity
$x_0=s^{-2}a^{-1}$, direct evaluation gives
\[
c(x_0)=
 -(S^{-1}+S^{-2})c(s)
 -S^{-2}A_0^{-1}c(a),
\]
which is exactly
\eqref{eq:dyadic-strict-nielsen-s1}.
The inverse word
$a=x_0^{-1}\sigma^{-2}$
gives \eqref{eq:dyadic-strict-nielsen-inverse}.
For a principal degree-zero coboundary
$c(g)=(\rho_M(g)-1)v$,
both sides evaluate on the same group
elements, proving the first identity of
\eqref{eq:dyadic-strict-nielsen-chain-square}.

The relation words satisfy the exact
\emph{free-group} identity
\begin{equation}\label{eq:dyadic-nielsen-relator-conjugacy}
r_{\rm std}=a^2s^4[s,y]
          =(Q_0P_0)^{-1},\qquad
r_{\rm mark}=P_0Q_0
          =Q_0^{-1}r_{\rm std}^{-1}Q_0
\end{equation}
after replacing $a=x_0^{-1}\sigma^{-2}$,
$s=\sigma$, and $y=x_1$.
Because $\rho_M(r_{\rm std})=1$,
the crossed-derivation identity yields
\[
c(r_{\rm mark})
=-\rho_M(Q_0)^{-1}c(r_{\rm std}).
\]
Evaluation of arbitrary generator
$1$-cochains on the corresponding
free relator is precisely the completed
Fox differential $d^1$.
This proves the second identity of
\eqref{eq:dyadic-strict-nielsen-chain-square}
without requiring the generator values
to be cocycles.
Both $\mathsf N^1$ and $\mathsf N^2$
are invertible integral matrices
over $B$ and involve only transport
matrices and their inverses; therefore
base change is strict.
The invariant cup product is that of
$R\Gamma(P,M)$ (or its paired coefficient
objects); identifying both Fox complexes
with the same intrinsic cochains
shows that cup classes agree.
\end{proof}

\begin{theorem}[Tame-block contraction of the full marked Jacobian]
\label{thm:dyadic-tame-jacobian-fox-splitting}
Let $B$ and $M$ be as in
Theorem~\ref{thm:dyadic-strict-nielsen-fox-chain},
with the coefficient action factoring through $P$.
On the genuine pro-$2$ residual chart,
let $\mathcal J^\bullet_{\rm full}(M)$ be
the three-term \emph{linearized marked
relation complex}, with
\[
\mathcal J^0=M,\qquad
\mathcal J^1=M^4
\ \text{in coordinates }(\sigma,\tau,x_0,x_1),
\qquad
\mathcal J^2=M^2
\ \text{in the tame/wild relation order}.
\]
The degree-one differential evaluates
left crossed derivations on the
two full marked relators, and the
degree-zero differential is the
usual principal-cocycle map.
Writing $\Sigma=\rho_M(\sigma)$,
put
\begin{equation}\label{eq:dyadic-tame-L-unit}
\mathsf L=\Sigma^{-1}-2I_M
\in\operatorname{Aut}_B(M).
\end{equation}
After reordering degree-one coordinates as
$(x_0,\sigma,x_1;\tau)$, the marked
relation differential has block matrix
\begin{equation}\label{eq:dyadic-full-jacobian-block}
d^1_{\rm full}
=
\begin{pmatrix}
 0&\mathsf L\\
 d^1_{\rm mark}&\mathsf B_\tau
\end{pmatrix},
\end{equation}
where $\mathsf B_\tau:M\to M$ is the
finite-word derivative of the wild
relation in the tame direction at
$\tau=1$.  It also has the following
\emph{closed transport formula}.  Put
\[
X=\rho_M(x_0),\quad
G=\Sigma^2,\quad
Z=\Sigma^{-1}X\Sigma,\quad
P_0^\rho=G^{-1}XGX.
\]
Then
\begin{equation}\label{eq:dyadic-wild-tau-jacobian-closed}
\boxed{\quad
\mathsf B_\tau
=P_0^\rho(G^{-1}X+3X-I_M)
 +(Z^{-1}-I_M)X.
\quad}
\end{equation}
The output change
\begin{equation}\label{eq:dyadic-tame-row-elimination}
(u_{\rm tame},u_{\rm wild})
\longmapsto
\bigl(u_{\rm tame},
u_{\rm wild}-\mathsf B_\tau\mathsf L^{-1}u_{\rm tame}\bigr)
\end{equation}
is integral and invertible.
It yields an actual chain isomorphism
\begin{equation}\label{eq:dyadic-full-marked-chain-splitting}
\boxed{\quad
\mathcal J^\bullet_{\rm full}(M)
\cong
C^\bullet_{\rm Fox}(P,M)
\oplus
\bigl(0\to M
  \xrightarrow{\ \mathsf L\ }M\to0\bigr),
\quad}
\end{equation}
with the acyclic summand supported in
cohomological degrees $1$ and $2$.
Composing with
Theorem~\ref{thm:dyadic-strict-nielsen-fox-chain}
gives a \emph{direct integral chain-level
comparison} between the marked
full-$G_{\mathbf Q_2}$ Jacobian and the
standard dyadic Demu\v{s}kin Fox--Spencer
coefficient complex.

In this pro-$2$ image sector the full
two-relator marked Jacobian is therefore
cohomologically Fox-complete in degrees
$0,1,2$ and computes
$R\Gamma(G_{\mathbf Q_2},M)$.
This conclusion is \emph{not} extended to
arbitrary residual images, for which
the fixed full-group marked relator
Jacobian may retain higher syzygies
and the all-Sylow derived construction
remains necessary.
\end{theorem}

\begin{proof}
At a genuine pro-$2$ representation the
tame generator $\tau$ acts as the identity,
and by Theorem~\ref{thm:dyadic-full-marked-pro2-compression}
every occurrence of $\omega_2$ in the
marked relations is the identity
operation on the ambient pro-$2$
deformation sector, including first-order
semidirect extensions by $M$.
The tame relation can be written as
\[
 r_t=\sigma^{-1}\tau\sigma\tau^{-2}=1.
\]
For an arbitrary crossed derivation $c$
with $c(\sigma)=v_\sigma$ and
$c(\tau)=v_\tau$, and
$\rho_M(\tau)=I$, one calculates
\[
 c(r_t)=(\Sigma^{-1}-2I_M)v_\tau;
\]
all $v_\sigma$ terms cancel.
Modulo the maximal ideal of $B$,
the reduction of $\mathsf L$ is
$\bar\Sigma^{-1}$ since $2=0$ there,
hence $\mathsf L$ is invertible over
the complete local ring $B$.
For completeness, the closed
$\tau$-derivative in
\eqref{eq:dyadic-wild-tau-jacobian-closed}
is obtained directly from the literal
Roe--Turturean auxiliary words.
At $\tau=1$ and with $c(\tau)=v_\tau$
and all other generator variations zero,
\[
c(d_0)=Xv_\tau,\qquad
c(d_g)=G^{-1}Xv_\tau,\qquad
c(h_c)=0,\qquad
c(c_0)=(Z^{-1}-I_M)Xv_\tau.
\]
Since $h_0=(x_0^{g_0}x_0)d_gd_0^3h_c$
and $u_1=x_1\tau$, the first two
contributions to $c(r_{\rm wild})$ are
$P_0^\rho(G^{-1}X+3X)v_\tau$
and $-P_0^\rho v_\tau$, followed by
the displayed contribution from $c_0$.
This proves the closed formula
without an unspecified Fox derivative.

The derivative of the wild relation
with $\tau$ fixed equal to $1$ is
the Fox derivative of $P_0Q_0$;
its derivative with respect to
$v_\tau$ is by definition
$\mathsf B_\tau$.  This proves the
block display
\eqref{eq:dyadic-full-jacobian-block}.
The row operation
\eqref{eq:dyadic-tame-row-elimination}
kills $\mathsf B_\tau$ without changing
the first Fox block.

Since $\tau$ acts trivially on $M$,
the degree-zero principal cocycle has
$v_\tau=(\rho_M(\tau)-1)v=0$;
hence the row/column decomposition
also splits the entire complex,
not only its degree-one differential.
The one-block complex with differential
$\mathsf L$ is contractible, yielding
\eqref{eq:dyadic-full-marked-chain-splitting}.
The Nielsen chain maps above then give
the explicit comparison in the standard
$a,s,y$ coordinates.
The standard Fox resolution of the
Demu\v{s}kin group $P$ computes
$R\Gamma(P,M)$ by
Theorem~\ref{thm:fox-spencer-demushkin},
and inflation from $P$ computes
$R\Gamma(G_{\mathbf Q_2},M)$ for
pro-$2$ coefficient actions by
Theorem~\ref{thm:sylow-fox-retract}.
No claim about the naive marked complex
on other residual sectors is used.
\end{proof}

\begin{corollary}[Exact full-group rank-one hypersurface elimination]
\label{cor:dyadic-full-rankone-tame-elimination}
For the trivial residual rank-one
representation $\bar\rho=\mathbf1$,
write the four marked scalar generator
matrices as
\[
\Sigma=1+z_\sigma,\quad
T=1+z_\tau,\quad
X_0=1+z_0,\quad
X_1=1+z_1.
\]
The completed full-group tame and wild
marked relation ideal is
\begin{equation}\label{eq:dyadic-full-rankone-two-relator-ideal}
\boxed{\quad
I_{\rm marked}
=(z_\tau,\ z_0(z_0+2))
\subseteq
\mathcal O_E[[z_\sigma,z_\tau,z_0,z_1]].
\quad}
\end{equation}
Thus the full-group marked
complete local presentation is already
a regular sequence, with a contractible
tame variable and one genuine
dyadic hypersurface:
\[
R^\square_{\mathbf1}
\cong
\mathcal O_E[[z_\sigma,z_0,z_1]]/
(z_0(z_0+2)).
\]
Under the Nielsen coordinate transform
\begin{equation}\label{eq:dyadic-rankone-u-z0-change}
 1+u=X_0^{-1},\qquad
 u=-\frac{z_0}{1+z_0},
\end{equation}
the identity of ideals is
\begin{equation}\label{eq:dyadic-rankone-hypersurface-comparison}
u(u+2)=-\frac{z_0(z_0+2)}{(1+z_0)^2}.
\end{equation}
Consequently the two full-group marked
relations and the one-relator
Demu\v{s}kin deformation hypersurface
of Theorem~\ref{thm:dyadic-q2-rankone-integral-ring}
agree \emph{integrally and scheme-theoretically},
not merely on their tangent spaces
or after inverting $2$.
\end{corollary}

\begin{proof}
Every scalar lift of $\mathbf1$ takes
values in the pro-$2$ principal units.
Conjugation is trivial, and the tame
relation $\tau^\sigma=\tau^2$
therefore becomes $T=T^2$,
equivalently $z_\tau=0$ since
$T\in1+\mathfrak m_A$ is invertible.
All $\omega_2$ powers act as the identity
on these principal units.  More strongly,
for arbitrary scalar $T$ near $1$ the
\emph{literal} full marked tame and wild
relations simplify exactly to
\[
r_t=T^{-1},\qquad
r_{\rm wild}=X_0^2T^3.
\]
Indeed $d_0=d_g=T$, $h_c=c_0=1$,
$h_0=X_0^2T^4$, and
$u_1^{-1}x_1^\sigma=T^{-1}$.
With $T=1$, the wild relation
reduces to $P_0Q_0$ and, for scalar
commuting matrices, to $X_0^2=1$.
Thus its equation is exactly
$z_0(z_0+2)=0$ modulo $z_\tau$.
Elementary ideal elimination gives
\eqref{eq:dyadic-full-rankone-two-relator-ideal}
in the complete local ambient
power-series ring; the tame relation
has the explicit unit multiple
$T^{-2}(T-T^2)$ of $z_\tau$,
and the wild relation differs from
$X_0^2-1$ by an element of
$(z_\tau)$.  Both generators form
a regular sequence.
Finally
$a=x_0^{-1}\sigma^{-2}$ and $s=\sigma$
give $a s^2=x_0^{-1}$ in the
scalar quotient, so
$1+u=\rho(a)\rho(s)^2=X_0^{-1}$.
The displayed rational power-series
identity
\eqref{eq:dyadic-rankone-hypersurface-comparison}
is integral because $1+z_0$ is a
unit in the completed local ring.
This proves the exact identification.
\end{proof}

\begin{corollary}[The literal full-group rank-one Jacobian]
\label{cor:dyadic-full-rankone-literal-jacobian}
At the trivial rank-one integral character
with coefficient module $M=\mathcal O_E$,
the full marked two-relator cochain
Jacobian is the exact integer matrix
\begin{equation}\label{eq:dyadic-full-rankone-jacobian-matrix}
d^1_{\rm full}:
\mathcal O_E^4\longrightarrow\mathcal O_E^2,
\qquad
\boxed{\quad
\begin{pmatrix}
0&-1&0&0\\
0&3&2&0
\end{pmatrix},
\quad}
\end{equation}
in $(\sigma,\tau,x_0,x_1)$
source coordinates and
$(r_t,r_{\rm wild})$ target order.
Its preceding coboundary differential
$d^0:\mathcal O_E\to\mathcal O_E^4$
is zero.  Hence the full marked complex
has cohomology
\begin{equation}\label{eq:dyadic-full-rankone-jacobian-cohomology}
H^0\simeq\mathcal O_E,\qquad
H^1\simeq\mathcal O_E^2,\qquad
H^2\simeq\mathcal O_E/2\mathcal O_E.
\end{equation}
After reduction to $k_E$ its first
cohomology has dimension $3$ and
its second has dimension $1$;
the extra mod-$2$ tangent direction
is exactly the integral hypersurface
square obstruction.
Thus a single $2\times4$ matrix
simultaneously checks the full-group
relation derivative, the
Demu\v{s}kin Fox cohomology, and the
singularity of the rank-one
framed deformation ring.
\end{corollary}

\begin{proof}
At the trivial scalar point the
tame and wild relators are
$r_t=T^{-1}$ and $r_{\rm wild}=X_0^2T^3$
by the preceding corollary.
Their first derivatives with
respect to $(\sigma,\tau,x_0,x_1)$
are exactly $(-1)$ in the tame
$\tau$ column and $(3,2)$ in the
wild $(\tau,x_0)$ columns;
all other derivatives vanish.
A rank-one adjoint action is trivial,
so its principal coboundary $d^0$
vanishes.
The matrix has a unit pivot $-1$
and a remaining invariant factor $2$.
Its integral kernel has free rank $2$
and its cokernel is $\mathcal O_E/2$.
Modulo the residue uniformizer,
the remaining factor $2$ vanishes,
giving the stated degree-one and
degree-two dimensions.  These agree
with \eqref{eq:dyadic-q2-rankone-special-fiber}
and \eqref{eq:dyadic-untwisted-tate-groups}.
\end{proof}

\begin{remark}[Sharp scope of the direct integral bridge]
\label{rem:dyadic-full-chain-bridge-scope}
If $\operatorname{im}\bar\rho$ is a
$2$-group, its Sylow preimage
in Theorem~\ref{thm:sylow-fox-retract}
is all of $G_{\mathbf Q_2}$.
The present chain equivalence therefore
makes the pointwise Sylow--Fox
comparison completely native on
these sectors, without choosing
a chain homotopy inverse of a raw
Herr representative.
For a general residual image,
the Sylow preimage is a proper
finite-index subgroup; the
full-group marked relation
Jacobian need not be cohomologically
complete, and an explicit
degree-two integral comparison
requires actual Schreier
relations and syzygy control.
The all-Sylow/derived construction
retains the unconditional result.
Neither this restriction nor the
finite-jet profinite-power formula
of fixed126 is a claim of a uniform
four-generator full-group
cohomological Fox resolution
for all residual sectors.
\end{remark}

% Version fixed128 standalone insertion, fully integrated into the mother TeX.
% New non-pro-2 residual S_3 sector, integral Schreier transfer, Hecke
% projector, and a coefficient-specific full marked Jacobian Smith certificate.

\subsection{The first non-pro-$2$ residual sector: tame $S_3$, Schreier,
Hecke, and a native integral Jacobian certificate}
\label{subsec:dyadic-s3-schreier-hecke}

The strict Nielsen bridge of
Theorem~\ref{thm:dyadic-tame-jacobian-fox-splitting}
assumes that the \emph{entire} residual image is pro-$2$.
The smallest honest test outside that domain is the tame
$S_3=\mathrm{GL}_2(\mathbf F_2)$ quotient of $G_{\mathbf Q_2}$.
For this quotient the Sylow preimage has odd index three
but is nonnormal.  We give its actual coset words,
integral transfer and Hecke identities, and a nontrivial
second-degree coefficient calculation.  The comparison is
at the coefficient level; no uniform full-group
resolution is inferred from the two marked relators.

Put $\Gamma=G_{\mathbf Q_2}$, choose
$\alpha^3=2$, and define
\[
F=\mathbf Q_2(\alpha),\qquad
L=F(\zeta_3).
\]
Let $\pi:\Gamma\twoheadrightarrow S_3$ be the
permutation action on the three cubic roots,
and write $N=\ker\pi$.  Choose compatible tame
Roe--Turturean generators $\sigma,\tau$ whose
images are a transposition $s$ and a $3$-cycle $u$:
$sus^{-1}=u^{-1}=u^2$.
The wild marked generators $x_0,x_1$
map to $1$ in this tame quotient.
Take the Sylow subgroup $S=\langle s\rangle$
and $H=\pi^{-1}(S)=G_F$, so that
\[
[\Gamma:H]=3,\qquad [H:N]=2.
\]

\begin{theorem}[A concrete non-pro-$2$ integral residual carrier]
\label{thm:dyadic-s3-integral-sector}
The extension $F/\mathbf Q_2$ is totally
and tamely ramified of degree $3$;
$L/F$ is unramified quadratic, and
$\operatorname{Gal}(L/\mathbf Q_2)\simeq S_3$.
On the standard integral rank-two
representation $T=\mathbf Z_2^2$, the
chosen images act by
\begin{equation}\label{eq:dyadic-s3-integral-matrices}
\boxed{\quad
\rho_T(\sigma)=
S_0=\begin{pmatrix}0&1\\1&0\end{pmatrix},
\qquad
\rho_T(\tau)=
U_0=\begin{pmatrix}-1&-1\\1&0\end{pmatrix},
\qquad
\rho_T(x_0)=\rho_T(x_1)=I_2.
\quad}
\end{equation}
They satisfy $S_0^2=U_0^3=I_2$ and
$S_0U_0S_0^{-1}=U_0^2$.
Reduction modulo $2$ gives the faithful
natural representation
$\bar\rho_T:\Gamma\twoheadrightarrow
\mathrm{GL}_2(\mathbf F_2)$,
whose image is \emph{not} a $2$-group,
while $\rho_T(H)=\langle S_0\rangle$
is a $2$-group.
The maximal pro-$2$ local quotient
$P_H=G_F(2)$ is a Demu\v{s}kin group
of rank five with $q(P_H)=2$ and
full cyclotomic orientation image
$\mathbf Z_2^\times$.
It therefore admits a (not yet
Schreier-marked) five-generator
one-relator model
\begin{equation}\label{eq:dyadic-s3-h-demushkin-five-relator}
P_H\cong\bigl\langle z_1,\ldots,z_5
\,\big|\,
z_1^2z_2^4[z_2,z_3][z_4,z_5]
\bigr\rangle_{\mathrm{pro}-2}.
\end{equation}

The complete integral cohomology ledgers
in degrees zero through two are
\begin{equation}\label{eq:dyadic-s3-natural-cohom-ledger}
\boxed{
\begin{array}{c|ccc}
 &H^0(-,T)&H^1(-,T)&H^2(-,T)\\ \hline
 \Gamma&0&\mathbf Z_2^{\,2}&0\\
 H&\mathbf Z_2&\mathbf Z_2^{\,7}&\mathbf F_2 .
\end{array}}
\end{equation}
Consequently the integral Sylow transfer
idempotent
\begin{equation}\label{eq:dyadic-s3-sylow-idempotent}
e_H=\tfrac13\,\operatorname{res}_H^\Gamma
\operatorname{cor}_H^\Gamma
\end{equation}
annihilates $H^0(H,T)$ and $H^2(H,T)$.
On $H^1(H,T)\cong\mathbf Z_2^7$
it is an actual integral projection
onto a free rank-two summand, with
free rank-five kernel.
\end{theorem}

\begin{proof}
The polynomial $X^3-2$ is Eisenstein
at $2$, so $F/\mathbf Q_2$ has degree
three and ramification index three.
Because this index is odd, the
ramification is tame.  The polynomial
$X^2+X+1$ is irreducible modulo $2$;
its root $\zeta_3$ generates the
unramified quadratic extension.
The Galois closure is thus a
six-element tame extension with the
usual $S_3$ action.
The matrices in
\eqref{eq:dyadic-s3-integral-matrices}
satisfy the displayed finite-group
relations, and their reductions
generate $\mathrm{GL}_2(\mathbf F_2)$.
They define the standard integral
permutation-difference lattice.
The odd-degree extension $F/\mathbf Q_2$
contains no $\mu_4$, so the exact
$2$-primary roots of unity in $F$
form $\mu_2$.  Moreover the
cyclotomic pro-$2$ tower of
$\mathbf Q_2$ has trivial
intersection with $F$, since
their finite intersection would
have degree dividing both $3$
and a power of $2$.
The orientation image for $G_F(2)$
is therefore all of $\mathbf Z_2^\times$.
The local Demu\v{s}kin classification
then gives rank $[F:\mathbf Q_2]+2=5$,
$q=2$, and the odd-rank $f=2$
standard relation
\eqref{eq:dyadic-s3-h-demushkin-five-relator}.

For the cohomology computation put
$V=T\otimes_{\mathbf Z_2}\mathbf Q_2$ and
$D=T^\vee(1)\otimes_{\mathbf Z_2}
(\mathbf Q_2/\mathbf Z_2)$.
The wild inertia $W$ acts trivially
on $T$ because $\rho_T$ is tame.
Its cyclotomic image contains $-1$
(and in fact the entire cyclotomic
pro-$2$ image), so
$D^W=D[2]=T^\vee/2T^\vee$.
The finite residual $S_3$ action
on $T^\vee/2$ has no nonzero
invariants: the $3$-cycle $U_0$
has no eigenvalue $1$ modulo $2$.
Integral local Tate duality
therefore gives $H^2(\Gamma,T)=0$.
For $H$, the representation is
unramified quadratic, and
$(T^\vee/2)^{\langle S_0\rangle}$
is one-dimensional, giving
$H^2(H,T)\cong\mathbf F_2$.

The fixed lattices are
$T^\Gamma=0$ and
$T^H=\mathbf Z_2(1,1)^t$.
The local Euler--Poincar\'e rank
formula \cite{NSW} gives
\[
\operatorname{rank}_{\mathbf Z_2}H^1(\Gamma,T)=2,
\qquad
\operatorname{rank}_{\mathbf Z_2}H^1(H,T)=2[ F:\mathbf Q_2]+1=7.
\]
These groups have no $2$-primary
torsion.  Indeed the invariants in
$V/T$ are zero for $\Gamma$, since
$(T/2T)^\Gamma=0$; for $H$ the
invariants $(V/T)^H$ are the full
diagonal $\mathbf Q_2/\mathbf Z_2$
and are already the image of $V^H$.
The long exact sequence for
$0\to T\to V\to V/T\to0$
therefore introduces no torsion in
$H^1$.
This proves
\eqref{eq:dyadic-s3-natural-cohom-ledger}.
Finally $3$ is a unit in
$\mathbf Z_2$, and the standard
restriction--corestriction identity
makes \eqref{eq:dyadic-s3-sylow-idempotent}
a derived idempotent whose image is
$R\Gamma(\Gamma,T)$.
Its actions in degrees zero, one
and two follow from the ledger
and the splitting of finitely
generated free $\mathbf Z_2$-modules.
\end{proof}

\begin{theorem}[Finite nonnormal Schreier table and corestriction]
\label{thm:dyadic-s3-schreier-corestriction}
Choose right-coset representatives
$t_0=1,t_1=\tau,t_2=\tau^2$
for $H\backslash\Gamma$.
Write
\[
t_i g=h_{i,g}\,t_{j(i,g)},\qquad
h_{i,g}\in H,
\]
where $j(i,g)$ is determined by the
right $S_3$-coset action.
Then the complete Schreier data
for the four marked generators are
\begin{equation}\label{eq:dyadic-s3-schreier-table}
\boxed{
\begin{array}{c|ccc}
g& i=0&i=1&i=2\\ \hline
\sigma&
(0,\sigma)&(2,\sigma)&(1,\sigma\tau^3)\\
\tau&
(1,1)&(2,1)&(0,\tau^3)\\
x_j&
(0,x_j)&(1,\tau x_j\tau^{-1})&
(2,\tau^2x_j\tau^{-2}),\quad j=0,1 .
\end{array}}
\end{equation}
An entry $(k,h)$ means $j(i,g)=k$
and $h_{i,g}=h$.
In particular a finite profinite
Schreier generating system for $H$
consists of
\[
\sigma,\ \tau^3,\quad
\tau^i x_j\tau^{-i}
\quad(i=0,1,2,\ j=0,1).
\]
There are eight displayed
generators after removal of trivial
and repeated Schreier words.
The six wild conjugates have
pro-$2$ normal closure, while the
tame element $\tau^3$ has pro-odd
order.  Thus $\tau^3$ maps to $1$
under $H\twoheadrightarrow P_H=H(2)$.

For any continuous $H$-cocycle
$c\in Z^1(H,T)$, the corestriction
is represented on marked generators
by the finite coset formula
\begin{equation}\label{eq:dyadic-s3-transfer-one}
(\operatorname{cor}_H^\Gamma c)(g)
=\sum_{i=0}^{2}
  \rho_T(t_i)^{-1}c(h_{i,g}).
\end{equation}
For the specific rank-two $S_3$
module \eqref{eq:dyadic-s3-integral-matrices}
this simplifies \emph{exactly} to
\begin{equation}\label{eq:dyadic-s3-transfer-simplified}
\boxed{
\begin{aligned}
(\operatorname{cor}c)(\sigma)&=0,\\
(\operatorname{cor}c)(\tau)&=0,\\
(\operatorname{cor}c)(x_j)
 &=\sum_{i=0}^{2} U_0^{-i}
   c(\tau^i x_j\tau^{-i}),
 \qquad j=0,1 .
\end{aligned}}
\end{equation}
This is an integral formula, with
no average or division by $2$;
the normalized Sylow projector
divides by the odd index $3$ only
\emph{after} corestriction.

The analogous higher-degree bar
corestriction is obtained by tracking
successive coset indices:
\begin{equation}\label{eq:dyadic-s3-transfer-all-arity}
(\operatorname{cor}c)(g_1,\dots,g_n)
=\sum_{i=0}^2\rho_T(t_i)^{-1}
c\!\left(h_{i,g_1},
h_{j(i,g_1),g_2},\ldots,
h_{j(i,g_1\cdots g_{n-1}),g_n}\right).
\end{equation}
It is compatible with the normalized
bar differential and hence realizes
the usual corestriction on
cohomology; this is \emph{not}
by itself a finite Demu\v{s}kin
Fox change-of-basis formula.
\end{theorem}

\begin{proof}
The tame relation is
$\sigma^{-1}\tau\sigma=\tau^2$,
so $\tau^i\sigma=\sigma\tau^{2i}$
as an exact profinite word identity.
Reducing $2i$ modulo $3$ against
the chosen coset representatives
gives the first line of
\eqref{eq:dyadic-s3-schreier-table};
the second and third lines follow
from right multiplication by
$\tau$ and by a wild generator
with trivial residual $S_3$ image.
The Schreier rewriting theorem for
finite-index open profinite subgroups
gives the displayed generating
system; the normal closure of the
wild conjugates is contained in
the marked wild pro-$2$ subgroup,
and it is pro-$2$.
The element $\tau^3$ lies in tame
inertia, of pro-odd order, so its
image in the maximal pro-$2$
quotient of $H$ is trivial.

The standard bar corestriction for
right-coset representatives is
\eqref{eq:dyadic-s3-transfer-all-arity}.
For degree one the crossed-cocycle
identity and
$h_{i,g_1g_2}=
h_{i,g_1}h_{j(i,g_1),g_2}$
show directly that its output is
a $\Gamma$-cocycle, proving
\eqref{eq:dyadic-s3-transfer-one}.
The same identity gives the bar
chain-map formula in every degree.

Because the action of $\tau^3$
on $T$ is trivial and its closed
cyclic group has pro-odd order,
every continuous crossed
homomorphism $H\to T$ vanishes on
$\tau^3$.  Therefore the
$\sigma$-column of the table gives
\[
(\operatorname{cor}c)(\sigma)
=(I+U_0^{-1}+U_0^{-2})c(\sigma)=0,
\]
using $I+U_0+U_0^2=0$.
The $\tau$ column gives zero, and
the two wild columns give exactly
\eqref{eq:dyadic-s3-transfer-simplified}.
\end{proof}

\begin{theorem}[The index-three integral Hecke--Sylow polynomial]
\label{thm:dyadic-s3-hecke-sylow-polynomial}
On $R\Gamma(H,T)$ define
\begin{equation}\label{eq:dyadic-s3-hecke-operator}
\mathsf A=\operatorname{res}_H^\Gamma
\operatorname{cor}_H^\Gamma,
\qquad
\mathsf T_{\rm Hk}=\mathsf A-I.
\end{equation}
These operators obey, as derived
endomorphisms (and hence in every
cohomological degree),
\begin{equation}\label{eq:dyadic-s3-hecke-minpoly}
\boxed{\qquad
\mathsf A^2=3\mathsf A,\qquad
\mathsf T_{\rm Hk}^2=
\mathsf T_{\rm Hk}+2I,\qquad
e_H=\frac{I+\mathsf T_{\rm Hk}}3.
\qquad}
\end{equation}
The nonidentity term also has
the exact Mackey double-coset
description
\begin{equation}\label{eq:dyadic-s3-double-coset-hecke}
\mathsf T_{\rm Hk}=
\operatorname{cor}_N^H\circ
(\tau)_*\circ
\operatorname{res}_N^H,
\qquad N=\ker\pi,
\end{equation}
where $(\tau)_*$ is conjugation on
the normal subgroup $N$, together
with its prescribed action on $T$.
Both the restriction and the
corestriction through $N$ are
integrally defined: no division
by $[H:N]=2$ occurs.

For the coefficient lattice $T$,
the two eigensummands selected by
$e_H$ have the following precise
meaning:
\[
\begin{array}{c|ccc}
& H^0(H,T)&H^1(H,T)&H^2(H,T)\\ \hline
\operatorname{im}e_H&0&\mathbf Z_2^2&0\\
\ker e_H&\mathbf Z_2&\mathbf Z_2^5&\mathbf F_2.
\end{array}
\]
Thus the genuine non-pro-$2$
descent kills a torsion top Fox
class even though the index-three
operator itself remains integral.
\end{theorem}

\begin{proof}
The standard cohomological identity
$\operatorname{cor}\operatorname{res}
=3I$ implies
$\mathsf A^2=3\mathsf A$.
Substituting $\mathsf T_{\rm Hk}=\mathsf A-I$
gives the displayed quadratic
polynomial and projector formula.
In the finite residual quotient,
$S=\langle s\rangle$ has precisely
two double cosets in $S_3$:
$S$ and $SuS$.  The latter
has intersection
$S\cap uSu^{-1}=\{1\}$,
whose inverse image in $\Gamma$
is exactly $N$.
The Mackey restriction--corestriction
formula therefore yields
\eqref{eq:dyadic-s3-double-coset-hecke}.
All maps are defined before
inverting $2$ and require no
normality of $H$.
The final table follows from
Theorem~\ref{thm:dyadic-s3-integral-sector}
and the idempotent's identification
with $R\Gamma(\Gamma,T)$.
\end{proof}

\begin{theorem}[The first non-pro-$2$ integral adjoint cohomology ledger]
\label{thm:dyadic-s3-adjoint-local-cohomology}
Retain the integral $S_3$ lattice $T$
and put
\[
M=\operatorname{ad}T
=\operatorname{End}_{\mathbf Z_2}(T),
\qquad \bar M=M/2M.
\]
The cohomology of $\Gamma=G_{\mathbf Q_2}$
is
\begin{equation}\label{eq:dyadic-s3-adjoint-integral-ledger}
\boxed{\quad
H^0(\Gamma,M)=\mathbf Z_2,\qquad
H^1(\Gamma,M)=\mathbf Z_2^{\,5},\qquad
H^2(\Gamma,M)=\mathbf F_2.
\quad}
\end{equation}
Its mod-$2$ cohomology dimensions are
\begin{equation}\label{eq:dyadic-s3-adjoint-mod2-ledger}
\boxed{\quad
\bigl(\dim H^0,\dim H^1,\dim H^2\bigr)
(\Gamma,\bar M)=(1,6,1).
\quad}
\end{equation}
The tame order-three averaging operator
is integral on both $T$ and $M$:
\begin{equation}\label{eq:dyadic-s3-tame-averaging}
\Pi_{\rm tame}^T
=\tfrac13(I+U_0+U_0^2)=0,\qquad
\Pi_{\rm tame}^M(A)
=\tfrac13\sum_{j=0}^2U_0^jAU_0^{-j}.
\end{equation}
The second projector has image
$\mathbf Z_2[U_0]$ of rank two;
conjugation by $S_0$ fixes exactly
the scalar rank-one submodule.
\end{theorem}

\begin{proof}
Over $\mathbf F_2$, the natural
representation of $S_3=\mathrm{GL}_2(\mathbf F_2)$
is absolutely irreducible and its
endomorphism centralizer consists
only of scalars.  Consequently
$M^\Gamma=\mathbf Z_2 I$ and
$\bar M^\Gamma=\mathbf F_2 I$.
As in Theorem~\ref{thm:dyadic-s3-integral-sector},
wild inertia acts trivially on $M$
but has cyclotomic image containing
$-1$.  Therefore the invariants in
\[
M^\vee(1)\otimes_{\mathbf Z_2}
(\mathbf Q_2/\mathbf Z_2)
\]
are precisely the $S_3$-invariants
of $M^\vee/2M^\vee$.
The trace pairing identifies the
latter with $\bar M$;
it has one-dimensional invariant
space.  Integral local Tate duality
gives $H^2(\Gamma,M)\cong\mathbf F_2$.
Rational local Euler--Poincar\'e
gives
$\operatorname{rank}H^1(\Gamma,M)
=4+\operatorname{rank}H^0(\Gamma,M)=5$.

There is no torsion in $H^1$:
if an invariant class in
$(M\otimes\mathbf Q_2)/M$
has denominator $2^n$, its
numerator modulo $2$ is
$S_3$-invariant and hence scalar.
Subtracting an appropriate
rational scalar matrix reduces
the denominator; iteration shows
that all these invariants come
from $(M\otimes\mathbf Q_2)^\Gamma$.
The coefficient long exact
sequence then forces
$H^1(\Gamma,M)$ to be torsion-free.
The mod-$2$ Euler formula with
$\dim H^0=\dim H^2=1$ gives
$\dim H^1=4+1+1=6$.

Finally $U_0$ has
$U_0^2+U_0+I=0$, and
$3\in\mathbf Z_2^\times$,
so the displayed tame averages
are genuine integral projectors.
The centralizer of $U_0$ in
$M_2(\mathbf Z_2)$ is
$\mathbf Z_2[U_0]$ because its
minimal polynomial is irreducible
modulo $2$.  Conjugation by
$S_0$ sends $U_0$ to $U_0^2$;
its fixed part is exactly
$\mathbf Z_2 I$.
\end{proof}

\begin{theorem}[Explicit full-marked $8\times16$ integral Jacobian certificate]
\label{thm:dyadic-s3-marked-adjoint-jacobian}
Use the full Roe--Turturean
$\Gamma$-presentation of
Theorem~\ref{thm:dyadic-q2-marked-word-jet-solver}
and the integral residual lift
\eqref{eq:dyadic-s3-integral-matrices}.
Let
\begin{equation}\label{eq:dyadic-s3-marked-adjoint-complex}
\mathcal J^\bullet_{S_3}(M):
\quad
M\xrightarrow{d^0} M^{4}
\xrightarrow{d^1} M^{2}
\end{equation}
be the completed marked two-relation
Jacobian complex at this lift.
Its differentials are the principal
cocycle differential and the first
derivatives of the literal tame/wild
marked matrix relations.
In the row-major basis
$(E_{11},E_{12},E_{21},E_{22})$
of $M$, order source variables by
$(\sigma,\tau,x_0,x_1)$
and target relation matrices by
(tame, wild).  Then the reduction
of $d^1$ modulo $2$ is the
\emph{actual} following $8\times16$
matrix:
\begin{equation}\label{eq:dyadic-s3-jacobian-mod2-matrix}
\boxed{\ \!
{\renewcommand{\arraystretch}{0.95}
\setlength{\arraycolsep}{2.5pt}
\begin{array}{rrrr|rrrr|rrrr|rrrr}
1&1&0&1&1&1&0&0&0&0&0&0&0&0&0&0\\
0&1&1&0&1&0&0&1&0&0&0&0&0&0&0&0\\
1&0&1&1&0&0&1&0&0&0&0&0&0&0&0&0\\
1&1&0&1&0&1&0&1&0&0&0&0&0&0&0&0\\ \hline
0&0&0&0&1&1&1&0&0&0&0&0&1&1&1&1\\
0&0&0&0&1&0&0&1&0&0&0&0&1&0&1&1\\
0&0&0&0&1&0&0&1&0&0&0&0&1&1&0&1\\
0&0&0&0&0&1&1&1&0&0&0&0&1&1&1&1
\end{array}}\ \! } .
\end{equation}
Its rank is $7$, and its
left nullspace is generated by
\begin{equation}\label{eq:dyadic-s3-jacobian-left-null}
\ell=(1,0,0,1\mid1,0,0,1),
\end{equation}
the sum of the two matrix traces.

More strongly, the genuine
\emph{integral} Jacobian has
Smith invariant factors
\begin{equation}\label{eq:dyadic-s3-jacobian-smith}
\boxed{\quad
\operatorname{SNF}_{\mathbf Z_2}(d^1)
=\operatorname{diag}
(1,1,1,1,1,1,1,2)
\quad}
\end{equation}
with eight additional zero
source columns.
An explicit finite-precision
certificate is obtained as follows:
work in
$\bigl(\mathbf Z/4[\varepsilon]/
(\varepsilon^2)\bigr)$, use the
left variations
\[
\rho_\varepsilon(g)
=(I+\varepsilon A_g)\rho_T(g),
\quad g\in\{\sigma,\tau,x_0,x_1\},
\]
set $\sigma^{\omega_2}=\sigma$,
and evaluate
$(x_i\tau)^{\omega_2}=(x_i\tau)^9$.
In the resulting $8\times16$
Jacobian modulo $4$, the
$8\times8$ minor formed by the
one-based columns
\begin{equation}\label{eq:dyadic-s3-jacobian-minor-columns}
\boxed{\qquad
\{1,2,5,6,9,13,14,15\}
\qquad}
\end{equation}
gives the following fully displayed
minor, with those columns in
increasing order:
\begin{equation}\label{eq:dyadic-s3-jacobian-mod4-minor}
\mathsf D_4=
\begin{pmatrix}
1&3&3&3&0&0&0&0\\
0&3&1&2&0&0&0&0\\
3&0&0&2&0&0&0&0\\
3&1&0&1&0&0&0&0\\
0&0&1&1&2&1&1&3\\
0&0&3&2&0&3&2&3\\
0&0&1&2&0&1&3&2\\
0&0&2&3&0&3&3&1
\end{pmatrix}
\pmod4,\qquad
\boxed{\det(\mathsf D_4)\equiv2\pmod4}.
\end{equation}
Together with
the rank-seven reduction
\eqref{eq:dyadic-s3-jacobian-mod2-matrix},
this proves
\eqref{eq:dyadic-s3-jacobian-smith}.

The preceding degree-zero map
has mod-$2$ rank three.
Consequently the entire marked
complex has the integral cohomology
\begin{equation}\label{eq:dyadic-s3-marked-cochain-ledger}
H^0(\mathcal J_{S_3})=\mathbf Z_2,\qquad
H^1(\mathcal J_{S_3})=\mathbf Z_2^5,\qquad
H^2(\mathcal J_{S_3})=\mathbf F_2.
\end{equation}
It is homotopy equivalent, after
finite integral row/column reduction,
to the finite elementary complex
\begin{equation}\label{eq:dyadic-s3-marked-smith-model}
\boxed{\quad
\mathbf Z_2[0]\ \oplus\
\mathbf Z_2^{\,5}[-1]\ \oplus\
\bigl[\mathbf Z_2
\xrightarrow{\ 2\ }\mathbf Z_2\bigr]_{[1,2]} .
\quad}
\end{equation}
In particular, the full marked
two-relator coefficient Jacobian
is \emph{cohomologically complete}
for this specified integral
$S_3$ \emph{adjoint} representation.
It is not asserted to be a
projective group resolution over
$\mathbf Z_2[[\Gamma]]$.
\end{theorem}

\begin{proof}
The residual group action on
$\operatorname{ad}\bar T$ is
conjugation by the two matrices
$S_0,U_0$.  In the dual-number
ring modulo $2$, every lift
of a matrix with residual order
three has $2$-primary exponent
at most two; its $\omega_2$
projection is therefore its third
power.  Modulo $4$ the corresponding
exponent divides four: every congruence
element in the dual-number ring has form
$I+K$ with $K=2A+\varepsilon B$;
one has $K^2=2\varepsilon(AB+BA)$
and $K^3=0$, hence $(I+K)^4=I$.
Thus the CRT exponent $9$, congruent
to $1$ modulo $4$ and $0$ modulo $3$,
computes its $\omega_2$ projection.
For the residual order-two matrix
$S_0$, the whole lift is pro-$2$,
so its $\omega_2$ projection is
itself.  These are finite instances
of the integral binomial theorem
\ref{thm:dyadic-omega2-power-formula}.

Substitute these powers into the
literal auxiliary words
$\sigma_2,u_i,d_0,z_0,c_0,g_0,
d_g,h_c,h_0$
of Roe--Turturean, and differentiate
the two relation matrices with
respect to the sixteen standard
left tangent directions
$A_g=E_{ij}$.
This gives exactly
\eqref{eq:dyadic-s3-jacobian-mod2-matrix};
a direct row elimination yields
rank seven.
The displayed $\ell$ annihilates
every column, so the left
nullspace is exactly one-dimensional.
Performing the same calculation
modulo $4$ with exponent $9$
gives the minor specified by
\eqref{eq:dyadic-s3-jacobian-minor-columns},
of determinant $2$ modulo $4$.
Since all $8\times8$ minors
vanish modulo $2$ but at least one
has valuation exactly one,
the invariant factors are seven
units followed by a single factor
of valuation one.  This proves
\eqref{eq:dyadic-s3-jacobian-smith}.

The principal-cocycle map has
kernel equal to the scalar
centralizer $M^\Gamma=\mathbf Z_2$;
its mod-$2$ reduction has kernel
$\mathbf F_2$, so $d^0$ has
three unit invariant factors.
By $d^1d^0=0$ and integral
elementary divisors, the image
of $d^0$ is a saturated rank-three
submodule of $\ker d^1$.
Therefore $H^1$ is free of
rank $16-8-3=5$, while $H^2$
is the cokernel $\mathbf F_2$.
The degree-zero invariant
space is $\mathbf Z_2$.
Removing the seven unit
degree-one--two summands and
three unit degree-zero--one
summands proves the homotopy
normal form
\eqref{eq:dyadic-s3-marked-smith-model}.
\end{proof}

\begin{corollary}[Coefficient-specific comparison with the Sylow--Fox carrier]
\label{cor:dyadic-s3-adjoint-smith-sylow-comparison}
For the rank-four coefficient
$M=\operatorname{ad}T$ of
Theorem~\ref{thm:dyadic-s3-adjoint-local-cohomology},
both the full marked coefficient
Jacobian
$\mathcal J^\bullet_{S_3}(M)$ and
the all-Sylow Fox retract
$\operatorname{Im}(e_H)$ are bounded
perfect integral complexes
with cohomology
$(\mathbf Z_2,\mathbf Z_2^5,\mathbf F_2)$
in degrees $(0,1,2)$.
They are therefore isomorphic in
$D(\mathbf Z_2)$ and can be
connected by an integral
chain-homotopy equivalence
\emph{after choices of Smith
reduction and a top-degree
Tate-dual generator}.
The degree-zero and degree-one
identifications are canonically
fixed by invariants and evaluation
of actual framed deformations.
The top-degree identification
may be normalized by the
mod-$2$ trace obstruction
\eqref{eq:dyadic-s3-jacobian-left-null}
and local Tate duality.

This removes the \emph{existence}
and integer-divisor gap for
a direct finite coefficient
comparison in the first
non-pro-$2$ adjoint sector.
It does \emph{not} supply a
canonical strict marked-word
chain map to the chosen five-
generator Fox resolution of $P_H$,
nor does it establish a
presentation-independent
second-syzygy theorem for
all non-pro-$2$ residual sectors.
\end{corollary}

\begin{proof}
Theorem~\ref{thm:dyadic-s3-marked-adjoint-jacobian}
gives the explicit integral Smith
model of the marked complex.
The all-Sylow retract computes
$R\Gamma(\Gamma,M)$ by
Theorem~\ref{thm:all-sylow-transfer}.
Its cohomology is computed in
Theorem~\ref{thm:dyadic-s3-adjoint-local-cohomology}.
Over the discrete valuation
ring $\mathbf Z_2$, any bounded
perfect complex admits finite
elementary-divisor decompositions,
with its homotopy type determined
by its cohomology modules in each
degree.  Hence its finite
perfect representatives may be
reduced to the same elementary
complex
\eqref{eq:dyadic-s3-marked-smith-model},
which gives an integral chain
homotopy equivalence after the
chosen reductions.
The actual cocycle relation
follows from the marked
presentation and residual
pro-$2$ automaticity
(Theorem~\ref{thm:marked-residual-strictification}),
while local Tate duality selects
a nonzero generator of the
one-dimensional top torsion.
No explicit universal
Schreier-to-minimal-Fox basis is
constructed by this argument,
explaining the stated limitation.
\end{proof}

\begin{warning}[What the $S_3$ computation does not prove]
\label{warn:dyadic-s3-syzygy-scope}
The Schreier table gives finite
words in the full group, and
the all-arity bar transfer is
an exact integral operation.
Turning those ten raw
Schreier-word entries into a
canonical minimal dyadic
Demu\v{s}kin basis for
$P_H=G_{\mathbf Q_2(\sqrt[3]2)}(2)$,
with the necessary
integral relation syzygies and
an explicit multiplicative
chain comparison, is a
strictly stronger problem.
The $S_3$ natural and adjoint
coefficient results close their
own cohomology, Hecke and
Smith-normal-form ledgers,
not the general non-pro-$2$
full-word resolution problem.
The existence of the
all-Sylow derived atlas
remains unconditional and is
not upgraded here to a
uniform full-group projective
Fox resolution.
\end{warning}

% Fixed129 independent substantive insertion.
% A strict integral natural-coefficient marked Jacobian contraction,
% Schreier evaluation projector and derived Hecke equivalence
% on the genuine non-pro-2 S3 residual sector of fixed128.

\subsection{Strict natural-coefficient descent in the non-pro-$2$ $S_3$ sector}
\label{subsec:dyadic-s3-natural-strict-transfer}

The coefficient-specific adjoint Smith certificate of
Theorem~\ref{thm:dyadic-s3-marked-adjoint-jacobian}
does not, by itself, give an explicit map to the
Sylow-preimage cochains.  We now obtain a stronger,
\emph{strict chain-level} result for the distinct
natural rank-two coefficient lattice $T$.
Unlike the adjoint module, the natural $S_3$ module
satisfies the integral norm-zero relation
$I+U_0+U_0^2=0$.  This kills the pro-$2$
projection of every affine lift of the tame
order-three generator, and is precisely the
extra algebra needed for an integral closed
formula at all chain degrees.

Continue with $\Gamma=G_{\mathbf Q_2}$,
$H=G_{\mathbf Q_2(\sqrt[3]{2})}$,
$T=\mathbf Z_2^2$, and the fixed matrices
\[
S=S_0=\begin{pmatrix}0&1\\1&0\end{pmatrix},
\qquad
U=U_0=\begin{pmatrix}-1&-1\\1&0\end{pmatrix}
\]
of \eqref{eq:dyadic-s3-integral-matrices}.
Write $I=I_2$ and put
\begin{equation}\label{eq:dyadic-s3-natural-tame-blocks}
\begin{aligned}
R_\sigma&=S-I=
\begin{pmatrix}-1&1\\1&-1\end{pmatrix},
&
R_\tau&=U-I=
\begin{pmatrix}-2&-1\\1&-1\end{pmatrix},
\\
L_\sigma&=S^{-1}(U-I)
=\begin{pmatrix}1&-1\\-2&-1\end{pmatrix},
&
L_\tau&=S^{-1}-I-U
=\begin{pmatrix}0&2\\0&-1\end{pmatrix}.
\end{aligned}
\end{equation}
The integral determinants are
\begin{equation}\label{eq:dyadic-s3-natural-unit-dets}
\det(R_\tau)=3,\qquad
\det(L_\sigma)=-3,\qquad
\det(S)=-1.
\end{equation}
All these are units in $\mathbf Z_2$.

\begin{theorem}[Exact natural-coefficient full-marked Fox matrix]
\label{thm:dyadic-s3-natural-full-marked-matrix}
Let
$\mathcal J^\bullet_{\Gamma,\mathrm{nat}}(T)$
be the linear \emph{coefficient}
crossed-derivation complex obtained
from the full Roe--Turturean
marked relations in the module $T$
(not the adjoint tangent module
$\operatorname{ad}T$):
\begin{equation}\label{eq:dyadic-s3-natural-three-term-complex}
T\xrightarrow{d^0}T^{\oplus4}
\xrightarrow{d^1}T^{\oplus2},
\end{equation}
with degree-one source coordinates
$(a,b,c,d)$, meaning values at
$(\sigma,\tau,x_0,x_1)$, and target
coordinates in the (tame,wild)
relation order.  The \emph{complete
integral} matrices are
\begin{equation}\label{eq:dyadic-s3-natural-differentials}
\boxed{
\begin{aligned}
d^0(v)&=(R_\sigma v,\ R_\tau v,\ 0,\ 0),
\\
d^1(a,b,c,d)
&=(L_\sigma a+L_\tau b,\ -2c+S^{-1}d).
\end{aligned}}
\end{equation}
Equivalently, in the ordinary
coordinate basis the $4\times8$
degree-one matrix is
\begin{equation}\label{eq:dyadic-s3-natural-4x8-jacobian}
\boxed{\quad
d^1=
\begin{pmatrix}
 1&-1&0&2&0&0&0&0\\
 -2&-1&0&-1&0&0&0&0\\
 0&0&0&0&-2&0&0&1\\
 0&0&0&0&0&-2&1&0
\end{pmatrix}.
\quad}
\end{equation}
Both the tame block $L_\sigma$
and the wild $x_1$ block
$S^{-1}$ are invertible
integrally, while the $\tau$
component of $d^0$ is
$R_\tau$, another integral unit.
Thus $d^1$ has four integral
unit elementary divisors,
and $d^0$ has two further
unit pivots.  After cancelling
these four degree-one--two and
two degree-zero--one contractible
summands, its remaining degree-one
cohomology is free of rank two:
\begin{equation}\label{eq:dyadic-s3-natural-marked-cohomology}
H^0(\mathcal J_{\Gamma,\mathrm{nat}})=0,
\qquad
H^1(\mathcal J_{\Gamma,\mathrm{nat}})
\cong T,\qquad
H^2(\mathcal J_{\Gamma,\mathrm{nat}})=0.
\end{equation}
These are the genuine integral
$G_{\mathbf Q_2}$ cohomology groups
of $T$ in
Theorem~\ref{thm:dyadic-s3-integral-sector}.
\end{theorem}

\begin{proof}
For the coefficient module $T$,
a crossed derivation is a section
of the affine semidirect group
$T\rtimes_\rho S_3$, whose elements
are written $(A,v)$ with law
\[
(A,v)(B,w)=(AB,\ v+Aw).
\]
On the four marked generators
take $(S,a),(U,b),(I,c),(I,d)$.
The tame relation is
$r_t=\sigma^{-1}\tau\sigma\tau^{-2}$;
the affine calculation gives
\[
c(r_t)
=S^{-1}(U-I)a+(S^{-1}-I-U)b,
\]
exactly the first row block of
\eqref{eq:dyadic-s3-natural-differentials}.

The key cancellation is the exact
integral identity
\[
I+U+U^2=0.
\]
For every $v\in T$,
$(U,v)^3=(I,(I+U+U^2)v)=(I,0)$.
Thus every affine $2$-primary
projection $(U,v)^{\omega_2}$
is exactly the identity, while
the projection of $(S,a)$ is
$(S,a)$ because its image
is pro-$2$.
Consequently, in the literal
Roe--Turturean auxiliary words
\eqref{eq:dyadic-omega2-binomial-power}
and Theorem~\ref{thm:dyadic-q2-marked-word-jet-solver},
\[
\sigma_2=\sigma,\quad
u_0=u_1=1,\quad
d_0=x_0^{-1},\quad c_0=h_c=1.
\]
The wild generators act by pure
translations; those translations
commute, and $g_0=\sigma^2$
acts trivially on the translation
subgroup because its linear part
is $S^2=I$.
It follows exactly, not only
modulo $2$, that
\[
h_0=(I,-2c),\qquad
x_1^\sigma=(I,S^{-1}d),
\]
and hence
$c(r_{\rm wild})=-2c+S^{-1}d$.
This proves the full wild row.

The expression for $d^0$
is the principal crossed
derivation $(\rho(g)-I)v$.
The group relator identity implies
$d^1d^0=0$; it also follows
directly from
\eqref{eq:dyadic-s3-natural-tame-blocks}.
By \eqref{eq:dyadic-s3-natural-unit-dets}
the tame $L_\sigma$ and wild
$S^{-1}$ blocks give four
independent unit pivots in
$d^1$, hence $\operatorname{coker}d^1=0$.
The $\tau$ block $R_\tau$
of $d^0$ has two independent
unit pivots, so the image
of $d^0$ is a saturated rank-two
submodule of $\ker d^1$.
Thus $H^0=H^2=0$ and
$H^1$ is free of rank
$8-4-2=2$, as claimed.
The identification with intrinsic
cohomology is made explicit in
the next theorem rather than
assuming the marked relator
complex is a full universal
group-algebra resolution.
\end{proof}

\begin{theorem}[A strict integral deformation retract of the full marked complex]
\label{thm:dyadic-s3-natural-strict-retract}
Define maps of graded
$\mathbf Z_2$-modules between
$\mathcal J^\bullet_{\Gamma,\mathrm{nat}}(T)$
and the single finite module
$T[-1]$:
\begin{equation}\label{eq:dyadic-s3-natural-pi-maps}
\begin{aligned}
p^1(a,b,c,d)&=c,
 &p^0=p^2&=0,\\
i^1(t)&=(0,0,t,2St),
 &i^0=i^2&=0.
\end{aligned}
\end{equation}
Define degree-minus-one
maps $h^1:T^{\oplus4}\to T$ and
$h^2:T^{\oplus2}\to T^{\oplus4}$
by the \emph{closed integral}
matrices
\begin{equation}\label{eq:dyadic-s3-natural-explicit-homotopy}
\boxed{
\begin{aligned}
h^1(a,b,c,d)&=(U-I)^{-1}b,
\\
h^2(u,w)&=(L_\sigma^{-1}u,\ 0,\ 0,\ Sw).
\end{aligned}}
\end{equation}
Then $i,p$ are strict chain
maps and
\begin{equation}\label{eq:dyadic-s3-natural-contraction-identities}
\boxed{\qquad
p i=I_T,\qquad
i p=I_{\mathcal J}-(d h+h d),
\qquad
h^1d^0=I_T,\quad d^1h^2=I_{T^2}.
\qquad}
\end{equation}
In particular
\begin{equation}\label{eq:dyadic-s3-natural-complex-retract}
\mathcal J^\bullet_{\Gamma,\mathrm{nat}}(T)
\ \simeq_{\mathrm{chain},\,\mathbf Z_2}\
T[-1]
\end{equation}
is a \emph{strict, explicit
integral chain-homotopy equivalence},
not just a derived equivalence
from coincident cohomology ranks.

The class $t\in T$ corresponds
to the unique normalized continuous
$\Gamma$-cocycle $c_t$ satisfying
\begin{equation}\label{eq:dyadic-s3-natural-normalized-cocycle}
\boxed{
c_t(\sigma)=c_t(\tau)=0,\qquad
c_t(x_0)=t,\qquad
c_t(x_1)=2St.
}
\end{equation}
Every continuous $\Gamma$-cocycle
with coefficients in $T$ is
cohomologous to exactly one
$c_t$.  Thus evaluation at the
wild source generator $x_0$
gives an intrinsic-to-the-marked-source
isomorphism
\begin{equation}\label{eq:dyadic-s3-natural-cohom-coordinate}
H^1(\Gamma,T)\xrightarrow{\;\sim\;}T,
\qquad[c]\longmapsto c(x_0),
\end{equation}
which is integral and requires
no Fox-to-Herr parametrix.
\end{theorem}

\begin{proof}
The $x_0$ action on $T$ is
trivial, so $p^1d^0=0$.
Equation
\eqref{eq:dyadic-s3-natural-differentials}
gives $d^1i^1=0$, while
$p^1i^1=I_T$.
Hence $i,p$ are chain maps
and one-sided inverses.

The tame block identity
\[
L_\sigma R_\sigma+
L_\tau R_\tau=0
\]
follows from $d^1d^0=0$.
Both $L_\sigma$ and $R_\tau$
are invertible over $\mathbf Z_2$,
so
\[
R_\sigma R_\tau^{-1}
=-L_\sigma^{-1}L_\tau.
\]
Using this identity and
\eqref{eq:dyadic-s3-natural-differentials},
one checks for every
$(a,b,c,d)\in T^4$:
\[
\begin{aligned}
d^0h^1(a,b,c,d)&=
(R_\sigma R_\tau^{-1}b,\ b,\ 0,\ 0),\\
h^2d^1(a,b,c,d)&=
(a+L_\sigma^{-1}L_\tau b,\
 0,\ 0,\ d-2Sc).
\end{aligned}
\]
Their sum is exactly
$(a,b,c,d)-i p(a,b,c,d)$.
Likewise $h^1d^0=I_T$ and
$d^1h^2=I_{T^2}$ because
$L_\sigma^{-1}$, $(U-I)^{-1}$,
and $S$ have integral entries.
This proves all identities in
\eqref{eq:dyadic-s3-natural-contraction-identities}.

The generator values in
\eqref{eq:dyadic-s3-natural-normalized-cocycle}
satisfy both marked word
relations by the exact affine
calculation in
Theorem~\ref{thm:dyadic-s3-natural-full-marked-matrix}.
In each finite quotient
$T/2^nT\rtimes S_3$, the normal
closure of the images of the
wild generators lies in the
$2$-group $T/2^nT$;
the pro-$2$ marked condition is
therefore satisfied.  The
Roe--Turturean presentation
\cite{RoeTurtureanGQ2} supplies
a compatible family of
finite representations, whose
inverse limit defines the
continuous cocycle $c_t$.
Conversely, every continuous
cocycle determines a generator
tuple satisfying the linearized
marked relations, and the
contraction above subtracts
a unique principal cocycle
to reach the displayed
normalized tuple.
As $x_0$ acts trivially,
$c(x_0)$ is unchanged by a
principal coboundary.
This proves the uniqueness
and the actual cohomology
isomorphism.
\end{proof}

\begin{theorem}[Explicit index-three Sylow cocycle projector]
\label{thm:dyadic-s3-explicit-cocycle-projector}
Let
$t_i=\tau^i$ for $i=0,1,2$
be the right-coset representatives
of Theorem~\ref{thm:dyadic-s3-schreier-corestriction},
and write
\[
w_{ij}=\tau^i x_j\tau^{-i}\in H
\qquad(i=0,1,2;\ j=0,1).
\]
For \emph{every} continuous
$H$-cocycle $c\in Z^1(H,T)$
define the finite integral
Schreier-coordinate functional
\begin{equation}\label{eq:dyadic-s3-strict-phi-functional}
\boxed{\quad
\Phi_H(c)=\frac13
\sum_{i=0}^{2}U^{-i}c(w_{i0})\in T.
\quad}
\end{equation}
This descends to an
integral surjection on
$H^1(H,T)$ because each
$w_{i0}$ acts trivially
on the coefficient lattice.
The normalized global cocycles
of \eqref{eq:dyadic-s3-natural-normalized-cocycle}
satisfy
\begin{equation}\label{eq:dyadic-s3-strict-res-values}
\boxed{\quad
\begin{gathered}
(\operatorname{res}c_t)(\sigma)=0,\qquad
(\operatorname{res}c_t)(\tau^3)=0,\\
(\operatorname{res}c_t)(w_{i0})=U^i t,
\qquad
(\operatorname{res}c_t)(w_{i1})=2U^iS t,
\quad i=0,1,2.
\end{gathered}\quad}
\end{equation}
The restriction map is
injective in degree one,
and its left inverse is
$\Phi_H$:
\begin{equation}\label{eq:dyadic-s3-strict-res-phi}
\Phi_H(\operatorname{res}c_t)=t.
\end{equation}
More strongly, for every
$H$-\emph{cocycle}, not just
its cohomology class,
\begin{equation}\label{eq:dyadic-s3-strict-e-cocycle}
\boxed{\quad
\frac13\operatorname{res}
\operatorname{cor}(c)
=\operatorname{res}c_{\Phi_H(c)}
\qquad\text{as literal }H
\text{-cocycles}.\quad}
\end{equation}

In the eight displayed
Schreier generator coordinates
\[
(\sigma,\tau^3,\
w_{00},w_{10},w_{20},\
w_{01},w_{11},w_{21}),
\]
the following \emph{entirely
finite} block matrices describe
the degree-one projector:
\begin{equation}\label{eq:dyadic-s3-schreier-16x16-projector}
\boxed{
\begin{aligned}
\mathsf I_{\rm Sch}&=
\begin{pmatrix}
0\\0\\I\\U\\U^2\\2S\\2US\\2U^2S
\end{pmatrix}:T\longrightarrow T^8,\\[3pt]
\mathsf P_{\rm Sch}&=\frac13
\begin{pmatrix}
0&0&I&U^{-1}&U^{-2}&0&0&0
\end{pmatrix}:T^8\longrightarrow T,\\
\mathsf E_{\rm Sch}
&=\mathsf I_{\rm Sch}\mathsf P_{\rm Sch},
\qquad
\mathsf E_{\rm Sch}^2=\mathsf E_{\rm Sch}.
\end{aligned}}
\end{equation}
Every block is an integral
$2\times2$ matrix; the resulting
$\mathsf E_{\rm Sch}$ is a
strict rank-two $16\times16$
projection on the ambient
Schreier evaluation space.
On actual cocycles it is
\emph{exactly} the normalized
restriction--corestriction
operator.  On cohomology
\begin{equation}\label{eq:dyadic-s3-h1-primitive-decomposition}
\boxed{\quad
H^1(H,T)
=\operatorname{res}H^1(\Gamma,T)
\oplus\ker\Phi_H
\cong\mathbf Z_2^2
\oplus\mathbf Z_2^5.
\quad}
\end{equation}
The formula does not presume
that these eight displayed
Schreier words form a
\emph{minimal} pro-$2$
Demu\v{s}kin presentation.
\end{theorem}

\begin{proof}
Every $w_{ij}$ lies in
the wild normal subgroup
of the marked Galois
presentation and maps
to the identity under $\rho_T$.
Thus $\Phi_H$ kills
$H$-coboundaries.
The crossed-cocycle formula
and $\rho_T(\tau)=U$ give
\[
c_t(\tau^i x_j\tau^{-i})
=U^i c_t(x_j)
\]
because $c_t(\tau)=0$ and
$\rho_T(x_j)=I$.
This gives
\eqref{eq:dyadic-s3-strict-res-values}
and
\eqref{eq:dyadic-s3-strict-res-phi}.

For arbitrary $c\in Z^1(H,T)$,
the element $\tau^3\in H$
has pro-odd order and
acts trivially on $T$.
Continuity forces
$c(\tau^3)=0$.
The coset formula
\eqref{eq:dyadic-s3-transfer-one}
then gives
\[
(\operatorname{cor}c)(\sigma)
=(I+U^{-1}+U^{-2})c(\sigma)=0,
\quad
(\operatorname{cor}c)(\tau)=0,
\]
and
\[
(\operatorname{cor}c)(x_0)
=\sum_{i=0}^{2}U^{-i}c(w_{i0})
=3\Phi_H(c).
\]
By the exact normalized-cocycle
description in
Theorem~\ref{thm:dyadic-s3-natural-strict-retract},
a $\Gamma$-cocycle with
zero tame-generator values
is determined uniquely by
its value at $x_0$ and
is $c_t$ with that value.
Thus $\operatorname{cor}c
=c_{3\Phi_H(c)}$,
and dividing by the odd
unit $3$ gives the exact
cocycle identity
\eqref{eq:dyadic-s3-strict-e-cocycle}.

Matrix multiplication shows
\[
\mathsf P_{\rm Sch}
\mathsf I_{\rm Sch}
=\tfrac13(I+U^{-1}U+U^{-2}U^2)
=I,
\]
hence $\mathsf E_{\rm Sch}$
is a strict integral
idempotent of rank two.
For actual $H$-cocycles,
the matrix is the evaluation
of \eqref{eq:dyadic-s3-strict-e-cocycle}.
The rank-seven freeness of
$H^1(H,T)$ from
Theorem~\ref{thm:dyadic-s3-integral-sector}
now gives the rank-five kernel
and direct sum
\eqref{eq:dyadic-s3-h1-primitive-decomposition}.
\end{proof}

\begin{theorem}[A strict integral bar-cochain Hecke representative]
\label{thm:dyadic-s3-strict-bar-hecke-projector}
Let $C^\bullet(H,T)$ be the
normalized continuous $H$-bar
cochain complex.  Extend
\eqref{eq:dyadic-s3-strict-phi-functional}
to \emph{arbitrary} continuous
degree-one cochains
$f:H\to T$, and set
\[
p_H^1(f)=\frac13\sum_{i=0}^{2}
U^{-i}f(w_{i0}),\quad
p_H^k=0\ (k\ne1).
\]
Define the chain map
$i_H:T[-1]\to C^\bullet(H,T)$
by $i_H^1(t)=\operatorname{res}c_t$,
and $i_H^k=0$ for $k\ne1$.
Then $p_H:C^\bullet(H,T)\to T[-1]$
is a \emph{strict cochain map},
$p_Hi_H=I_T$, and
\begin{equation}\label{eq:dyadic-s3-strict-bar-projector}
\boxed{\quad
\Pi_H=i_Hp_H,\qquad
\Pi_H^2=\Pi_H,\qquad
\mathsf T_{\rm Hk}^{\rm str}=3\Pi_H-I.
\quad}
\end{equation}
Both the idempotent identity
and the Hecke polynomial hold
as \emph{literal identities of
integral cochain maps}:
\begin{equation}\label{eq:dyadic-s3-strict-bar-hecke-polynomial}
\boxed{\quad
(\mathsf T_{\rm Hk}^{\rm str})^2
=\mathsf T_{\rm Hk}^{\rm str}+2I.
\quad}
\end{equation}
The strict operator
$\Pi_H$ represents the
\emph{canonical derived}
Sylow idempotent:
\begin{equation}\label{eq:dyadic-s3-bar-derived-equality}
[\Pi_H]=
\left[
\tfrac13\operatorname{res}_H^\Gamma
\operatorname{cor}_H^\Gamma
\right]
\quad
\text{in }
\operatorname{End}_{D(\mathbf Z_2)}
R\Gamma(H,T).
\end{equation}

There is also an explicit
integral chain-level comparison
between the \emph{finite}
full-marked coefficient
complex and the \emph{image}
of the Sylow bar-cochain
projector.  With $i,p$ from
Theorem~\ref{thm:dyadic-s3-natural-strict-retract},
set
\begin{equation}\label{eq:dyadic-s3-marked-to-sylow-chain-maps}
\mathsf F=i_Hp:
\mathcal J^\bullet_{\Gamma,\mathrm{nat}}
\longrightarrow C^\bullet(H,T),
\qquad
\mathsf G=i p_H:
C^\bullet(H,T)
\longrightarrow\mathcal J^\bullet_{\Gamma,\mathrm{nat}}.
\end{equation}
These are \emph{strict integral}
cochain maps and satisfy
\begin{equation}\label{eq:dyadic-s3-marked-to-sylow-composites}
\boxed{\quad
\mathsf G\mathsf F=i p
\simeq I_{\mathcal J}
\ \text{by the explicit }h,
\qquad
\mathsf F\mathsf G=\Pi_H
\ \text{strictly}.
\quad}
\end{equation}
In particular the entire
degree-one Hecke/Sylow
summand is finite and
directly reconstructed from
the marked $x_0$-coordinate.
This closes the natural-$T$
\emph{integral descent chain}
without selecting a
five-generator Demu\v{s}kin
Fox basis for $P_H$.

The statement does \emph{not}
upgrade the actual infinite
bar complex to a globally
finite free group-algebra
resolution.  It identifies
a \emph{strict finite
cochain summand} of it.
\end{theorem}

\begin{proof}
For a $0$-cochain $v\in T$,
the continuous bar coboundary
has
\[
(d^0v)(w_{i0})
=(\rho_T(w_{i0})-I)v=0.
\]
Thus $p_H^1d^0=0$,
and all other cochain
map conditions for
$p_H:C^\bullet(H,T)\to T[-1]$
are automatic.
The chosen $i_H^1(t)$
is a continuous cocycle,
so $i_H$ is a chain map.
Equation
\eqref{eq:dyadic-s3-strict-res-phi}
gives $p_Hi_H=I_T$.
Consequently $\Pi_H$
is strictly idempotent.
Expanding
$(3\Pi_H-I)^2$
using $\Pi_H^2=\Pi_H$
gives
\eqref{eq:dyadic-s3-strict-bar-hecke-polynomial}.

To compare with the canonical
derived transfer, define
$p_\Gamma:C^\bullet(\Gamma,T)\to T[-1]$
by $p_\Gamma^1(f)=f(x_0)$
and zero in other degrees.
It is a cochain map because
$x_0$ acts trivially on $T$.
Let $i_\Gamma^1(t)=c_t$.
Theorem~\ref{thm:dyadic-s3-natural-strict-retract}
proves $p_\Gamma i_\Gamma=I$
and that $p_\Gamma$ and
$i_\Gamma$ induce inverse
isomorphisms on intrinsic
$G_{\mathbf Q_2}$ cohomology.
Hence their composites are
identities in $D(\mathbf Z_2)$.

The explicit degree-one
corestriction formula
\eqref{eq:dyadic-s3-transfer-one}
holds for arbitrary
continuous cochains as
well as cocycles, and
gives the \emph{strict}
identity of cochain maps
\[
p_H=\tfrac13 p_\Gamma
\,\operatorname{cor}_H^\Gamma,
\qquad
i_H=\operatorname{res}_H^\Gamma
\,i_\Gamma.
\]
Consequently
\[
[\Pi_H]=
\left[\operatorname{res}\,
i_\Gamma\,\tfrac13p_\Gamma
\operatorname{cor}\right]
=\left[\tfrac13\operatorname{res}
\operatorname{cor}\right]
\]
in the derived category,
proving
\eqref{eq:dyadic-s3-bar-derived-equality}.
Finally all maps in
\eqref{eq:dyadic-s3-marked-to-sylow-chain-maps}
are composites of strict
cochain maps, and
$p_Hi_H=I$, $pi=I$
imply
\[
\mathsf G\mathsf F
=i(p_Hi_H)p=ip,\qquad
\mathsf F\mathsf G
=i_H(pi)p_H=i_Hp_H.
\]
The homotopy in
\eqref{eq:dyadic-s3-natural-explicit-homotopy}
satisfies $ip\simeq I_{\mathcal J}$
over $\mathbf Z_2$,
giving the claimed chain-level
comparison.
\end{proof}

\begin{warning}[What is and is not native in fixed129]
\label{warn:dyadic-s3-natural-minimal-fox-boundary}
The natural-coefficient descent
projector is now strict on
the integral bar complex,
is an actual finite
$16\times16$ block matrix
on the displayed Schreier
generator-evaluation space,
and is a strict chain
summand equivalent to the
finite full-marked
coefficient Jacobian.
Its image is the same
derived object as the
all-Sylow Fox retract.
Nevertheless an
\emph{explicit identification
of the eight Schreier words
with a specific minimal
five-generator Demu\v{s}kin
basis}, including the
integer relation module,
its Fox syzygies and the
direct finite
Schreier-to-minimal-Fox
second-degree chain map,
has not been constructed.
Nor does the norm-zero
identity hold for the
adjoint $S_3$ module:
$\operatorname{ad}T$ has
nonzero tame $C_3$ invariants
and requires the separate
eight-by-sixteen Smith
certificate of fixed128.
The theorem is therefore
a sharp result for the
fixed natural coefficient
carrier, not an assertion
of an all-representation
full-group Fox resolution.
\end{warning}

% Fixed130 complete independent insertion.
% Analytic 2-integral splitting of ad(T) on the genuine tame S3
% sector; strict elementary chain reductions and Hecke top obstruction.

\subsection{Tame-idempotent splitting and an analytic integral
adjoint Fox--Smith certificate in the $S_3$ sector}
\label{subsec:dyadic-s3-adjoint-tame-splitting}

The coefficient-specific matrix computation in
Theorem~\ref{thm:dyadic-s3-marked-adjoint-jacobian}
admits a stronger structural proof.  For the fixed
integral $S_3$ representation $T=\mathbf Z_2^2$,
the adjoint module $\operatorname{End}(T)$
has a canonical order-three tame idempotent.
Its invariant and noninvariant summands are
both explicitly rank two, and the latter
is \emph{integrally isomorphic} to the natural
coefficient representation treated by
Theorem~\ref{thm:dyadic-s3-natural-strict-retract}.
The invariant summand is an induced
unramified-quadratic permutation lattice.
This splits the full $8\times16$ marked
Jacobian \emph{before} elementary reduction,
giving a symbolic proof of its unique
factor $2$ and an actual generator
of the top obstruction.

Write, as in \eqref{eq:dyadic-s3-integral-matrices},
\[
S=\begin{pmatrix}0&1\\1&0\end{pmatrix},
\qquad
U=\begin{pmatrix}-1&-1\\1&0\end{pmatrix},
\quad
M=\operatorname{End}_{\mathbf Z_2}(T).
\]
Let $\operatorname{Ad}_g(A)=gAg^{-1}$.
The order-three averaging operator is
\begin{equation}\label{eq:dyadic-s3-adjoint-tame-projector}
\boxed{\quad
\Pi_0=\tfrac13
\bigl(I+\operatorname{Ad}_U+
\operatorname{Ad}_{U^2}\bigr),\quad
M_0=\operatorname{im}\Pi_0,\quad
M_1=\ker\Pi_0.
\quad}
\end{equation}
The idempotent commutes with $\operatorname{Ad}_S$,
because $SUS^{-1}=U^2$; it also commutes
with the whole fixed $S_3$ action since
$\langle U\rangle\triangleleft S_3$.

\begin{theorem}[A unimodular tame splitting of $\operatorname{ad}T$]
\label{thm:dyadic-s3-adjoint-unimodular-splitting}
Define two $2$-adic integral matrices
\begin{equation}\label{eq:dyadic-s3-adjoint-vw-matrices}
 V=(I-\Pi_0)E_{12}
 =\frac13\begin{pmatrix}-1&1\\2&1\end{pmatrix},
 \qquad
 W=SVS^{-1}
 =\frac13\begin{pmatrix}1&2\\1&-1\end{pmatrix}.
\end{equation}
Then
\begin{equation}\label{eq:dyadic-s3-adjoint-basis}
\boxed{\quad
M_0=\mathbf Z_2 U\oplus\mathbf Z_2 U^2,
\qquad
M_1=\mathbf Z_2 V\oplus\mathbf Z_2 W,
\qquad
M=M_0\oplus M_1.
\quad}
\end{equation}
In the row-major basis
$(E_{11},E_{12},E_{21},E_{22})$,
the change-of-basis matrix with
columns $(U,U^2,V,W)$ is
\begin{equation}\label{eq:dyadic-s3-adjoint-unimodular-basis}
\mathsf B=
\begin{pmatrix}
-1&0&-1/3&1/3\\
-1&1&1/3&2/3\\
1&-1&2/3&1/3\\
0&-1&1/3&-1/3
\end{pmatrix},
\qquad
\boxed{\det\mathsf B=1\in\mathbf Z_2^\times}.
\end{equation}
The $S_3$ action is, in these
two ordered rank-two blocks,
\begin{equation}\label{eq:dyadic-s3-adjoint-s3-block-actions}
\begin{array}{c|cc}
& M_0:(U,U^2)&M_1:(V,W)\\ \hline
\operatorname{Ad}_U
&I_2&
\begin{pmatrix}-1&-1\\1&0\end{pmatrix}\\[6pt]
\operatorname{Ad}_S
&\begin{pmatrix}0&1\\1&0\end{pmatrix}&
\begin{pmatrix}0&1\\1&0\end{pmatrix}.
\end{array}
\end{equation}
Consequently, as \emph{actual
integral} $\mathbf Z_2[S_3]$-modules,
\begin{equation}\label{eq:dyadic-s3-adjoint-induced-plus-natural}
\boxed{\quad
\operatorname{ad}T
\cong M_0\oplus M_1
\cong
\operatorname{Ind}_{G_{K_2}}^\Gamma\mathbf Z_2
\oplus T,
\qquad K_2=\mathbf Q_2(\zeta_3).
\quad}
\end{equation}
Here the induction is from the
index-two normal subgroup $G_{K_2}$
and uses its two cosets as a
\emph{permutation lattice};
it is not induction from the
index-three subgroup $H$.
\end{theorem}

\begin{proof}
The operator $\Pi_0$ is the
standard Reynolds idempotent
for the normal order-three
subgroup; $3$ is invertible in
$\mathbf Z_2$.  Its invariant
part is the centralizer
$\mathbf Z_2[U]$, which is free
of rank two with basis $(U,U^2)$
because $U^2+U+I=0$.
Evaluating $(I-\Pi_0)E_{12}$
gives the displayed $V$ and $W$.
The row-major matrix of these
four vectors is
\eqref{eq:dyadic-s3-adjoint-unimodular-basis}.
Multiplying its last two
columns by $3$ gives an
integer matrix of determinant $9$;
hence $\det\mathsf B=9/9=1$.
Thus $(U,U^2,V,W)$ is a
$\mathbf Z_2$ basis, and the
projector proves the direct sum
\eqref{eq:dyadic-s3-adjoint-basis}.

Conjugation by $S$ swaps
$U$ and $U^2$, and swaps $V,W$
by definition.  Conjugation by
$U$ fixes its centralizer
and, by direct multiplication,
\[
 UVU^{-1}=-V+W,\qquad
 UWU^{-1}=-V.
\]
These are exactly the two
matrices in
\eqref{eq:dyadic-s3-adjoint-s3-block-actions}.
The $M_1$ matrices are precisely
the natural representation
\eqref{eq:dyadic-s3-integral-matrices};
the $M_0$ matrices describe the
permutation action on the two
cosets of the normal subgroup
$\langle U\rangle$.
The preimage of $\langle U\rangle$
in $\Gamma$ is $G_{K_2}$,
so this permutation lattice is
$\operatorname{Ind}_{G_{K_2}}^\Gamma
\mathbf Z_2$.
This proves the integral
module isomorphism.
\end{proof}

\begin{theorem}[An analytic proof of the fixed128 adjoint Smith factors]
\label{thm:dyadic-s3-adjoint-analytic-smith}
In the coefficient basis
\eqref{eq:dyadic-s3-adjoint-basis}
the complete full-marked
three-term adjoint \emph{coefficient}
complex of fixed128 is the direct
sum of a $M_0$ complex and the
norm-zero natural complex on $M_1$.
The latter is the strict
$M_1[-1]$ retract of
Theorem~\ref{thm:dyadic-s3-natural-strict-retract}
via the canonical identification
$M_1\cong T$.

On the tame-invariant rank-two block
$M_0$ choose $S_M=\begin{psmallmatrix}
0&1\\1&0\end{psmallmatrix}$ in
the $(U,U^2)$ basis.
Since $\tau$ acts trivially
and $x_0,x_1$ act trivially,
the full marked differentials are
\begin{equation}\label{eq:dyadic-s3-tame-invariant-differentials}
\boxed{
\begin{aligned}
d_0^0(v)
&=((S_M-I)v,\ 0,\ 0,\ 0),\\
d_0^1(a,b,c,d)
&=\bigl((S_M-2I)b,\
  3b+2c+(S_M-I)d\bigr).
\end{aligned}}
\end{equation}
The operator $L=S_M-2I$ has
determinant $3$ and is a
$\mathbf Z_2$-unit.
The strict invertible target
change
\begin{equation}\label{eq:dyadic-s3-tame-invariant-row-change}
(u,w)\longmapsto
(u,w-3L^{-1}u)
\end{equation}
changes the degree-one
differential to
\begin{equation}\label{eq:dyadic-s3-tame-invariant-split-d1}
d_0^1(a,b,c,d)
\longmapsto
\bigl(Lb,\ 2c+(S_M-I)d\bigr).
\end{equation}
It therefore splits off the
contractible tame block
$M_0\xrightarrow{L}M_0$
in degrees $1$ and $2$,
leaving the three-generator
integral Fox complex
\begin{equation}\label{eq:dyadic-s3-tame-invariant-fox-core}
\boxed{\quad
M_0\xrightarrow{(S_M-I,0,0)^t}M_0^3
\xrightarrow{(a,c,d)\mapsto
2c+(S_M-I)d}M_0.
\quad}
\end{equation}

Put
\[
e=(1,0)^t,\qquad
e_+=(1,1)^t,\qquad
f=(-1,1)^t=(S_M-I)e.
\]
Both $(e_+,e)$ and $(f,e)$
are unimodular bases of $M_0$.
The following elementary
unit operations are sufficient
to reduce
\eqref{eq:dyadic-s3-tame-invariant-fox-core}:
\begin{enumerate}[label=\textnormal{(\roman*)}]
\item $e\mapsto f$ in $d^0$
splits off a degree-zero--one
unit pair, leaving $e_+$
as a generator of $H^0$;
\item the input $d=e$ maps
to $f$ in degree two,
giving one degree-one--two
unit pair;
\item $c=e$ maps to $2e$
modulo the image of $f$,
giving one surviving
two-primary elementary
pair $\mathbf Z_2\xrightarrow{2}\mathbf Z_2$;
\item the remaining three
degree-one coordinates
are freely generated
by a $\sigma$-value $e$,
by $(c,d)=(f,-2e)$,
and by $(c,d)=(0,e_+)$.
\end{enumerate}
For a literal diagonal normal form,
write
\begin{equation}\label{eq:dyadic-s3-adjoint-exact-fox-input-basis}
\begin{aligned}
v&=r e_++s e,&
a&=\alpha e+\beta f,\\
c&=c_e e+c_f f,&
d&=d_e e+d_+e_+,
\qquad \widetilde d_e=d_e+2c_f.
\end{aligned}
\end{equation}
All displayed basis changes are
unimodular over $\mathbf Z_2$.
In the degree-one coordinate order
$(\alpha,\beta,c_f,c_e,
\widetilde d_e,d_+)$ and
degree-two target order $(f,e)$,
the reduced Fox differentials
are \emph{exactly}
\begin{equation}\label{eq:dyadic-s3-adjoint-exact-fox-diagonal}
\boxed{
\begin{aligned}
d^0(r,s)&=(0,s,0,0,0,0),\\
d^1(\alpha,\beta,c_f,c_e,
   \widetilde d_e,d_+)
&=(\widetilde d_e,2c_e).
\end{aligned}}
\end{equation}
Thus this is a concrete
$\mathbf Z_2$ basis-level
chain diagonalization,
not merely a Smith-invariant
existence statement.

Thus the tame-invariant block has
the strict finite homotopy type
\begin{equation}\label{eq:dyadic-s3-tame-invariant-elementary}
\boxed{\quad
\mathcal J^\bullet(M_0)
\simeq_{\mathbf Z_2}
\mathbf Z_2[0]\oplus
\mathbf Z_2^3[-1]\oplus
[\mathbf Z_2\xrightarrow{2}
 \mathbf Z_2]_{[1,2]}.
\quad}
\end{equation}
Together with
$\mathcal J^\bullet(M_1)\simeq M_1[-1]$,
the full adjoint marked complex has
the \emph{analytically derived}
integral elementary form
\begin{equation}\label{eq:dyadic-s3-adjoint-elementary-from-tame}
\boxed{\quad
\mathcal J^\bullet_\Gamma(M)
\simeq_{\mathbf Z_2}
\mathbf Z_2[0]\oplus
\mathbf Z_2^5[-1]\oplus
[\mathbf Z_2\xrightarrow{2}
 \mathbf Z_2]_{[1,2]}.
\quad}
\end{equation}
In particular its degree-one
Jacobian has \emph{seven unit
Smith factors and exactly one
factor of $2$}, recovering
\eqref{eq:dyadic-s3-jacobian-smith}
without extracting a numerical
mod-$4$ minor.  Its degree-zero
Jacobian has three unit factors,
and its cohomology is precisely
\[
(H^0,H^1,H^2)(\Gamma,M)=
(\mathbf Z_2,\mathbf Z_2^5,\mathbf F_2).
\]
\end{theorem}

\begin{proof}
All marked differential matrices
are linear in the coefficient
representation, so the
$\mathbf Z_2[S_3]$ splitting
of Theorem~\ref{thm:dyadic-s3-adjoint-unimodular-splitting}
gives a strict splitting of
their cochain complex.
On $M_1$, the exact identity
$I+\operatorname{Ad}_U+
\operatorname{Ad}_{U^2}=0$
gives the natural norm-zero
proof of fixed129, with its
explicit strong contraction.

On $M_0$, the affine coefficient
action of $\tau$ is trivial.
For arbitrary crossed-derivation
values $(a,b,c,d)$,
the tame relation gives
$(S_M-2I)b$.  In the wild
relation the pro-$2$ projections
of $x_i\tau$ act as the identity
on each translation, so
$u_i=x_i\tau$,
$d_0=d_g=\tau$, and
$c_0=h_c=1$.  Since
$g_0=\sigma^2$ has trivial
linear part on $M_0$,
conjugation of a translation
by $g_0$ is trivial.
Thus
$h_0=2c+4b$,
while $u_1^{-1}x_1^\sigma$
adds $-d-b+S_M^{-1}d$.
As $S_M^{-1}=S_M$,
the wild differential is
$3b+2c+(S_M-I)d$,
proving
\eqref{eq:dyadic-s3-tame-invariant-differentials}.
The row change
\eqref{eq:dyadic-s3-tame-invariant-row-change}
then cancels the tame unit.
The remaining three-generator
complex is literally the
pro-$2$ marked Fox cochain
complex of
Theorem~\ref{thm:dyadic-strict-nielsen-fox-chain}
in these coefficients.

For the elementary reduction,
$(S_M-I)e=f$ and
$(S_M-I)e_+=0$.
The image of $d^0$ is the
primitive rank-one line
$\mathbf Z_2 f$ in the
$\sigma$ coordinate, so
one integral unit pair splits.
In degree one the $d=e$
coordinate maps to the
primitive target vector $f$.
Eliminating that row changes
the image of $c=e$ to
$2e$ in the complementary
free target quotient.
The other input coordinates
give exactly the three free
cocycles listed in (iv),
as can be checked by
$2c+(S_M-I)d=0$.
All changes use unimodular
$\mathbf Z_2$ bases; this
proves
\eqref{eq:dyadic-s3-tame-invariant-elementary}.
The direct sum with $M_1[-1]$
then gives
\eqref{eq:dyadic-s3-adjoint-elementary-from-tame}.
Counting its cancelled
degree-one--two blocks gives
two tame units plus one
Fox unit from $M_0$, and
four unit degree-one--two
blocks from $M_1$, hence
seven units and the
remaining factor $2$.
The degree-zero--one units
are one from $M_0$ and
two from $M_1$, hence three.
\end{proof}

\begin{corollary}[Complete marked adjoint cocycle and obstruction coordinates]
\label{cor:dyadic-s3-adjoint-explicit-h1-h2}
In the marked $S_3$ coefficient basis
$(U,U^2,V,W)$, every class in
$H^1(\Gamma,\operatorname{ad}T)$ has
a unique representative modulo
coboundaries of the form
\begin{equation}\label{eq:dyadic-s3-adjoint-normalized-cocycles}
\boxed{
\begin{aligned}
c(\sigma)&=\lambda e\in M_0,\\
c(\tau)&=0,\\
c(x_0)&=\mu f+t,\\
c(x_1)&=-2\mu e+\nu e_++2S_M t,
\end{aligned}
\qquad
\lambda,\mu,\nu\in\mathbf Z_2,\quad
t\in M_1\cong\mathbf Z_2^2.
}
\end{equation}
Here $e,e_+,f$ are the vectors
in the $(U,U^2)$ basis from
Theorem~\ref{thm:dyadic-s3-adjoint-analytic-smith},
while $S_M$ acts by the same swap
matrix on either block.
The parameters
$(\lambda,\mu,\nu,t)$ form
an exact integral
$\mathbf Z_2^5$ basis
of the global adjoint
$H^1$.

The top adjoint cohomology
has exactly one nonzero
$\mathbf F_2$ class:
in the reduced $M_0$
Fox target it is represented
by $e$ modulo
$2M_0+(S_M-I)M_0$.
The ordinary matrix trace
maps this class to
$\operatorname{tr}(U)=-1$
modulo two, hence
gives its canonical
nonzero scalar detector.
The trace annihilates
$M_1$ and therefore
locates the top
obstruction \emph{entirely}
in the tame-invariant
summand $M_0$.
\end{corollary}

\begin{proof}
The unit tame row forces
$c(\tau)=0$ for an actual
cocycle in either summand.
On $M_1$, the strong
normalization of
Theorem~\ref{thm:dyadic-s3-natural-strict-retract}
gives $(c(\sigma),c(x_0),c(x_1))
=(0,t,2S_Mt)$.
On $M_0$, the reduced Fox
equation $2c+(S_M-I)d=0$
forces $c=\mu f$ and
$d=-2\mu e+\nu e_+$.
The $\sigma$ coordinate is
defined modulo the primitive
coboundary line $\mathbf Z_2 f$,
and so has the unique
representative $\lambda e$.
These descriptions produce
the displayed five parameters.
The top Fox quotient is
\[
M_0/(2M_0+(S_M-I)M_0)
\cong\mathbf F_2,
\]
generated by $e$.
On $M_0=\mathbf Z_2[U]$,
matrix trace sends $U$ and
$U^2$ to $-1$, hence
annihilates $(S_M-I)M_0$,
and also annihilates
$2M_0$ modulo two.
Its value on $e=U$
is nonzero.  Since the
matrices $V,W$ are traceless,
the detector vanishes on
$M_1$ and proves the
final assertion.
\end{proof}

\begin{theorem}[Integral Shapiro decomposition and top-degree Hecke selection]
\label{thm:dyadic-s3-adjoint-shapiro-hecke}
Put $K_2=\mathbf Q_2(\zeta_3)$,
$F=\mathbf Q_2(\sqrt[3]2)$,
$L=FK_2$, and let
$N=G_L\subset H=G_F\subset\Gamma$.
For the fixed $S_3$ action on
$M=\operatorname{ad}T$,
the module decomposition of
Theorem~\ref{thm:dyadic-s3-adjoint-unimodular-splitting}
gives the following
\emph{integral} Shapiro and
restriction identifications:
\begin{equation}\label{eq:dyadic-s3-adjoint-shapiro-global-h}
\boxed{
\begin{aligned}
R\Gamma(\Gamma,M_0)&\simeq
R\Gamma(G_{K_2},\mathbf Z_2),\\
R\Gamma(\Gamma,M_1)&\simeq
R\Gamma(\Gamma,T),\\
R\Gamma(H,M_0)&\simeq
R\Gamma(N,\mathbf Z_2),\\
R\Gamma(H,M_1)&\simeq
R\Gamma(N,\mathbf Z_2).
\end{aligned}}
\end{equation}
The cohomology ledgers are
\begin{equation}\label{eq:dyadic-s3-adjoint-shapiro-ledgers}
\begin{array}{c|ccc}
& H^0&H^1&H^2\\ \hline
\Gamma,M_0&\mathbf Z_2&\mathbf Z_2^3&\mathbf F_2\\
\Gamma,M_1&0&\mathbf Z_2^2&0\\
H,M_0&\mathbf Z_2&\mathbf Z_2^7&\mathbf F_2\\
H,M_1&\mathbf Z_2&\mathbf Z_2^7&\mathbf F_2 .
\end{array}
\end{equation}
The index-three Sylow idempotent
$e_H=\tfrac13\operatorname{res}\operatorname{cor}$
preserves $M_0\oplus M_1$ and
acts on $H^2(H,M)$ as
\begin{equation}\label{eq:dyadic-s3-adjoint-top-hecke-projection}
\boxed{\quad
H^2(H,M)=\mathbf F_2\oplus\mathbf F_2,
\qquad
e_H^{(2)}=
\begin{pmatrix}1&0\\0&0\end{pmatrix}
\quad}
\end{equation}
with the first coordinate
the tame-invariant component.
Likewise its degree-one
image/kernel decomposition is
\[
\begin{array}{c|cc}
&\operatorname{rank}\operatorname{im}e_H&
\operatorname{rank}\ker e_H\\ \hline
H^1(H,M_0)&3&4\\
H^1(H,M_1)&2&5.
\end{array}
\]
Thus the one surviving global
two-primary obstruction in
fixed128 is the
\emph{odd-index restriction}
of the trivial-coefficient
obstruction over
$K_2=\mathbf Q_2(\zeta_3)$,
while the entire top class
from the natural $M_1$
summand is killed by
the Hecke/Sylow projector.
\end{theorem}

\begin{proof}
The first two Shapiro identities
follow directly from
\eqref{eq:dyadic-s3-adjoint-induced-plus-natural}.
On restriction to $H$,
the generator $\sigma$ acts
on both $M_0$ and $M_1$
by the swap matrix, while
the wild generators act
trivially.  The index-two
kernel of this action
inside $H$ is $N=G_L$.
Thus each restricted module
is the permutation-induced
lattice $\operatorname{Ind}_N^H
\mathbf Z_2$, proving the
remaining two Shapiro identities.
The degree-one
and degree-two cohomology
of a trivial integral
$\mathbf Z_2$ lattice
over a dyadic local
absolute Galois group
is given by local class
field theory and Tate
duality:
\[
H^0(G_D,\mathbf Z_2)=\mathbf Z_2,
\quad H^1(G_D,\mathbf Z_2)
=\mathbf Z_2^{[D:\mathbf Q_2]+1},
\quad H^2(G_D,\mathbf Z_2)=\mathbf F_2
\]
for $D=K_2,L$, since
both contain exactly $\mu_2$
among the $2$-power roots
of unity.  Combined with
the natural-$T$ calculation
of fixed128, these give
\eqref{eq:dyadic-s3-adjoint-shapiro-ledgers}.

For $M_0$ the ordinary Mackey
identification
\[
\operatorname{Res}_H^\Gamma
\operatorname{Ind}_{G_{K_2}}^\Gamma
\mathbf Z_2
\simeq
\operatorname{Ind}_N^H\mathbf Z_2
\]
holds because $H G_{K_2}=\Gamma$
and $H\cap G_{K_2}=N$.
Under Shapiro, restriction
$\Gamma\to H$ becomes
the usual restriction
$G_{K_2}\to N=G_L$.
Since $[L:K_2]=3$ and
$3$ is a unit modulo $2$,
this restriction is an
isomorphism on the
one-dimensional
$H^2(-,\mathbf Z_2)$:
its local Tate-dual
effect is multiplication
by the odd degree three
on the $\mu_2$ invariant.
Thus the idempotent
acts as identity on
the first $H^2$ component.
For $M_1$, global
$H^2(\Gamma,M_1)=0$ by
the norm-zero calculation,
so the same idempotent
annihilates the second
$H^2$ component.
This gives
\eqref{eq:dyadic-s3-adjoint-top-hecke-projection}.
The degree-one ranks follow
from the Shapiro ledger
and the usual
$\operatorname{cor}\operatorname{res}=3$
splitting, separately on
each summand.
\end{proof}

\begin{warning}[Scope of the analytic tame decomposition]
\label{warn:dyadic-s3-adjoint-tame-splitting-scope}
The tame projector $\Pi_0$
splits the \emph{fixed}
$S_3$ adjoint coefficient
lattice.  The stronger
family-level question is
resolved in
Theorems~\ref{thm:dyadic-f131-canonical-tame-extraction}
and~\ref{thm:dyadic-f131-universal-split-locus}:
the tame relation rigidifies
$\rho(\tau)$ to exact order
three, and the polynomial
Reynolds idempotent extends
to every framed deformation;
however, its image is
Galois-stable exactly when
the wild matrices commute
with the tame order-three
transport.  The explicit
dual-number counterexample
of Theorem~\ref{thm:dyadic-f131-first-order-nosplit}
shows that no Galois-equivariant
lift of the fixed projector
exists on the full
unrestricted deformation
ring.

The analytic computation
upgrades fixed128 from
a numerical Smith
certificate to explicit
integral cochain
elementary divisors
and a top-degree Hecke
selection \emph{for this
specific adjoint lattice}.
It does not select
canonical minimal generators
of the rank-five
Demu\v{s}kin group
$G_F(2)$, or give
the universal Schreier
relation-syzygy chain
isomorphism required
for arbitrary non-pro-$2$
residual families.
\end{warning}

% Fixed131: family-level tame C3 projector, its exact wild-descent ideal,
% a first-order nonextendability witness, and curved two-block transport.
% This subsection is integrated verbatim into the full fixed131 manuscript.

\subsection{The universal tame projector and its wild descent obstruction}
\label{subsec:dyadic-s3-universal-tame-projector}

Version fixed130 proves a canonical decomposition for the
\emph{fixed} tame $S_3$ adjoint lattice.  It does not imply
that the resulting projector is a map of Galois
representations after arbitrary framed deformation:
the wild generators need not preserve its image.
We now separate two questions:
\emph{(a)} whether a canonical $C_3$ projector exists
integrally over the full completed residual ring, and
\emph{(b)} exactly where that projector descends to
an invariant coefficient decomposition.
Both questions admit finite residual-jet formulas.
The answer to (a) is positive; the answer to (b)
is an explicit closed proper locus.

Let
\[
\bar\rho:\Gamma=G_{\mathbf Q_2}
\twoheadrightarrow S_3
\simeq\mathrm{GL}_2(\mathbf F_2)
\]
be the residual representation of
Theorem~\ref{thm:dyadic-s3-integral-sector}.
Let $(A,\mathfrak m,k)$ be a complete
Noetherian local $\mathcal O_E$-algebra,
with a framed lift
$\rho_A:\Gamma\to\mathrm{GL}_2(A)$
whose reduction is $\bar\rho$.
Write
\[
X=\rho_A(\tau),\quad
Y=\rho_A(\sigma),\quad
X_j=\rho_A(x_j)\ (j=0,1),
\qquad
\bar X=\bar U,\quad
\bar Y=\bar S,\quad
\bar X_j=I_2.
\]
The marked tame relation remains
\begin{equation}\label{eq:dyadic-f131-tame-relation}
Y^{-1}XY=X^2.
\end{equation}

\begin{theorem}[Canonical integral order-three extraction in every residual family]
\label{thm:dyadic-f131-canonical-tame-extraction}
Put $Q=X^3-I_2\in M_2(\mathfrak m)$.
The binomial series
\begin{equation}\label{eq:dyadic-f131-tame-root}
u_A:=(X^3)^{1/3}
=\sum_{j\ge0}\binom{1/3}{j}Q^j
\in I_2+M_2(\mathfrak m)
\end{equation}
converges $\mathfrak m$-adically
and is integral over $A$.
The matrix
\begin{equation}\label{eq:dyadic-f131-odd-tame-matrix}
\boxed{\qquad
z_A=Xu_A^{-1}
\qquad}
\end{equation}
is the unique order-three
component of $X$ in its
procyclic closure.  Moreover the
\emph{exact marked tame relation}
forces $u_A=I_2$ and $z_A=X$.
In particular the full universal
rank-two tame matrix already
has exact order three; it satisfies
\begin{equation}\label{eq:dyadic-f131-cubic-etale}
\boxed{\qquad
z_A^3=I_2,\qquad
z_A^2+z_A+I_2=0,\qquad
z_A\bmod\mathfrak m=\bar U.
\qquad}
\end{equation}
The algebra
\begin{equation}\label{eq:dyadic-f131-canonical-tame-etale-algebra}
\mathcal C_A:=A[z_A]\subseteq M_2(A)
\simeq A[T]/(T^2+T+1)
\end{equation}
is a finite étale $A$-algebra
of rank two and is exactly the
centralizer of $z_A$ in $M_2(A)$.

On $M_A=\operatorname{End}_A(A^2)$
define the intrinsic tame Reynolds operator
\begin{equation}\label{eq:dyadic-f131-universal-tame-projector}
\boxed{\qquad
\Pi_A=\frac13
\left(I+\operatorname{Ad}_{z_A}
+\operatorname{Ad}_{z_A^2}\right).
\qquad}
\end{equation}
It is a strict integral idempotent,
with finite projective direct summands
\begin{equation}\label{eq:dyadic-f131-universal-tame-summands}
M_A=M_{0,A}\oplus M_{1,A},
\qquad M_{0,A}=\mathcal C_A=\operatorname{im}\Pi_A,
\qquad M_{1,A}=\ker\Pi_A,
\end{equation}
both of rank two.
This construction commutes with
continuous local coefficient base
change.  On the genuine
marked residual deformation ring
$z_A=X$, so the tame
projector is already a
\emph{finite polynomial} in
$X$ and its adjoint action,
with denominator the unit $3$
only.  The binomial formula
\eqref{eq:dyadic-f131-tame-root}
remains a compatible calculation
of the pro-$2$ part, which is
identically $I_2$ here.
There is no tame power-series
solve, even before passing
to finite precision.

Furthermore the full tame relation
\eqref{eq:dyadic-f131-tame-relation}
forces
\begin{equation}\label{eq:dyadic-f131-sigma-tame-normalization}
\boxed{\qquad
Y^{-1}z_A Y=z_A^2,
\quad
[\operatorname{Ad}_Y,\Pi_A]=0,
\quad
[\operatorname{Ad}_X,\Pi_A]=0.
\qquad}
\end{equation}
Thus the potential failure of
Galois-equivariant splitting comes
\emph{only} from wild transport.
\end{theorem}

\begin{proof}
The matrix $Q=X^3-I_2$ lies in
$M_2(\mathfrak m)$ because the
residual tame matrix has order
three.  The coefficients
$\binom{1/3}{j}\in\mathbf Z_2$
are integral and the $j$th
summand belongs to
$M_2(\mathfrak m^j)$.
The usual binomial identity,
valid for powers of the single
commuting matrix $Q$, proves
$u_A^3=X^3$.
Since $u_A$ is a power series
in $X$, it commutes with $X$,
so $z_A^3=I_2$.
Its reduction is $\bar U$.
The matrix $z_A-I_2$ is invertible:
its reduction $\bar U-I_2$
has determinant one.
Multiplying
$(z_A-I_2)(z_A^2+z_A+I_2)=0$
by its inverse proves the
quadratic polynomial in
\eqref{eq:dyadic-f131-cubic-etale}.

The polynomial $T^2+T+1$
is separable in residue
characteristic two; its
derivative $2T+1$ is a unit.
The map
$A[T]/(T^2+T+1)\to M_2(A)$
sending $T$ to $z_A$
is injective, since
$I_2,z_A$ are linearly
independent after reduction.
It is a finite étale rank-two
algebra.
More concretely, for the
first framed basis vector $e_1$,
the pair $(e_1,z_A e_1)$
is an $A$-basis because its
reduction is a basis for
$\bar U$.  In this cyclic basis,
every matrix commuting with $z_A$
is a unique $A$-linear
polynomial of degree at most
one in $z_A$.
This proves the centralizer
claim and its rank.

The conjugation
$\operatorname{Ad}_{z_A}$
has order three, so its
Reynolds average is an
idempotent over $A$.
Its image is precisely the
centralizer $\mathcal C_A$.
Its reduction is the rank-two
idempotent of
Theorem~\ref{thm:dyadic-s3-adjoint-unimodular-splitting};
the complementary summand is
therefore finite free of rank
two by Nakayama and local
projectivity.
All constructions use integral
convergent power series and
unit matrices and hence commute
with completed local base
change.  Terms of order
$j\ge t$ vanish modulo
$\mathfrak m^t$.

Finally the exact tame relation
also proves the stronger
\emph{rigidity} claim.
Let $d=\det X\in A^\times$.
Conjugating and taking
determinants in
\eqref{eq:dyadic-f131-tame-relation}
gives $d=d^2$, hence $d=1$.
Put $t=\operatorname{tr}X$.
The two sides of the tame
relation have equal trace,
so Cayley--Hamilton yields
\[
t=\operatorname{tr}(X^2)=t^2-2d=t^2-2,
\qquad (t-2)(t+1)=0.
\]
Since $\bar t=\operatorname{tr}(\bar U)=1$
in residue characteristic two,
$t-2$ is a unit.
Therefore $t=-1$.
Cayley--Hamilton now gives
$X^2+X+I_2=0$ and $X^3=I_2$.
Consequently the unique
pro-$2$ cubic root of $X^3$
in \eqref{eq:dyadic-f131-tame-root}
is $u_A=I_2$, and
$z_A=X$ identically.
In particular the claimed
tame projector is \emph{finite
polynomial} on the full
marked deformation ring.
The tame relation directly gives
$Y^{-1}z_A Y=z_A^2$.
Conjugation by $Y$ therefore
normalizes the cyclic
order-three adjoint action,
and commutes with its
Reynolds projector.
The matrix $X$ commutes with
its own odd component $z_A$,
so its adjoint action does as
well.  This proves all
assertions.
\end{proof}

\begin{theorem}[Exact wild equations for the universal split locus]
\label{thm:dyadic-f131-universal-split-locus}
Let $R_{\bar\rho}^{\square}$ be
the framed complete local
Galois deformation ring
and form $z^\square,\Pi^\square$
using Theorem~\ref{thm:dyadic-f131-canonical-tame-extraction}.
Let $\mathfrak I_{\rm split}$
be the closed ideal of
$R_{\bar\rho}^{\square}$
generated by the matrix entries of
\begin{equation}\label{eq:dyadic-f131-wild-split-equations}
\boxed{\qquad
[X_0,z^\square],\qquad
[X_1,z^\square].
\qquad}
\end{equation}
Equivalently, by
\eqref{eq:dyadic-f131-universal-tame-summands},
the ideal is generated by
the two rank-two block
components
\begin{equation}\label{eq:dyadic-f131-wild-split-projector-equations}
\boxed{\qquad
(1-\Pi^\square)X_0,\qquad
(1-\Pi^\square)X_1,
\qquad}
\end{equation}
where $\Pi^\square$ acts
on the \emph{matrix} algebra.
On any framed deformation
over a complete local $A$,
the following are equivalent:
\begin{enumerate}[label=\textnormal{(\roman*)}]
\item $\operatorname{Ad}_{\rho_A(g)}$
commutes with $\Pi_A$
for every $g\in\Gamma$;
\item both blocks $M_{0,A}$ and
$M_{1,A}$ of
\eqref{eq:dyadic-f131-universal-tame-summands}
are $\Gamma$-subrepresentations;
\item $X_0,X_1$ commute with $z_A$;
\item every marked generator
normalizes the finite étale
algebra $\mathcal C_A$,
hence
\[
\rho_A(\Gamma)\subseteq
N_{\mathrm{GL}_2(A)}(\mathcal C_A).
\]
\end{enumerate}
Therefore
\begin{equation}\label{eq:dyadic-f131-split-ring}
R_{\rm split}:=
R_{\bar\rho}^{\square}/\mathfrak I_{\rm split}
\end{equation}
represents the \emph{exact}
tame-isotypic split
subfunctor of the framed
residual deformation functor.

Over $R_{\rm split}$ there
is an integral Galois-equivariant
coefficient decomposition,
stable under all derived
base changes within the
split subfunctor:
\begin{equation}\label{eq:dyadic-f131-split-cochains}
R\Gamma(\Gamma,M_{R_{\rm split}})
\simeq
R\Gamma(\Gamma,M_{0,R_{\rm split}})
\oplus
R\Gamma(\Gamma,M_{1,R_{\rm split}}).
\end{equation}
The same splitting applies
to the all-Sylow finite
perfect coefficient carrier.
This does \emph{not} assert
the fixed130 constant Smith
normal form on the entire
split family.

The étale algebra
$\mathcal C_A$ makes $A^2$
a free rank-one
$\mathcal C_A$-module.
On the split locus the
index-two Galois subgroup
$G_{\mathbf Q_2(\zeta_3)}$
acts $\mathcal C_A$-linearly
and the other coset
acts through the
nontrivial involution
$z_A\mapsto z_A^2$.
Thus the split locus is
precisely the associated
rank-one étale-toric/
semilinear normalizer
realization, not an
arbitrary extension of
the fixed $S_3$ matrix frame.
\end{theorem}

\begin{proof}
The equivalence of (i)
and (ii) is idempotent
linear algebra.
By
\eqref{eq:dyadic-f131-sigma-tame-normalization},
$\sigma$ and $\tau$
already commute with
$\Pi_A$ on the adjoint
module.
Suppose that an individual
wild matrix $X_j$, whose
residue is $I_2$, satisfies
$[\operatorname{Ad}_{X_j},\Pi_A]=0$.
Then $\operatorname{Ad}_{X_j}$
preserves
$\mathcal C_A=\operatorname{im}\Pi_A$,
so $X_j z_A X_j^{-1}$
lies in $\mathcal C_A$,
is a cubic root of one,
and reduces to $\bar U$.
As $3z_A^2$ is a unit
in the commutative
$\mathfrak m$-complete
algebra $\mathcal C_A$,
the root lifting the given
residue is unique.
Consequently
$X_jz_AX_j^{-1}=z_A$.
Conversely, if $X_j$
commutes with $z_A$,
its adjoint action
commutes with the
Reynolds projector.
This proves (i)--(iii).

Every marked generator
normalizes $\mathcal C_A$
if (iii) holds, because
$\sigma$ acts on it by
$z_A\mapsto z_A^2$,
while $\tau,x_0,x_1$
centralize it.
Conversely, if a wild
$X_j$ normalizes
$\mathcal C_A$, the
same unique-root argument
shows it centralizes $z_A$.
This proves (iv) and the
functorial vanishing
criterion.
Since all equations are
entries of matrices over
the complete universal
deformation ring,
they define the stated
closed represented subfunctor.
The equivalent projected
form follows from
$\ker\operatorname{ad}(z_A)
=\mathcal C_A=\operatorname{im}\Pi_A$.

Over that quotient the two
idempotent summands are
preserved by $\Gamma$.
Ordinary continuous cochains
and the perfect all-Sylow
carrier are additive in
the coefficient module,
giving the displayed
derived splitting.
Its base-change
compatibility follows
because the idempotent and
all actions are defined
already over $R_{\rm split}$.

The cyclic basis
$(e_1,z_Ae_1)$ makes
$A^2$ a free rank-one
module over $\mathcal C_A$.
The normalizer acts either
linearly over $\mathcal C_A$
or semilinearly under its
unique nontrivial étale
involution.  The tame
generator $\sigma$ acts
semilinearly, while the
even-parity index-two
subgroup is generated
inside the normalizer
by $\tau$, the wild
generators and
$\sigma^2$ and therefore
acts linearly.
This proves the final
description without
claiming a global
minimal Fox basis.
\end{proof}

\begin{theorem}[A genuine wild first-order obstruction and failure of universal splitting]
\label{thm:dyadic-f131-first-order-nosplit}
Let
$k=\mathbf F_2$ and
$A_\epsilon=k[\epsilon]/(\epsilon^2)$.
Reduce the fixed matrices
$S,U,V$ of
Theorem~\ref{thm:dyadic-s3-adjoint-unimodular-splitting}
modulo two:
\begin{equation}\label{eq:dyadic-f131-wild-counterexample-matrices}
\bar S=
\begin{pmatrix}0&1\\1&0\end{pmatrix},
\quad
\bar U=
\begin{pmatrix}1&1\\1&0\end{pmatrix},
\quad
\bar V=
\begin{pmatrix}1&1\\0&1\end{pmatrix}.
\end{equation}
The marked generator assignment
\begin{equation}\label{eq:dyadic-f131-wild-counterexample-rep}
\boxed{\quad
\rho_\epsilon(\sigma)=\bar S,\quad
\rho_\epsilon(\tau)=\bar U,\quad
\rho_\epsilon(x_0)=I_2+\epsilon\bar V,\quad
\rho_\epsilon(x_1)=I_2
\quad}
\end{equation}
satisfies the literal
Roe--Turturean tame and
wild relations, including
the marked pro-$2$ wild
normal-closure condition.
It therefore defines a
genuine continuous
$G_{\mathbf Q_2}$ residual
deformation.

Its canonical odd tame
component is
$z_\epsilon=\bar U$.
However
\begin{equation}\label{eq:dyadic-f131-nonzero-wild-commutator}
\boxed{\quad
[\bar V,\bar U]=\bar V\ne0,\qquad
[\rho_\epsilon(x_0),z_\epsilon]
=\epsilon\bar V\ne0.
\quad}
\end{equation}
Thus $\rho_\epsilon$ lies
outside the split locus of
Theorem~\ref{thm:dyadic-f131-universal-split-locus}.

Moreover this failure
cannot be repaired by
replacing the canonical
projector with \emph{any}
Galois-equivariant
idempotent
$E_\epsilon$ that reduces
to the fixed130 tame
projector $\Pi_0$.
Hence the fixed130
$\operatorname{ad}T=M_0\oplus M_1$
splitting does \emph{not}
extend to the unrestricted
universal framed
$S_3$ deformation space,
even after deforming the
projector itself.
\end{theorem}

\begin{proof}
Over the dual-number ring
one has
$\bar U^3=I_2$ and
$\bar S^{-1}\bar U\bar S
=\bar U^2$.
Also
\[
\bar V+\operatorname{Ad}_{\bar U}\bar V
+\operatorname{Ad}_{\bar U^2}\bar V=0.
\]
Therefore
\[
\bigl((I_2+\epsilon\bar V)\bar U\bigr)^3
=I_2
\]
and the $\omega_2$
component of
$\rho_\epsilon(x_0\tau)$
is the identity.
Similarly the
$\omega_2$ component
of $\bar U$ is the
identity, whereas
$\bar S^{\omega_2}=\bar S$.
Consequently the literal
marked auxiliary words
satisfy
\[
u_0=u_1=1,\quad
d_0=x_0^{-1}=x_0,\quad
g_0=\sigma^2=1,\quad
c_0=h_c=1.
\]
Because $x_0^2=1$ and
$d_0=x_0$, the remaining
word $h_0$ is also one.
The full wild relator
$h_0u_1^{-1}x_1^\sigma c_0$
is therefore one, and the
tame relator was checked
above.
The wild images lie in the
normal subgroup
$1+\epsilon M_2(k)$,
which is an elementary
abelian $2$-group.
The marked pro-$2$
condition is automatic.
The universal property
of the marked
Roe--Turturean presentation
then gives an actual
continuous Galois
representation.

Since $\rho_\epsilon(\tau)$
has exact order three,
its pro-$2$ component
is one and
$z_\epsilon=\bar U$.
Direct characteristic-two
matrix multiplication gives
$[\bar V,\bar U]=\bar V$,
proving
\eqref{eq:dyadic-f131-nonzero-wild-commutator}.

Suppose that
$E_\epsilon=\Pi_0+\epsilon D$
were a Galois-equivariant
idempotent reducing to
$\Pi_0$.
As an action on the
adjoint module,
\[
\operatorname{Ad}_{\rho_\epsilon(x_0)}
=I+\epsilon\operatorname{ad}_{\bar V}.
\]
Commutation with $E_\epsilon$
would imply, already
at order $\epsilon$,
\[
[\operatorname{ad}_{\bar V},\Pi_0]=0.
\]
But $\Pi_0(\bar U)=\bar U$,
$\operatorname{ad}_{\bar V}(\bar U)=\bar V$,
and $\Pi_0(\bar V)=0$,
so this commutator maps
$\bar U$ to the nonzero
matrix $\bar V$.
This contradiction is
independent of $D$ and
proves the strong
nonextendability.
\end{proof}

\begin{corollary}[Exact first-order codimension of the split subfunctor]
\label{cor:dyadic-f131-split-tangent-codimension}
At the residual $S_3$
point over $\mathbf F_2$,
the first-order split
condition on a framed
Galois cocycle $c$ is
\begin{equation}\label{eq:dyadic-f131-tangent-wild-character}
\boxed{\quad
\kappa_{\rm wild}([c])
:=(I-\bar\Pi_0)c(x_0)=0
\quad\text{in }\bar M_1.
\quad}
\end{equation}
The map
\[
\kappa_{\rm wild}:
H^1(\Gamma,\operatorname{ad}\bar T)
\longrightarrow\bar M_1
\cong\mathbf F_2^2
\]
is a \emph{surjection}
with kernel
$H^1(\Gamma,\bar M_0)$.
Thus there is a short
exact sequence
\begin{equation}\label{eq:dyadic-f131-tangent-exact-sequence}
\boxed{
0\to H^1(\Gamma,\bar M_0)
\longrightarrow
H^1(\Gamma,\operatorname{ad}\bar T)
\xrightarrow{\;\kappa_{\rm wild}\;}
\bar M_1\to0.
}
\end{equation}
The unframed residual
tangent dimensions are
\[
\dim_{\mathbf F_2}H^1(\Gamma,\operatorname{ad}\bar T)=6,
\qquad
\dim_{\mathbf F_2}\ker\kappa_{\rm wild}=4.
\]
The corresponding framed
tangent dimensions are
$9$ and $7$, respectively.
Hence the tame-isotypic
split subfunctor is a
\emph{proper} closed
subfunctor, with tangent
codimension exactly two.
No assertion that its
global defining ideal is
a regular sequence is made.
\end{corollary}

\begin{proof}
At the residual point
the wild generators
act trivially, so
$c(x_j)$ is unchanged
under infinitesimal
conjugation.
Differentiating the
commutator equations
\eqref{eq:dyadic-f131-wild-split-equations}
gives
$[c(x_j),\bar U]=0$,
equivalently
$(I-\bar\Pi_0)c(x_j)=0$.
The fixed130 blockwise
linearized relations show
that the $\bar M_1$
component of $c(x_1)$
is already zero for
\emph{every} global
cocycle, whereas the
$\bar M_1$ component of
$c(x_0)$ is unrestricted.
Indeed the $\bar M_1$
wild differential reduces
to $\bar S^{-1}c(x_1)$
in characteristic two.
The norm-zero natural
block of fixed129
identifies its $H^1$
with $\bar M_1$ by
evaluation at $x_0$.
Consequently
$\kappa_{\rm wild}$
is the canonical
projection onto the
$\bar M_1$ summand of
the fixed130
decomposition, giving
the short exact
sequence.

The fixed130 $M_0$
Fox complex has
$H^1$ of dimension
four after mod-two
reduction: its three
integral free directions
and its one top-degree
$2$-torsion Bockstein
direction.
The $\bar M_1$
complex contributes two
additional dimensions.
Thus the full unframed
$H^1$ dimension is six
and the kernel has
dimension four.
The residual adjoint
invariants are scalars
of dimension one, so
the degree-zero
coboundary image has
dimension $4-1=3$.
Adding these gauge
directions gives the
framed tangent
dimensions $9$ and $7$.
The explicit deformation
of
Theorem~\ref{thm:dyadic-f131-first-order-nosplit}
realizes a nonzero
obstruction class, so
the ideal is genuinely
proper.
\end{proof}

\begin{theorem}[Exact curved two-block Cartan multiplication]
\label{thm:dyadic-f131-curved-tame-transport}
For an arbitrary framed
deformation over $A$,
retain its canonical tame
projector $\Pi=\Pi_A$
and put
$P_0=\Pi$, $P_1=I-\Pi$.
For $g\in\Gamma$ define
\[
\mathcal R(g)=\operatorname{Ad}_{\rho_A(g)}
\quad\text{on }M_A,
\qquad
\mathcal R_{ab}(g)=P_a\mathcal R(g)P_b,
\quad a,b\in\{0,1\}.
\]
Write
\begin{equation}\label{eq:dyadic-f131-curved-two-blocks}
\mathcal R(g)=
\begin{pmatrix}
A_g&B_g\\ C_g&D_g
\end{pmatrix}
\quad\text{on }M_{0,A}\oplus M_{1,A}.
\end{equation}
The true profinite
transport remains
multiplicatively flat,
but its two truncated
diagonal carriers obey
the \emph{exact} formulas
\begin{equation}\label{eq:dyadic-f131-curved-block-products}
\boxed{
\begin{aligned}
A_{gh}&=A_gA_h+B_gC_h,\\
B_{gh}&=A_gB_h+B_gD_h,\\
C_{gh}&=C_gA_h+D_gC_h,\\
D_{gh}&=C_gB_h+D_gD_h.
\end{aligned}}
\end{equation}
In particular
\begin{equation}\label{eq:dyadic-f131-two-step-return-curvature}
\boxed{\quad
\mathscr F_{00}(g,h)
:=A_{gh}-A_gA_h=B_gC_h,
\quad
\mathscr F_{11}(g,h)
:=D_{gh}-D_gD_h=C_gB_h.
\quad}
\end{equation}
Both diagonal defects
vanish on the split
locus, but outside it
the off-diagonal
wild transport is a
genuine first-order
source of possible
two-step return defects.
The off-diagonal blocks
vanish identically for
$\sigma$ and $\tau$,
and all their nonzero
sources are carried by
wild words.

For any finite residual
precision $\mathfrak m^t$,
each block matrix,
its product formula,
and the projectors
are finite computations
from the marked
transport matrices
and unit-matrix inverses.
The tame projector itself
is already polynomial
on this marked residual
sector, without any
binomial truncation;
only other marked
profinite word evaluations
may require the fixed126
finite-jet procedure.
These formulas define
a \emph{source-faithful
curved two-block
coefficient chart} on
the unrestricted
$S_3$ residual family.
They are a
finite arithmetic
transport identity,
not an assertion of
additional analytic
differential equations.
\end{theorem}

\begin{proof}
The group law is
$\mathcal R(gh)
=\mathcal R(g)\mathcal R(h)$.
Insert $P_0+P_1=I$
between the two factors
and project to each of
the four pairs of block
degrees.  This gives
\eqref{eq:dyadic-f131-curved-block-products}
and
\eqref{eq:dyadic-f131-two-step-return-curvature}
exactly, with no
approximation.
By
\eqref{eq:dyadic-f131-sigma-tame-normalization},
the $\sigma,\tau$
transport preserves
both tame summands.
Their off-diagonal
blocks therefore vanish.
Vanishing of all wild
off-diagonal blocks is
equivalent to the split
criterion of
Theorem~\ref{thm:dyadic-f131-universal-split-locus}.
The stated finite-jet
effectivity follows
from the binomial
tail estimate and
ordinary finite
matrix multiplication.
\end{proof}

\begin{remark}[Two distinct meanings of closure]
\label{rem:dyadic-f131-split-versus-global-atlas}
The universal tame
\emph{matrix projector}
exists over the full
residual deformation
ring.  A
\emph{Galois-equivariant
decomposition} by that
projector exists
precisely on the
wild-centralizer
subfunctor, and a
genuine first-order
counterexample shows
that the latter is
proper.
By contrast the
all-Sylow perfect
cochain carrier and
the Cayley--Hamilton
full-transport
realization exist
over the whole
framed residual
sector regardless
of this failure.
Therefore failure of
a strict tame
eigensplitting does
not reopen the
derived faithful
finite-atlas theorem.
Nor does the curved
block identity
provide the still
missing canonical
Schreier-to-minimal-
Demu\v{s}kin Fox
syzygy comparison
for arbitrary
non-pro-$2$
deformation families.
\end{remark}

\subsection{Potential Fox closure}

After a finite residual base change every representation enters a maximal pro-$p$ chart.  Choose a stable lattice $T\subset V$ and a finite Galois extension $L/K$ killing $\bar\rho$ and containing $\mu_p$; then
\[
\rho(G_L)\subset 1+\varpi_E\End_{\mathcal O_E}(T),
\]
which is pro-$p$, so the restricted action factors through the Demu\v{s}kin quotient $P_L=G_L(p)$.

\begin{proposition}[Potential finite Fox closure]\label{prop:potential-fox-closure}
With $L$ as above, $R\Gamma(G_L,V)$ is computed by the finite Fox--Spencer complex of $P_L$.  Moreover, if $\Delta=\Gal(L/K)$, then
\[
R\Gamma(G_K,V)
\simeq
R\Gamma\bigl(\Delta,R\Gamma(G_L,V)\bigr).
\]
Since $E$ has characteristic zero and $\Delta$ is finite, derived $\Delta$-invariants are represented by ordinary invariants after a homotopy-coherent $\Delta$-action has been transferred to a finite Fox model.
\end{proposition}

\begin{proof}
Let
\[
N:=\ker(G_L\twoheadrightarrow P_L).
\]
Because $\rho(G_L)$ is pro-$p$, the universal property of $P_L$ shows that $N$ acts trivially on $T$ and on every $T/\varpi_E^nT$.  We use the standard \v{S}afarevi\v{c}--Demu\v{s}kin cohomology comparison for a $p$-adic local field containing $\mu_p$: inflation
\[
H^i(P_L,\mathbf F_p)\longrightarrow H^i(G_L,\mathbf F_p)
\]
is an isomorphism for $i=0,1,2$ \cite{NSW}.  The Hochschild--Serre five-term sequence therefore gives
\[
H^1(N,\mathbf F_p)^{P_L}=0.
\]
The module $H^1(N,\mathbf F_p)$ is a discrete $p$-primary module for the pro-$p$ group $P_L$; any nonzero such module contains nonzero invariants.  Hence
\[
H^1(N,\mathbf F_p)=0.
\]
To justify the degree-two comparison one must use the full Hochschild--Serre spectral sequence, not only its five-term truncation.  Since the $q=1$ row vanishes,
\[
d_2:E_2^{0,2}=H^2(N,\mathbf F_p)^{P_L}
   \longrightarrow
   E_2^{2,1}=H^2\bigl(P_L,H^1(N,\mathbf F_p)\bigr)
\]
has zero target.  The only remaining possible outgoing differential is
\[
d_3:E_3^{0,2}\longrightarrow E_3^{3,0},
\]
and its target is zero because $P_L$ has $p$-cohomological dimension two.  Hence
\[
E_\infty^{0,2}=E_2^{0,2}=H^2(N,\mathbf F_p)^{P_L}.
\]
But inflation in degree two is already surjective by the displayed local comparison above.  In the Hochschild--Serre filtration of $H^2(G_L,\mathbf F_p)$ this forces the complementary graded piece $E_\infty^{0,2}$ to vanish.  Thus
\[
H^2(N,\mathbf F_p)^{P_L}=0.
\]
Again $H^2(N,\mathbf F_p)$ is a discrete $p$-primary $P_L$-module, so the same pro-$p$ invariant argument gives
\[
H^2(N,\mathbf F_p)=0.
\]
Devissage through the $\varpi_E$-adic filtration therefore yields
\[
H^j(N,T/\varpi_E^nT)=0
\qquad (j=1,2,\ n\ge1).
\]
The Hochschild--Serre spectral sequence for $T/\varpi_E^nT$ consequently has only the $q=0$ row in total degrees at most two, so inflation is an isomorphism in degrees $0,1,2$.  Both $G_L$ and the Demu\v{s}kin group $P_L$ have $p$-cohomological dimension two, so there is no higher cohomology.  The transition maps in the finite coefficient systems are surjective, so the Mittag--Leffler condition removes derived inverse-limit terms.  Passing to the inverse limit over $n$ and then inverting $\varpi_E$ gives
\[
R\Gamma(P_L,V)\simeq R\Gamma(G_L,V).
\]
The finite Fox--Spencer complex computes the left-hand side by Theorem~\ref{thm:fox-spencer-demushkin}.

The second assertion is Hochschild--Serre for
\[
1\longrightarrow G_L\longrightarrow G_K\longrightarrow\Delta\longrightarrow1.
\]
Averaging by $|\Delta|^{-1}\in E$ makes invariants exact on $E[\Delta]$-modules.  A strict action on a chosen Fox presentation need not be canonical, so the descent action on the finite model is naturally homotopy coherent.
\end{proof}

\subsection{Rational descent and the integral defect}

Keep the notation of Proposition~\ref{prop:potential-fox-closure}, and put
\[
Q_{L/K}:=G_K/N,
\qquad
1\longrightarrow P_L\longrightarrow Q_{L/K}\longrightarrow\Delta\longrightarrow1.
\]
The representation $V$ factors through $Q_{L/K}$.  Conjugation in $Q_{L/K}$ gives the intrinsic object $R\Gamma(P_L,V)$ a homotopy-coherent $\Delta$-action; a choice of section $\Delta\to Q_{L/K}$ turns this into explicit action maps together with a factor-set homotopy.

\begin{theorem}[Canonical rational finite descent from a potential Fox chart]\label{thm:rational-fox-descent}
Let $C^\bullet_{\mathrm{Fox}}(P_L,V)$ be any finite Fox--Spencer model of $R\Gamma(P_L,V)$.  The intrinsic homotopy-coherent $\Delta$-object $R\Gamma(P_L,V)$ admits a bounded representative
\[
C^\bullet_{L/K}(V)\in K^b\bigl(E[\Delta]\text{-}\mathrm{proj}\bigr).
\]
For this representative there are canonical equivalences
\begin{equation}\label{eq:rational-fox-descent}
R\Gamma(G_K,V)
\simeq
R\Gamma\bigl(\Delta,R\Gamma(P_L,V)\bigr)
\simeq
C^\bullet_{L/K}(V)^\Delta.
\end{equation}
The quasi-isomorphism class of the invariant complex is independent of the stable lattice, the section of $Q_{L/K}\to\Delta$, the Demu\v{s}kin presentation, and the chosen strictification.
\end{theorem}

\begin{proof}
The first equivalence is Hochschild--Serre together with Proposition~\ref{prop:potential-fox-closure}.  The intrinsic $\Delta$-action is an action in the derived category: changing a lift of $\delta\in\Delta$ modifies the induced automorphism of $P_L$ by an inner automorphism, and inner automorphisms act chain-homotopically trivially on group cochains.  Hence the resulting action is independent of a section up to coherent homotopy.

Because $E$ has characteristic zero, Maschke's theorem makes $E[\Delta]$ semisimple.  A perfect $E$-complex with homotopy-coherent $\Delta$-action therefore admits a bounded strict representative by finite projective $E[\Delta]$-modules.  Moreover, the invariants functor is exact, with Reynolds operator
\[
\operatorname{Av}_\Delta
=\frac1{|\Delta|}\sum_{\delta\in\Delta}\delta.
\]
Consequently derived invariants are computed termwise, giving the second equivalence in \eqref{eq:rational-fox-descent}.  Every choice made above represents the same object in $\operatorname{Perf}(E[\Delta])$, so its invariant complex has a choice-independent quasi-isomorphism class.
\end{proof}

\begin{corollary}[Prime-to-$p$ integral descent]\label{cor:prime-to-p-integral-descent}
If $p\nmid |\Delta|$, then the same conclusion holds over $\mathcal O_E$: derived $\Delta$-invariants are ordinary invariants, and any finite $\Delta$-equivariant integral Fox model yields a finite integral complex computing $R\Gamma(G_K,T)$.
\end{corollary}

\begin{proof}
Now $|\Delta|\in\mathcal O_E^\times$, so the same averaging idempotent is defined integrally and makes invariants exact.
\end{proof}

\begin{warning}[Naive integral invariants fail for $p$-primary descent]\label{warn:integral-delta-invariants}
If $p\mid|\Delta|$, ordinary invariants do not compute derived invariants over $\mathcal O_E$.  Already for $\Delta=C_p$ acting trivially on $\mathcal O_E$, the standard periodic resolution gives
\[
H^{2j}(C_p,\mathcal O_E)\cong \mathcal O_E/p\mathcal O_E
\quad(j\ge1),
\qquad
H^{2j+1}(C_p,\mathcal O_E)=0.
\]
Thus replacing a potential integral Fox complex by its termwise $\Delta$-invariants loses genuine descent classes.  In the Hochschild--Serre spectral sequence
\[
E_2^{a,b}=H^a\bigl(\Delta,H^b(P_L,T)\bigr)
\Longrightarrow H^{a+b}(G_K,T),
\]
the finite-extension class and its transgressions must cancel the apparently unbounded $p$-primary quotient cohomology, because the abutment vanishes above degree two.  This cancellation is part of the arithmetic transport data; it is not visible in ordinary invariants of the Fox complex.
\end{warning}


\subsection{Sylow absorption and a prime-to-p Fox retract}

Pointwise integral closure avoids this fixed normal quotient: absorb the $p$-primary residual image into an open pro-$p$ chart and descend only across prime-to-$p$ index.

\begin{theorem}[Sylow absorption and prime-to-$p$ Fox retract]\label{thm:sylow-fox-retract}
Let $T$ be a $G_K$-stable $\mathcal O_E$-lattice, let
\[
\overline\rho_T:G_K\longrightarrow \mathrm{GL}(T/\varpi_ET)
\]
be its residual representation, and put $\overline G_T:=\operatorname{Im}(\overline\rho_T)$.  Choose a Sylow $p$-subgroup $S\subset\overline G_T$ and let
\[
H:=\overline\rho_T^{-1}(S),
\qquad
m:=[G_K:H]=[\overline G_T:S].
\]
Then:
\begin{enumerate}[label=\textnormal{(\roman*)}]
\item $m$ is prime to $p$, and the image of $H$ in $\mathrm{GL}_{\mathcal O_E}(T)$ is pro-$p$;
\item the restricted representation factors through the maximal pro-$p$ quotient $P_H:=H(p)$;
\item if $K_H$ is the finite extension fixed by $H$, then $P_H$ is free pro-$p$ when $\mu_p\not\subset K_H$ and is Demu\v{s}kin when $\mu_p\subset K_H$; hence $R\Gamma(H,T)$ is computed by a finite free Fox--Spencer complex $C^\bullet_{\mathrm{Fox}}(P_H,T)$ of amplitude contained in $[0,2]$;
\item for the chosen $S$, restriction identifies $R\Gamma(G_K,T)$ with the image of the homotopy idempotent
\begin{equation}\label{eq:sylow-transfer-idempotent}
 e_H
 :=
 \frac1m\,\operatorname{res}_{H}^{G_K}
 \circ
 \operatorname{cor}_{H}^{G_K}
 \quad\in\quad
 \operatorname{End}_{D(\mathcal O_E)}R\Gamma(H,T).
\end{equation}
Consequently $R\Gamma(G_K,T)$ is represented by a bounded finite-projective direct summand of the Fox chart of $P_H$.
\end{enumerate}
\end{theorem}

\begin{proof}
The index $m$ is prime to $p$ by the Sylow theorem.  The kernel of reduction
\[
\mathrm{GL}_r(\mathcal O_E)\longrightarrow \mathrm{GL}_r(k_E)
\]
is pro-$p$.  The image $\rho(H)$ is an extension of a closed subgroup of this kernel by the finite $p$-group $S$, and is therefore pro-$p$.  Hence the restricted action factors through $P_H=H(p)$.

The group $H$ is the absolute Galois group of the finite extension $K_H/K$.  For every prime $p$, the maximal pro-$p$ quotient of the absolute Galois group of a $p$-adic local field is free pro-$p$ when $\mu_p$ is not contained in the field and is a finite-rank Demu\v{s}kin group when $\mu_p$ is contained in the field; in particular the dyadic case is Demu\v{s}kin and requires only the standard pro-$2$ relation case distinction \cite{NSW,JardenShusterman}.  In all cases there is a finite free Fox resolution of amplitude contained in $[0,2]$.

If $\mu_p\subset K_H$, the inflation comparison is exactly the Demu\v{s}kin argument of Proposition~\ref{prop:potential-fox-closure}.  If $\mu_p\not\subset K_H$, then $P_H$ is free pro-$p$, so
\[
H^2(P_H,\mathbf F_p)=0,
\]
while local Tate duality gives
\[
H^2(H,\mathbf F_p)
\cong H^0(H,\mu_p)^\vee
=0.
\]
Thus mod-$p$ inflation is also an isomorphism in degree two in the free branch; degree one is unchanged by passage to the maximal pro-$p$ quotient.  For $N_H:=\ker(H\twoheadrightarrow P_H)$, the five-term sequence and the full Hochschild--Serre degree-two argument used in Proposition~\ref{prop:potential-fox-closure} therefore give
\[
H^1(N_H,\mathbf F_p)=H^2(N_H,\mathbf F_p)=0.
\]
Since $N_H$ acts trivially on $T$, d\'evissage through $T/\varpi_E^nT$ followed by Mittag--Leffler passage gives in both the free and Demu\v{s}kin branches
\[
R\Gamma(P_H,T)\simeq R\Gamma(H,T).
\]

There are restriction and corestriction morphisms in continuous cohomology, satisfying
\[
\operatorname{cor}_{H}^{G_K}\circ
\operatorname{res}_{H}^{G_K}
=m\,\operatorname{id}_{R\Gamma(G_K,T)}.
\]
Since $m\in\mathcal O_E^\times$, the endomorphism $e_H$ in \eqref{eq:sylow-transfer-idempotent} satisfies
\[
e_H^2
=\frac1{m^2}\operatorname{res}\operatorname{cor}
\operatorname{res}\operatorname{cor}
=\frac1m\operatorname{res}\operatorname{cor}
=e_H.
\]
The restriction map identifies $R\Gamma(G_K,T)$ with $\operatorname{Im}(e_H)$, with inverse $m^{-1}\operatorname{cor}$.  The category of perfect $\mathcal O_E$-complexes is idempotent complete, so this image is represented by a bounded finite-projective direct summand of any finite Fox model of $R\Gamma(H,T)$.
\end{proof}

\begin{corollary}[Pointwise integral Fox closure without $p$-primary compression]\label{cor:sylow-pointwise-closure}
Every continuous integral $p$-adic representation admits a finite Fox--Spencer coefficient model after restriction to an open subgroup of prime-to-$p$ index, and its original $G_K$-cohomology is a transfer-idempotent summand of that model.  Thus a separate extension-resolution compression of the $p$-primary quotient is not required for the existence of pointwise finite integral closure.
\end{corollary}


\subsection{Residual-family Sylow--Fox atlases}

Fixing the residual framed representation makes the construction uniform; only descent to the residual-determinant stack remains.

\begin{proposition}[Universal pro-$p$ restriction on a residual framed chart]\label{prop:universal-sylow-prop}
Let
\[
\bar\rho:G_K\longrightarrow \mathrm{GL}_r(k_E)
\]
be a continuous residual representation, let $R_{\bar\rho}^{\square}$ be its universal framed deformation ring, and let
\[
\rho^{\square}:G_K\longrightarrow
\mathrm{GL}_r(R_{\bar\rho}^{\square})
\]
be the universal framed representation.  Put $\bar G:=\operatorname{Im}(\bar\rho)$, choose a Sylow $p$-subgroup $S\subset\bar G$, and set
\[
H_{\bar\rho,S}:=\bar\rho^{-1}(S),
\qquad
m_{\bar\rho}:=[G_K:H_{\bar\rho,S}]=[\bar G:S].
\]
Then $m_{\bar\rho}$ is prime to $p$, the image
\[
\rho^{\square}(H_{\bar\rho,S})
\subset
\mathrm{GL}_r(R_{\bar\rho}^{\square})
\]
is pro-$p$ for the maximal-ideal topology, and the universal restricted representation factors through
\[
P_{\bar\rho,S}:=H_{\bar\rho,S}(p).
\]
\end{proposition}

\begin{proof}
The index statement is Sylow's theorem.  Reduction modulo the maximal ideal of $R_{\bar\rho}^{\square}$ maps the universal restricted image into $S$.  The kernel
\[
1+M_r(\mathfrak m_{R_{\bar\rho}^{\square}})
\]
is pro-$p$: its quotients by the congruence filtration are successive extensions of finite additive $k_E$-vector groups.  Hence the universal image of $H_{\bar\rho,S}$ is an extension of a finite $p$-group by a pro-$p$ group and is therefore pro-$p$.  The universal property of the maximal pro-$p$ quotient gives the factorization through $P_{\bar\rho,S}$.
\end{proof}

\begin{theorem}[Family Sylow--Fox coefficient atlas]\label{thm:family-sylow-fox}
Retain the notation of Proposition~\ref{prop:universal-sylow-prop}.  Let
\[
\mathcal T^{\square}:=(R_{\bar\rho}^{\square})^r
\]
with its universal $G_K$-action.  Then:
\begin{enumerate}[label=\textnormal{(\roman*)}]
\item the group $P_{\bar\rho,S}$ is free pro-$p$ or Demu\v{s}kin, and a finite Fox resolution gives a bounded finite-free complex
\[
C^\bullet_{\mathrm{Fox}}
(P_{\bar\rho,S},\mathcal T^{\square})
\]
over $R_{\bar\rho}^{\square}$ computing
$R\Gamma(H_{\bar\rho,S},\mathcal T^{\square})$;
\item the restriction--corestriction endomorphism
\begin{equation}\label{eq:family-sylow-idempotent}
 e_{\bar\rho,S}^{\square}
 :=
 \frac1{m_{\bar\rho}}
 \operatorname{res}_{H_{\bar\rho,S}}^{G_K}
 \circ
 \operatorname{cor}_{H_{\bar\rho,S}}^{G_K}
\end{equation}
is a homotopy idempotent of that perfect complex, and its image
\[
\mathcal K_{\bar\rho,S}^{\mathrm{Syl}}
:=\operatorname{Im}(e_{\bar\rho,S}^{\square})
\in\Perf(R_{\bar\rho}^{\square})
\]
is naturally equivalent to
\[
R\Gamma(G_K,\mathcal T^{\square});
\]
\item for every derived morphism
$R_{\bar\rho}^{\square}\to B$, one has a canonical base-change equivalence
\begin{equation}\label{eq:family-sylow-base-change}
\mathcal K_{\bar\rho,S}^{\mathrm{Syl}}
\otimes_{R_{\bar\rho}^{\square}}^{\mathbf L}B
\simeq
R\Gamma
\left(
G_K,
\mathcal T^{\square}
\otimes_{R_{\bar\rho}^{\square}}^{\mathbf L}B
\right);
\end{equation}
\item the quasi-isomorphism class of
$\mathcal K_{\bar\rho,S}^{\mathrm{Syl}}$
is independent of the chosen Sylow subgroup $S$.
\end{enumerate}
Thus the pointwise Sylow--Fox retract extends to a finite family-level coefficient chart on the entire framed residual deformation space.
\end{theorem}

\begin{proof}
The fixed field of $H_{\bar\rho,S}$ is a finite extension of $K$.  Proposition~\ref{prop:universal-sylow-prop} makes the universal coefficient action factor through its maximal pro-$p$ Galois group, which is free pro-$p$ or Demu\v{s}kin for every $p$; at $p=2$ it lies in the standard dyadic Demu\v{s}kin cases.  Base-changing the corresponding finite free Fox resolution and applying continuous $\operatorname{Hom}$ gives the finite-free complex in (i).  The inflation comparison is the same finite-level d\'evissage and Mittag--Leffler argument used in Proposition~\ref{prop:potential-fox-closure}.

Restriction and corestriction are natural for coefficient families and satisfy
\[
\operatorname{cor}\circ\operatorname{res}
=m_{\bar\rho}\operatorname{id}.
\]
Because $m_{\bar\rho}$ is a unit in $R_{\bar\rho}^{\square}$, Equation~\eqref{eq:family-sylow-idempotent} is a homotopy idempotent.  Idempotent completeness of $\Perf(R_{\bar\rho}^{\square})$ gives its image, and the usual restriction--corestriction argument identifies that image with the universal $G_K$-cohomology object.  All maps are defined over $R_{\bar\rho}^{\square}$ and commute with derived coefficient change, which proves (iii).

If $S'=\bar gS\bar g^{-1}$ is another Sylow subgroup, choose a lift $g\in G_K$ of $\bar g$.  Conjugation by $g$ identifies $H_{\bar\rho,S}$ with $H_{\bar\rho,S'}$, their maximal pro-$p$ quotients, Fox resolutions, and transfer idempotents.  Two lifts differ by an element of $\ker(\bar\rho)\subset H_{\bar\rho,S}$; the resulting inner automorphisms act chain-homotopically trivially on continuous cochains.  Hence the induced equivalence is independent of the lift in the derived category.
\end{proof}


\begin{proposition}[Universal formal Fox matrices on a residual framed chart]\label{prop:universal-formal-fox-matrices}
Retain the notation of Theorem~\ref{thm:family-sylow-fox}, and choose a free or Demu\v{s}kin presentation
\[
 P_{\bar\rho,S}
 =\langle x_1,\ldots,x_d\mid r_1,\ldots,r_m\rangle_{\mathrm{pro}\text{-}p},
 \qquad m\in\{0,1\}.
\]
Let
\[
 \rho_S^{\square}:\mathcal O_E[[P_{\bar\rho,S}]]
 \longrightarrow
 M_r(R_{\bar\rho}^{\square})
\]
be the continuous algebra homomorphism induced by the universal restricted representation.  With the left-module conventions of Section~\ref{sec:fox-spencer}, the universal Fox--Spencer differential is the finite matrix complex
\begin{equation}\label{eq:universal-formal-fox-complex}
 \mathcal T^{\square}
 \xrightarrow{\ d^0_{\square}\ }
 (\mathcal T^{\square})^d
 \xrightarrow{\ d^1_{\square}\ }
 (\mathcal T^{\square})^m,
\end{equation}
where
\begin{align}
 d^0_{\square}(v)
 &=\bigl((\rho_S^{\square}(x_i)-I)v\bigr)_{i=1}^d,\label{eq:universal-fox-d0}\\
 d^1_{\square}(v_1,\ldots,v_d)_j
 &=\sum_{i=1}^d
 \rho_S^{\square}\!\left(
 \overline{\frac{\partial r_j}{\partial x_i}}
 \right)v_i.\label{eq:universal-fox-d1}
\end{align}
The entries are convergent elements of
$M_r(R_{\bar\rho}^{\square})$, one has
$d^1_{\square}d^0_{\square}=0$, and the construction commutes strictly with every continuous base change
$R_{\bar\rho}^{\square}\to B$.
Consequently the Fox differential is already an explicit formal full-transport matrix on the complete residual framed chart.  Theorem~\ref{thm:ch-fox-effectivity} below shows that these matrices algebraize on a finite-type pro-$p$ factorization locus of the Cayley--Hamilton generic-matrix chart.  Coherent comparison of the resulting algebraic loci is organized later by the stratified atlas; finite-type effectivity itself is already closed.
\end{proposition}

\begin{proof}
The factorization in Proposition~\ref{prop:universal-sylow-prop} extends continuously from the group to the completed group algebra and therefore evaluates every completed Fox derivative in the universal matrix algebra.  Equations~\eqref{eq:universal-fox-d0} and \eqref{eq:universal-fox-d1} are the dual of the completed Fox resolution.  The fundamental identity
\[
 r_j-1
 =\sum_i
 \frac{\partial r_j}{\partial x_i}(x_i-1)
\]
becomes zero after passing to the quotient group, and hence gives
$d^1_{\square}d^0_{\square}=0$.  Continuous evaluation of completed group-ring elements commutes with continuous coefficient change, proving the base-change assertion.
\end{proof}

\begin{corollary}[Finite residual Sylow--Fox atlas]\label{cor:finite-residual-sylow-atlas}
For fixed $K$, $E$, and rank $r$, the coefficient cohomology of every continuous $r$-dimensional representation is covered by finitely many family Sylow--Fox charts
\[
\left
\{
\mathcal K_{\bar\rho,S}^{\mathrm{Syl}}
\right\}_{[\bar\rho]}.
\]
Each chart is integral, finite perfect of amplitude contained in $[0,2]$, compatible with derived base change, and retains the universal framed $G_K$-representation as its coefficient transport.
\end{corollary}

\begin{proof}
There are only finitely many residual representation classes in fixed rank, because $G_K$ is topologically finitely generated and $\mathrm{GL}_r(k_E)$ is finite.  Apply Theorem~\ref{thm:family-sylow-fox} to each class.
\end{proof}

\subsection{Canonical all-Sylow averaging and Morita descent}\label{subsec:all-sylow-descent}

Keeping all Sylow subgroups simultaneously gives a symmetric finite model and removes the choice on the residual representation groupoid.

\begin{definition}[All-Sylow Fox carrier]\label{def:all-sylow-carrier}
Retain a residual representation
\[
 \bar\rho:G_K\longrightarrow\mathrm{GL}_r(k_E),
 \qquad
 \bar G:=\operatorname{Im}(\bar\rho),
\]
and let $\operatorname{Syl}_p(\bar G)$ be the finite set of Sylow $p$-subgroups of $\bar G$.  Put
\[
 \nu_{\bar\rho}:=\#\operatorname{Syl}_p(\bar G),
 \qquad
 m_{\bar\rho}:=[\bar G:S]
\]
for any $S\in\operatorname{Syl}_p(\bar G)$.  For each $S$, let
\[
 H_S:=\bar\rho^{-1}(S),
 \qquad
 P_S:=H_S(p),
\]
and let $\mathcal T^{\square}=(R_{\bar\rho}^{\square})^r$ be the universal framed coefficient module.  The intrinsic all-Sylow carrier is
\begin{equation}\label{eq:all-sylow-carrier}
 \mathcal C_{\bar\rho}^{\mathrm{allSyl}}
 :=
 \bigoplus_{S\in\operatorname{Syl}_p(\bar G)}
 R\Gamma(H_S,\mathcal T^{\square})
 \in\operatorname{Perf}(R_{\bar\rho}^{\square}).
\end{equation}
After choosing free or Demu\v{s}kin presentations of the groups $P_S$, it is represented by a finite direct sum of Fox--Spencer complexes.
\end{definition}

\begin{theorem}[Canonical all-Sylow transfer retract]\label{thm:all-sylow-transfer}
With the notation of Definition~\ref{def:all-sylow-carrier}, define
\[
 \operatorname{res}_{\mathrm{all}}
 :=
 (\operatorname{res}_{H_S}^{G_K})_S,
 \qquad
 \operatorname{cor}_{\mathrm{all}}
 :=
 \sum_S\operatorname{cor}_{H_S}^{G_K}.
\]
Then:
\begin{enumerate}[label=\textnormal{(\roman*)}]
\item both $m_{\bar\rho}$ and $\nu_{\bar\rho}$ are prime to $p$;
\item
\begin{equation}\label{eq:all-sylow-cor-res}
 \operatorname{cor}_{\mathrm{all}}
 \circ
 \operatorname{res}_{\mathrm{all}}
 =
 \nu_{\bar\rho}m_{\bar\rho}\,\operatorname{id}
\end{equation}
on $R\Gamma(G_K,\mathcal T^{\square})$;
\item the endomorphism
\begin{equation}\label{eq:all-sylow-idempotent}
 e_{\bar\rho}^{\mathrm{allSyl}}
 :=
 \frac1{\nu_{\bar\rho}m_{\bar\rho}}
 \operatorname{res}_{\mathrm{all}}
 \circ
 \operatorname{cor}_{\mathrm{all}}
\end{equation}
is a canonical homotopy idempotent of
$\mathcal C_{\bar\rho}^{\mathrm{allSyl}}$, and
\[
 \operatorname{Im}(e_{\bar\rho}^{\mathrm{allSyl}})
 \simeq
 R\Gamma(G_K,\mathcal T^{\square});
\]
\item the pair
$\bigl(\mathcal C_{\bar\rho}^{\mathrm{allSyl}},
 e_{\bar\rho}^{\mathrm{allSyl}}\bigr)$
is equivariant under conjugation of $\bar\rho$ and under automorphisms of the residual representation.  It therefore descends from the framed deformation chart to the corresponding unframed residual representation stack.
\end{enumerate}
The construction commutes with derived base change on $R_{\bar\rho}^{\square}$.
\end{theorem}

\begin{proof}
The index $m_{\bar\rho}$ is prime to $p$ by definition of a Sylow subgroup, while Sylow's congruence gives
$\nu_{\bar\rho}\equiv1\pmod p$; hence
$\nu_{\bar\rho}m_{\bar\rho}\in\mathcal O_E^\times$.
For every $S$, restriction and corestriction satisfy
\[
 \operatorname{cor}_{H_S}^{G_K}
 \operatorname{res}_{H_S}^{G_K}
 =m_{\bar\rho}\operatorname{id}
\]
in continuous cohomology \cite{BrownCohomology}.  Summing over all Sylow subgroups gives Equation~\eqref{eq:all-sylow-cor-res}; the same calculation as in Theorem~\ref{thm:sylow-fox-retract} proves that Equation~\eqref{eq:all-sylow-idempotent} is a homotopy idempotent whose image is the global cohomology object.

Conjugation in $\bar G$ permutes $\operatorname{Syl}_p(\bar G)$ transitively and identifies the corresponding open subgroups, maximal pro-$p$ quotients, restriction maps, and transfer maps.  Hence the direct sum and the averaged idempotent are invariant under the whole residual representation groupoid, not merely independent of one chosen Sylow subgroup.  Standard descent for perfect complexes on quotient stacks gives (iv).  Naturality of restriction, transfer, finite direct sums, and idempotent images gives the base-change statement.
\end{proof}

\begin{corollary}[Canonical formal descent on a residual determinant sector]\label{cor:all-sylow-det-formal}
Fix a residual determinant $\bar D$.  Taking the finite disjoint union of Theorem~\ref{thm:all-sylow-transfer} over the residual representation classes with determinant $\bar D$ produces a canonical perfect coefficient carrier on the formal residual-representation refinement of the Cayley--Hamilton stack $\mathfrak X_{\bar D}$.  Its pullback to every residual framed chart is represented by a finite all-Sylow Fox complex and its canonical idempotent.
\end{corollary}

\begin{proof}
There are finitely many residual representation classes in fixed rank.  Each class carries an equivariant all-Sylow object by Theorem~\ref{thm:all-sylow-transfer}; the conjugation-equivariant descent data glue on the corresponding quotient substacks.  Their finite disjoint union is precisely the formal refinement of the Cayley--Hamilton representation stack by residual isomorphism class.
\end{proof}

\begin{proposition}[Gerbe weight and Morita descent]\label{prop:all-sylow-morita}
On the absolutely irreducible locus of a residual determinant sector, the all-Sylow coefficient carrier descends canonically to the Cayley--Hamilton splitting $\Gm$-gerbe, and hence as a perfect module over the universal Azumaya Cayley--Hamilton algebra.  In general it need not descend as an ordinary perfect complex on the coarse pseudodeformation base.  By contrast, its adjoint coefficient complex has scalar weight zero and descends to the base.
\end{proposition}

\begin{proof}
On an Azumaya splitting chart the universal representation bundle has scalar weight one under the stabilizer $\Gm$.  Applying $R\Gamma(G_K,-)$ preserves this weight, so the coefficient complex is naturally an object on the splitting gerbe, equivalently a Morita module for the Azumaya algebra.  A weight-one object on a nontrivial $\Gm$-gerbe is not the pullback of an ordinary perfect complex from the coarse base.  The adjoint bundle is an endomorphism bundle, so scalar automorphisms act trivially; its cohomology complex has weight zero and therefore descends.  This is the same Morita distinction recorded in Proposition~\ref{prop:azumaya-gerbe} and Definition~\ref{def:morita-carrier}.
\end{proof}

\begin{warning}[Center-valued descent is too small]\label{warn:center-sylow-descent}
The all-Sylow construction removes the choice of a Sylow subgroup and gives descent to the representation stack.  It does not turn the commutative pseudodeformation ring $R_{\bar D}$ into a faithful coefficient carrier.  Reducible representations with the same determinant may have different extension data and residual images, while irreducible coefficient objects can carry nonzero gerbe weight.  The correct target is the Cayley--Hamilton representation stack or its Morita-valued algebra carrier, not the center alone.
\end{warning}

\begin{remark}[Status after all-Sylow descent]\label{rem:sylow-not-global}
The all-Sylow construction removes the choice of Sylow subgroup and gives formal unframed descent.  On residual framed charts, Proposition~\ref{prop:universal-formal-fox-matrices} gives universal Fox matrices; Theorem~\ref{thm:ch-fox-effectivity} algebraizes them on finite-type pro-$p$ factorization loci; and Proposition~\ref{prop:determinantal-contraction} gives explicit contractions on constant-rank strata.  Multiplicative and Morita descent are supplied by Theorem~\ref{thm:all-sylow-ainfty} and Corollary~\ref{cor:abstract-morita-multiplicative}, while Theorems~\ref{thm:fox-kuranishi-comparison} and~\ref{thm:cech-stasheff-gauges} control the deformation comparison and its overlap maps.

The all-Sylow step already removes the Sylow choice and gives the canonical perfect object.  Rank-jump specialization of its full perfect matrices is given by derived base change.  The later stratified-globalization and raw-Herr Tate-dual theorems make the cross-gauge comparison finite-native and recursively explicit; what survives is only the optional choice of a particular strict raw-period representative.
\end{remark}


\subsection{Multiplicative all-Sylow descent}\label{subsec:all-sylow-multiplicative}

The all-Sylow perfect object is canonical, but normalized transfer is not multiplicative.  Strict projection fails; homotopy-coherent multiplicative transfer remains available.

\begin{proposition}[Arithmetic normalized transfer is not multiplicative]\label{prop:transfer-not-multiplicative}
Let $p$ be odd, let $K/\Qp$ be a finite extension containing $\mu_p$, and let $L/K$ be a quadratic extension.  Put
\[
 G=G_K,
 \qquad
 H=G_L,
 \qquad
 \Delta=G/H=\langle\tau\rangle.
\]
After choosing a $G$-equivariant identification $\mathbf F_p\simeq\mu_p$, the normalized restriction--corestriction projector
\[
 e=\frac12\operatorname{res}_{H}^{G}
 \operatorname{cor}_{H}^{G}
 \quad\text{on}\quad H^*(H,\mathbf F_p)
\]
is not a graded-ring homomorphism.  The failure occurs already on
$H^1(H,\mathbf F_p)\times H^1(H,\mathbf F_p)\to H^2(H,\mathbf F_p)$.
\end{proposition}

\begin{proof}
Because $H$ is normal of index two, the Mackey formula identifies $e$ with the Reynolds operator
\[
 e=\frac12(1+\tau).
\]
Write
\[
 H^1(H,\mathbf F_p)=H^1_+\oplus H^1_-
\]
for the $\pm1$ eigenspace decomposition under $\tau$.  Since $2$ is invertible in $\mathbf F_p$, inflation--restriction gives
\[
 H^1_+=H^1(H,\mathbf F_p)^\Delta
 \cong H^1(G,\mathbf F_p).
\]
Let $d=[K:\Qp]$.  Kummer theory and the standard local dimension formula give
\[
 \dim_{\mathbf F_p}H^1(G,\mathbf F_p)=d+2,
 \qquad
 \dim_{\mathbf F_p}H^1(H,\mathbf F_p)=2d+2,
\]
because both $K$ and $L$ contain $\mu_p$ \cite{NSW}.  Hence
\[
 \dim_{\mathbf F_p}H^1_-=d>0.
\]
Because $\mu_p\subset K$, the cyclotomic extension $\Qp(\mu_p)$ is contained in $K$, so $p-1$ divides $d=[K:\Qp]$.  Since $p$ is odd, $d$ is even; thus the nondegenerate alternating form on $H^1_-$ obtained below has the required even dimension.

The local invariant map identifies $H^2(H,\mathbf F_p)$ with $\mathbf F_p$ and is $\Delta$-invariant.  Local Tate duality therefore makes the cup product
\[
 H^1(H,\mathbf F_p)\times H^1(H,\mathbf F_p)
 \longrightarrow H^2(H,\mathbf F_p)\cong\mathbf F_p
\]
a perfect alternating $\Delta$-invariant pairing \cite{NSW}.  If
$a\in H^1_+$ and $b\in H^1_-$, then
\[
 \langle a,b\rangle
 =\langle\tau a,\tau b\rangle
 =-\langle a,b\rangle,
\]
so $H^1_+\perp H^1_-$.  Perfectness then implies that the restriction of the pairing to $H^1_-$ is nondegenerate.  Thus there exist
$a,b\in H^1_-$ with
\[
 a\smile b\ne0.
\]
For these classes,
\[
 e(a)=e(b)=0,
\]
whereas $\tau$ acts trivially on $H^2(H,\mathbf F_p)$, so
\[
 e(a\smile b)=a\smile b\ne0=e(a)\smile e(b).
\]
Thus $e$ is not multiplicative, and the failure lies entirely in the cohomological range of the local Galois group.
\end{proof}

\begin{warning}[No strict multiplicative transfer projection]\label{warn:no-strict-multiplicative-transfer}
Proposition~\ref{prop:transfer-not-multiplicative} gives a local arithmetic counterexample in the exact cohomological range $H^1\times H^1\to H^2$.  It therefore rules out using the normalized transfer idempotent itself as a strict dga projection in the arithmetic Sylow construction.  The obstruction persists even though restriction identifies global cohomology with a direct summand and the image on cohomology is a subring.  A finite multiplicative all-Sylow chart must therefore retain coherent homotopies; strict multiplicativity cannot be imposed on the transfer projector as an additional axiom.
\end{warning}

The construction is coefficient-functorial: the same restricted-cochain sum and restriction--corestriction formulas apply to any perfect coefficient complex $\mathcal M$ with continuous $G_K$-action.

\begin{theorem}[Multiplicative all-Sylow closure]\label{thm:all-sylow-ainfty}
Let $\mathcal A$ be a perfect continuous coefficient dga in a fixed residual framed sector, equipped with a continuous $G_K$-action by dg-algebra automorphisms; for example
\[
 \mathcal A=\operatorname{End}_{R_{\bar\rho}^{\square}}(\mathcal T^{\square}).
\]
Let $\mathcal K_{\bar\rho}^{\mathrm{allSyl}}(\mathcal A)$ denote the perfect image of the all-Sylow transfer idempotent.  Then:
\begin{enumerate}[label=\textnormal{(\roman*)}]
\item the specified restriction--normalized-corestriction equivalence
\[
 u:\mathcal K_{\bar\rho}^{\mathrm{allSyl}}(\mathcal A)
 \xrightarrow{\ \sim\ }
 R\Gamma(G_K,\mathcal A)
\]
transports the intrinsic $E_1$-algebra structure on continuous cochains to an $E_1$-algebra structure on
$\mathcal K_{\bar\rho}^{\mathrm{allSyl}}(\mathcal A)$.
On every bounded finite-projective representative this transported object is represented by an $A_\infty$-algebra and an $A_\infty$ quasi-isomorphism
\[
 \mathcal K_{\bar\rho}^{\mathrm{allSyl}}(\mathcal A)
 \simeq_{A_\infty}
 C^\bullet_{\mathrm{cont}}(G_K,\mathcal A);
\]
\item to write chain-level formulas, choose a $K$-projective dg-algebra model $\mathcal B\to C^\bullet_{\mathrm{cont}}(G_K,\mathcal A)$, a bounded projective representative $\mathcal K$ of the all-Sylow image, and, after replacing these representatives by homotopy-equivalent models and adding contractible summands if necessary, contraction data
\[
 \left(
 \mathcal K
 \mathop{\rightleftarrows}^{i}_{p}
 \mathcal B,
 h
 \right),
 \qquad
 p i=\operatorname{id},
 \quad
 \operatorname{id}-i p=dh+hd.
\]
The first transferred operations are
\begin{align}
 m_2&=p\,\mu\,(i\otimes i),\label{eq:all-sylow-m2}\\
 m_3&=p\,\mu\bigl(h\mu\otimes1-1\otimes h\mu\bigr)i^{\otimes3},\label{eq:all-sylow-m3}
\end{align}
with the standard Koszul signs, and the higher $m_n$ are the usual finite sums over planar rooted trees;
\item different choices of dg replacements, splittings, comparison maps, and contracting homotopies produce $A_\infty$-isomorphic models compatible with the specified equivalence $u$.  The intrinsic transported $E_1$ object commutes with derived base change.  A particular rooted-tree formula need not commute strictly with base change unless the contraction data have themselves been chosen coherently;
\item on the Cayley--Hamilton splitting gerbe the adjoint $E_1$ algebra has scalar weight zero and descends Morita-covariantly.  The universal coefficient carrier is an $E_1$ module over this algebra and retains its natural gerbe weight; finite projective representatives yield corresponding $A_\infty$-module models.
\end{enumerate}
\end{theorem}

\begin{proof}
The continuous bar cochains form an $E_1$ algebra, represented by their usual cup-product dga.  The all-Sylow restriction map lands in the image of the averaged idempotent, and normalized corestriction restricts to its inverse.  Thus Theorem~\ref{thm:all-sylow-transfer} supplies the specified equivalence $u$ in the \emph{underlying} stable $\infty$-category of coefficient modules.  It is not asserted to be a monoidal map: Proposition~\ref{prop:transfer-not-multiplicative} excludes that interpretation for normalized transfer.  An algebra structure can nevertheless be transported along an equivalence of the underlying objects, by requiring $u$ to become an equivalence of $E_1$-algebras after the transport.  This gives the intrinsic $E_1$ structure in (i) and makes its derived base-change compatibility formal at the level of algebra objects.

For explicit formulas one must choose chain representatives.  Replace the continuous cochain dga by a $K$-projective, for example semifree, dg-algebra model and represent the perfect all-Sylow image by a bounded projective complex.  After the standard mapping-cylinder stabilization, the chosen derived equivalence is represented by contraction data as in (ii).  The homotopy transfer theorem then gives Equations~\eqref{eq:all-sylow-m2}--\eqref{eq:all-sylow-m3} and the higher rooted-tree operations \cite{BurkeTransfer,ManettiPerturbation,VokrinekHPT}.  Its comparison theorem gives an $A_\infty$ isomorphism between the structures arising from different replacements and contractions.  This proves the intrinsic assertion in (iii), while also explaining why strict chain-level base-change naturality requires a coherent choice of contractions rather than following from perfectness alone.

The scalar-weight statement is the Morita descent of Proposition~\ref{prop:all-sylow-morita}: endomorphism coefficients have weight zero, while the universal representation has weight one.  Transporting the algebra and module structures along the corresponding all-Sylow equivalences proves (iv).
\end{proof}

\begin{corollary}[Abstract Morita-multiplicative descent]\label{cor:abstract-morita-multiplicative}
The all-Sylow coefficient carrier has a canonical $E_1$-algebra/module structure relative to its specified equivalence with intrinsic continuous cochains, and this object descends to the residual representation stack and Morita-valuedly to the Cayley--Hamilton splitting gerbe.  Every bounded finite-projective all-Sylow representative therefore admits an $A_\infty$ model, while its rooted-tree formulas depend on dg replacements and contraction data.  The later common-source, presentation, Fitting, and globalization theorems make those choices recursively explicit and coherently comparable; no canonical rooted-tree representative is asserted or required.
\end{corollary}

\begin{remark}[Dyadic closure of the Sylow coefficient atlas]\label{rem:sylow-two}
The prime $2$ is now included in Theorems~\ref{thm:sylow-fox-retract}, \ref{thm:family-sylow-fox}, and \ref{thm:all-sylow-transfer}.  Nothing in the restriction--corestriction or residual-family argument requires odd residue characteristic: the Sylow index and the number of Sylow $2$-subgroups are odd and hence units in every $2$-adic coefficient ring.  The extra input is the standard dyadic Demu\v{s}kin presentation of the maximal pro-$2$ quotient, made explicit in Theorems~\ref{thm:dyadic-relator-cases} and~\ref{thm:dyadic-twisted-fox-orientation}.  Its different relation types give the finite Fox matrices of Proposition~\ref{prop:universal-formal-fox-matrices}, and Corollary~\ref{thm:dyadic-classification-certificate} supplies a faithful orientation-enriched symbol.  Finite Fox perfectness and derived base change use one-relator projectivity, not a universal full-$G_K$ word normal form.  The integral cyclotomic torsion boundary is separately recorded in Proposition~\ref{prop:dyadic-integral-herr-derived}.  The orientation-twisted cochain row and its transfer-normalized local Tate pairing are computed in Theorem~\ref{thm:dyadic-fox-kummer-tate} and Proposition~\ref{prop:dyadic-sylow-fox-tate-pairing}; the pairing does not treat normalized transfer as a multiplicative map.
\end{remark}

\begin{corollary}[Normal-chart $p$-primary comparison is derivedly closed]\label{cor:normal-p-primary-derived-comparison}
Keep a fixed normal extension
\[
1\to P_L\to Q_{L/K}\to\Delta\to1,
\qquad p\mid|\Delta|,
\]
and a Demu\v{s}kin Fox resolution of $P_L$.  For every perfect coefficient object on the normal chart, the Hochschild--Serre total complex, the Sylow--Fox transfer image, and any finite projective full-transport resolution represent the same intrinsic continuous cochain object.  Hence they are connected by a zigzag of quasi-isomorphisms, unique up to the usual chain homotopies after projective replacement.  The $p$-primary nonsemisimplicity of $\Delta$ obstructs treating normalized finite-quotient averaging as a strict multiplicative projection, but it creates no derived comparison or finite-atlas obstruction.
\end{corollary}

\begin{proof}
The Hochschild--Serre construction computes the continuous cochains of the extension.  The Sylow--Fox image is identified with the same intrinsic cochain object by Theorems~\ref{thm:sylow-fox-retract} and~\ref{thm:all-sylow-transfer}, while Theorem~\ref{thm:abstract-full-transport-resolution} computes it from a finite projective full-transport resolution.  Standard comparison of projective resolutions gives the stated zigzag and homotopy uniqueness.  Proposition~\ref{prop:transfer-not-multiplicative} records the separate strict-multiplicativity obstruction.
\end{proof}

Thus the former normal-chart comparison problem is no longer part of the mainline.  A closed factor-set formula for one chosen Hochschild--Serre representative remains a coordinate refinement only.  The next argument supplies the finite full-transport representative used by the atlas.

\subsection{Abstract integral full-transport perfectness}\label{subsec:abstract-full-transport}

Let
\[
\Lambda_{K,E}^{\circ}:=\mathcal O_E[[G_K]]
\]
and work in the abelian category of pseudocompact left $\Lambda_{K,E}^{\circ}$-modules.  Recall that a profinite group is of type $\mathrm{FP}_2$ over $\mathcal O_E$ when the trivial module admits a projective resolution whose terms in homological degrees $0,1,2$ are finitely generated \cite{CookFP}.

\begin{lemma}[Finite presentation gives $\mathrm{FP}_2$]\label{lem:finite-presentation-fp2}
The trivial pseudocompact $\Lambda_{K,E}^{\circ}$-module $\mathcal O_E$ is of type $\mathrm{FP}_2$.
\end{lemma}

\begin{proof}
Choose a finite profinite presentation of $G_K$ with $d$ generators and $m$ relations.  Completed cellular chains, or equivalently the first two stages of completed Fox calculus, give an exact partial resolution
\[
(\Lambda_{K,E}^{\circ})^m
\longrightarrow
(\Lambda_{K,E}^{\circ})^d
\longrightarrow
\Lambda_{K,E}^{\circ}
\longrightarrow
\mathcal O_E
\longrightarrow0.
\]
Exactness is asserted at the last three terms; no claim is made that the leftmost map is injective.  This is precisely the $\mathrm{FP}_2$ condition.  Finite presentability of $G_K$ is supplied by \cite{JannsenWingberg,JardenShusterman,NeukirchLocal}.
\end{proof}

\begin{theorem}[Abstract integral full-transport resolution]\label{thm:abstract-full-transport-resolution}
There exists an exact sequence of pseudocompact $\Lambda_{K,E}^{\circ}$-modules
\begin{equation}\label{eq:abstract-full-transport-resolution}
0\longrightarrow P_2^{\circ}
\longrightarrow P_1^{\circ}
\longrightarrow P_0^{\circ}
\longrightarrow\mathcal O_E
\longrightarrow0,
\end{equation}
where each $P_i^{\circ}$ is finitely generated projective.  Equivalently,
\[
\mathcal O_E\in\operatorname{Perf}(\Lambda_{K,E}^{\circ})
\]
has Tor amplitude contained in $[0,2]$.
\end{theorem}

\begin{proof}
By Lemma~\ref{lem:finite-presentation-fp2}, choose a projective resolution with finitely generated terms through degree two.  Put
\[
Z_1:=\ker(P_1^{\circ}\to P_0^{\circ}).
\]
Brumer's comparison theorem for pseudocompact group algebras identifies the projective dimension of the trivial compact module with the $p$-cohomological dimension: after extension of coefficients from $\mathbf Z_p$ to $\mathcal O_E$, one has
\[
\operatorname{pd}_{\mathcal O_E[[G_K]]}(\mathcal O_E)
=
\operatorname{cd}_p(G_K)
=2.
\]
See \cite{Brumer,NSW}.  Hence the second syzygy $Z_1$ is projective.  It is finitely generated because it is the image of the finitely generated degree-two term in the chosen $\mathrm{FP}_2$ resolution.  Replacing that term by $Z_1$ gives \eqref{eq:abstract-full-transport-resolution}.
\end{proof}

\begin{theorem}[Projective row reduction of a full completed Fox matrix]\label{thm:full-fox-projector}
Choose a finite profinite presentation of $G_K$ with $d$ generators and $m$
relations, and write the completed Fox partial resolution as
\begin{equation}\label{eq:full-fox-partial-resolution}
 (\Lambda_{K,E}^{\circ})^m
 \xrightarrow{d_2^{\mathrm{Fox}}}
 (\Lambda_{K,E}^{\circ})^d
 \xrightarrow{d_1^{\mathrm{Fox}}}
 \Lambda_{K,E}^{\circ}
 \longrightarrow\mathcal O_E\longrightarrow0.
\end{equation}
Put $Z_1:=\ker(d_1^{\mathrm{Fox}})=\operatorname{im}(d_2^{\mathrm{Fox}})$.
Then there exists an idempotent
\[
 e_{\mathrm{Fox}}\in
 M_m(\Lambda_{K,E}^{\circ})
\]
such that
\begin{equation}\label{eq:fox-projector-identities}
 d_2^{\mathrm{Fox}}e_{\mathrm{Fox}}=d_2^{\mathrm{Fox}},
 \qquad
 \ker(d_2^{\mathrm{Fox}})
 =(1-e_{\mathrm{Fox}})(\Lambda_{K,E}^{\circ})^m,
\end{equation}
and restriction induces an isomorphism
\[
 d_2^{\mathrm{Fox}}:
 e_{\mathrm{Fox}}(\Lambda_{K,E}^{\circ})^m
 \xrightarrow{\sim} Z_1.
\]
Consequently
\begin{equation}\label{eq:reduced-full-fox-resolution}
0\longrightarrow
 e_{\mathrm{Fox}}(\Lambda_{K,E}^{\circ})^m
 \xrightarrow{d_2^{\mathrm{Fox}}}
 (\Lambda_{K,E}^{\circ})^d
 \xrightarrow{d_1^{\mathrm{Fox}}}
 \Lambda_{K,E}^{\circ}
 \longrightarrow\mathcal O_E\longrightarrow0
\end{equation}
is a length-two finite projective full-transport resolution.

For any finite projective coefficient module $\mathcal W$, dualizing
\eqref{eq:reduced-full-fox-resolution} gives a three-term cochain complex in
which the obstruction term is the direct summand cut out by the dual
idempotent $e_{\mathrm{Fox}}^{\vee}$ on $\mathcal W^m$.  Thus every finite presentation admits a genuine projective Fox row reduction.  Computing a particularly useful matrix representative $e_{\mathrm{Fox}}$ in marked or Cayley--Hamilton coordinates is only an optional local gauge task: the projector-free syzygy-augmented resolution below is already canonical enough for the invariant construction.
\end{theorem}

\begin{proof}
By Theorem~\ref{thm:abstract-full-transport-resolution}, the syzygy $Z_1$ is a
finitely generated projective $\Lambda_{K,E}^{\circ}$-module.  The Fox map
$d_2^{\mathrm{Fox}}:(\Lambda_{K,E}^{\circ})^m\twoheadrightarrow Z_1$ is
surjective, hence admits a continuous module section
\[
 s:Z_1\longrightarrow(\Lambda_{K,E}^{\circ})^m.
\]
Set
\[
 e_{\mathrm{Fox}}:=s\,d_2^{\mathrm{Fox}}.
\]
Then $e_{\mathrm{Fox}}^2=e_{\mathrm{Fox}}$ and
$d_2^{\mathrm{Fox}}e_{\mathrm{Fox}}=d_2^{\mathrm{Fox}}$.  If
$e_{\mathrm{Fox}}x=0$, then $s(d_2^{\mathrm{Fox}}x)=0$; injectivity of the
section gives $d_2^{\mathrm{Fox}}x=0$.  The reverse inclusion is immediate,
so the kernel identity in \eqref{eq:fox-projector-identities} follows.  The
restriction of $d_2^{\mathrm{Fox}}$ to the image of $e_{\mathrm{Fox}}$ is
therefore bijective onto $Z_1$.  This proves exactness of
\eqref{eq:reduced-full-fox-resolution}.  Dualization identifies the final
cochain term with the corresponding direct summand of $\mathcal W^m$.
\end{proof}

\begin{warning}[The full Fox projector is not canonical]\label{warn:fox-projector-noncanonical}
The idempotent $e_{\mathrm{Fox}}$ depends on a splitting of the relation map.
Different choices yield isomorphic projective syzygies and chain-homotopy
equivalent full-transport resolutions, but not a preferred matrix idempotent.
Proposition~\ref{prop:no-functorial-projection} rules out expecting a globally functorial strict choice.  The invariant object is the projective syzygy $Z_1$ and its class in the perfect derived category.  A useful projector may still be chosen or computed in marked or Cayley--Hamilton coordinates, but this is a chart gauge rather than a missing invariant datum.
\end{warning}

\begin{theorem}[Projector-free syzygy-augmented Fox resolution]\label{thm:syzygy-augmented-fox}
In the setting of Theorem~\ref{thm:full-fox-projector}, put
\[
 K_2^{\mathrm{Fox}}:=\ker(d_2^{\mathrm{Fox}}).
\]
Then $K_2^{\mathrm{Fox}}$ is a finitely generated projective
$\Lambda_{K,E}^{\circ}$-module, and there is a canonical exact sequence,
relative only to the chosen profinite presentation,
\begin{equation}\label{eq:syzygy-augmented-fox-resolution}
0\longrightarrow K_2^{\mathrm{Fox}}
\longrightarrow(\Lambda_{K,E}^{\circ})^m
\xrightarrow{d_2^{\mathrm{Fox}}}
(\Lambda_{K,E}^{\circ})^d
\xrightarrow{d_1^{\mathrm{Fox}}}
\Lambda_{K,E}^{\circ}
\longrightarrow\mathcal O_E\longrightarrow0.
\end{equation}
No splitting or idempotent is required.

For every finite projective coefficient module $\mathcal W$, dualizing
\eqref{eq:syzygy-augmented-fox-resolution} gives the four-term finite
projective cochain complex
\begin{equation}\label{eq:syzygy-augmented-fox-cochain}
\mathcal W
\longrightarrow\mathcal W^d
\longrightarrow\mathcal W^m
\longrightarrow
\operatorname{Hom}^{\mathrm{cont}}_{\Lambda_{K,E}^{\circ}}
(K_2^{\mathrm{Fox}},\mathcal W),
\end{equation}
which computes $R\Gamma_{\mathrm{cont}}(G_K,\mathcal W)$.  Its degree-three
cohomology vanishes.  Choosing a section of $d_2^{\mathrm{Fox}}$ decomposes
\eqref{eq:syzygy-augmented-fox-resolution} as the reduced length-two Fox
resolution of Theorem~\ref{thm:full-fox-projector} plus a contractible
summand.  Thus the projector model is a chosen reduced gauge of the canonical
nonminimal syzygy-augmented model.
\end{theorem}

\begin{proof}
Write $Z_1=\operatorname{im}(d_2^{\mathrm{Fox}})$.  By
Theorem~\ref{thm:abstract-full-transport-resolution}, $Z_1$ is finitely
generated projective.  Hence the exact sequence
\[
0\longrightarrow K_2^{\mathrm{Fox}}
\longrightarrow(\Lambda_{K,E}^{\circ})^m
\longrightarrow Z_1\longrightarrow0
\]
splits, so $K_2^{\mathrm{Fox}}$ is a finitely generated projective direct
summand of a finite free module.  Exactness of
\eqref{eq:syzygy-augmented-fox-resolution} is then immediate from the Fox
partial resolution.  Since all terms are projective, applying continuous
$\operatorname{Hom}(-,\mathcal W)$ computes derived invariants and gives
\eqref{eq:syzygy-augmented-fox-cochain}.  A chosen splitting identifies
$(\Lambda_{K,E}^{\circ})^m$ with
$K_2^{\mathrm{Fox}}\oplus e_{\mathrm{Fox}}(\Lambda_{K,E}^{\circ})^m$.
Under this decomposition the first two terms contain the identity complex
$K_2^{\mathrm{Fox}}\xrightarrow{\mathrm{id}}K_2^{\mathrm{Fox}}$, which is
contractible, and the remaining summand is exactly
\eqref{eq:reduced-full-fox-resolution}.
\end{proof}

\begin{proposition}[No functorial Fox projector]\label{prop:no-functorial-fox-projector}
A projector onto the projective image of a surjection cannot in general be
chosen functorially from the surjection.  In particular,
$e_{\mathrm{Fox}}$ cannot be an intrinsic functor of the completed Fox
relation matrix alone.
\end{proposition}

\begin{proof}
Let $R$ be a nonzero ring and let $q:R^2\to R$ be $q(x,y)=x$.  Its splittings
are
\[
s_a(1)=(1,a),\qquad a\in R.
\]
The automorphism
\[
u(x,y)=(x,y+x)
\]
preserves $q$ and sends $s_a$ to $s_{a+1}$.  Hence no splitting is invariant
under all automorphisms of the surjection.  The associated idempotents are
therefore equally nonfunctorial.  Applying this elementary obstruction to a
split relation map shows that a canonical matrix $e_{\mathrm{Fox}}$ cannot
be extracted from $d_2^{\mathrm{Fox}}$ without additional gauge data.
\end{proof}

\begin{warning}[The projector problem is a gauge problem]\label{warn:fox-projector-gauge}
Computing a canonical full Fox projector is therefore not an intrinsic mathematical target.  A useful projector may be
chosen on a marked, determinantal, or Cayley--Hamilton chart, but it is only
a local frame for the projective module $K_2^{\mathrm{Fox}}$ or its
complement.  The invariant object is the syzygy-augmented resolution
\eqref{eq:syzygy-augmented-fox-resolution}.
\end{warning}

\begin{theorem}[Pro-$p$ Fox syzygies are free and admit column-reduced frames]\label{thm:pro-p-fox-syzygy-free}
Let $P$ be a pro-$p$ group and put
\[
 \Lambda_P^{\circ}:=\mathcal O_E[[P]],
 \qquad
 \mathfrak m_P:=(\varpi_E,\,g-1\mid g\in P).
\]
Then $\Lambda_P^{\circ}$ is a pseudocompact local ring with Jacobson
radical $\mathfrak m_P$, and every finitely generated projective
pseudocompact $\Lambda_P^{\circ}$-module is finite free
\cite{Brumer}.

Suppose a finite pro-$p$ presentation gives a completed Fox map
\[
 d_2^{\mathrm{Fox}}:(\Lambda_P^{\circ})^m
 \longrightarrow(\Lambda_P^{\circ})^d
\]
whose image $Z_1$ is finitely generated projective, and put
$K_2^{\mathrm{Fox}}:=\ker(d_2^{\mathrm{Fox}})$.  Then there are integers
$s,t\ge0$, with $s+t=m$, an injective matrix
\[
 B:(\Lambda_P^{\circ})^s\hookrightarrow(\Lambda_P^{\circ})^d,
\]
and a change of relation coordinates
\[
 U\in\mathrm{GL}_m(\Lambda_P^{\circ})
\]
such that
\begin{equation}\label{eq:pro-p-column-reduction}
 d_2^{\mathrm{Fox}}U=(B\ \ 0),
 \qquad
 K_2^{\mathrm{Fox}}
 =U\bigl(0\oplus(\Lambda_P^{\circ})^t\bigr).
\end{equation}
In particular, $Z_1\cong(\Lambda_P^{\circ})^s$ and
$K_2^{\mathrm{Fox}}\cong(\Lambda_P^{\circ})^t$.

If $U'$ is a second reduction, then on the two decompositions the
transition matrix has parabolic form
\begin{equation}\label{eq:pro-p-frame-transition}
 U^{-1}U'
 =
 \begin{pmatrix}
  a&0\\ c&d
 \end{pmatrix},
 \qquad
 a\in\mathrm{GL}_s(\Lambda_P^{\circ}),
 \quad
 d\in\mathrm{GL}_t(\Lambda_P^{\circ}),
\end{equation}
and $d$ is precisely the transition matrix between the two syzygy
frames.  For every cofinal system of open two-sided ideals
$I_\nu\subset\mathfrak m_P$, the frame torsor satisfies
\begin{equation}\label{eq:pro-p-frame-limit}
 \operatorname{Isom}_{\Lambda_P^{\circ}}
 \bigl((\Lambda_P^{\circ})^t,K_2^{\mathrm{Fox}}\bigr)
 \cong
 \varprojlim_{\nu}
 \operatorname{Isom}_{\Lambda_P^{\circ}/I_\nu}
 \bigl((\Lambda_P^{\circ}/I_\nu)^t,
       K_2^{\mathrm{Fox}}/I_\nu K_2^{\mathrm{Fox}}\bigr).
\end{equation}
Thus a syzygy frame can be approximated, to arbitrary $p$-adic and
profinite precision, by compatible finite-quotient column reductions.
\end{theorem}

\begin{proof}
Brumer's pseudocompact local-ring theorem identifies
$\mathfrak m_P$ as the Jacobson radical of $\Lambda_P^{\circ}$.
Let $M$ be a finitely generated projective pseudocompact
$\Lambda_P^{\circ}$-module.  Choose a basis of
$M/\mathfrak m_PM$ over the residue field and lift it to $M$.
Topological Nakayama gives a surjection
$(\Lambda_P^{\circ})^r\twoheadrightarrow M$.  Since $M$ is projective,
the map splits.  Its kernel is finitely generated projective and has
zero reduction modulo $\mathfrak m_P$; topological Nakayama therefore
kills the kernel.  Hence $M$ is free.

Apply this to $Z_1$ and $K_2^{\mathrm{Fox}}$.  The exact sequence
\[
0\longrightarrow K_2^{\mathrm{Fox}}
\longrightarrow(\Lambda_P^{\circ})^m
\longrightarrow Z_1\longrightarrow0
\]
splits.  Choose bases of the two free summands and a section of the
surjection onto $Z_1$.  The resulting direct-sum isomorphism is the
matrix $U$, and in these coordinates the Fox map has the form
\eqref{eq:pro-p-column-reduction}.  A second direct-sum frame preserves
the standard kernel $0\oplus(\Lambda_P^{\circ})^t$; this is equivalent
to the parabolic form \eqref{eq:pro-p-frame-transition}.

Finally, finite projective pseudocompact modules and their continuous
homomorphisms are recovered from all finite quotients.  Completeness gives
injectivity in \eqref{eq:pro-p-frame-limit}; compatible quotient frames
lift to a continuous isomorphism, and invertibility can be tested modulo
the Jacobson radical.  This proves surjectivity.
\end{proof}

\begin{corollary}[Free local frames on Sylow--Fox charts]\label{cor:sylow-fox-free-frames}
Every free or Demu\v{s}kin Sylow--Fox chart used in the potential and
family-level closure theorems admits a global free syzygy frame over its
pro-$p$ completed group algebra.  For the standard free and minimal
Demu\v{s}kin resolutions one has $K_2^{\mathrm{Fox}}=0$; for a
nonminimal pro-$p$ presentation, Theorem~\ref{thm:pro-p-fox-syzygy-free}
provides a finite free redundant-syzygy block and explicit parabolic
transition matrices.  Hence the existence of local frames is not an open
problem in the pro-$p$ charts.
\end{corollary}

\begin{warning}[Full-group projectives need not admit a single frame]\label{warn:full-group-not-local}
The conclusion above uses that $P$ is pro-$p$.  The full completed algebra
$\mathcal O_E[[G_K]]$ is not local in general, and a finitely generated
projective full-transport module need not be free.  Therefore the global
object should not be forced into one matrix frame.  Its correct coordinates
are blockwise or Morita-valued, compatible with the residual determinant
and Cayley--Hamilton decompositions.  The pro-$p$ theorem solves the local
frame problem after Sylow reduction; it does not manufacture a canonical
full-$G_K$ basis.
\end{warning}

\begin{remark}[Non-rigid flag descent: final status]\label{rem:nonrigid-effective-descent-closed}
The non-rigid flag branch is already effective at the level required by the native atlas.  Theorems~\ref{thm:nonrigid-flag-hdescent} and~\ref{thm:nonrigid-strict-local} give derived descent and bounded strict local representatives; Theorems~\ref{thm:nonrigid-finite-stage-retract} and~\ref{thm:schur-koszul-flag-dg} reduce the descended carrier to a finite Schur--Koszul Cech totalization; and Theorems~\ref{thm:linear-descent-retraction} and~\ref{thm:descent-exponent-stratification} turn the first working descent stage and its idempotent into terminating finite matrix problems with a finite bound on every fixed quasi-compact chart.  Theorem~\ref{thm:no-uniform-descent-exponent-bound} shows that one cannot strengthen this to a bound depending only on the representation rank, flag type, and number of group generators under arbitrary nonreduced coefficient base change.  Thus there is no remaining descent-existence problem here; a prettier closed solver on a chosen GMA chart is only an optional coordinate optimization.
\end{remark}

\begin{corollary}[Integral coefficient closure and base change]\label{cor:integral-full-transport-coeff}
Let $A$ be a complete $\mathcal O_E$-flat coefficient algebra, let
\[
\Lambda_{A,K}:=A\widehat\otimes_{\mathcal O_E}\Lambda_{K,E}^{\circ}=A[[G_K]],
\]
and let $\mathcal W$ be a finite projective $A$-module with continuous $G_K$-action.  Define
\begin{equation}\label{eq:abstract-ft-cochain}
C_{\mathrm{FT}}^{\bullet}(G_K,\mathcal W)
:=
\operatorname{Hom}^{\mathrm{cont}}_{\Lambda_{A,K}}
\bigl(A\widehat\otimes_{\mathcal O_E}P_{\bullet}^{\circ},\mathcal W\bigr).
\end{equation}
Then $C_{\mathrm{FT}}^{\bullet}(G_K,\mathcal W)$ is a bounded complex of finite projective $A$-modules in degrees $0,1,2$, and there is a natural quasi-isomorphism
\[
C_{\mathrm{FT}}^{\bullet}(G_K,\mathcal W)
\simeq
R\Gamma_{\mathrm{cont}}(G_K,\mathcal W).
\]
For a morphism $A\to B$ of complete $\mathcal O_E$-flat coefficient algebras,
\[
C_{\mathrm{FT}}^{\bullet}(G_K,\mathcal W)
\widehat\otimes_A B
\cong
C_{\mathrm{FT}}^{\bullet}
\bigl(G_K,\mathcal W\widehat\otimes_A B\bigr)
\]
term by term.  Torsion or nonflat specializations are obtained by derived base change.
\end{corollary}

\begin{proof}
Completed base change of \eqref{eq:abstract-full-transport-resolution} remains exact because all syzygies are $\mathcal O_E$-torsion free, hence flat over the discrete valuation ring $\mathcal O_E$.  The resulting complex is a finite projective resolution of the trivial $A[[G_K]]$-module $A$.  Applying continuous $\operatorname{Hom}_{A[[G_K]]}(-,\mathcal W)$ computes continuous derived invariants.  Finite projectivity of the $P_i^{\circ}$ and of $\mathcal W$ gives finite projectivity of every term and the displayed base-change identity.
\end{proof}

\subsection{Abstract multiplicative enhancement}\label{subsec:abstract-multiplicative}

Perfectness also gives multiplicative structure abstractly; only explicit and covariant formulas remain.

\begin{theorem}[Diagonal approximation and transferred higher operations]\label{thm:abstract-diagonal}
Let $P_\bullet^{\circ}\to\mathcal O_E$ be a finite projective resolution as in Theorem~\ref{thm:abstract-full-transport-resolution}.  Give
\[
P_\bullet^{\circ}\widehat\otimes_{\mathcal O_E}P_\bullet^{\circ}
\]
the diagonal $G_K$-action.  Then:
\begin{enumerate}
\item the total completed tensor product is a projective pseudocompact $\Lambda_{K,E}^{\circ}$-resolution of $\mathcal O_E$;
\item the diagonal map $1\mapsto1\otimes1$ on $\mathcal O_E$ lifts to a continuous $\Lambda_{K,E}^{\circ}$-linear chain map
\begin{equation}\label{eq:abstract-diagonal}
\Delta_P:P_\bullet^{\circ}
\longrightarrow
P_\bullet^{\circ}\widehat\otimes_{\mathcal O_E}P_\bullet^{\circ},
\end{equation}
unique up to continuous chain homotopy;
\item for continuous coefficient modules $M,N,Q$ and a continuous $G_K$-equivariant pairing $\mu:M\widehat\otimes N\to Q$, the formula
\[
(f\smile_P g):=\mu\circ(f\widehat\otimes g)\circ\Delta_P
\]
defines a cochain-level cup product whose induced product on continuous cohomology is canonical;
\item after choosing a continuous chain-homotopy equivalence between the normalized continuous bar resolution and $P_\bullet^{\circ}$, homotopy transfer equips finite cochain complexes with coefficients in an algebra object with integral $A_\infty$-enhancements quasi-isomorphic to the usual bar-cochain structure.  After inverting $\varpi_E$, adjoint coefficients also inherit the corresponding $L_\infty$ enhancement.  Different choices determine equivalent objects in the relevant infinity-homotopy category.
\end{enumerate}
\end{theorem}

\begin{proof}
Every $P_i^{\circ}$ is a direct summand of a finite free pseudocompact $\Lambda_{K,E}^{\circ}$-module and is $\mathcal O_E$-flat.  With diagonal action, the completed tensor of a free transport module with an $\mathcal O_E$-projective pseudocompact module is topologically induced and hence projective; direct summands remain projective.  Since $P_\bullet^{\circ}$ is bounded and $\mathcal O_E$-flat, its completed tensor square is a resolution of $\mathcal O_E$.

The comparison theorem for projective resolutions lifts the diagonal of the trivial module to \eqref{eq:abstract-diagonal}, and any two lifts are chain homotopic.  The standard diagonal-approximation construction then gives the cup product and proves independence on cohomology \cite{BrownCohomology}.  Finally, choose comparison maps and a homotopy between the bar resolution and $P_\bullet^{\circ}$.  Homological perturbation transfers the integral bar cup product and its coherences to an $A_\infty$ structure on the finite model; after passing to $E$, it also transfers the cup--commutator deformation Lie structure.  Existence and homotopy invariance follow from the homotopy-transfer theorem over projective resolutions \cite{ManettiPerturbation,BurkeTransfer,VokrinekHPT}.
\end{proof}

\begin{corollary}[Abstract multiplicative full-transport closure]\label{cor:abstract-multiplicative-closure}
The finite full-transport complex of Corollary~\ref{cor:integral-full-transport-coeff} admits, after choices, a diagonal approximation and an integral $A_\infty$ multiplicative enhancement recovering cup products and Massey products; after passage to $E$, it also recovers brackets and the intrinsic deformation $L_\infty$ structure.  Thus neither binary multiplicativity nor abstract higher coherence is an existence problem.
\end{corollary}

\begin{theorem}[Infinite transport rank of the naive universal tensor diagonal]\label{thm:infinite-diagonal-transport-rank}
Let $A$ be a nonzero complete Noetherian coefficient ring, let $G$ be an
infinite profinite group, and put $\Lambda=A[[G]]$.  For $r\ge2$, give
\[
M_r:=\widehat{\bigotimes}_{A}^{\,r}\Lambda
\]
the diagonal left $G$-action.  Then $M_r$ is not finitely generated as a
pseudocompact left $\Lambda$-module.

More precisely, for every open normal subgroup $U\triangleleft G$, put
$H=G/U$.  The finite quotient of $M_r$ obtained by reducing every tensor
factor through $A[H]$ is, with diagonal $H$-action, a free left $A[H]$-module
of rank $|H|^{r-1}$:
\begin{equation}\label{eq:finite-quotient-diagonal-rank}
\bigotimes_A^{\,r}A[H]
\cong
A[H]^{\oplus |H|^{r-1}}.
\end{equation}
Consequently the generator number of the universal diagonal tensor target is
unbounded along finite quotients of $G$.
\end{theorem}

\begin{proof}
For $r=2$, define
\[
\Phi_H:A[H]\otimes_A A[H]\longrightarrow A[H]\otimes_A A[H],
\qquad
\Phi_H(g\otimes h)=g\otimes g^{-1}h.
\]
On the source let $H$ act diagonally; on the target let it act only on the
first tensor factor.  Then
\[
\Phi_H\bigl(xg\otimes xh\bigr)
=xg\otimes g^{-1}h
=x\,\Phi_H(g\otimes h),
\]
so $\Phi_H$ is $A[H]$-linear, with inverse
$g\otimes k\mapsto g\otimes gk$.  Hence the diagonal tensor square is free of
rank $|H|$ over $A[H]$.  Iterating the same change of variables gives
\eqref{eq:finite-quotient-diagonal-rank} for every $r$.

Suppose $M_r$ were generated by $N$ elements over $\Lambda$.  Every finite
quotient in \eqref{eq:finite-quotient-diagonal-rank} would then be generated by
at most $N$ elements over $A[H]$.  After reduction modulo a maximal ideal of
$A$, dimension over the residue field shows that the free
$k[H]$-module $k[H]^{\oplus |H|^{r-1}}$ requires at least
$|H|^{r-1}$ generators.  Since an infinite profinite group has finite
quotients of unbounded order, this is impossible.
\end{proof}

\begin{corollary}[The universal diagonal is projective but not a finite transport matrix]\label{cor:universal-diagonal-not-finite-matrix}
Take the standard normalization of a full-transport resolution with
$P_0=\Lambda$.  Then the degree-zero term of
$P_\bullet\widehat\otimes_A P_\bullet$ is
$\Lambda\widehat\otimes_A\Lambda$.  The coefficient-base-changed form of
Theorem~\ref{thm:abstract-diagonal} lifts the cup diagonal into a projective
pseudocompact resolution whose tensor target is already infinitely generated
in degree zero.  Therefore abstract diagonal approximation cannot, by itself,
be converted into a finite matrix problem over $A[[G]]$.

After coefficient evaluation and finite contraction, every fixed $A_\infty/L_\infty$ operation is a finite matrix.  The obstruction is specifically to a naive
coefficient-independent finite tensor-square transport resolution.
\end{corollary}

\begin{warning}[Multiplicative existence is not an explicit Cartan chart]\label{warn:abstract-resolution-not-explicit}
Theorems~\ref{thm:abstract-full-transport-resolution} and~\ref{thm:abstract-diagonal} establish multiplicative existence but do not choose a finite presentation of $G_K$, a basis of $P_2^{\circ}$, a Fox Jacobian, or explicit formulas for $\Delta_P$ and the transferred higher operations.  Theorem~\ref{thm:infinite-diagonal-transport-rank} explains why no naive strict Hopf-type replacement can supply these data: the universal tensor target is not finitely generated over the completed group algebra.  Moreover, the diagonal and contraction are noncanonical, and Proposition~\ref{prop:no-functorial-projection} excludes a globally natural strict choice.

At this stage the missing ingredient is explicit coefficient-level comparison with the finite Cayley--Hamilton/Morita tensor carrier, not existence of an integral multiplicative full-transport complex.  The later Cayley--Hamilton evaluation and recursive-globalization theorems supply precisely that finite comparison, so this warning is a boundary on naive universal tensor models rather than a final atlas gap.
\end{warning}

\subsection{Comparison boundary}

The finite complex computes cohomology and supports transferred products, but it does not by itself present a preferred computable full-$G_K$ dg algebra or dg Lie model.  A three-term complex may still carry infinitely many nonzero transferred operations \cite{WangErickson}.  We therefore pass to the derived representation stack, where the intrinsic deformation object is canonical; the later finite-atlas theorems realize it by explicit finite gauges without imposing a universal strict dg normal form.

\section{Derived representation stacks}\label{sec:derived-envelope}

\subsection{The formal moduli problem}

For a continuous representation $\rho:G\to\mathrm{GL}_E(V)$, let
\[
\mathfrak{Rep}_{n,E}(G):=\operatorname{Map}_{\mathrm{cont}}(BG,B\mathrm{GL}_{n,E})
\]
denote the derived representation problem, interpreted formally at $\rho$ when no global analytic stack is fixed \cite{GalatiusVenkatesh,AntonioPadic,AntonioProEtale}.

\begin{theorem}[Intrinsic derived Cartan envelope]\label{thm:derived-envelope}
The formal derived completion of $\mathfrak{Rep}_{n,E}(G)$ at $\rho$ has tangent complex
\begin{equation}\label{eq:derived-tangent}
\mathbb T_{\mathfrak{Rep}_{n,E}(G),\rho}
\simeq
C^\bullet_{\mathrm{cont}}(G,\operatorname{ad}\rho)[1].
\end{equation}
Its formal deformation problem is governed by the cup--commutator differential graded Lie algebra on $C^\bullet_{\mathrm{cont}}(G,\operatorname{ad}\rho)$.  In particular, the Maurer--Cartan equation of Section~\ref{sec:deformation} is not an added structure: it is the intrinsic formal geometry of the derived representation stack.
\end{theorem}

\begin{proof}
The tangent of a derived mapping stack at $f:X\to Y$ is $R\Gamma(X,f^*\mathbb T_Y)$.  With $X=BG$ and $Y=B\mathrm{GL}_n$, this is the shifted adjoint cochain complex in \eqref{eq:derived-tangent}; its cup--commutator bracket is the standard controlling Lie structure \cite{GalatiusVenkatesh,AntonioPadic}.
\end{proof}

Framing removes the infinitesimal conjugation term.  In the unframed complex, $H^0$, $H^1$, and $H^2$ record stabilizers, tangents, and obstructions.

\subsection{Naive relation charts and the syzygy correction}

For a finite profinite or pro-$p$ presentation
\[
P=\langle x_1,\ldots,x_d\mid r_1,\ldots,r_m\rangle,
\]
the relation words define $\mathbf r:\mathrm{GL}(V)^d\to\mathrm{GL}(V)^m$.  Its ordinary fiber is the framed representation locus; its homotopy fiber is only a presentation-dependent naive derived chart.

\begin{warning}[Redundant-relation derived no-go]\label{warn:redundant-derived}
The presentations $\langle x\mid x\rangle$ and $\langle x\mid x,x\rangle$ define the same trivial group, but their naive tangent complexes are respectively
\[
\mathfrak{gl}(V)\xrightarrow{\mathrm{id}}\mathfrak{gl}(V),
\qquad
\mathfrak{gl}(V)\xrightarrow{(\mathrm{id},\mathrm{id})}\mathfrak{gl}(V)^2.
\]
The second has a spurious cokernel.  Hence displayed relators alone do not define an intrinsic derived zero locus; higher syzygies are essential.
\end{warning}

Derived representation schemes are therefore defined by deriving the representation functor on a cofibrant simplicial or semi-free resolution, rather than by applying a Koszul construction to an arbitrary list of relators \cite{BKRDerived,BKRCyclic}.  The intrinsic formal object remains Theorem~\ref{thm:derived-envelope}.

\begin{definition}[Cohomologically Fox-complete presentation]\label{def:fox-complete}
A finite transport presentation is \emph{cohomologically Fox-complete over $E$ at $\rho$} if its completed cellular/Fox complex is a finite projective resolution through the full amplitude of $C^\bullet_{\mathrm{cont}}(P,\operatorname{ad}\rho)$.  This is a statement about the intrinsic tangent/cohomology complex; it does not by itself assert that the displayed relators form a cofibrant presentation of the full derived representation functor.
\end{definition}

Free pro-$p$ groups and Demu\v{s}kin groups have cohomologically Fox-complete presentations: the former have cohomological dimension one, while the latter possess the length-two completed Fox resolution of Theorem~\ref{thm:fox-spencer-demushkin}.

\begin{theorem}[Intrinsic Fox tangent in the cohomologically complete sector]\label{thm:intrinsic-fox-chart}
For any finite presentation, the tangent complex of its naive framed relation-derived fiber is
\begin{equation}\label{eq:framed-fox-tangent}
(\operatorname{ad}V)^d
\xrightarrow{\ J^{\mathrm{Fox}}_\rho\ }
(\operatorname{ad}V)^m,
\end{equation}
where $J^{\mathrm{Fox}}_\rho$ is the evaluated Fox Jacobian.  After quotienting by conjugation, the tangent complex is, in degrees $[-1,0,1]$,
\begin{equation}\label{eq:unframed-fox-tangent}
\underset{-1}{\operatorname{ad}V}
\xrightarrow{d^0_\rho}
\underset{0}{(\operatorname{ad}V)^d}
\xrightarrow{J^{\mathrm{Fox}}_\rho}
\underset{1}{(\operatorname{ad}V)^m}.
\end{equation}
If the presentation is cohomologically Fox-complete at $\rho$, then \eqref{eq:unframed-fox-tangent} is quasi-isomorphic to the intrinsic tangent complex
\[
C^\bullet_{\mathrm{cont}}(P,\operatorname{ad}\rho)[1].
\]
For a Demu\v{s}kin presentation, $m=1$ and this is exactly the Fox--Spencer complex of Theorem~\ref{thm:fox-spencer-demushkin} with coefficients in $\operatorname{ad}V$.
\end{theorem}

\begin{proof}
Fox differentiation computes the linearization of the word maps, and the Fox fundamental identity gives $J^{\mathrm{Fox}}_\rho d^0_\rho=0$.  Cohomological Fox-completeness is precisely the additional hypothesis that this finite cellular complex computes the intrinsic continuous cochains in the relevant amplitude.  The free and Demu\v{s}kin cases follow from their explicit finite resolutions.
\end{proof}

\begin{remark}[Tangent completeness is not full derived completeness]\label{rem:tangent-vs-derived-complete}
A Fox quasi-isomorphism identifies tangent objects, not automatically formal moduli problems.  Full equivalence requires a cofibrant transport model, or a quasi-isomorphism of the controlling dg- or $L_\infty$ algebras, including nonlinear higher syzygies.
\end{remark}

\begin{corollary}[Presentation independence after derived syzygy completion]\label{cor:derived-presentation-independence}
Any two cofibrant resolutions of the same profinite group determine equivalent formal derived representation stacks.  If two finite presentations are cohomologically Fox-complete, their Fox--Spencer tangent complexes model the same intrinsic tangent object
\[
C^\bullet_{\mathrm{cont}}(P,\operatorname{ad}\rho)[1].
\]
For arbitrary finite presentations, full presentation independence holds only after adjoining the necessary higher syzygies or replacing the presentation by a cofibrant simplicial resolution.
\end{corollary}

\begin{remark}[Finite presentation versus finite derived presentation]\label{rem:finite-vs-derived-presentation}
Finite generators and relators control degree zero.  A finite derived chart must also control their higher syzygies.  The free and Demu\v{s}kin resolutions do so; an abstract finite presentation of $G_K$ need not.
\end{remark}

\subsection{Finite recursive control}

\begin{proposition}[Finite recursive source from a finite cofibrant transport model]\label{prop:finite-recursive-source}
Suppose the formal representation problem near $\rho$ is presented by a finite cofibrant simplicial transport model, or equivalently by finitely many generators in each of finitely many homological degrees with differential given by finite transport words.  Then every Taylor coefficient and every transferred $A_\infty/L_\infty$ operation is recursively determined by that finite syntax.  A finite cohomological Fox resolution supplies the tangent complex, but a finite all-order source additionally requires such a cofibrant enhancement or an explicit finite recursive contraction from the bar construction.
\end{proposition}

\begin{proof}
Applying the derived representation functor to the finite cofibrant model gives a finite semi-free coordinate dga.  Formal expansion and homotopy transfer then generate all Taylor coefficients recursively.  Without the syzygy data, Warning~\ref{warn:redundant-derived} shows that the relators control only the classical locus and its naive Koszul enhancement.
\end{proof}

\begin{warning}[Finite recursion does not imply formality]\label{warn:no-formality}
A finite recursive source need not kill higher operations.  Demu\v{s}kin cochains already show $q$-dependent $A_3$-formality behavior \cite{PalQuickOdd,PalQuickTwo}.  Those results concern mod-$p$ trivial coefficients, not the characteristic-zero adjoint complex; here they serve only to separate recursive finite syntax from formality.
\end{warning}

\subsection{Scope of the derived envelope}

\begin{warning}[A canonical derived closure is not a classification]\label{warn:derived-not-classification}
The derived stack is a canonical faithful formal envelope, but not a finite list of slopes, filtrations, or normal forms.  The matrix-pair embedding of Theorem~\ref{thm:matrix-pair-embedding} makes this wildness intrinsic; merely renaming the derived stack would add no classification content.
\end{warning}

\begin{remark}[Derived transport-carrier comparison: final status]\label{rem:derived-transport-comparison-closed}
The comparison anticipated here is carried out by the later Cayley--Hamilton, all-Sylow, lci/Kuranishi, and finite-atlas theorems.  The Cayley--Hamilton layer supplies the finite universal full-transport algebra, the intrinsic derived envelope supplies the cotangent and Maurer--Cartan deformation object, and the projector-free all-Sylow carrier gives a finite perfect representative of continuous cochains.  Theorem~\ref{thm:derived-faithful-finite-closure} pairs these layers into a faithful derived/Morita chart, while Theorem~\ref{thm:recursive-native-globalization} supplies coordinate independence, base change, tensor compatibility, and coherent Fox/Herr/Koszul comparison.  Thus this is a completed structural step, not a remaining normal-form problem.
\end{remark}

\section{Local complete intersections and Kuranishi carriers}\label{sec:lci-carrier}

For framed deformations of local Galois representations, finite syzygy-complete formal models exist in every residual sector.

\subsection{Framed residual deformation sectors}

Let $k$ be the residue field of a finite extension $E/\Qp$, let $\mathcal O$ be its ring of integers, and let
\[
\bar\rho:G_K\longrightarrow \mathrm{GL}_n(k)
\]
be continuous.  Write $R_{\bar\rho}^{\square}$ for the framed universal deformation ring.

\begin{theorem}[Local complete-intersection strictification]\label{thm:local-lci}
There is a presentation
\begin{equation}\label{eq:lci-presentation}
R_{\bar\rho}^{\square}
\cong
\mathcal O[[x_1,\ldots,x_r]]/(f_1,\ldots,f_s)
\end{equation}
in which $\varpi,f_1,\ldots,f_s$ is a regular sequence in $\mathcal O[[x_1,\ldots,x_r]]$, $r$ is the dimension of the framed tangent space, and
\[
s=\dim_k H^2(G_K,\operatorname{ad}\bar\rho).
\]
In particular, $R_{\bar\rho}^{\square}$ is flat over $\mathcal O$ and is a local complete intersection of the expected relative dimension.  Moreover, the Galatius--Venkatesh derived framed deformation ring is homotopy discrete.
\end{theorem}

\begin{proof}
B\"ockle--Iyengar--Pa\v{s}k\=unas prove flatness, the expected-dimensional complete-intersection presentation, and regularity of $\varpi,f_1,\ldots,f_s$ \cite{BIPLocal}.  Galatius--Venkatesh's comparison criterion then gives homotopy discreteness \cite{GalatiusVenkatesh}.  The corrigendum changes only an affiliation \cite{BIPCorrigendum}.
\end{proof}

\begin{corollary}[Finite Koszul Cartan carrier]\label{cor:koszul-carrier}
For every residual framed sector, the Koszul differential graded algebra
\[
K_{\bar\rho}^{\square}
:=K_{\mathcal O[[x_1,\ldots,x_r]]}(f_1,\ldots,f_s)
\]
is a finite free semi-free differential graded resolution of $R_{\bar\rho}^{\square}$.  It is therefore a finite syzygy-complete formal carrier for the framed deformation problem.  The same conclusion holds after completing at any characteristic-zero point of the generic fiber.
\end{corollary}

\begin{proof}
The Koszul complex of a regular sequence is a finite free resolution.  Localizations and completions of a local complete intersection remain local complete intersections, so the construction persists at characteristic-zero points.
\end{proof}

\subsection{The finite Kuranishi map}

Over $E$, continuous local Galois cohomology with finite-dimensional coefficients is finite-dimensional and is concentrated in degrees $0,1,2$.  After choosing a contraction, homotopy transfer places a minimal $L_\infty$ structure $\{\ell_m\}_{m\ge2}$ on
\[
H^\bullet(G_K,\operatorname{ad}\rho).
\]
Modulo gauge, the formal deformation equations are represented by the Kuranishi map
\begin{equation}\label{eq:kuranishi-map}
\kappa_\rho:H^1(G_K,\operatorname{ad}\rho)
\longrightarrow
H^2(G_K,\operatorname{ad}\rho),
\qquad
\kappa_\rho(x)=\sum_{m\ge2}\frac{1}{m!}\ell_m(x,\ldots,x).
\end{equation}

\begin{proposition}[Finite Kuranishi carrier]\label{prop:finite-kuranishi}
The formal unframed deformation problem at $\rho$ is determined, up to formal gauge equivalence, by the finite-dimensional graded vector space
\[
H^0(G_K,\operatorname{ad}\rho)\oplus
H^1(G_K,\operatorname{ad}\rho)\oplus
H^2(G_K,\operatorname{ad}\rho)
\]
together with its transferred minimal $L_\infty$ structure.  Equivalently, after a formal slice is chosen, all obstruction equations are compressed into the finitely many scalar formal power series comprising \eqref{eq:kuranishi-map}.
\end{proposition}

\begin{proof}
Homotopy transfer gives a minimal $L_\infty$ model on the finite-dimensional cohomology in degrees $0,1,2$, and filtered $L_\infty$ quasi-isomorphisms preserve Maurer--Cartan functors.  The equivalent formal-power-series presentation is the standard transferred description of Galois deformation rings \cite{WangErickson}.
\end{proof}

\begin{warning}[Finite carrier is not finite Taylor truncation]\label{warn:lci-not-formal}
Finite variables and equations do not imply finitely many nonzero $\ell_m$ or bounded polynomial degree: the Kuranishi equations may be formal power series with infinitely many Taylor coefficients.
\end{warning}

\subsection{The lci comparison boundary}

Local lci strictification settles finite pointwise existence, but not transport-explicit covariance across residual sectors.  The chosen regular sequence is noncanonical and sector dependent; the Koszul model does not itself retain the full $\Lambda_{K,E}$-action or supply monoidal descent; and wild sectors admit no transparent normal form.  The comparison problem is therefore to make these local carriers transport-explicit and coherently compatible, not to establish finite local syzygies anew.


\section{Cayley--Hamilton transport}\label{sec:ch-carrier}

Local lci carriers encode deformation equations but not full transport; the completed group algebra has full transport but is not finite over a commutative base.  Cayley--Hamilton theory supplies the finite intermediate carrier.

\subsection{Residual determinants are independent of the stable lattice}

Let $V$ be an $n$-dimensional continuous $E$-representation and let $T\subset V$ be a $G_K$-stable $\mathcal O_E$-lattice.  Reduction gives
\[
\bar\rho_T:G_K\longrightarrow \mathrm{GL}_n(k_E).
\]
Write $\bar D_T$ for the associated $n$-dimensional determinant, in the sense of Chenevier.

\begin{proposition}[Canonical residual determinant]\label{prop:residual-determinant-independent}
The determinant $\bar D_T$ is independent of the choice of stable lattice $T$.  Its associated semisimple residual representation is the Brauer--Nesbitt semisimplification of $T/\varpi_ET$.  We denote the resulting lattice-independent residual determinant by $\bar D_V$.
\end{proposition}

\begin{proof}
For each $g\in G_K$, the characteristic polynomial $\det(X-\rho(g))$ is independent of a basis.  Stability of $T$ implies that its coefficients lie in $\mathcal O_E$, and reduction modulo $\varpi_E$ gives the characteristic polynomial of $g$ acting on $T/\varpi_ET$.  Hence every stable lattice produces the same reduced characteristic-polynomial law, and therefore the same determinant $\bar D_V$.  Brauer--Nesbitt identifies the corresponding semisimple residual representation.
\end{proof}

The determinant therefore gives a canonical sector index, while the stack retains lattice-dependent extension data.

\begin{corollary}[Finitely many canonical determinant sectors]\label{cor:finite-determinant-sectors}
For fixed $K$, $E$, and $n$, only finitely many residual determinants $\bar D_V$ occur among continuous $n$-dimensional $E$-representations of $G_K$.
\end{corollary}

\begin{proof}
There are only finitely many continuous homomorphisms $G_K\to\mathrm{GL}_n(k_E)$ by Lemma~\ref{lem:finite-residual-sectors}.  Passing from a residual representation to its determinant has finite image, so only finitely many determinant sectors occur.
\end{proof}

\subsection{The universal Cayley--Hamilton quotient}

Fix a residual determinant
\[
\bar D:G_K\longrightarrow k
\]
of dimension $n$, with universal pseudodeformation ring $R_{\bar D}$ and universal determinant $D^{\mathrm u}$.  The universal determinant extends continuously to $R_{\bar D}[[G_K]]$.  Let $\mathrm{CH}(D^{\mathrm u})$ be the closed two-sided ideal generated by the evaluations of the universal characteristic polynomials, and set
\begin{equation}\label{eq:universal-ch-algebra}
\mathcal E^{\mathrm{CH}}_{\bar D}
:=R_{\bar D}[[G_K]]/\mathrm{CH}(D^{\mathrm u}).
\end{equation}
The quotient carries a universal continuous transport map
\begin{equation}\label{eq:universal-ch-transport}
\rho^{\mathrm{CH}}_{\bar D}:G_K
\longrightarrow
\bigl(\mathcal E^{\mathrm{CH}}_{\bar D}\bigr)^\times.
\end{equation}

\begin{theorem}[Finite full-transport Cayley--Hamilton carrier]\label{thm:ch-full-transport}
For $G_K$ and every residual determinant $\bar D$ of dimension $n$:
\begin{enumerate}
\item $R_{\bar D}$ is complete Noetherian;
\item $\mathcal E^{\mathrm{CH}}_{\bar D}$ is finite as an $R_{\bar D}$-module;
\item the $\mathfrak m_{\bar D}$-adic topology, the quotient topology from $R_{\bar D}[[G_K]]$, and the topology induced by the image of $G_K$ on $\mathcal E^{\mathrm{CH}}_{\bar D}$ are equivalent;
\item every continuous rank-$n$ representation with residual determinant $\bar D$ factors through the universal map \eqref{eq:universal-ch-transport}, and the corresponding representation stack is the stack of rank-$n$ representations of the finite $R_{\bar D}$-algebra $\mathcal E^{\mathrm{CH}}_{\bar D}$.
\end{enumerate}
\end{theorem}

\begin{proof}
Mazur finiteness gives the Noetherian base; module-finiteness, topology, factorization, and algebraization are the universal Cayley--Hamilton results of Wang--Erickson, built on Chenevier's determinant theory \cite{WangEricksonFamilies,ChenevierDeterminants}.
\end{proof}

\begin{corollary}[Finite canonical residual-determinant atlas]\label{cor:ch-atlas}
For fixed $K$, $E$, and $n$, the continuous representation functor is partitioned into finitely many canonical residual determinant sectors.  Each sector has a formal chart
\[
\widehat{\operatorname{Rep}}_n
 \bigl(\mathcal E^{\mathrm{CH}}_{\bar D_i}\bigr)
 \longrightarrow \operatorname{Spf}R_{\bar D_i}
\]
and an algebraic chart
\[
\operatorname{Rep}_n
 \bigl(\mathcal E^{\mathrm{CH}}_{\bar D_i}\bigr)
 \longrightarrow \operatorname{Spec}R_{\bar D_i}.
\]
The first is the $\mathfrak m_{\bar D_i}$-adic continuous-representation stack.  The second is its finite-type algebraization over the indicated base.  Each sector carries a module-finite universal full-transport algebra and contains every residual extension class with the prescribed semisimplified determinant.
\end{corollary}

\begin{proof}
Use Corollary~\ref{cor:finite-determinant-sectors}, Theorem~\ref{thm:ch-full-transport}, and Wang--Erickson's algebraization theorem.  Proposition~\ref{prop:residual-determinant-independent} shows that the index $\bar D_i$ is independent of the stable lattice.
\end{proof}

\subsection{Irreducible and multiplicity-free sectors}

The finite carrier becomes substantially more rigid on standard open sectors.

\begin{corollary}[Azumaya full-transport carrier on the irreducible locus]\label{cor:ch-azumaya}
On the absolutely irreducible locus of the rigid analytic pseudorepresentation space, the universal Cayley--Hamilton algebra is an Azumaya algebra of rank $n^2$ over its center and carries the universal continuous transport.  Consequently it is finite projective, and after an etale cover it becomes a matrix algebra with a universal rank-$n$ representation.
\end{corollary}

\begin{proof}
This is the universal Cayley--Hamilton representation on Chenevier's absolutely irreducible locus \cite{ChenevierDeterminants}.  An Azumaya algebra is finite locally free over its center and is etale-locally a matrix algebra.
\end{proof}

\begin{proposition}[Multiplicity-free GMA coordinates]\label{prop:ch-gma}
Assume that $\bar D$ is split and multiplicity-free.  Then the algebra
$\mathcal E^{\mathrm{CH}}_{\bar D}$ admits a generalized matrix algebra structure over
$R_{\bar D}$ compatible with the universal determinant.  In the two-character case it has the form
\[
\mathcal E^{\mathrm{CH}}_{\bar D}
\cong
\begin{pmatrix}
R_{\bar D} & \mathcal B\\
\mathcal C & R_{\bar D}
\end{pmatrix},
\]
where the off-diagonal modules $\mathcal B$ and $\mathcal C$ carry the two extension directions and their products define the reducibility ideals.
\end{proposition}

\begin{proof}
The generalized matrix algebra structure is the multiplicity-free part of the universal Cayley--Hamilton theorem of Wang--Erickson \cite{WangEricksonFamilies}, extending the standard GMA construction of Bellaiche--Chenevier.  The block decomposition follows from lifting the orthogonal idempotents associated to the distinct residual simple factors.  In dimension two the off-diagonal blocks are precisely the two extension bimodules.
\end{proof}

The off-diagonal GMA modules are exactly the extension data lost by the center.

\subsection{Flagged GMA charts and finite blockwise reconstruction}

A GMA decomposition does not by itself make the individual Peirce blocks stable under conjugation by the universal transport.  A blockwise cochain model therefore requires an \emph{oriented flag chart}: the representation, not merely its Cayley--Hamilton algebra, must preserve a chosen filtration.  On such a chart the entire interaction among the blocks is controlled by one nilpotent Maurer--Cartan element.

\begin{theorem}[Finite blockwise reconstruction on a flagged chart]\label{thm:flagged-gma-perturbation}
Let $A$ be a complete Noetherian $\mathcal O_E$-algebra or a $\Qp$-affinoid algebra, and let $T$ be a finite projective $A$-module with a continuous $G_K$-action preserving a finite filtration by direct summands
\[
0=T_0\subset T_1\subset\cdots\subset T_s=T.
\]
Assume that the graded pieces $V_a=T_a/T_{a-1}$ are finite projective.  After an affine refinement choose an $A$-linear splitting
\[
T\cong\bigoplus_{a=1}^s V_a
\]
and write $\rho_{\mathrm{gr}}=\bigoplus_a\rho_a$ for the diagonal graded transport.  Put
\[
\mathfrak n
:=\bigoplus_{a<b}\operatorname{Hom}_A(V_b,V_a)
\subset\operatorname{End}_A(T).
\]
Then the following hold.

\begin{enumerate}[label=\textnormal{(\roman*)}]
\item There is a unique normalized cochain
\[
\eta\in C^1_{\mathrm{cont}}(G_K,\mathfrak n)
\]
such that
\begin{equation}\label{eq:flagged-transport-mc-coordinate}
\rho(g)=(1+\eta(g))\rho_{\mathrm{gr}}(g).
\end{equation}
It satisfies the Maurer--Cartan equation
\begin{equation}\label{eq:flagged-extension-mc}
 d_{\mathrm{gr}}\eta+\eta\smile\eta=0.
\end{equation}
Thus the linear off-diagonal entries are extension cocycles, while the higher upper blocks record their finite Yoneda-composition corrections.

\item Give
$\operatorname{Hom}_A(V_b,V_a)$ block weight $b-a$ and let $F^q$ be the sum of blocks of weight at least $q$.  The adjoint cochain dga of $\rho$ is identified, after the chosen splitting, with
\begin{equation}\label{eq:flagged-adjoint-twist}
\left(
 C^\bullet_{\mathrm{cont}}
 \bigl(G_K,\operatorname{End}_A(\textstyle\bigoplus_aV_a)\bigr),
 D_\eta=d_{\mathrm{gr}}+[\eta,-],
 \smile
\right).
\end{equation}
The operator $[\eta,-]$ strictly raises block weight.  Hence $D_\eta$ preserves the finite filtration, and the associated graded complex is
\begin{equation}\label{eq:flagged-adjoint-associated-graded}
\operatorname{gr}^q_F C^\bullet_{\mathrm{cont}}(G_K,\operatorname{ad}T)
\cong
\bigoplus_{b-a=q}
 C^\bullet_{\mathrm{cont}}
 \bigl(G_K,\operatorname{Hom}_A(V_b,V_a)\bigr)
\end{equation}
with its ordinary blockwise Galois differential.

\item For every block $\operatorname{Hom}_A(V_b,V_a)$, choose a bounded finite-projective chart.  Replace the untwisted block cochain dga by a filtered $K$-projective dg-algebra model
\[
\mathcal B_{\mathrm{gr}}
\xrightarrow{\ \sim\ }
C^\bullet_{\mathrm{cont}}
\bigl(G_K,\operatorname{End}_A(\textstyle\bigoplus_aV_a)\bigr)
\]
in which the Maurer--Cartan twist is represented by a filtration-raising perturbation $\delta_\eta$.  After replacing the finite chart and $\mathcal B_{\mathrm{gr}}$ by homotopy-equivalent models and adding contractible summands if necessary, choose a block-preserving contraction
\[
(C_{\mathrm{blk}}^\bullet
 \mathop{\rightleftarrows}^{i}_{p}
 \mathcal B_{\mathrm{gr}},h).
\]
Such data may be chosen locally by filtered cofibrant replacement; they are gauges, not intrinsic structure.  Since $\delta_\eta h$ strictly raises the finite block filtration, the operator
\begin{equation}\label{eq:flagged-finite-perturbation}
\mathcal A_\eta
:=(1-\delta_\eta h)^{-1}\delta_\eta
=\sum_{m\ge0}(\delta_\eta h)^m\delta_\eta
\end{equation}
is a finite sum.  The basic perturbation formulas
\begin{align}
 d_{\mathrm{blk},\eta}
 &=d_{\mathrm{blk}}+p\mathcal A_\eta i,
 \label{eq:flagged-block-differential}\\
 i_\eta&=i+h\mathcal A_\eta i,
 &p_\eta&=p+p\mathcal A_\eta h,
 &h_\eta&=h+h\mathcal A_\eta h
 \label{eq:flagged-block-contraction}
\end{align}
produce a bounded finite-projective complex
$C_{\mathrm{blk},\eta}^\bullet$ and a contraction
\[
C_{\mathrm{blk},\eta}^\bullet
\mathop{\rightleftarrows}^{i_\eta}_{p_\eta}
\bigl(\mathcal B_{\mathrm{gr}},d_{\mathrm{gr}}+\delta_\eta\bigr)
\simeq
C^\bullet_{\mathrm{cont}}(G_K,\operatorname{ad}T).
\]
Every matrix entry is a finite expression in the blockwise comparison matrices, the extension coordinates $\eta$, and the inverse minors used by the chosen contractions.  In particular, the construction is integral whenever the input charts and contractions are integral.

\item Filtered colored homotopy transfer applies simultaneously to block multiplication
\[
\operatorname{Hom}(V_b,V_a)
\otimes
\operatorname{Hom}(V_c,V_b)
\longrightarrow
\operatorname{Hom}(V_c,V_a),
\]
to the cup-commutator bracket, and to coefficient modules.  It gives homotopy-coherent $A_\infty/L_\infty$ block transitions whose associated graded maps are the common-source chart transitions for the individual block coefficients.  Different splittings, block contractions, or minor charts produce filtered $A_\infty/L_\infty$-equivalent models.
\end{enumerate}

In the two-step case $0\subset V_1\subset T$, one has
$\eta\in Z^1(G_K,\operatorname{Hom}(V_2,V_1))$ and
$\eta\smile\eta=0$.  The block filtration is
\[
C^\bullet\!\operatorname{Hom}(V_2,V_1)
\subset
C^\bullet\!\begin{pmatrix}
\operatorname{End}(V_1)&\operatorname{Hom}(V_2,V_1)\\
0&\operatorname{End}(V_2)
\end{pmatrix}
\subset
C^\bullet\!\operatorname{End}(T),
\]
and the perturbation series reduces to
\begin{equation}\label{eq:two-block-finite-perturbation}
\mathcal A_\eta
=\delta_\eta+\delta_\eta h\delta_\eta.
\end{equation}
\end{theorem}

\begin{proof}
Equation~\eqref{eq:flagged-transport-mc-coordinate} is the unique upper-unipotent factorization of the filtered transport relative to the chosen splitting.  Expanding the representation identity gives Equation~\eqref{eq:flagged-extension-mc}, exactly as in the nonlinear cocycle calculation of Theorem~\ref{thm:MC}.  Twisting the adjoint cochain dga by this Maurer--Cartan element replaces $d_{\mathrm{gr}}$ by $d_{\mathrm{gr}}+[\eta,-]$, proving~\eqref{eq:flagged-adjoint-twist}.

Block composition adds weights.  Every component of $\eta$ has positive weight, so $[\eta,-]$ strictly raises the finite filtration and disappears on the associated graded.  This proves~\eqref{eq:flagged-adjoint-associated-graded}.  Choose the filtered $K$-projective replacement and stabilized contraction stated in (iii).  Its homotopy $h$ has weight zero, while the lifted twist $\delta_\eta$ has positive weight; hence $\delta_\eta h$ is nilpotent.  The basic perturbation lemma gives Equations~\eqref{eq:flagged-finite-perturbation}--\eqref{eq:flagged-block-contraction}.  No convergence or completion beyond that already present in the coefficient rings is involved because the series terminates by weight.

The multiplication and bracket respect block-weight addition, so the filtered colored homotopy-transfer theorem supplies the higher operations and comparison maps.  Choice-independence is its comparison theorem.  For two blocks the possible weights are $-1,0,1$, so every term containing three copies of $\delta_\eta$ vanishes, giving Equation~\eqref{eq:two-block-finite-perturbation} \cite{BurkeTransfer,ManettiPerturbation,VokrinekHPT}.
\end{proof}

\begin{corollary}[Oriented multiplicity-free GMA block assembly]\label{cor:oriented-gma-block-assembly}
Assume that $\bar D$ is split and multiplicity-free.  Refine the Cayley--Hamilton GMA locus by Fitting strata on which the Peirce blocks are finite projective, and pass to an oriented framed flag chart on which the universal representation is block upper triangular.  Then the lifted GMA idempotents give the graded pieces $V_a$, the off-diagonal universal matrix entries give the Maurer--Cartan coordinate $\eta$, and the GMA products
\[
e_a\mathcal E^{\mathrm{CH}}_{\bar D}e_b
\otimes
e_b\mathcal E^{\mathrm{CH}}_{\bar D}e_c
\longrightarrow
e_a\mathcal E^{\mathrm{CH}}_{\bar D}e_c
\]
give the quadratic terms of Equation~\eqref{eq:flagged-extension-mc}.  Consequently the adjoint Fox, Herr, Koszul, and Cayley--Hamilton carriers on that chart admit an integral finite blockwise reconstruction by Equations~\eqref{eq:flagged-block-differential}--\eqref{eq:flagged-block-contraction}, together with homotopy-coherent multiplicative and internal-Hom transitions.

Thus blockwise assembly is no longer an existence problem on a fixed oriented flag chart.  Theorem~\ref{thm:flag-change-gma-descent} supplies the comparison across changes of flag, idempotent splitting, and opposite triangular orientation.  Theorem~\ref{thm:rigid-flag-etale-descent} gives finite etale effectivity on the finite rigid-flag locus.  At this stage only the non-rigid flag strata remain; Theorems~\ref{thm:nonrigid-flag-hdescent}--\ref{thm:descent-exponent-stratification} below close their derived descent and finite chartwise effectivity.
\end{corollary}

\begin{proof}
Proposition~\ref{prop:ch-gma} supplies the Peirce decomposition and its multiplication maps.  On the oriented flag refinement the universal transport preserves the selected filtration, so Theorem~\ref{thm:flagged-gma-perturbation} applies.  Fitting and minor refinement make the block modules and their finite representatives projective and permit the local filtered contractions used in the theorem.  The final statement records precisely which choices are removed by the theorem and which are not.
\end{proof}

\begin{theorem}[Flag-change and idempotent-splitting comparison]\label{thm:flag-change-gma-descent}
Let $U$ be an overlap of two oriented flagged GMA charts for the same family $T$,
with filtered splittings $(F_\bullet,\sigma)$ and $(F'_\bullet,\sigma')$ and
finite block carriers $C_{F,\eta}^\bullet$ and $C_{F',\eta'}^\bullet$ supplied by
Theorem~\ref{thm:flagged-gma-perturbation}.  After a Fitting/minor refinement,
the following statements hold.

\begin{enumerate}[label=\textnormal{(\roman*)}]
\item Suppose first that $F_\bullet=F'_\bullet$ and only the splitting, or the
chosen lifts of the GMA idempotents, changes.  The two split frames differ by an
upper-unipotent automorphism
\[
 u=1+q,
 \qquad q\in\bigoplus_{a<b}\operatorname{Hom}(V_b,V_a).
\]
The two Maurer--Cartan coordinates are related by the finite gauge formula
\begin{equation}\label{eq:flag-splitting-gauge}
 1+\eta'(g)
 =u^{-1}(1+\eta(g))\rho_{\mathrm{gr}}(g)
   u\rho_{\mathrm{gr}}(g)^{-1}.
\end{equation}
Since $q$ is nilpotent, $u^{-1}=\sum_{m\ge0}(-q)^m$ is a finite sum.  Conjugation
therefore gives a strict filtered dga isomorphism
\begin{equation}\label{eq:flag-splitting-strict-dga}
 \operatorname{Ad}(u^{-1}):
 (C^\bullet,D_\eta,\smile)
 \xrightarrow{\ \cong\ }
 (C^\bullet,D_{\eta'},\smile).
\end{equation}
Any two complete systems of lifted orthogonal GMA idempotents reducing to the
same residual system are, locally on the complete GMA chart, conjugate by such
a unit in $1+\operatorname{Jac}(\mathcal E^{\mathrm{CH}}_{\bar D})$.  Thus
splitting and idempotent-lift dependence is removed by the same finite formula.

\item For arbitrary stable flags on the overlap, choose the transition matrix
$g_{F'F}$ between the two split frames.  On the full adjoint cochain dga there
is a strict isomorphism
\begin{equation}\label{eq:flag-change-ambient-conjugation}
 \operatorname{Ad}(g_{F'F}^{-1}):
 C^\bullet_{\mathrm{cont}}(G_K,\operatorname{ad}T)_F
 \xrightarrow{\ \cong\ }
 C^\bullet_{\mathrm{cont}}(G_K,\operatorname{ad}T)_{F'}.
\end{equation}
Extend the perturbed contractions of Theorem~\ref{thm:flagged-gma-perturbation}
to filtered $A_\infty$ quasi-isomorphisms
$I_F^\infty,P_F^\infty$ and $I_{F'}^\infty,P_{F'}^\infty$.  Then
\begin{equation}\label{eq:flag-change-ainfty-transition}
 \mathfrak F_{F'F}
 :=P_{F'}^\infty\circ
   \operatorname{Ad}(g_{F'F}^{-1})\circ I_F^\infty
 :C_{F,\eta}^\bullet\rightsquigarrow C_{F',\eta'}^\bullet
\end{equation}
is a filtered $A_\infty$ quasi-isomorphism, and likewise an $L_\infty$
quasi-isomorphism for the cup-commutator brackets.  Its linear term is
\begin{equation}\label{eq:flag-change-linear-term}
 (\mathfrak F_{F'F})_1
 =p_{F',\eta'}\operatorname{Ad}(g_{F'F}^{-1})i_{F,\eta}.
\end{equation}
Every fixed Taylor component is a finite rooted-tree expression in the
transition matrix, the extension coordinates, the block multiplications, and
the inverse minors defining the contractions.

\item On a triple overlap, the strict ambient identities
$g_{F''F}=g_{F''F'}g_{F'F}$ imply that
$\mathfrak F_{F''F'}\mathfrak F_{F'F}$ and $\mathfrak F_{F''F}$ are connected
by the explicit $A_\infty/L_\infty$ homotopy obtained by inserting
$ I_{F'}^\infty P_{F'}^\infty\simeq\operatorname{id}$ in the common ambient
dga.  A coherent transfer convention extends these homotopies over the entire
nerve of the flag groupoid.  Changes of minor contractions or split frames
change the result by a coherent filtered $A_\infty/L_\infty$ equivalence.

\item Opposite triangular orientations are included in (ii)--(iii): on the
locus where both oriented flags exist, the transition matrix may include the
corresponding block permutation.  No comparison is asserted away from this
overlap, where one of the two flags is not preserved by the transport.
\end{enumerate}
\end{theorem}

\begin{proof}
For a fixed filtration, two splittings differ by a filtration-raising
endomorphism $q$.  Rewriting the same representation in the new frame gives
Equation~\eqref{eq:flag-splitting-gauge}; conjugation intertwines the two
adjoint actions and hence the two twisted differentials.  Nilpotence of the
finite block filtration makes every inverse polynomial.  The corresponding
conjugacy statement for lifted idempotents is the standard idempotent-conjugacy
argument in a complete semiperfect algebra, applied successively to the
orthogonal residual idempotents.

For two different flags, Equation~\eqref{eq:flag-change-ambient-conjugation} is
ordinary change of frame.  Composing it with the filtered perturbation
inclusions and projections gives Equation~\eqref{eq:flag-change-ainfty-transition}.
The homotopy-transfer comparison theorem gives the rooted-tree formulas and
shows that the linear term is Equation~\eqref{eq:flag-change-linear-term}.
On triple overlaps, strict composition in the ambient cochain dga leaves only
the middle factor $I_{F'}^\infty P_{F'}^\infty$; its chosen homotopy to the
identity gives the stated coherence.  The full nerve is treated by the same
cofibrant/Boardman--Vogt convention used in
Theorem~\ref{thm:cech-stasheff-gauges}.
\end{proof}

\begin{corollary}[Homotopy-coherent descent on the flag groupoid]\label{cor:flag-groupoid-descent}
On every split multiplicity-free GMA Fitting stratum, the oriented flagged
block assemblies form a homotopy-coherent $A_\infty/L_\infty$ and lax
monoidal atlas on the groupoid of invariant finite-projective flags.  This
atlas presents the pullback of the intrinsic adjoint cochain carrier, so
changes of stable flag, lifted idempotents, splitting, and triangular
orientation are no longer comparison obstructions.

At this stage the unflagged effectivity has not yet been constructed.  The
next sequence of theorems supplies it: rigid flags descend etale-locally,
non-rigid flags descend by derived $h$-descent, finite strict local models exist,
and finite-stage Schur--Koszul extraction with a terminating linear retraction
solver removes the infinite Cech totalization from the final native atlas.
\end{corollary}


\begin{theorem}[Rigid-flag finite etale effectivity]\label{thm:rigid-flag-etale-descent}
Let $X$ be a Noetherian finite-type framed representation chart carrying a
finite-projective universal $G_K$-module $\mathcal T$, and fix flag ranks
$\mathbf d=(d_1<\cdots<d_{s-1})$.  Let
\[
 \pi:\operatorname{Flag}^{G_K}_{\mathbf d}(\mathcal T)\longrightarrow X
\]
be the functor of $G_K$-stable flags by direct summands.  Then:

\begin{enumerate}[label=\textnormal{(\roman*)}]
\item The invariant-flag functor is represented by a closed subscheme of the
relative flag variety $\operatorname{Flag}_{\mathbf d}(\mathcal T)$; hence
$\pi$ is projective and of finite presentation.  If
$g_1,\ldots,g_e$ topologically generate $G_K$, the closed equations are the
vanishing of
\begin{equation}\label{eq:invariant-flag-equations}
 \mathcal F_a\xrightarrow{\rho(g_j)}\mathcal T
 \longrightarrow\mathcal T/\mathcal F_a,
 \qquad 1\le a<s,\quad 1\le j\le e,
\end{equation}
for the universal flag $\mathcal F_\bullet$.

\item At a geometric invariant flag $(x,F_\bullet)$, let
$\mathfrak p_F\subset\operatorname{End}(\mathcal T_x)$ be its stabilizer
parabolic.  The relative Zariski tangent space is
\begin{equation}\label{eq:invariant-flag-relative-tangent}
 T_{(x,F_\bullet)}
 \bigl(\operatorname{Flag}^{G_K}_{\mathbf d}(\mathcal T)/X\bigr)
 \cong
 H^0\!\left(G_K,
 \operatorname{End}(\mathcal T_x)/\mathfrak p_F\right).
\end{equation}
Call the flag \emph{rigid} when this space vanishes.

\item Let $X^{\mathrm{fin,rig}}\subset X$ be the maximal locally constructible
locus over which every geometric fiber of $\pi$ is finite and every invariant
flag is rigid.  It admits a finite locally closed stratification
\[
 X^{\mathrm{fin,rig}}=\coprod_\lambda X_\lambda
\]
such that
\begin{equation}\label{eq:rigid-flag-finite-etale}
 \pi_\lambda:
 \operatorname{Flag}^{G_K}_{\mathbf d}(\mathcal T)\times_X X_\lambda
 \longrightarrow X_\lambda
\end{equation}
is finite etale of constant degree.  On a degree-one stratum,
$\pi_\lambda$ is an isomorphism.

\item On every $X_\lambda$, the flagged Fox/Herr/Koszul/Cayley--Hamilton
carriers, their projective syzygy modules, and every fixed finite level of the
$A_\infty/L_\infty$ and lax monoidal structure descend effectively along
\eqref{eq:rigid-flag-finite-etale}.  The descended datum is independent of the
chosen flag and agrees with the intrinsic unflagged perfect carrier.  A single
global matrix is not asserted unless the descended projective module is
trivialized; the invariant output is a finite projective or Morita-valued
object with descended structure maps.
\end{enumerate}
\end{theorem}

\begin{proof}
The relative flag variety is projective over $X$.  Stability under the finite
set of topological generators is expressed by
\eqref{eq:invariant-flag-equations}; a direct summand stable under the dense
subgroup they generate is stable under $G_K$ by continuity.  This proves~(i).

The tangent space of the ordinary flag variety at $F_\bullet$ is
$\operatorname{End}(\mathcal T_x)/\mathfrak p_F$.  Linearizing the invariance
equations selects the vectors fixed by the adjoint $G_K$-action, giving
\eqref{eq:invariant-flag-relative-tangent}.

Over the base locus on which every flag fiber is finite, projectivity and
quasi-finiteness make $\pi$ finite.  Remove the image of the non-unramified
locus; Equation~\eqref{eq:invariant-flag-relative-tangent} shows that the
remaining finite morphism is unramified.  Apply the flattening stratification
to its finite direct-image algebra.  On each flat stratum the morphism is
finite locally free and unramified, hence finite etale
\cite{StacksFitting,StacksEtale}.  A finite etale morphism of rank one is an
isomorphism, proving~(iii).

Finite etale morphisms are fpqc covers.  Finite projective modules, perfect
complexes, algebra/module objects, and their morphisms satisfy fpqc descent.
The flag-groupoid transition maps of
Corollary~\ref{cor:flag-groupoid-descent} provide the descent datum; applying
descent to each finite arity and to a fixed coherent-transfer convention gives
(iv).  Different conventions yield equivalent descended homotopy-coherent
objects.
\end{proof}

\begin{corollary}[Two-character nonsplit descent]\label{cor:two-character-rigid-flag}
Consider a two-character multiplicity-free chart with a family of nonsplit
extensions
\[
 0\longrightarrow\chi_1\longrightarrow\mathcal T
 \longrightarrow\chi_2\longrightarrow0
\]
of rank-one characters.  On any locally closed locus on which
$\chi_1\ne\chi_2$ at every geometric point and the extension remains
nonsplit, the oriented invariant line of type $\chi_1$ is unique and rigid.
After refinement, the oriented flag chart is therefore isomorphic to the
unflagged nonsplit locus, and the finite blockwise Fox/Herr/Koszul and
Cayley--Hamilton formulas of Theorem~\ref{thm:flagged-gma-perturbation}
descend there without a remaining flag choice.
\end{corollary}

\begin{proof}
A second invariant line different from $\chi_1$ would project nontrivially to
the rank-one quotient $\chi_2$, hence isomorphically, and would split the
extension.  Thus the oriented invariant line is unique.  Its relative tangent
space is
\[
 H^0\!\left(G_K,\operatorname{Hom}(\chi_1,\chi_2)\right),
\]
which vanishes when the two characters are distinct.  The degree-one case of
Theorem~\ref{thm:rigid-flag-etale-descent} applies.
\end{proof}

\begin{theorem}[Derived $h$-effectivity on non-rigid flag loci]\label{thm:nonrigid-flag-hdescent}
Retain the notation of Theorem~\ref{thm:rigid-flag-etale-descent}.  Let
$X_{\mathbf d}^{\mathrm{red}}\subset X$ be a locally closed stratum whose
geometric points admit a $G_K$-stable flag of type $\mathbf d$, and let
$\mathfrak X_{\mathbf d}$ be the locally Noetherian derived enhancement
induced from the intrinsic derived representation problem
$\mathfrak{Rep}_{n,E}(G_K)$.  Write $P_{\mathbf d}\subset\mathrm{GL}_n$
for the corresponding parabolic and define the derived stack of reductions
\begin{equation}\label{eq:derived-invariant-flag-stack}
 \mathfrak F_{\mathbf d}^{\mathrm{der}}
 :=
 \mathfrak X_{\mathbf d}
 \mathop{\times}^{\mathbf R}_{\mathfrak{Rep}_{n,E}(G_K)}
 \operatorname{Map}_{\mathrm{cont}}(BG_K,BP_{\mathbf d}).
\end{equation}
On the finite residual word/lci charts of this paper, this fiber product is
interpreted as the corresponding locally Noetherian formal derived stack.
Then:

\begin{enumerate}[label=\textnormal{(\roman*)}]
\item The classical truncation of
$\mathfrak F_{\mathbf d}^{\mathrm{der}}$ is the invariant-flag scheme
$\operatorname{Flag}^{G_K}_{\mathbf d}(\mathcal T)$ over
$X_{\mathbf d}^{\mathrm{red}}$.  The forgetful morphism
\[
 \pi^{\mathrm{der}}:
 \mathfrak F_{\mathbf d}^{\mathrm{der}}
 \longrightarrow \mathfrak X_{\mathbf d}
\]
is almost finitely presented, and its classical truncation is projective and
surjective.  Hence $\pi^{\mathrm{der}}$ is a derived $h$-cover.

\item At a geometric invariant flag $(x,F_\bullet)$, with stabilizer
parabolic Lie algebra $\mathfrak p_F$, the relative tangent complex is
\begin{equation}\label{eq:derived-invariant-flag-tangent}
 \mathbb T_{\mathfrak F_{\mathbf d}^{\mathrm{der}}/
                  \mathfrak X_{\mathbf d},(x,F_\bullet)}
 \simeq
 R\Gamma\left(
 G_K,\operatorname{End}(\mathcal T_x)/\mathfrak p_F
 \right).
\end{equation}
It is perfect of amplitude contained in $[0,2]$.  Its $H^0$ is the tangent
space in Equation~\eqref{eq:invariant-flag-relative-tangent}; $H^1$ is the
first obstruction space for deforming the invariant flag with the ambient
representation fixed, and $H^2$ is the terminal higher obstruction.

\item Let $\mathfrak F^{\mathrm{der}}_{\mathbf d,\bullet}$ be the
derived Cech nerve of $\pi^{\mathrm{der}}$.  Pullback induces symmetric
monoidal equivalences
\begin{equation}\label{eq:nonrigid-flag-hdescent}
 \operatorname{Perf}(\mathfrak X_{\mathbf d})
 \simeq
 \operatorname*{Tot}_{[q]\in\Delta}
 \operatorname{Perf}
 \left(\mathfrak F^{\mathrm{der}}_{\mathbf d,q}\right),
\end{equation}
and likewise for almost perfect complexes.

\item The flagged Fox/Herr/Koszul/Cayley--Hamilton carriers and the
flag-change coherences of Theorem~\ref{thm:flag-change-gma-descent} define an
object of the totalization in
\eqref{eq:nonrigid-flag-hdescent}.  They therefore descend uniquely to a
perfect unflagged carrier on $\mathfrak X_{\mathbf d}$, agreeing with the
intrinsic adjoint cochain carrier.  Algebra, module, $A_\infty/L_\infty$,
lax monoidal, and Morita-valued structures descend in the same way.
\end{enumerate}

Thus positive-dimensional, nonreduced, and cardinality-jumping flag fibers do
not obstruct derived categorical descent.  The theorem does not produce a
bounded strict matrix for the descended object: the full derived Cech
totalization and its higher coherences may still be required.
\end{theorem}

\begin{proof}
A reduction of the universal $\mathrm{GL}_n$-transport to
$P_{\mathbf d}$ is exactly a stable flag of type $\mathbf d$, proving the
truncation statement.  The classical forgetful map is the projective invariant
flag morphism of Theorem~\ref{thm:rigid-flag-etale-descent}; it is surjective
onto $X_{\mathbf d}^{\mathrm{red}}$ by definition.  This makes the derived
map an $h$-cover.

The tangent complex of a derived mapping stack is obtained by continuous
cochains in the pulled-back tangent representation.  The fiber sequence
$\mathfrak p_F\to\operatorname{End}(\mathcal T_x)\to
\operatorname{End}(\mathcal T_x)/\mathfrak p_F$ therefore gives
\eqref{eq:derived-invariant-flag-tangent}; local $p$-cohomological dimension
two gives its perfect amplitude.

Halpern--Leistner and Preygel prove derived $h$-descent for
$\operatorname{APerf}$ and $\operatorname{Perf}$ on locally Noetherian
derived stacks \cite[Theorem~3.3.1]{HalpernLeistnerPreygel}.  Applying their
theorem to $\pi^{\mathrm{der}}$ gives
\eqref{eq:nonrigid-flag-hdescent}.  The equivalence is symmetric monoidal, so
it also descends algebra objects, modules, and their homotopy-coherent
structure maps.  The flag-change theorem supplies the required Cech datum,
and its comparison with the common ambient adjoint cochain dga identifies the
descent with the intrinsic unflagged carrier.
\end{proof}


\begin{theorem}[Finite strict local effectivity of non-rigid descent]\label{thm:nonrigid-strict-local}
Retain the notation of Theorem~\ref{thm:nonrigid-flag-hdescent}, and assume
that the locally closed derived stratum $\mathfrak X_{\mathbf d}$ is
quasi-compact.  Let $\mathcal P$ be any perfect carrier descended by
Equation~\eqref{eq:nonrigid-flag-hdescent}; this includes the adjoint
cochain carrier and the underlying perfect complexes of the descended
algebra, module, and Morita-valued objects.  Then:

\begin{enumerate}[label=\textnormal{(\roman*)}]
\item There is a finite affine smooth cover
\[
 U_\lambda=\operatorname{Spec}A_\lambda
 \longrightarrow \mathfrak X_{\mathbf d}
\]
such that $\mathcal P|_{U_\lambda}$ is represented by a bounded complex
$P_\lambda^\bullet$ of finite projective $A_\lambda$-modules.  If
$\mathcal P$ has Tor amplitude in $[a,b]$, the representative may be chosen
with $P_\lambda^j=0$ for $j\notin[a,b]$.  In particular, the local Galois
cochain carriers of amplitude $[0,2]$ admit three-term strict matrices, and
the framed evaluation fibers of amplitude $[0,1]$ admit two-term matrices.

\item A common finite Fitting stratification and minor refinement may be
chosen for any finite collection of descended carriers and structure maps.
On every resulting minor chart the ranks of the differentials and cohomology
modules are constant, Proposition~\ref{prop:determinantal-contraction}
gives explicit contractions, and Theorem~\ref{thm:cech-stasheff-gauges}
gives computable $A_\infty/L_\infty$, module, lax monoidal, and Morita
transition data.

\item Pulling $P_\lambda^\bullet$ back to the derived flag Cech nerve
gives the same object as the flagged block complexes of
Theorem~\ref{thm:flagged-gma-perturbation}.  After choosing contractions,
the comparison maps may be taken to be the common-source chart maps of
Theorems~\ref{thm:common-source-chart-change} and
\ref{thm:common-source-multiplicative-tower}; their higher coherences are
supplied by the common ambient adjoint cochain object.  Thus the strict local
matrices realize, rather than replace, the $h$-descended unflagged carrier.

\item A single finite-projective syzygy module exists on a stratum only when
the relevant descended perfect object has Tor amplitude $[0,0]$.  In general
the invariant strict output is a bounded finite-projective complex, or a
Morita-valued perfect object, and not one globally free matrix module.
\end{enumerate}

Consequently bounded strict matrix \emph{existence} on non-rigid strata is
not an open problem.  This theorem alone does not produce a flag-native
finite-stage extraction or explicit gauge-to-transport comparison; the next
theorem supplies a noncanonical finite-stage retract.
\end{theorem}

\begin{proof}
A perfect complex on an affine derived chart is represented by a bounded
complex of finite projective modules.  If its Tor amplitude is $[a,b]$, the
representative may be chosen in the same interval
\cite[Section~15.76, especially Lemma~15.76.2]{StacksPerfect}.  Quasi-
compactness gives a finite affine smooth cover, proving~(i).

Apply the Fitting stratification simultaneously to the finitely many matrices
in the chosen representatives and refine by the nonvanishing minors used for
row reduction.  The determinantal contraction and Cech--Stasheff theorems
then prove~(ii).  Part~(iii) follows because both the pulled-back strict model
and the flagged block model represent the same object of the totalization in
Equation~\eqref{eq:nonrigid-flag-hdescent}; the common-source comparison
theorems provide strict linear maps and homotopy-coherent higher maps.
Finally, a perfect object is represented by a finite projective module in one
degree precisely on the locus where its Tor amplitude is $[0,0]$.  This proves
(iv) and the final distinction between local strict existence and closed-form
effectivity.
\end{proof}


\begin{theorem}[Finite-stage Cech--Koszul extraction on non-rigid strata]\label{thm:nonrigid-finite-stage-retract}
Retain the notation of Theorems~\ref{thm:nonrigid-flag-hdescent} and
\ref{thm:nonrigid-strict-local}.  After replacing a finite affine smooth chart
of $\mathfrak X_{\mathbf d}$ by a finite affine refinement, write
$U=\operatorname{Spec}A$ and choose finitely many affine opens $W_i$ covering
the relative flag variety $\operatorname{Flag}_{\mathbf d}(\mathcal T_U)$.
Let $\mathcal F_a$ be the universal subbundles and let
\[
 \mathcal V
 :=\bigoplus_{a,j}
 \underline{\operatorname{Hom}}
 \left(\mathcal F_a,\mathcal T_U/\mathcal F_a\right)
\]
on the flag variety, where $j$ ranges over fixed topological generators
$g_1,\ldots,g_e$ of $G_K$.  The stability equations of
Equation~\eqref{eq:invariant-flag-equations} define a section
$s_\rho\in\Gamma(\mathcal V)$.  Then:

\begin{enumerate}[label=\textnormal{(\roman*)}]
\item On every $W_i$, the derived invariant-flag chart is the derived zero
locus of $s_\rho$ and has the finite commutative Koszul model
\begin{equation}\label{eq:invariant-flag-koszul-cdga}
 B_i
 :=\Gamma\!\left(
 W_i,
 \left(\bigwedge{}^\bullet\mathcal V^\vee,
 d=\iota_{s_\rho}\right)
 \right).
\end{equation}
Let $B_i^{\mathrm{cl}}:=\pi_0(B_i)$ be the coordinate ring of the
classical invariant-flag chart, and put
$B^{\mathrm{cl}}=\prod_iB_i^{\mathrm{cl}}$.  The map of Noetherian rings
$A\to B^{\mathrm{cl}}$ is a finitely presented $h$-cover.  Its derived
Amitsur terms
$(B^{\mathrm{cl}})^{\otimes_A^{\mathbf L}(q+1)}$ are finite-type, while
the derived flag directions on each term are represented by finite tensor
products of the explicit Koszul cdgas $B_i$.

\item The classical algebra $B^{\mathrm{cl}}$ is descendable over $A$.
Consequently there is an integer $N\ge0$ and maps in
$\operatorname{Perf}(A)$
\begin{equation}\label{eq:finite-stage-unit-retract}
 A\xrightarrow{\ \alpha_N\ }
 \operatorname{Tot}_{\le N}
 \left((B^{\mathrm{cl}})^{\otimes_A^{\mathbf L}(\bullet+1)}\right)
 \xrightarrow{\ \beta_N\ }A,
 \qquad
 \beta_N\alpha_N=\operatorname{id}_A.
\end{equation}
After a finite affine cover of the quasi-compact stratum, one may choose one
common value of $N$ by taking the maximum of the finitely many local
exponents.

\item Let $\mathcal Q^\bullet$ be the explicit flagged carrier with its Cech
descent datum, and let $\mathcal P$ be the descended unflagged carrier.  Tensoring
Equation~\eqref{eq:finite-stage-unit-retract} with $\mathcal P$, or equivalently
using the flagged descent datum under
Equation~\eqref{eq:nonrigid-flag-hdescent}, gives a retract
\begin{equation}\label{eq:finite-stage-carrier-retract}
 \mathcal P
 \longrightarrow
 \operatorname{Tot}_{\le N}
 \left(\mathcal Q^\bullet|_{[0,N]}\right)
 \longrightarrow
 \mathcal P.
\end{equation}
Thus the infinite derived Cech totalization is not required to obtain a finite
presentation of the descended perfect object: a finite partial
Cech--Koszul totalization, followed by one idempotent splitting, suffices.

\item On every affine minor chart, the middle term of
Equation~\eqref{eq:finite-stage-carrier-retract} is represented by a bounded
finite-projective double complex whose entries are finite expressions in the
stability equations, the Peirce/GMA multiplication matrices, the flagged
perturbation matrices of Theorem~\ref{thm:flagged-gma-perturbation}, and the
chosen Cech face maps.  Splitting the retract idempotent in
$\operatorname{Perf}(A)$ gives a bounded finite-projective representative of
$\mathcal P$.  The same finite-stage construction applies simultaneously to
algebra, module, $A_\infty/L_\infty$, lax monoidal, and Morita-valued
structures.
\end{enumerate}

The integer $N$, the maps $\alpha_N,\beta_N$, and the resulting matrix
idempotent are gauges.  Descendability proves that finite-stage extraction
exists, but it does not supply a canonical polynomial formula for this
idempotent or a uniform optimal exponent on every GMA stratum.
\end{theorem}

\begin{proof}
A derived zero locus of a section of a finite locally free sheaf is represented
by its Koszul commutative dg algebra, which proves the first assertion of~(i).
The finite affine open cover of the classical invariant-flag scheme may be
replaced by its disjoint union; composing it with the proper surjective
invariant-flag morphism gives the finitely presented classical $h$-cover
$\operatorname{Spec}B^{\mathrm{cl}}\to U$.

Bhatt and Scholze prove that a finitely presented $h$-cover of Noetherian
affines is descendable \cite[Proposition~11.25]{BhattScholzeWitt}.  Mathew's
characterization of descendability identifies this with the condition that
the tensor unit belongs to the thick tensor ideal generated by $B$; equivalently,
the unit is a retract of a finite partial totalization of the Amitsur complex
\cite[Section~3.3, especially Propositions~3.20--3.22]{MathewGalois}.  This
gives Equation~\eqref{eq:finite-stage-unit-retract} and proves~(ii).

Tensoring the retract with the strict classical perfect $A$-module
$\mathcal P$ gives Equation~\eqref{eq:finite-stage-carrier-retract}.  Under
the classical restriction of the $h$-descent equivalence, its Amitsur terms are
precisely the finite levels of the flagged Cech datum, proving~(iii).  The full
derived flag directions remain encoded by the finite Koszul enhancements
$B_i$ and compare with this strict extraction through
Theorem~\ref{thm:nonrigid-flag-hdescent}.  Finally, each affine
invariant-flag term has a finite Koszul enhancement and each flagged carrier
already has finite block matrices.  A
finite partial totalization is therefore a bounded finite double complex.
Idempotent completeness of $\operatorname{Perf}(A)$ splits the retract, while
the symmetric monoidal form of descendability carries all the higher algebra,
module, and Morita structures.  This proves~(iv).
\end{proof}


\begin{theorem}[Schur--Koszul dg presentation of the flag Cech terms]\label{thm:schur-koszul-flag-dg}
Retain the notation of Theorem~\ref{thm:nonrigid-finite-stage-retract} and
write
\[
 F_{\mathbf d}:=\operatorname{Flag}_{\mathbf d}(\mathcal T_U),
 \qquad
 Z^{\mathrm{der}}\hookrightarrow F_{\mathbf d}
\]
for the derived invariant-flag zero locus.  For $q\ge0$, put
\[
 F_{\mathbf d,q}:=F_{\mathbf d}^{\times_U(q+1)},
 \qquad
 Z_q^{\mathrm{der}}:=(Z^{\mathrm{der}})^{\times_U(q+1)}.
\]
Then the following hold.

\begin{enumerate}[label=\textnormal{(\roman*)}]
\item The relative flag variety is an iterated relative Grassmannian.  There
is a finite vector bundle
\begin{equation}\label{eq:flag-schur-generator}
 \mathcal G_{\mathbf d}
 =\bigoplus_{\boldsymbol\lambda}
 \bigotimes_a
 \mathbf S_{\lambda_a}(\mathcal Q_a)
\end{equation}
on $F_{\mathbf d}$, where the partitions $\lambda_a$ range over finite
rectangles determined only by $n$ and $\mathbf d$, such that
$\mathcal G_{\mathbf d}$ is a compact generator of
$\operatorname{Perf}(F_{\mathbf d})$.  Its formation commutes with arbitrary
base change on $U$.  One may replace the Schur functors in
\eqref{eq:flag-schur-generator} by the characteristic-free exterior-power
generators of Buchweitz--Leuschke--Van den Bergh.

For the product $F_{\mathbf d,q}$, the external tensor product
\[
 \mathcal G_{\mathbf d,q}
 :=\mathcal G_{\mathbf d}^{\boxtimes(q+1)}
\]
is a compact generator.

\item Let $\mathcal V_\ell$ and $s_{\rho,\ell}$ denote the pullbacks of the
stability bundle and stability section from the $\ell$-th flag factor.  Put
\begin{equation}\label{eq:schur-koszul-flag-complex}
 \mathcal K_q
 :=\bigotimes_{\ell=0}^{q}
 \left(\bigwedge{}^\bullet\mathcal V_\ell^\vee,
 d=\iota_{s_{\rho,\ell}}\right).
\end{equation}
The pullback of $\mathcal G_{\mathbf d,q}$ to $Z_q^{\mathrm{der}}$ is a
compact generator, and its derived endomorphism algebra is
\begin{equation}\label{eq:schur-koszul-endomorphism-dga}
 \mathcal A_q^{\mathrm{SK}}
 :=R\operatorname{End}_{Z_q^{\mathrm{der}}}
 \left(\mathcal G_{\mathbf d,q}|_{Z_q^{\mathrm{der}}}\right)
 \simeq
 R\Gamma\!\left(
 F_{\mathbf d,q},
 \mathcal End(\mathcal G_{\mathbf d,q})\otimes\mathcal K_q
 \right).
\end{equation}
It is a perfect $A$-linear dg algebra, compatible with arbitrary base change,
and there is a Morita equivalence
\begin{equation}\label{eq:schur-koszul-morita-equivalence}
 \operatorname{Perf}(Z_q^{\mathrm{der}})
 \simeq
 \operatorname{Perf}(\mathcal A_q^{\mathrm{SK}}).
\end{equation}
Thus each derived Cech term is represented by one finite dg algebra, rather
than by a separately chosen affine cover.

\item Every face or degeneracy map of the flag Cech nerve induces an explicit
perfect dg bimodule between the corresponding algebras
$\mathcal A_q^{\mathrm{SK}}$.  These bimodules form a cosimplicial Morita
diagram.  After choosing bounded finite-projective representatives, the
partial totalization through level $N$ in
Equation~\eqref{eq:finite-stage-carrier-retract} is computed by a finite
matrix totalization built from
\[
 \mathcal A_0^{\mathrm{SK}},\ldots,\mathcal A_N^{\mathrm{SK}}
\]
and the finitely many face bimodules.  The same diagram carries the algebra,
module, $A_\infty/L_\infty$, lax monoidal, and Morita structure maps.

\item If $e$ topologically generates $G_K$, set
\begin{equation}\label{eq:flag-stability-koszul-rank}
 R_{\mathbf d}
 :=e\sum_{a=1}^{s-1}d_a(n-d_a).
\end{equation}
Then $\mathcal K_q$ has Koszul length at most
$(q+1)R_{\mathbf d}$.  The number and ranks of the Schur summands in
$\mathcal G_{\mathbf d,q}$ depend only on $(n,\mathbf d,q)$.  Consequently,
for every fixed partial-totalization level $N$, the size and cohomological
length of the Schur--Koszul matrix problem admit an a priori bound depending
only on $(n,\mathbf d,e,N)$ and on the ranks of the selected coefficient
carriers.
\end{enumerate}

The commutative affine $h$-cover used in
Theorem~\ref{thm:nonrigid-finite-stage-retract} is still needed to produce a
descendability retract.  The present theorem removes the additional choice of
an affine cover for each derived flag term: after Morita replacement, the
remaining unknown is one finite-stage retract idempotent in an explicit
finite Schur--Koszul dg diagram.
\end{theorem}

\begin{proof}
A relative flag variety is a tower of relative Grassmannians.  Efimov's
integral Grassmannian generators, equivalently the characteristic-free
tilting generators of Buchweitz--Leuschke--Van den Bergh, commute with base
change to an arbitrary commutative ring.  Applying them successively in the
Grassmannian tower and tensoring their pullbacks gives the finite Schur
generator~\eqref{eq:flag-schur-generator}; generation is local on the base and
therefore descends from a trivialization of the successive vector bundles
\cite{EfimovGrassmannians,BuchweitzLeuschkeVdBGrass}.  External tensor products
of compact generators generate the product flag variety.

The morphism $Z_q^{\mathrm{der}}\to F_{\mathbf d,q}$ is affine and perfect,
with direct-image algebra the finite Koszul complex
\eqref{eq:schur-koszul-flag-complex}.  If an object of
$\operatorname{Perf}(Z_q^{\mathrm{der}})$ is right-orthogonal to the pullback
of $\mathcal G_{\mathbf d,q}$, its affine pushforward is right-orthogonal to
the generator on $F_{\mathbf d,q}$ and hence vanishes.  The pullback is
therefore a compact generator.  Projection formula gives
\eqref{eq:schur-koszul-endomorphism-dga}, and derived Morita theory gives
\eqref{eq:schur-koszul-morita-equivalence}.  Properness of the flag variety
and perfectness of every displayed bundle make the endomorphism dg algebra
perfect over $A$; all constructions commute with base change.

A morphism between Cech terms acts on their compact generators by a perfect
Fourier--Mukai bimodule.  Composition of the maps gives the cosimplicial
bimodule identities, and choosing finite-projective representatives converts
every fixed truncation into a finite matrix totalization.  The colored
homotopy-transfer constructions already used in
Theorems~\ref{thm:common-source-multiplicative-tower} and
\ref{thm:common-source-lax-monoidal-tower} transport the higher structures
through the same bimodules.  Finally,
$\operatorname{rank}\underline{\operatorname{Hom}}
(\mathcal F_a,\mathcal T/\mathcal F_a)=d_a(n-d_a)$, which proves
\eqref{eq:flag-stability-koszul-rank} and the stated finite bounds.
\end{proof}


\begin{theorem}[Finite linear retraction equations and terminating exponent search]\label{thm:linear-descent-retraction}
Retain the notation of Theorems~\ref{thm:nonrigid-finite-stage-retract} and
\ref{thm:schur-koszul-flag-dg}.  For every $N\geq0$, choose a bounded
finite-projective $A$-complex $T_N^\bullet$ representing the partial
Schur--Koszul Cech totalization through level $N$, and let
\[
 \alpha_N:A[0]\longrightarrow T_N^\bullet
\]
be the canonical augmentation.  Put
\[
 \mathcal H_N^\bullet
 :=\operatorname{Hom}_A^\bullet(T_N^\bullet,A[0])
\]
with its ordinary Hom differential, and define
\begin{equation}\label{eq:descent-retraction-evaluation}
 \varepsilon_N:Z^0(\mathcal H_N^\bullet)\longrightarrow A,
 \qquad
 \beta\longmapsto (\beta\alpha_N)(1).
\end{equation}
Then:

\begin{enumerate}[label=\textnormal{(\roman*)}]
\item The unit $A[0]$ is a retract of $T_N^\bullet$ in
$\operatorname{Perf}(A)$ if and only if
\begin{equation}\label{eq:descent-stage-obstruction}
 1\in\operatorname{im}(\varepsilon_N).
\end{equation}
Equivalently, the obstruction
\[
 \omega_N:=[1]\in\operatorname{coker}(\varepsilon_N)
\]
vanishes.  When it vanishes, every preimage
$\beta_N\in Z^0(\mathcal H_N^\bullet)$ with
$\varepsilon_N(\beta_N)=1$ is a strict chain retraction
\[
 \beta_N\alpha_N=\operatorname{id}_{A[0]}.
\]

\item After choosing finite free frames on an affine minor chart, write
$D_N$ for the matrix of the Hom differential
$\mathcal H_N^0\to\mathcal H_N^1$ and $v_N$ for the row matrix of
$\varepsilon_N$.  A retraction is exactly a solution of the finite linear
system
\begin{equation}\label{eq:finite-descent-linear-system}
 \begin{pmatrix}D_N\\ v_N\end{pmatrix}b
 =
 \begin{pmatrix}0\\1\end{pmatrix}.
\end{equation}
Thus $N$ is a valid descent stage if and only if the displayed right-hand
side lies in the image of one explicit finite matrix whose entries are
Schur--Koszul, Cech-face, and coefficient-carrier matrices.

\item Descendability implies that $\omega_N=0$ for some finite $N$.
Consequently the sequence
\[
 \omega_0,\omega_1,\omega_2,\ldots
\]
is a terminating recognition procedure for a working descent stage.  Over
an effectively presented coefficient ring for which membership in finitely
presented submodules is decidable, successively constructing
\eqref{eq:finite-descent-linear-system} computes a valid $N$ and a retraction
$\beta_N$.  No a priori optimal uniform bound is asserted.

\item A solution gives the strict chain idempotent
\begin{equation}\label{eq:explicit-descent-idempotent}
 e_N:=\alpha_N\beta_N\in
 Z^0\operatorname{End}_A^\bullet(T_N^\bullet),
 \qquad e_N^2=e_N.
\end{equation}
On the finite Fitting stratification on which every degreewise rank of $e_N$
is constant, a minor cover gives explicit matrices
\[
 \operatorname{im}(e_N^j)
 \mathop{\rightleftarrows}^{\,i_N^j}_{\,p_N^j}
 T_N^j,
 \qquad p_N^ji_N^j=1,
 \qquad i_N^jp_N^j=e_N^j.
\]
The restricted differentials form a bounded finite-projective complex
$\operatorname{im}(e_N)$ representing the unit.  Tensoring this splitting
with any descended perfect carrier gives its finite-stage matrix retract.
Different solutions of~\eqref{eq:finite-descent-linear-system} give equivalent
perfect retracts, not a canonical matrix.
\end{enumerate}

Hence, after the Schur--Koszul replacement, neither the finite descent stage nor the retract idempotent is an unspecified higher-categorical datum: both are detected by a sequence of finite module-membership problems.  Theorem~\ref{thm:descent-exponent-stratification} below makes the working stage constructible and chartwise bounded, while Theorem~\ref{thm:no-uniform-descent-exponent-bound} shows that no bound depending only on the fixed representation/flag combinatorics can exist under arbitrary nonreduced coefficient base change.  A prettier closed solution on one chosen GMA chart is therefore only an optional coordinate optimization.
\end{theorem}

\begin{proof}
Because $T_N^\bullet$ is a bounded complex of finite projective modules, it is
$K$-projective.  Therefore every morphism
$T_N^\bullet\to A[0]$ in the derived category is represented by a degree-zero
cocycle of $\mathcal H_N^\bullet$.  Moreover
$\operatorname{Hom}_{D(A)}(A[0],A[0])=A$, and a homotopy between two
endomorphisms of the complex $A[0]$ is necessarily zero.  Thus a derived
retraction is equivalent to a strict cocycle $\beta_N$ satisfying
$\beta_N\alpha_N=1$, proving~(i).

Choosing finite free frames turns the cocycle and normalization conditions
into~\eqref{eq:finite-descent-linear-system}; this proves~(ii).
Theorem~\ref{thm:nonrigid-finite-stage-retract} gives a retraction for some
$N$, so the corresponding obstruction eventually vanishes.  Module
membership in the image of a finite matrix is precisely the test in~(ii),
proving~(iii).

Equation~\eqref{eq:explicit-descent-idempotent} is immediate from
$\beta_N\alpha_N=1$.  A degreewise idempotent on a finite projective module
has finite-projective image.  On a constant-rank minor chart, selecting an
invertible rank minor gives the standard explicit inclusion and projection
matrices; because $e_N$ is a chain map, the differential preserves its image.
Tensoring the unit retract with a perfect carrier is exactly the finite-stage
carrier retract of Equation~\eqref{eq:finite-stage-carrier-retract}.  This
proves~(iv).
\end{proof}


\begin{theorem}[Constructible descent exponent and specialization law]\label{thm:descent-exponent-stratification}
Retain the notation of Theorem~\ref{thm:linear-descent-retraction}.  For each
$N\geq0$, combine the cocycle and normalization equations into the finite
$A$-linear map
\begin{equation}\label{eq:descent-augmented-linear-map}
 L_N:\mathcal H_N^0\longrightarrow \mathcal H_N^1\oplus A,
 \qquad
 b\longmapsto\bigl(db,\varepsilon_N(b)\bigr),
\end{equation}
and put
\[
 r_N:=(0,1),\qquad
 Q_N:=\operatorname{coker}(L_N),\qquad
 \bar\omega_N:=[r_N]\in Q_N.
\]
Let $M_N:=A\bar\omega_N\subset Q_N$, let
$J_N:=\operatorname{Ann}_A(M_N)$, and define
\begin{equation}\label{eq:descent-exponent-open}
 U_N:=D(J_N)=\operatorname{Spec}A\setminus\operatorname{Supp}(M_N).
\end{equation}
Then:

\begin{enumerate}[label=\textnormal{(\roman*)}]
\item For every prime $\mathfrak p\subset A$, one has
\begin{equation}\label{eq:local-descent-stage-criterion}
 \mathfrak p\in U_N
 \quad\Longleftrightarrow\quad
 \bar\omega_N|_{A_{\mathfrak p}}=0
 \quad\Longleftrightarrow\quad
 \text{Equation~\eqref{eq:finite-descent-linear-system} has a solution over }
 A_{\mathfrak p}.
\end{equation}
Thus $U_N$ is exactly the locus on which level $N$ admits a strict local
retraction.

\item The opens are increasing:
\begin{equation}\label{eq:descent-exponent-open-filtration}
 U_0\subseteq U_1\subseteq U_2\subseteq\cdots.
\end{equation}
Moreover, there is a finite $N_*$ with $U_{N_*}=\operatorname{Spec}A$.
Consequently the local minimal descent exponent
\begin{equation}\label{eq:minimal-local-descent-exponent}
 \nu(\mathfrak p):=\min\{N:\mathfrak p\in U_N\}
\end{equation}
is a finite upper-semicontinuous integer-valued function.  In particular,
with $U_{-1}=\varnothing$, the finitely many locally closed strata
\begin{equation}\label{eq:descent-exponent-strata}
 S_N:=U_N\setminus U_{N-1},
 \qquad 0\leq N\leq N_*,
\end{equation}
are precisely the loci of constant minimal exponent $N$.

\item Each $S_N$ admits a finite Fitting/flattening stratification and a finite
principal minor cover on which a solution $b_N$ of
Equation~\eqref{eq:finite-descent-linear-system} is obtained by ordinary block
row reduction.  The resulting retraction $\beta_N$ and idempotent
$e_N=\alpha_N\beta_N$ are therefore rational expressions in one selected
invertible minor and in the entries of the Schur--Koszul descent matrix.
The set of all local solutions is an affine torsor under
$\ker(L_N)$; different choices give gauge-equivalent strict representatives
of the same descended perfect object.

\item The matrices $L_N$, the classes $\bar\omega_N$, and the partial
Schur--Koszul totalizations commute with coefficient base change.  Flat base
change pulls back the opens $U_N$ exactly; an arbitrary base change may make
$\bar\omega_N$ vanish at an earlier stage but cannot destroy an existing
retraction.  Along specialization inside $\operatorname{Spec}A$, the minimal
exponent can therefore increase, exactly as allowed by upper
semicontinuity.  The full finite totalization and all its differential matrices
still specialize entrywise.  Hence a rank jump may invalidate a chosen
minor or retraction formula, but it creates no additional derived
``rank-jump differential''; one passes to the appropriate exponent and
Fitting stratum.
\end{enumerate}

The integer $N_*$ is the first level for which $J_{N_*}=A$.  It is a uniform
bound on the chosen quasi-compact affine chart and is computable by the
finite module-membership search of Theorem~\ref{thm:linear-descent-retraction}.
This does not give a bound depending only on $(n,\mathbf d,e)$, nor a canonical
closed formula for the row-reduction solutions.
\end{theorem}

\begin{proof}
The map $L_N$ is a morphism between finite $A$-modules, so $Q_N$ and the cyclic
submodule $M_N$ are finite.  Localization is exact and commutes with cokernels.
Thus $\bar\omega_N$ vanishes over $A_{\mathfrak p}$ exactly when
$(M_N)_{\mathfrak p}=0$, which is equivalent to
$\mathfrak p\notin\operatorname{Supp}(M_N)$ and proves~(i).

The truncation map from the $(N+1)$-st partial totalization to the $N$-th one
carries $\alpha_{N+1}$ to $\alpha_N$.  A retraction at level $N$ therefore
composes with this truncation to give a retraction at level $N+1$, proving
\eqref{eq:descent-exponent-open-filtration}.  Descendability, in the finite
form of Theorem~\ref{thm:nonrigid-finite-stage-retract}, gives one global
working level; hence $U_{N_*}=\operatorname{Spec}A$ for some finite $N_*$.  The
identity $\{\nu\leq N\}=U_N$ proves upper semicontinuity and the finite
stratification in~(ii).

On $S_N$ the section $r_N$ lies locally in the image of $L_N$.  Simultaneous
Fitting stratification of the source, image, and cokernel makes the relevant
ranks locally constant.  Refining by a nonvanishing rank minor gives the
usual row-reduction solution and proves~(iii).  Two solutions differ by an
element of $\ker(L_N)$, while both split the same augmentation and hence
represent the same retract in $\operatorname{Perf}(A)$.

Every object entering~\eqref{eq:descent-augmented-linear-map} is obtained by
finite tensor products, totalizations, Hom complexes, and cokernels of finite
projective matrices.  These constructions commute with flat base change, and
the equations themselves pull back under arbitrary base change.  A local flat
map is faithfully flat, so the support of the finite cyclic module $M_N$ pulls
back exactly.  This proves the base-change statement.  Upper semicontinuity
allows $\nu$ to increase under specialization; entrywise specialization of the
full matrices is the perfect-complex base-change statement already used in
Proposition~\ref{prop:rank-jump-specialization}.  This proves~(iv).
\end{proof}


\begin{theorem}[No combinatorial uniform bound for the non-rigid descent exponent]\label{thm:no-uniform-descent-exponent-bound}
Assume $p$ is odd.  There is no bound for the finite-stage non-rigid flag
descent exponent depending only on the representation rank, the flag ranks, and
the number of ambient topological group generators.  More precisely, fix once
and for all a finite extension $L/K$ for which
Proposition~\ref{prop:free-two-quotient} supplies a quotient
$G_L\twoheadrightarrow F_2^{(p)}$.  The ambient group $G_L$ then has one fixed
finite generator number, independent of the family below.  Already for rank $2$
and one-step flags of rank $1$, there are coefficient charts factoring through
this same two-generator quotient for which the least working Cech/Amitsur stage
tends to infinity.

For every $N\ge2$, put
\[
 A_N^\circ:=\mathcal O_E[\varepsilon]/(\varepsilon^N),
 \qquad
 A_N:=A_N^\circ[1/p].
\]
Choose the free generators $a,b$ of the fixed quotient $F_2^{(p)}$ and fix
$\lambda=\varpi_E^m$ with $m\ge1$ large enough that
$1+\lambda M_2(A_N^\circ)$ is contained in the standard pro-$p$ congruence
subgroup (the same $m$ works for every $N$).  Define:
\begin{equation}\label{eq:unbounded-descent-matrix-family}
 \rho_N(a)
 =I+\lambda
 \begin{pmatrix}0&0\\0&1\end{pmatrix},
 \qquad
 \rho_N(b)
 =I+\lambda
 \begin{pmatrix}0&1\\ \varepsilon&0\end{pmatrix}.
\end{equation}
The universal property of $F_2^{(p)}$ makes~\eqref{eq:unbounded-descent-matrix-family}
a continuous $A_N$-linear Galois-representation family.

Its invariant-line scheme is
\begin{equation}\label{eq:unbounded-descent-flag-scheme}
 \operatorname{Flag}^{G_L}_{1}(\rho_N)
 \cong
 \operatorname{Spec}(A_N/(\varepsilon))
 \longrightarrow \operatorname{Spec}A_N.
\end{equation}
In particular this is a finite proper surjective $h$-cover.  If $I_N=(\varepsilon)$,
then the tensor-nilpotence exponent of the fiber of
$A_N\to A_N/I_N$ is at least $N$: for every $q<N$ the multiplication map
\begin{equation}\label{eq:nilpotent-descent-power-nonzero}
 I_N^{\otimes_{A_N}q}\longrightarrow A_N
\end{equation}
is nonzero, since it sends $\varepsilon^{\otimes q}$ to
$\varepsilon^q\ne0$.  Hence the descendability exponent, and therefore the
least partial-Amitsur totalization level required for a retract up to the usual
one-step indexing convention, is unbounded as $N\to\infty$.

Consequently no uniform exponent bound depending only on the fixed ambient
generator number, representation rank, and flag type can be imposed under
arbitrary nonreduced coefficient base change.  The correct general effectivity
statement is chart-specific: the finite linear retraction search terminates on
each fixed quasi-compact chart, while stronger bounds require additional
hypotheses controlling nilpotent thickness.
\end{theorem}

\begin{proof}
The matrices in~\eqref{eq:unbounded-descent-matrix-family} lie in the chosen
pro-$p$ congruence subgroup, so the universal property of the fixed quotient
$F_2^{(p)}$ gives the representation of $G_L$.  Since the action factors through
that quotient, invariance under the two displayed generator matrices is
equivalent to invariance under the whole image of $G_L$.  A line with homogeneous vector
$(x,y)$ is stable under $\rho_N(a)$ exactly when
\[
 xy=0,
\]
because the two diagonal eigenvalues differ by the unit $\lambda\in A_N^\times$.
Thus the invariant-line scheme is supported on the two coordinate sections of
$\mathbf P^1_{A_N}$.  Stability under $\rho_N(b)$ gives
\[
 \varepsilon x^2-y^2=0.
\]
On the section $y=0$ this is precisely $\varepsilon=0$, while on $x=0$ it is
$1=0$.  This proves~\eqref{eq:unbounded-descent-flag-scheme}.  Since both source
and target have one underlying point, the finite closed morphism is proper and
surjective, hence an $h$-cover.

For a quotient $A\to A/I$, the fiber in $D(A)$ is represented by $I$, and the
iterated tensor map from the fiber to the tensor unit is the multiplication
$I^{\otimes q}\to A$.  The tensor-nilpotence criterion for descendability says
that a bounded Amitsur retract forces one of these iterated fiber maps to vanish
\cite[Section~3.3]{MathewGalois}.  In the present family the map is nonzero for
every $q<N$ by~\eqref{eq:nilpotent-descent-power-nonzero}.  Hence no bound that sees only the fixed rank $2$, flag type $(1)$, and the
fixed ambient generator number of $G_L$ can control the required stage.
\end{proof}

\begin{corollary}[The descent frontier is chart-specific, not uniform]\label{cor:descent-frontier-chart-specific}
Theorems~\ref{thm:linear-descent-retraction} and
\ref{thm:descent-exponent-stratification} already give a terminating finite
algorithm and a finite bound on each fixed quasi-compact chart.  Theorem~\ref{thm:no-uniform-descent-exponent-bound}
shows that these chartwise bounds cannot be replaced, in the stated generality,
by a function of the representation/flag combinatorics alone.  Thus failure of a
universal exponent formula is structural and does not represent an unfinished
finite-atlas globalization argument.
\end{corollary}

\begin{warning}[Explicit non-rigid boundary]\label{warn:nonrigid-flag-locus}
Finite etale descent gives ordinary matrices on rigid finite-flag strata,
Theorem~\ref{thm:nonrigid-flag-hdescent} gives derived categorical descent,
Theorem~\ref{thm:nonrigid-strict-local} gives bounded local matrices, and
Theorem~\ref{thm:nonrigid-finite-stage-retract} replaces the infinite Cech
nerve by a finite partial Cech--Koszul totalization and one retract, while
Theorem~\ref{thm:schur-koszul-flag-dg} replaces every derived flag term by one
finite Schur--Koszul dg algebra with explicit size bounds, and
Theorem~\ref{thm:linear-descent-retraction} reduces the first working stage
and its retract idempotent to finite linear module-membership problems, while
Theorem~\ref{thm:descent-exponent-stratification} gives an increasing open
solvability filtration, a finite uniform bound on every quasi-compact affine
chart, inverse-minor formulas on a finite stratification, and the precise
upper-semicontinuous rank-jump law.  Theorem~\ref{thm:no-uniform-descent-exponent-bound} shows that the chartwise
finite bound cannot, under arbitrary nonreduced coefficient base change, be
replaced by a bound depending only on representation rank, flag type, and number
of group generators.  Thus the terminating finite linear solver is the correct
general effectivity statement; a prettier closed solution may exist on special
GMA charts but is not a missing theorem-level ingredient.  Theorem~\ref{thm:yoneda-gauge-transport}
shows that no additional representation-by-representation compatibility datum
is needed.  No canonical global matrix, optimal uniform exponent, or globally
minimal gauge should be imposed as a general axiom.
\end{warning}


\subsection{Transport reconstruction and the cochain-gauge boundary}

The universal transport algebra already reconstructs the representation.  What
remains at this stage is the explicit finite cochain gauge and its comparison
with the intrinsic cochains; transport reconstruction and cochain-gauge
comparison are therefore kept separate.

\begin{theorem}[Reconstruction factorization through the Cayley--Hamilton transport module]\label{thm:reconstruction-factorization}
Let $A$ be a complete Noetherian $\mathcal O_E$-algebra or a $\Qp$-affinoid
algebra equipped with a continuous $R_{\bar D}$-algebra structure, fix the
residual determinant sector $\bar D$, and let $M$ be a finite
projective $A$-module equipped with a continuous algebra homomorphism
\begin{equation}\label{eq:ch-module-action-reconstruction}
 \theta_M:
 \mathcal E^{\mathrm{CH}}_{\bar D}
 \widehat\otimes_{R_{\bar D}}A
 \longrightarrow \operatorname{End}_A(M).
\end{equation}
Then:

\begin{enumerate}[label=\textnormal{(\roman*)}]
\item The formula
\begin{equation}\label{eq:ch-transport-reconstruction}
 \rho_M(g):=\theta_M\!\left(\rho^{\mathrm{CH}}_{\bar D}(g)\right)
 \qquad(g\in G_K)
\end{equation}
reconstructs a continuous $A$-linear representation.  Every representation
in the residual determinant sector arises in this way, and morphisms are
exactly the $\mathcal E^{\mathrm{CH}}_{\bar D}$-linear maps between the
corresponding modules.  The construction commutes with coefficient base
change.  Because the quotient, $\mathfrak m_{\bar D}$-adic, and transport
image topologies on $\mathcal E^{\mathrm{CH}}_{\bar D}$ are equivalent, this
is a reconstruction in the completed transport category, not merely an
algebraic reconstruction of a dense subgroup.

\item Let $\mathcal F(M)$ be any perfect coefficient construction generated
from $M$ by finite sums, tensor products, duals, cones, and derived internal
Hom; in particular one may take $\mathcal F(M)=M$ or
$\mathcal F(M)=\operatorname{End}_A(M)$.  The finite full-transport complex
\begin{equation}\label{eq:reconstruction-ft-cochain}
 C_{\mathrm{FT}}^\bullet\bigl(G_K,\mathcal F(M)\bigr)
 =\operatorname{Hom}^{\mathrm{cont}}_{A[[G_K]]}
 \left(A\widehat\otimes_{\mathcal O_E}P_\bullet^\circ,
       \mathcal F(M)\right)
\end{equation}
is naturally quasi-isomorphic to
$R\Gamma_{\mathrm{cont}}(G_K,\mathcal F(M))$ and commutes with derived
coefficient base change.  A chosen comparison of the finite projective
resolution with the normalized bar resolution reconstructs the intrinsic
cochain object.  For algebra and adjoint coefficients, the diagonal and
homotopy-transfer constructions of Theorem~\ref{thm:abstract-diagonal}
reconstruct the intrinsic $A_\infty$ and $L_\infty$ operations up to coherent
equivalence.

\item The adjoint cochain carrier, even together with all of its transferred
higher operations, is not conservative on representations.  Indeed, already over $A=E$, let
$\chi_1\ne\chi_2:G_K\to E^\times$ be two rank-one characters with the same
residual reduction; for example, take two distinct sufficiently small
unramified lifts of one residual character.  Then
\begin{equation}\label{eq:rankone-adjoint-nonreconstruction}
 \operatorname{ad}A(\chi_1)
 \cong A
 \cong \operatorname{ad}A(\chi_2)
\end{equation}
with trivial $G_K$-action.  Hence their adjoint cochain dgas and controlling
$L_\infty$ algebras are identical, although the representations are not
isomorphic.

\item Consequently faithful reconstruction separates into two layers:
\begin{equation}\label{eq:two-layer-reconstruction}
\begin{gathered}
 \boxed{\text{finite full-transport module}}\\
 {}+{}\\
 \boxed{\text{finite cochain gauge with a bar/full-transport comparison}}.
\end{gathered}
\end{equation}
The first layer is supplied canonically by
Theorem~\ref{thm:ch-full-transport}; the second is supplied abstractly by
Corollary~\ref{cor:integral-full-transport-coeff} and locally by the Fox,
Herr, Koszul, and Schur--Koszul strictifications.  On non-rigid flag strata,
the descent exponent and retract idempotent strictify only the cochain-gauge
layer: the unflagged Cayley--Hamilton transport module already exists before
choosing a flag.
\end{enumerate}
\end{theorem}

\begin{proof}
Part~(i) is the universal factorization statement of
Theorem~\ref{thm:ch-full-transport}: applying the algebra action to the
universal units gives the representation, and the equivalence of the three
topologies proves continuity and completedness.  The same universal property
identifies intertwiners with module maps and is stable under coefficient
change \cite{WangEricksonFamilies,ChenevierDeterminants}.

Part~(ii) is Corollary~\ref{cor:integral-full-transport-coeff}, applied to the
indicated perfect coefficient constructions, together with the projective
resolution comparison and Theorem~\ref{thm:abstract-diagonal}.  For~(iii), a
rank-one character acts on its endomorphism module by conjugation, which is
trivial; therefore every adjoint cochain operation is the one attached to the
trivial coefficient module.  Part~(iv) follows by combining the preceding
three statements with the flag-descent and strictification theorems.
\end{proof}

\begin{corollary}[Two-layer criterion for a faithful local closure]\label{cor:two-layer-faithful-closure}
A local finite-atlas chart is faithful provided that it retains:
\begin{enumerate}[label=\textnormal{(\alph*)}]
\item the Cayley--Hamilton transport module $(M,\theta_M)$;
\item a bounded finite-projective cochain gauge $C_h^\bullet$;
\item a quasi-isomorphism
\[
 C_h^\bullet\xrightarrow{\sim}
 C_{\mathrm{FT}}^\bullet(G_K,\operatorname{End}_A(M))
\]
compatible, up to the stated coherent homotopies, with the transferred
products, brackets, tensor maps, and Morita structure.
\end{enumerate}
The completed prolongation first reconstructs $\rho_M$ by
Equation~\eqref{eq:ch-transport-reconstruction} and then reconstructs its
intrinsic adjoint cochains through Equation~\eqref{eq:reconstruction-ft-cochain}.
Thus the only issue at this stage is explicit comparison and coherence of the
finite gauge, not existence of the underlying completed transport
representation.  The universal-resolution, common-source, and recursive-native
theorems below close that comparison at the finite-atlas level.
\end{corollary}


\begin{theorem}[Universal Yoneda reduction of gauge-to-transport comparison]\label{thm:yoneda-gauge-transport}
Let $A$ be as in Theorem~\ref{thm:reconstruction-factorization}, put
$\Lambda=\Lambda_{A,K}=A[[G_K]]$, and work in the pseudocompact derived
category.  Let $P,Q\in\operatorname{Perf}(\Lambda)$ be represented by bounded
complexes of finite projective pseudocompact $\Lambda$-modules.  Define exact
$A$-linear coefficient functors
\[
 \mathcal C_P(W):=R\operatorname{Hom}^{\mathrm{cont}}_{\Lambda}(P,W),
 \qquad
 \mathcal C_Q(W):=R\operatorname{Hom}^{\mathrm{cont}}_{\Lambda}(Q,W).
\]
Then:
\begin{enumerate}[label=\textnormal{(\roman*)}]
\item Enriched derived Yoneda gives a natural equivalence
\begin{equation}\label{eq:yoneda-gauge-transport}
 \operatorname{Map}_{\operatorname{Perf}(\Lambda)}(Q,P)
 \simeq
 \operatorname{Nat}_{A}^{\mathrm{ex}}(\mathcal C_P,\mathcal C_Q).
\end{equation}
A morphism $f:Q\to P$ corresponds to precomposition
\begin{equation}\label{eq:universal-resolution-comparison}
 f_W^*:
 R\operatorname{Hom}_{\Lambda}(P,W)
 \longrightarrow
 R\operatorname{Hom}_{\Lambda}(Q,W).
\end{equation}
Thus a coefficient-natural comparison for every continuous perfect
$\Lambda$-module is determined by one universal completed-transport map.
The comparison is an equivalence for every $W$ if and only if $f$ is an
equivalence.

\item Choose local projective frames or idempotent presentations for $P$ and
$Q$.  If their differentials and $f$ are represented by matrices with entries
in $\Lambda$, then the differentials and comparison matrices of
\eqref{eq:universal-resolution-comparison} are obtained by applying the
continuous coefficient action
\[
 \theta_W:\Lambda\longrightarrow\operatorname{End}_A(W)
\]
entrywise, with the transpose or standard involution dictated by the chosen
left-module Hom convention.  This construction commutes term by term with
flat coefficient base change and by derived tensor product with arbitrary
coefficient change.

\item For $W=\operatorname{End}_A(M)$, the map $\theta_W$ is the conjugation
action induced by the Cayley--Hamilton module action $\theta_M$ of
Equation~\eqref{eq:ch-module-action-reconstruction}.  Hence every universal
resolution map evaluates to a comparison of finite adjoint gauges on the
reconstructed representation.  In particular, both finite gauges are
obtained by evaluating the same completed transport cells, and the resulting
comparison commutes with the reconstruction formula
\eqref{eq:ch-transport-reconstruction}.

\item The same statement is colored and coefficient-functorial.  Universal
maps among transport resolutions automatically induce compatible comparisons
for tensor products, duals, internal Homs, coefficient algebra objects, and
adjoint dg-Lie objects.  Diagonals and higher homotopies on the universal
resolution transport to the corresponding $A_\infty/L_\infty$ and lax
monoidal structure maps after evaluation.
\end{enumerate}
\end{theorem}

\begin{proof}
The functor $\mathcal C_P$ is the corepresentable exact functor represented by
$P$.  Enriched Yoneda identifies natural transformations from
$R\operatorname{Hom}_{\Lambda}(P,-)$ to
$R\operatorname{Hom}_{\Lambda}(Q,-)$ with maps $Q\to P$, proving~(i).
Representing all perfect objects by bounded projective complexes turns
precomposition into ordinary matrix substitution.  A continuous
$\Lambda$-action evaluates every completed-group-algebra entry, proving~(ii).
For adjoint coefficients this action is conjugation through $\theta_M$, which
proves~(iii).  Part~(iv) follows by applying the same enriched Yoneda argument
in the colored symmetric-monoidal coefficient category and then using the
diagonal and homotopy-transfer constructions already established above.
\end{proof}

\begin{corollary}[Gauge-to-transport compatibility is a universal resolution problem]\label{cor:gauge-transport-universal-problem}
Suppose a finite coefficient gauge is presented functorially as
$\mathcal C_{P_h}$ and the abstract full-transport gauge as
$\mathcal C_{P_{\mathrm{FT}}}$.  To compare them for every representation and
every perfect coefficient construction, it is enough to construct one map
\[
 P_{\mathrm{FT}}\longrightarrow P_h
 \quad\text{in }\operatorname{Perf}(A[[G_K]]).
\]
Once this map is written in completed-transport coordinates, evaluation on
the Cayley--Hamilton module produces all coefficient comparison matrices and
their base-change, tensor, dual, internal-Hom, and Morita compatibilities.
Hence one does not need a separate comparison for each representation:
transport-resolved gauges reduce to universal resolution maps in
completed-group-algebra coordinates.  Fox and Hochschild--Serre gauges are
already of this form.  Raw cyclotomic Herr gauges are handled later by the
Tate-dual finite-output perfectification and its direct Fox map, so no
Frobenius-parametrix replacement is required for finite-atlas closure; period
realizations beyond Herr remain downstream comparison functors.
\end{corollary}


\begin{theorem}[Universal comparison of finite full-transport resolutions]\label{thm:universal-transport-resolution-comparison}
Let $A$ be a complete Noetherian $\mathcal O_E$-algebra or a $\Qp$-affinoid algebra, put
$\Lambda=A[[G_K]]$, and let
\[
0\longrightarrow P_2\xrightarrow{d_2^P}P_1
\xrightarrow{d_1^P}P_0\xrightarrow{\varepsilon_P}A\longrightarrow0,
\qquad
0\longrightarrow Q_2\xrightarrow{d_2^Q}Q_1
\xrightarrow{d_1^Q}Q_0\xrightarrow{\varepsilon_Q}A\longrightarrow0
\]
be augmented length-two resolutions by finite projective pseudocompact
$\Lambda$-modules.  Then:
\begin{enumerate}[label=\textnormal{(\roman*)}]
\item There are augmentation-preserving chain maps
\[
 f:Q_\bullet\longrightarrow P_\bullet,
 \qquad
 g:P_\bullet\longrightarrow Q_\bullet
\]
which are inverse up to chain homotopy.  Their augmentation-preserving
chain-homotopy classes are unique.
\item A comparison $f$ is obtained recursively from the three finite lifting
systems
\begin{equation}\label{eq:transport-resolution-lifting-system}
 \varepsilon_Pf_0=\varepsilon_Q,
 \qquad
 d_1^Pf_1=f_0d_1^Q,
 \qquad
 d_2^Pf_2=f_1d_2^Q.
\end{equation}
If $f'$ is another solution, a homotopy $h:f\simeq f'$ is obtained from
\begin{equation}\label{eq:transport-resolution-homotopy-system}
 f_0-f_0'=d_1^Ph_0,
 \qquad
 f_1-f_1'=d_2^Ph_1+h_0d_1^Q,
 \qquad
 f_2-f_2'=h_1d_2^Q.
\end{equation}
Thus, after choosing finite free frames or idempotent presentations, both the
comparison and its choice-independence reduce to finite module equations over
$A[[G_K]]$.
\item The chain-homotopy class commutes with derived coefficient base change.
Precomposition with $f$ gives the coefficient-natural gauge comparison of
Theorem~\ref{thm:yoneda-gauge-transport}; evaluation on a Cayley--Hamilton
module gives its explicit compatibility with reconstructed transport.
\end{enumerate}
Consequently, comparison among finite projective full-transport resolutions
has no further existence obstruction.  What may remain difficult is a useful
closed solution of the lifting systems in a specified Fox, Herr, or
Cayley--Hamilton coordinate chart.
\end{theorem}

\begin{proof}
Because $Q_0$ is projective, $\varepsilon_Q$ lifts through the surjection
$\varepsilon_P$, giving $f_0$.  The map $f_0d_1^Q$ lands in
$\ker(\varepsilon_P)=\operatorname{im}(d_1^P)$; projectivity of $Q_1$ gives
$f_1$.  Likewise $f_1d_2^Q$ lands in
$\ker(d_1^P)=\operatorname{im}(d_2^P)$, and projectivity of $Q_2$ gives $f_2$.
This proves Equation~\eqref{eq:transport-resolution-lifting-system}.  The same
construction gives $g$.

The standard comparison theorem for projective resolutions constructs
homotopies $gf\simeq\operatorname{id}_Q$ and
$fg\simeq\operatorname{id}_P$.  Applied to two lifts of the identity on $A$,
it also gives Equation~\eqref{eq:transport-resolution-homotopy-system} and
uniqueness up to homotopy.  All constructions are finite lifting problems
between finite projectives, hence become finite matrix equations after free or
idempotent presentation.  Tensoring the maps and homotopies gives the
base-change assertion; the final statement is Theorem~\ref{thm:yoneda-gauge-transport}.
\end{proof}

\begin{corollary}[Strict Fox--full-transport deformation retract]\label{cor:fox-full-transport-retract}
Retain Theorems~\ref{thm:full-fox-projector} and
\ref{thm:syzygy-augmented-fox}, and abbreviate
$\Lambda=\Lambda_{K,E}^{\circ}$,
$K=K_2^{\mathrm{Fox}}$, and $e=e_{\mathrm{Fox}}$.  Let
$F_\bullet^{\mathrm{aug}}$ be the augmented Fox resolution
\[
0\longrightarrow K\longrightarrow\Lambda^m
\xrightarrow{d_2^{\mathrm{Fox}}}\Lambda^d
\xrightarrow{d_1^{\mathrm{Fox}}}\Lambda
\longrightarrow\mathcal O_E\longrightarrow0,
\]
and let $P_\bullet^{\mathrm{red}}$ be the reduced resolution with degree-two
term $e\Lambda^m$.  Define
\begin{equation}\label{eq:fox-full-retract-maps}
 r_3=0,
 \quad r_2=e,
 \quad r_1=r_0=\operatorname{id},
 \qquad
 j_2:e\Lambda^m\hookrightarrow\Lambda^m,
 \quad j_1=j_0=\operatorname{id}.
\end{equation}
Then $r:F_\bullet^{\mathrm{aug}}\to P_\bullet^{\mathrm{red}}$ and
$j:P_\bullet^{\mathrm{red}}\to F_\bullet^{\mathrm{aug}}$ are chain maps,
$rj=\operatorname{id}$, and
\begin{equation}\label{eq:fox-augmented-explicit-homotopy}
 \operatorname{id}_{F^{\mathrm{aug}}}-jr=dH+Hd,
 \qquad
 H_2=(1-e):\Lambda^m\longrightarrow K=(1-e)\Lambda^m,
 \quad H_i=0\ (i\ne2).
\end{equation}
Thus the syzygy-augmented Fox resolution is the strict direct sum of the
reduced full-transport resolution and the contractible pair
$K\xrightarrow{\operatorname{id}}K$.  After applying continuous
$\operatorname{Hom}_\Lambda(-,W)$, Equations~\eqref{eq:fox-full-retract-maps}--
\eqref{eq:fox-augmented-explicit-homotopy} give a coefficient-natural strict
deformation retract for every perfect coefficient object $W$; its matrices are
obtained by evaluating $e$ on the coefficient action.
\end{corollary}

\begin{proof}
The identity $d_2^{\mathrm{Fox}}e=d_2^{\mathrm{Fox}}$ makes $r$ a chain map,
and $j$ is the evident inclusion.  The splitting
$\Lambda^m=e\Lambda^m\oplus(1-e)\Lambda^m$ identifies
$K$ with the second summand.  Equation~\eqref{eq:fox-augmented-explicit-homotopy}
then contracts the identity complex in degrees three and two.  Continuous Hom
reverses the displayed maps and preserves the homotopy identities.
\end{proof}

\begin{corollary}[The Fox/full-transport comparison is closed]\label{cor:fox-full-comparison-closed}
The comparison between the completed Fox gauge and a finite full-transport
resolution has no remaining chain-level existence problem.  One may either
use the canonical syzygy-augmented model, or choose $e_{\mathrm{Fox}}$ and use
the strict retract above.  Finite-type projective frames and their higher
coherences are supplied later on the pro-$p$, determinantal, and
Cayley--Hamilton charts.  Computing one especially convenient projector is
therefore a local gauge optimization, not a Fox-side atlas obstruction.
\end{corollary}

\begin{warning}[Cohomological compression alone is not faithful]\label{warn:cochain-not-transport-reconstruction}
No finite adjoint cochain complex, Kuranishi map, or transferred
$A_\infty/L_\infty$ algebra can serve by itself as a faithful representation
carrier.  Equation~\eqref{eq:rankone-adjoint-nonreconstruction} already gives
a rank-one counterexample.  The transport module or an equivalent holonomy
object must be retained as an independent layer.
\end{warning}

\subsection{Why the pseudodeformation ring alone is not faithful}

\begin{warning}[Center-only closure loses extension data]\label{warn:pseudorep-not-faithful}
The center $R_{\bar D}$ forgets nonsplit extensions.  In a two-block GMA it loses the individual modules $\mathcal B$ and $\mathcal C$, retaining only coarse products such as the reducibility ideal.  The faithful transport carrier is the noncommutative algebra $\mathcal E^{\mathrm{CH}}_{\bar D}$.
\end{warning}

This leaves a narrower raw-period comparison boundary, not an atlas-existence problem.

\begin{remark}[Raw-Herr multiplicative comparison boundary]\label{rem:ch-kuranishi-boundary}
The linear raw Herr--to--native-core comparison is closed on constant-rank Tate-duality charts by Theorem~\ref{thm:raw-herr-all-degree-perfectification} and Corollary~\ref{cor:raw-herr-all-degree-fox}.  Theorem~\ref{thm:raw-herr-abstract-ainfty-existence} closes representative-level multiplicative existence, and Theorem~\ref{thm:raw-herr-phantom-reduced-normalization} canonically identifies its finite-native linear class with the direct Tate-dual map.  Thus no globalization, projector, normal-form, additional Herr perfectification, or finite-dimensional compatibility condition remains.  Remark~\ref{rem:raw-herr-exact-gauge-boundary} isolates only the optional demand for literal equality on one unreduced Robba representative.
\end{remark}

Theorem~\ref{thm:finite-lci-ch-atlas} below establishes finite algebraization and pointwise finite-perfect carriers.  The global recursive assembly is closed later by Theorem~\ref{thm:recursive-native-globalization}; the stronger exact raw-period gauge is recorded only once, in Remark~\ref{rem:optional-exact-raw-herr-normalization}; it is not part of the globalization theorem.


\section{Cayley--Hamilton atlases}\label{sec:lci-ch-atlas}

The generic-matrix representation scheme places the finite Cayley--Hamilton transport and the pointwise lci deformation theory on one relative finite-type atlas.  The comparison is coefficient-relative and has a one-parameter correction at characteristic-$p$ local-field points.

Fix a residual determinant $\bar D$ of dimension $n$.  Let
\[
X_{\bar D}^{\mathrm{gen}}
:=\operatorname{Rep}^{\square}_n
\bigl(\mathcal E^{\mathrm{CH}}_{\bar D}\bigr)
=\operatorname{Spec}A_{\bar D}^{\mathrm{gen}}
\]
be the generic-matrix framed representation scheme of the universal Cayley--Hamilton algebra, and put
\[
\mathfrak X_{\bar D}
:=\bigl[X_{\bar D}^{\mathrm{gen}}/\mathrm{GL}_n\bigr].
\]
The universal algebra map
\[
\mathcal E^{\mathrm{CH}}_{\bar D}
\longrightarrow M_n(A_{\bar D}^{\mathrm{gen}})
\]
composed with \eqref{eq:universal-ch-transport} gives a universal continuous $G_K$-transport on the framed atlas.

\paragraph{Pointwise formal notation.}
Let $x\in X_{\bar D}^{\mathrm{gen}}$ be a closed framed point corresponding to $\rho_x$.  Let $\Lambda_x$ be the coefficient ring used by B\"ockle--Iyengar--Pa\v{s}k\=unas: the unramified coefficient DVR in the finite-residue case, $\kappa(x)$ in the characteristic-zero local-field case, and an $\mathcal O_E$-Cohen ring in the characteristic-$p$ local-field case.  Put
\[
\mathfrak q_x
:=\ker\!\left(
\Lambda_x\otimes_{\mathcal O_E}A_{\bar D}^{\mathrm{gen}}
\longrightarrow\kappa(x)
\right)
\]
and define
\begin{equation}\label{eq:coefficient-expanded-framed-completion}
\widehat X^{\Lambda_x}_{\bar D,x}
:=\operatorname{Spf}
\widehat{\left(\Lambda_x\otimes_{\mathcal O_E}
A_{\bar D}^{\mathrm{gen}}\right)_{\mathfrak q_x}}.
\end{equation}
For a local Artinian $\Lambda_x$-algebra $A$ with residue field $\kappa(x)$, set
\begin{equation}\label{eq:strict-frame-group}
\widehat{\mathrm{GL}}_{n,x}(A)
:=\ker\!\left(
\mathrm{GL}_n(A)\longrightarrow\mathrm{GL}_n(\kappa(x))
\right)
\end{equation}
and
\begin{equation}\label{eq:strict-gauge-formal-quotient}
\widehat{\mathfrak D}^{\mathrm{cl}}_{\bar D,x}
:=\left[
\widehat X^{\Lambda_x}_{\bar D,x}/
\widehat{\mathrm{GL}}_{n,x}
\right].
\end{equation}
If $\rho_x^{\mathrm u}$ is the universal framed deformation, define
\begin{equation}\label{eq:intrinsic-formal-cochain-carrier}
\mathcal C^{\mathrm{for},\square}_{\bar D,x}
:=R\Gamma_{\mathrm{cont}}
\bigl(G_K,\operatorname{ad}\rho_x^{\mathrm u}\bigr)
\end{equation}
and, after strict-gauge descent,
\begin{equation}\label{eq:classical-lci-formal-carrier}
\mathcal K^{\mathrm{lci}}_{\bar D,x}
:=\mathbb T_{
\widehat{\mathfrak D}^{\mathrm{cl}}_{\bar D,x}/
\operatorname{Spf}\Lambda_x}[-1].
\end{equation}

\begin{lemma}[$p$-power truncations and the nilcompletion boundary]\label{lem:min-p-power-nilcompletion}
Let $x$ be either a finite-residue point or a characteristic-$p$ local-field
point after the Cohen-ring coefficient expansion of
\eqref{eq:coefficient-expanded-framed-completion}, and write
\[
 R_x:=R_{\rho_x}^{\square}.
\]
Then $R_x$ is $p$-adically complete and separated, and for every $a\ge1$
\[
 R_{x,a}:=R_x/p^aR_x
\]
is a $p$-power-nilpotent coefficient ring.  Moreover:
\begin{enumerate}[label=\textnormal{(\roman*)}]
\item the coefficient-expanded formal neighborhood is recovered from these
thickenings by
\[
 \operatorname{Spf}R_x
 \simeq
 \varprojlim_a\operatorname{Spec}R_{x,a};
\]
\item in the characteristic-$p$ local-field case, if $\Lambda_x$ is the
chosen Cohen ring, reduction modulo $p^a$ is compatible with coefficient
expansion and completed localization: the $a$th thickening is obtained by
first replacing $\Lambda_x$ by $\Lambda_x/p^a$ and then forming the same
completed local chart;
\item after transport to the Emerton--Gee/Laurent $F$-crystal chart, Min's
classicality theorem applies separately to every $R_{x,a}$; compatible
comparison equivalences on these thickenings therefore pass, by
nilcompleteness, to the inverse-limit comparison on the $p$-nilpotent formal
locus.
\end{enumerate}
If $x$ is a characteristic-zero local-field point, $p$ is invertible in the
coefficient field and this truncation argument does not apply.
\end{lemma}

\begin{proof}
The first assertion is the defining adic completeness of the framed local
deformation ring in the two stated coefficient regimes.  Each quotient
$R_{x,a}$ is annihilated by $p^a$, so it lies in the $p$-nilpotent test
category.  For the Cohen-ring expansion, write the coefficient-expanded
local ring before completion as
$B_x=(\Lambda_x\otimes_{\mathcal O_E}A_{\bar D}^{\mathrm{gen}})_{\mathfrak q_x}$.
Because $B_x$ is Noetherian, adic completion commutes with quotient by the
finitely generated ideal $(p^a)$, giving
\[
 \widehat{B_x}/p^a\widehat{B_x}
 \cong
 \widehat{B_x/p^aB_x}.
\]
This is exactly the chart obtained from $\Lambda_x/p^a$.  Min's theorem can
therefore be applied levelwise after the stated transport.  The transition
maps are induced by reduction $R_{x,a+1}\to R_{x,a}$, so the levelwise
comparisons form a compatible inverse system.  Nilcomplete objects are
recovered from that system, and taking the inverse limit preserves the
comparison equivalence.  When $p$ is invertible there is no nontrivial
$p$-power-nilpotent tower recovering the formal neighborhood, which explains
the characteristic-zero boundary.
\end{proof}

\begin{theorem}[Finite algebraization and formal carriers]\label{thm:finite-lci-ch-atlas}
For every residual determinant sector $\bar D$ of $G_K$:
\begin{enumerate}[label=\textnormal{(\roman*)}]
\item \emph{Global algebraization.}
The scheme $X_{\bar D}^{\mathrm{gen}}$ is affine of finite type over $R_{\bar D}$ and carries the universal framed continuous representation.  The ring $A_{\bar D}^{\mathrm{gen}}$ is $\mathcal O_E$-torsion free and locally complete intersection, as is its special fiber.  The quotient
$\mathfrak X_{\bar D}=[X_{\bar D}^{\mathrm{gen}}/\mathrm{GL}_n]$
is a Noetherian algebraic stack of finite type over $\operatorname{Spec}R_{\bar D}$.  No claim is made that
$q_{\bar D}:\mathfrak X_{\bar D}\to\operatorname{Spec}R_{\bar D}$
is lci or that its ordinary global cotangent complex over $\mathcal O_E$ is perfect.

\item \emph{Pointwise deformation carriers.}
There is a natural isomorphism
\begin{equation}\label{eq:coefficient-expanded-equals-deformation}
\widehat X^{\Lambda_x}_{\bar D,x}
\cong\operatorname{Spf}R_{\rho_x}^{\square}.
\end{equation}
The quotient in \eqref{eq:strict-gauge-formal-quotient} is the pointed classical unframed deformation stack; it uses the strict change-of-frame group, not the whole algebraic group $\mathrm{GL}_n$.  Its relative tangent complex is perfect of amplitude contained in $[-1,1]$.

The complex $\mathcal C^{\mathrm{for},\square}_{\bar D,x}$ is perfect of amplitude $[0,2]$, is strict-gauge equivariant, and descends to
$\mathcal C^{\mathrm{for}}_{\bar D,x}\in\operatorname{Perf}(\widehat{\mathfrak D}^{\mathrm{cl}}_{\bar D,x})$.
For the closed-point map $i_x$,
\begin{equation}\label{eq:formal-tangent-closed-fiber}
i_x^*\mathcal C^{\mathrm{for}}_{\bar D,x}
\simeq C^\bullet_{\mathrm{cont}}(G_K,\operatorname{ad}\rho_x).
\end{equation}
The based evaluation fiber
\[
\operatorname{fib}\!\left(
\mathcal C^{\mathrm{for},\square}_{\bar D,x}
\longrightarrow\operatorname{ad}\rho_x^{\mathrm u}
\right)
\]
is the framed tangent carrier; after cancelling the degree-zero contractible pair it is represented by $[C^1\to C^2]$.  Strict-gauge descent restores the $C^0$ term exactly once.

\item \emph{Classicality comparison.}
The complex $\mathcal K^{\mathrm{lci}}_{\bar D,x}$ is perfect of amplitude $[0,2]$, and its closed fiber has automorphism, tangent, and obstruction spaces $H^0$, $H^1$, and $H^2$ of \eqref{eq:formal-tangent-closed-fiber}.  In finite-residue sectors, Galatius--Venkatesh gives
\begin{equation}\label{eq:residual-derived-discrete-comparison}
\mathcal K^{\mathrm{lci}}_{\bar D,x}
\simeq\mathcal C^{\mathrm{for}}_{\bar D,x}.
\end{equation}
More generally, Lemma~\ref{lem:min-p-power-nilcompletion} identifies the exact interface with Min's theorem: reduce the coefficient-expanded formal neighborhood modulo $p^a$, transport each $p$-power-nilpotent thickening to the equivalent Emerton--Gee/Laurent $F$-crystal formal chart, apply classicality levelwise, and then take the nilcomplete inverse limit.  This gives
\begin{equation}\label{eq:p-nilpotent-derived-classicality-comparison}
\mathcal K^{\mathrm{lci}}_{\bar D,x}
\simeq\mathcal C^{\mathrm{for}}_{\bar D,x}.
\end{equation}
This applies to finite-residue points and to characteristic-$p$ local-field points after Cohen-ring coefficient expansion.  At a characteristic-zero local-field point $p$ is invertible, so Min's $p$-nilpotent classicality theorem does not apply; there the comparison is retained as additional analytic derived data.

\item \emph{Bare completion and full transport.}
Write
\[
\widehat X^{\mathrm{alg}}_{\bar D,x}
:=\operatorname{Spf}\widehat{\mathcal O}_{X_{\bar D}^{\mathrm{gen}},x}.
\]
The coefficient-expanded and bare completions agree in the finite-residue and characteristic-zero cases.  In characteristic $p$, a choice of formal parameter gives a noncanonical smooth thickening
\begin{equation}\label{eq:charp-coefficient-thickening}
\widehat X^{\Lambda_x}_{\bar D,x}
\cong
\operatorname{Spf}\!\left(
\widehat{\mathcal O}_{X_{\bar D}^{\mathrm{gen}},x}[[T]]
\right).
\end{equation}
The bare completion need not carry a natural $\Lambda_x$-structure.

Finally, with
\[
\widetilde{\mathcal E}^{\mathrm{CH}}_{\bar D}
:=q_{\bar D}^{*}\mathcal E^{\mathrm{CH}}_{\bar D},
\qquad
\widehat{\widetilde{\mathcal E}}^{\mathrm{CH},\Lambda_x}_{\bar D,x}
:=\widetilde{\mathcal E}^{\mathrm{CH}}_{\bar D}
\big|_{\widehat{\mathfrak D}^{\mathrm{cl}}_{\bar D,x}},
\]
the pair
\begin{equation}\label{eq:pointwise-formal-carrier-pair}
\left(
\widehat{\widetilde{\mathcal E}}^{\mathrm{CH},\Lambda_x}_{\bar D,x},
\mathcal C^{\mathrm{for}}_{\bar D,x}
\right)
\end{equation}
is a finite full-transport and finite-perfect intrinsic carrier.  The lci complex is a second tangent--obstruction chart on the same formal neighborhood; on the absolutely irreducible locus the transport carrier descends Morita-theoretically.
\end{enumerate}
\end{theorem}

\begin{proof}
B\"ockle--Iyengar--Pa\v{s}k\=unas prove that the generic-matrix algebra is $\mathcal O_E$-torsion free and lci, with lci special fiber; Wang--Erickson supplies the module-finite Cayley--Hamilton algebra and its relative finite-type algebraization \cite{BIPLocal,WangEricksonFamilies}.  These are statements about the total ring over $R_{\bar D}$, not about an lci structure morphism or a globally perfect cotangent complex.  The corrigendum changes only an affiliation \cite{BIPCorrigendum}.

By Proposition~3.34 of \cite{BIPLocal}, the coefficient-expanded completion is $\operatorname{Spf}R_{\rho_x}^{\square}$.  Quotienting by strict frame changes gives the pointed unframed formal groupoid, changing cotangent amplitude from $[-1,0]$ in the framed lci chart to $[-1,1]$.  Applying the abstract full-transport resolution to the universal adjoint module gives the perfect equivariant cochain carrier and closed-fiber base change; its evaluation fiber is the framed complex, while strict gauge restores the $C^0$ term once.

Galatius--Venkatesh and the expected-dimension theorem identify the intrinsic and classical carriers in finite-residue sectors \cite{GalatiusVenkatesh,BIPLocal}.  For the larger $p$-nilpotent locus, transport the formal Galois problem through the equivalence with the Emerton--Gee/Laurent $F$-crystal chart.  Min proves classicality on truncated $p$-nilpotent animated rings and identifies the derived stack after nilcompletion \cite[Theorems~1.1 and~1.2]{MinDerivedEG}.  Lemma~\ref{lem:min-p-power-nilcompletion} makes the passage explicit: reduce the coefficient-expanded framed ring modulo $p^a$, apply Min levelwise, and recover the comparison by the compatible inverse limit.  In the characteristic-$p$ local-field case the Cohen-ring coefficient expansion commutes with these truncations; when $p$ is invertible at a characteristic-zero local-field point the argument has no $p$-nilpotent tower and yields no comparison.  BIP also identify the coefficient-expanded and bare completions in finite residue and characteristic zero, with a noncanonical one-variable thickening in characteristic $p$.  Pulling back and completing the Cayley--Hamilton algebra therefore places transport and both carriers on the same strict-gauge formal neighborhood.
\end{proof}

\begin{corollary}[Finite algebraization with pointwise intrinsic and lci carriers]\label{cor:lci-atlas-existence}
For fixed $K$, $E$, and rank $n$, finitely many residual-determinant algebraizations carry a module-finite universal transport algebra and an lci generic-matrix total space.  At every closed framed point, coefficient expansion identifies the completion with $\operatorname{Spf}R_{\rho_x}^{\square}$; the strict-gauge quotient carries the intrinsic cochain carrier and the classical lci carrier.  They agree in finite-residue sectors and, after Laurent $F$-crystal transport and nilcompletion, on the $p$-nilpotent formal locus.  The bare completion agrees in finite residue and characteristic zero, and is a noncanonical one-variable slice in characteristic $p$.

This proves finite algebraization and pointwise finite-perfect existence.  The stratified recursive globalization of these carriers is supplied by Theorem~\ref{thm:recursive-native-globalization}; the only stronger raw-period issue is the optional exact-representative normalization of Remark~\ref{rem:optional-exact-raw-herr-normalization}, which is already canonically normalized modulo phantoms by Theorem~\ref{thm:raw-herr-phantom-reduced-normalization}.
\end{corollary}

\begin{proof}
Combine Corollary~\ref{cor:finite-determinant-sectors}, Theorem~\ref{thm:finite-lci-ch-atlas}, and the universal factorization property of Theorem~\ref{thm:ch-full-transport}.
\end{proof}

\begin{warning}[Relative finite type and coefficient-relative completion]\label{warn:relative-finite-type-charp-slice}
``Finite type'' is relative to $\operatorname{Spec}R_{\bar D}$; it neither makes $q_{\bar D}$ lci nor yields a globally perfect cotangent complex over $\mathcal O_E$.  The two carriers live on the strict-gauge quotient of the coefficient-expanded completion.  In characteristic $p$ the bare completion is only a noncanonical codimension-one slice and has no canonical $\Lambda_x$-relative tangent complex.
\end{warning}

\subsection{Fox effectivity and contraction strata}\label{subsec:fox-effectivity}

After imposing the relevant pro-$p$ factorization condition, the universal formal Fox matrices algebraize on a finite-type closed Cayley--Hamilton chart.

\begin{theorem}[Cayley--Hamilton effectivity of the Fox differential]\label{thm:ch-fox-effectivity}
Let $x\in X_{\bar D}^{\mathrm{gen}}$ be a closed framed point represented by a residual homomorphism
\[
 \bar\rho:G_K\longrightarrow\mathrm{GL}_r(k_E).
\]
Choose a Sylow $p$-subgroup $S\subset\operatorname{Im}(\bar\rho)$ and put
\[
 H_S:=\bar\rho^{-1}(S),
 \qquad
 P_S:=H_S(p),
 \qquad
 N_S:=\ker(H_S\twoheadrightarrow P_S).
\]
Let $\rho^{\mathrm{CH}}_{\bar D}:G_K\to(\mathcal E_{\bar D}^{\mathrm{CH}})^\times$ be the universal Cayley--Hamilton transport and let
\begin{equation}\label{eq:pro-p-factorization-ideal}
 \mathcal J_S
 :=
 \overline{\left\langle
 \rho^{\mathrm{CH}}_{\bar D}(n)-1
 \;\middle|\;
 n\in N_S
 \right\rangle}
 \subseteq
 \mathcal E_{\bar D}^{\mathrm{CH}}
\end{equation}
be the closed two-sided ideal generated by the kernel transport.  Set
\[
 \mathcal E_{\bar D,S}^{\mathrm{CH}}
 :=
 \mathcal E_{\bar D}^{\mathrm{CH}}/\mathcal J_S,
 \qquad
 X_{\bar D,S}^{\mathrm{gen}}
 :=
 \operatorname{Rep}^{\square}_r
 (\mathcal E_{\bar D,S}^{\mathrm{CH}}).
\]
Then:
\begin{enumerate}[label=\textnormal{(\roman*)}]
\item $\mathcal J_S$ is finitely generated as a module and
$\mathcal E_{\bar D,S}^{\mathrm{CH}}$ is module-finite over $R_{\bar D}$; hence
$X_{\bar D,S}^{\mathrm{gen}}$ is a closed finite-type subscheme of
$X_{\bar D}^{\mathrm{gen}}$;
\item the formal completion of $X_{\bar D,S}^{\mathrm{gen}}$ at $x$ is the full framed deformation space
\[
 \widehat{X}_{\bar D,S,x}^{\mathrm{gen}}
 \simeq
 \operatorname{Spf}R_{\bar\rho}^{\square};
\]
\item the universal $H_S$-transport on
$X_{\bar D,S}^{\mathrm{gen}}$ factors through $P_S$;
\item for every free or Demu\v{s}kin presentation
\[
 P_S=\langle x_1,\ldots,x_d\mid r_1,\ldots,r_m\rangle_{\mathrm{pro}\text{-}p},
 \qquad m\in\{0,1\},
\]
the completed Fox derivatives evaluate under the universal matrix representation to an algebraic finite complex
\begin{equation}\label{eq:algebraic-ch-fox-complex}
 \mathcal V^{\mathrm u}
 \xrightarrow{\ d^0_{\mathrm{alg}}\ }
 (\mathcal V^{\mathrm u})^d
 \xrightarrow{\ d^1_{\mathrm{alg}}\ }
 (\mathcal V^{\mathrm u})^m
\end{equation}
on $X_{\bar D,S}^{\mathrm{gen}}$, with
\begin{align}
 d^0_{\mathrm{alg}}(v)
 &=\bigl((\rho^{\mathrm u}(x_i)-I)v\bigr)_i,\label{eq:algebraic-ch-fox-d0}\\
 d^1_{\mathrm{alg}}(v_1,\ldots,v_d)_j
 &=\sum_i
 \rho^{\mathrm u}\!\left(
 \overline{\frac{\partial r_j}{\partial x_i}}
 \right)v_i.\label{eq:algebraic-ch-fox-d1}
\end{align}
One has $d^1_{\mathrm{alg}}d^0_{\mathrm{alg}}=0$, and completion at $x$ recovers the universal formal Fox complex of Proposition~\ref{prop:universal-formal-fox-matrices}.
\end{enumerate}
Thus finite-type effectivity of the Fox differential is already available after passing to the finite pro-$p$ factorization locus; what is not canonical is the comparison among the different residual, Sylow, and presentation charts.
\end{theorem}

\begin{proof}
Module-finiteness over the Noetherian ring $R_{\bar D}$ makes $\mathcal J_S$ finitely generated and its quotient module-finite, proving the closed finite-type assertion.  On the completed deformation space the restricted image of $H_S$ is pro-$p$, so $N_S$ acts trivially and the completion factors through $P_S$.  The completed group-algebra factorization evaluates Fox derivatives algebraically; the Fox identity gives $d^1_{\mathrm{alg}}d^0_{\mathrm{alg}}=0$, and completion recovers the universal formal matrices \cite{WangEricksonFamilies}.
\end{proof}

The finite-type algebraization is also the correct place for the conormal-minor
stratification suppressed by Corollary~\ref{cor:local-transport-explicit-lci-row-reduction}.

\begin{theorem}[Finite-type transport-explicit lci row reduction]\label{thm:transport-explicit-lci-row-reduction}
Let $U=\operatorname{Spec}B$ be an affine finite-type transport chart in
$X_{\bar D}^{\mathrm{gen}}$, or in one of its finite-type closed transport
loci such as $X_{\bar D,S}^{\mathrm{gen}}$.  Suppose the chosen algebraic
transport presentation has the form
\[
 B=A/I,\qquad I=(F_1,\ldots,F_M),
\]
with $A$ Noetherian and with $I$ locally complete intersection.  Put
$\mathcal N=I/I^2$.  Then there is a finite locally closed Fitting
stratification of $U$ and, on each stratum, a finite cover by principal
conormal-minor charts $U_\alpha$ such that an $h_\alpha$-element subset
\[
 J_\alpha=\{j_1,\ldots,j_{h_\alpha}\}
\subset\{1,\ldots,M\}
\]
of the original transport equations has classes forming a basis of
$\mathcal N|_{U_\alpha}$.  On every such chart:
\begin{enumerate}[label=\textnormal{(\roman*)}]
\item the selected equations $g_{\alpha,a}=F_{j_a}$ generate the localized
ideal and form a regular sequence;
\item every redundant transport equation has an exact localized expression
\begin{equation}\label{eq:finite-type-redundant-row-expression}
 F_k=\sum_{a=1}^{h_\alpha}c_{ka}^{(\alpha)}g_{\alpha,a},
\end{equation}
where the coefficients are regular on the principal chart and are rational
expressions whose only denominators are powers of the chosen nonvanishing
conormal minor;
\item modulo $I$ the Jacobian rows satisfy
\begin{equation}\label{eq:transport-explicit-lci-row-reduction}
 dF_k=\sum_{a=1}^{h_\alpha}c_{ka}^{(\alpha)}\,dg_{\alpha,a};
\end{equation}
thus the lci Jacobian is obtained by selecting and row-reducing actual
algebraized transport rows;
\item on overlaps the selected bases differ by invertible matrices, so the
corresponding Koszul dg algebras are explicitly isomorphic.  When these
charts are compared with a finite perfect transport complex on a common
constant-rank stratum, the pairwise and higher finite-atlas transitions are
the \v{C}ech--Stasheff transitions of
Theorem~\ref{thm:cech-stasheff-gauges}.
\end{enumerate}
For the completion at a single residual point the stratification collapses
to the single-chart statement of
Corollary~\ref{cor:local-transport-explicit-lci-row-reduction}.
\end{theorem}

\begin{proof}
Stratify first by the Fitting ideals of the finite conormal module
$\mathcal N$; because the presentation is lci, on each resulting stratum it
is locally free of a fixed rank $h_\alpha$.  The finitely many
$h_\alpha\times h_\alpha$ minors of the matrix of the classes $[F_i]$
produce a finite principal cover on which one selected subset is a basis.
At every local ring of such a chart, Nakayama's lemma gives generation of
$I$ by the selected equations; since the ideal is lci of the same height,
they form a regular sequence.  The quotient $I/(g_{\alpha,1},\ldots,
g_{\alpha,h_\alpha})$ has zero stalk at every point of the chart and hence
vanishes, giving the exact expressions in (ii).  Cramer's rule on the chosen
conormal minor gives the stated denominator control modulo $I$, and any
lifts give the exact ideal identities.  Differentiation kills the terms
containing $g_{\alpha,a}$ after passage to $B$, proving (iii).  Changes of
conormal basis are invertible on overlaps, which gives the Koszul comparison;
the final coherent statement is then exactly the constant-rank transfer
formalism of Theorem~\ref{thm:cech-stasheff-gauges}.  In a complete local
residual chart one minor is already a unit by the basis selection in
Corollary~\ref{cor:local-transport-explicit-lci-row-reduction}, so no
stratification remains.
\end{proof}

Explicit transfer splittings exist on determinantal constant-rank strata, but not across rank jumps.

\begin{proposition}[Determinantal contraction strata]\label{prop:determinantal-contraction}
Let $A$ be Noetherian and let
\[
 C^\bullet:
 C^0\xrightarrow{d^0}C^1\xrightarrow{d^1}C^2
\]
be a complex of finite projective $A$-modules.  There is a finite locally closed stratification
\[
 \operatorname{Spec}A
 =
 \coprod_{\alpha}S_{\alpha}
\]
defined by Fitting ideals of the differentials and cohomology modules such that:
\begin{enumerate}[label=\textnormal{(\roman*)}]
\item on each $S_{\alpha}$, the ranks of $d^0,d^1$ are constant and every $H^i(C^\bullet)|_{S_{\alpha}}$ is finite locally free;
\item Zariski-locally on $S_{\alpha}$ there are decompositions
\[
 C^i
 \cong
 B^i\oplus H^i\oplus L^i
\]
for which $d:L^i\xrightarrow{\sim}B^{i+1}$ and $d$ vanishes on $B^i\oplus H^i$;
\item after choosing an invertible maximal minor of each nonzero differential, the inverse minor gives an explicit contraction
\[
 \left(
 H^\bullet(C)
 \mathop{\rightleftarrows}^{i}_{\pi}
 C^\bullet,
 h
 \right),
 \qquad
 \operatorname{id}-i\pi=dh+hd;
\]
\item if $C^\bullet$ is supplied with a dga or $A_\infty$ comparison to a multiplicative model, then the transferred operations on a stratum are explicit rational expressions in the matrix entries of the differential and the chosen inverse minors.  They commute with every base change that remains inside the same rank stratum.
\end{enumerate}
Applied to the algebraic Fox complex \eqref{eq:algebraic-ch-fox-complex}, this gives computable transfer homotopies and higher operations on a finite Fitting stratification of every pro-$p$ factorization chart.
\end{proposition}

\begin{proof}
Fitting stratification makes the relevant cokernels and cohomology modules locally free of constant rank \cite{StacksFitting}.  Inverting a nonzero maximal minor puts the differentials in split block form, yielding the stated decomposition and inverse-block homotopy.  Rooted-tree transfer then makes every higher operation rational in the matrix entries and chosen minors.
\end{proof}

\begin{warning}[No contraction across a rank jump]\label{warn:no-global-rank-jump-contraction}
For $C=[k[t]\xrightarrow{t}k[t]]$, the generic fiber is acyclic while $H^1(C)=k[t]/(t)$ at the jump.  A global contraction would split $k[t]\twoheadrightarrow k[t]/(t)$.  Hence contractions exist only stratumwise; the global object is the perfect complex, not a zero-differential minimal model.
\end{warning}

\begin{proposition}[Explicit specialization at rank jumps for the full perfect chart]\label{prop:rank-jump-specialization}
Let $A$ be Noetherian and let $C^\bullet$ be a bounded complex of finite projective $A$-modules, displayed by matrices for its differentials.  For every morphism $A\to B$ one has
\begin{equation}\label{eq:rank-jump-derived-specialization}
 C^\bullet\otimes_A B
 \simeq
 C^\bullet\otimes_A^{\mathbf L}B,
\end{equation}
and the specialized differential matrices are obtained by applying $A\to B$ entry by entry.  Thus specialization across a Fitting-rank boundary is completely explicit at the level of the intrinsic perfect complex.

The possible failure of cohomology to commute with specialization is measured by the standard hyper-Tor spectral sequence
\begin{equation}\label{eq:rank-jump-tor-spectral-sequence}
 \operatorname{Tor}^{A}_{-p}
 \bigl(H^q(C^\bullet),B\bigr)
 \Longrightarrow
 H^{p+q}\!\left(C^\bullet\otimes_A^{\mathbf L}B\right).
\end{equation}
Consequently no extra ``rank-jump differential'' is missing.  What cannot be specialized globally is a chosen contraction to the cohomology bundles or a zero-differential minimal gauge.
\end{proposition}

\begin{proof}
Termwise projectivity gives ordinary $=$ derived tensor product and entrywise specialization.  The usual hyper-Tor spectral sequence measures the failure of cohomology to commute with base change; Warning~\ref{warn:no-global-rank-jump-contraction} excludes a global minimal gauge.
\end{proof}

\begin{corollary}[Rank-jump specialization]\label{cor:rank-jump-specialization-solved}
On every finite Fox or Koszul chart, and on every completed Cayley--Hamilton formal chart carrying a finite perfect model, specialization across rank jumps is obtained by evaluating the universal matrices and applying derived base change.  The finite perfect complex already contains the specialization data; only the comparison of stratumwise minimal gauges remains noncanonical.
\end{corollary}

\begin{corollary}[Comparison after effectivity and contraction]\label{cor:after-effectivity-contraction}
Algebraization, perfect-complex specialization, and stratumwise contraction are explicit.  By Theorem~\ref{thm:yoneda-gauge-transport}, coefficient-natural compatibility with reconstructed transport reduces to universal maps among completed-group-algebra resolutions.  The universal resolution-comparison theorem, the common-source transition theorems, and the later residual/Morita globalization supply these maps at the recursively explicit finite-native level.  Direct closed Fox--Herr or Hochschild--Serre coordinate formulas may still be useful as chart optimizations, but they are not additional comparison or globalization data.
\end{corollary}

\subsection{Fox--Kuranishi and Stasheff descent}\label{subsec:fox-kuranishi-stasheff}

On a fixed residual framed sector, Fox and lci/Koszul complexes are finite gauges of the same formal moduli problem.

\begin{theorem}[Fox--Kuranishi comparison]\label{thm:fox-kuranishi-comparison}
Fix a residual framed representation $\bar\rho$ and a Sylow chart as in Theorem~\ref{thm:family-sylow-fox}.  Let
\[
 \mathrm{ev}:R\Gamma
 \bigl(G_K,\operatorname{ad}\mathcal T^{\square}\bigr)[1]
 \longrightarrow
 \operatorname{ad}\mathcal T^{\square}[1]
\]
be evaluation at the base point of $BG_K$, and define the intrinsic framed tangent object
\begin{equation}\label{eq:intrinsic-framed-tangent}
 \mathcal T_{\bar\rho}^{\square,\mathrm{fr}}
 :=
 \operatorname{fib}(\mathrm{ev}).
\end{equation}
The restriction--normalized-corestriction equivalence is compatible with evaluation in degree zero.  Hence the family Sylow--Fox retract supplies a bounded two-term representative
\[
 \mathcal F_{\bar\rho}^{\square,\mathrm{fr}}
 \in
 \operatorname{Perf}^{[0,1]}(R_{\bar\rho}^{\square})
\]
of \eqref{eq:intrinsic-framed-tangent}.  In normalized cochain coordinates this is the based complex
$[C^1\to C^2]$; it is not defined by deleting a term from an arbitrary quasi-isomorphic perfect representative.
Let
\[
 R_{\bar\rho}^{\square}
 \cong
 \mathcal O[[x_1,\ldots,x_r]]/(f_1,\ldots,f_s)
\]
be any local complete-intersection presentation from Theorem~\ref{thm:local-lci}, and let
\[
 \mathcal K_{\bar\rho}^{\mathrm{lci}}
 =
 \left[
 (R_{\bar\rho}^{\square})^r
 \xrightarrow{\ J_f\ }
 (R_{\bar\rho}^{\square})^s
 \right]
\]
be the dual cotangent/Jacobian complex in degrees $[0,1]$.  Then:
\begin{enumerate}[label=\textnormal{(\roman*)}]
\item both complexes identify with the intrinsic framed tangent object, and hence there is an equivalence in the perfect derived category
\begin{equation}\label{eq:fox-kuranishi-perfect-comparison}
 \mathcal F_{\bar\rho}^{\square,\mathrm{fr}}
 \simeq
 \mathcal K_{\bar\rho}^{\mathrm{lci}}
 \simeq
 \mathbb T_{\operatorname{Spf}R_{\bar\rho}^{\square}/\operatorname{Spf}\mathcal O};
\end{equation}
\item after completion at every characteristic-zero or residual point, the fibers of \eqref{eq:fox-kuranishi-perfect-comparison} identify the Fox deformation cocycles and obstructions with the tangent and obstruction spaces of the lci equations;
\item on every determinantal stratum on which a contraction of the adjoint Fox complex is chosen, homotopy transfer produces a finite-dimensional Kuranishi map
\[
 \kappa_{\mathrm{Fox}}:
 H^1(G_K,\operatorname{ad}\rho)
 \longrightarrow
 H^2(G_K,\operatorname{ad}\rho)
\]
whose Maurer--Cartan zero locus is formally isomorphic to the completion of the lci chart.  Consequently the regular sequence $(f_1,\ldots,f_s)$ and the components of $\kappa_{\mathrm{Fox}}$ differ only by formal changes of source coordinates and an invertible change of obstruction coordinates.
\end{enumerate}
The common intrinsic tangent object in (i) is canonical; a chosen equivalence between the two displayed finite complexes, and in particular a strict matrix quasi-isomorphism between the two displayed complexes requires choices and is not canonical.
\end{theorem}

\begin{proof}
The framed tangent object is the evaluation fiber, represented by $[C^1\to C^2]$.  The family Sylow--Fox equivalence respects evaluation, while homotopy discreteness of the lci deformation ring identifies its Jacobian complex with the same intrinsic object.  This proves the perfect comparison and its fibers.  On a constant-rank stratum, transferred dg-Lie structure yields the Kuranishi map; since it and the lci equations represent the same Maurer--Cartan functor, their zero loci differ by formal source and obstruction coordinates \cite{WangErickson}.
\end{proof}

\begin{warning}[No canonical equality between Fox and lci equations]\label{warn:no-canonical-fox-lci-equations}
Fox matrices and lci equations depend on presentations and coordinates.  The intrinsic perfect object is canonical; termwise matrix equality is only a gauge choice.
\end{warning}

\begin{theorem}[\v{C}ech--Stasheff transitions]\label{thm:cech-stasheff-gauges}
Let $S$ be a fixed constant-rank stratum of a finite algebraic Fox or perfect Cartan complex, and let $\{U_\alpha\}$ be a finite affine cover by minor charts carrying explicit contractions
\[
 (H_\alpha
 \mathop{\rightleftarrows}^{i_\alpha}_{p_\alpha}
 C|_{U_\alpha},h_\alpha).
\]
Transfer the dga or dg-Lie structure of $C$ to $H_\alpha$.  Then:
\begin{enumerate}[label=\textnormal{(\roman*)}]
\item on every overlap $U_{\alpha\beta}$ there is an explicit $A_\infty$ (respectively $L_\infty$) quasi-isomorphism
\[
 F^{\alpha\beta}:H_\alpha|_{U_{\alpha\beta}}
 \longrightarrow
 H_\beta|_{U_{\alpha\beta}}
\]
whose linear term is
\begin{equation}\label{eq:stashed-transition-linear}
 F^{\alpha\beta}_1=p_\beta i_\alpha;
\end{equation}
its higher components are finite rooted-tree expressions in the multiplication or bracket and in the matrices $i_\alpha,p_\alpha,h_\alpha,i_\beta,p_\beta,h_\beta$;
\item after choosing the transfer homotopies extending
$i_\beta p_\beta\simeq\operatorname{id}_C$, the two maps
$F^{\beta\gamma}F^{\alpha\beta}$ and $F^{\alpha\gamma}$ on a triple overlap are joined by an explicit $A_\infty$ or $L_\infty$ homotopy.  Each fixed higher coherence can likewise be written by finite leveled-tree formulas;
\item a full homotopy-coherent diagram on the entire \v{C}ech nerve exists, because all $H_\alpha$ are connected to the common ambient homotopy algebra $C$.  To obtain a specific chain-level system of all higher coherences one must, however, choose a coherent transfer convention, for example by using a cofibrant/Boardman--Vogt type resolution of the \v{C}ech indexing diagram.  The pairwise contractions alone do not canonically determine those fillers;
\item once such a coherent convention is fixed, every component is a rational expression in the matrix entries and the chosen inverse minors, and it commutes with base change on which those minors remain invertible.  Different coherent conventions give equivalent homotopy-coherent diagrams, not literally identical formulas.
\end{enumerate}
\end{theorem}

\begin{proof}
Each contraction connects $H_\alpha$ to the common ambient homotopy algebra $C$.  Composing transferred inclusion and projection maps gives the pairwise transition and its linear term; inserting the homotopy $i_\beta p_\beta\simeq\operatorname{id}_C$ gives the triple-overlap homotopy and leveled-tree formulas \cite{BurkeTransfer,ManettiPerturbation,VokrinekHPT}.  A cofibrant resolution of the \v{C}ech diagram supplies all higher fillers.  Because the contractions use only matrix operations and inverse minors, every chosen finite component is rational and stable under base change within the minor chart.
\end{proof}

\begin{corollary}[Constant-rank coherent existence and pairwise computability]\label{cor:constant-rank-gluing-solved}
On each Fitting stratum, minimal gauges form a homotopy-coherent atlas through the common perfect complex.  Pairwise maps and any fixed finite coherence level are computable after choosing a transfer convention; a strict global formula system is not canonical.
\end{corollary}

\begin{warning}[Rank-jump boundaries require a stratified atlas]\label{warn:stratified-not-global-minimal}
Across a rank jump no common minimal locally free complex exists.  The global carrier is the perfect complex with its Fitting-stratum gauges; only the zero-differential minimal model fails to globalize.
\end{warning}

\subsection{Azumaya gerbes and Morita-valued closure}

Let $U_{\bar D}^{\mathrm{irr}}$ be the absolutely irreducible locus in the generic fiber of $\operatorname{Spec}R_{\bar D}$.  On this locus the Cayley--Hamilton algebra is Azumaya of degree $n$.  Write
\[
\mathfrak X_{\bar D,E}:=\mathfrak X_{\bar D}\times_{\operatorname{Spec}\mathcal O_E}\operatorname{Spec}E
\]
for the coefficient-generic fiber.

\begin{proposition}[The irreducible representation stack is a splitting gerbe]\label{prop:azumaya-gerbe}
Over $U_{\bar D}^{\mathrm{irr}}$, the stack $\mathfrak X_{\bar D}$ is the stack of splittings of the Azumaya algebra $\mathcal E^{\mathrm{CH}}_{\bar D}$.  It is a $\Gm$-gerbe over $U_{\bar D}^{\mathrm{irr}}$ and is etale-locally equivalent to
\[
B\Gm\times U_{\bar D}^{\mathrm{irr}}.
\]
At a point $\rho$, after such an etale localization,
\begin{equation}\label{eq:azumaya-tangent-splitting}
\mathbb T_{\mathfrak X_{\bar D,E}/\operatorname{Spec}E,\rho}
\simeq
\kappa(\rho)[1]
\oplus
\mathbb T_{U_{\bar D}^{\mathrm{irr}}/\operatorname{Spec}E,\rho}.
\end{equation}
The first summand is the scalar automorphism direction; the second is the deformation direction of the determinant point.
\end{proposition}

\begin{proof}
The Azumaya statement is Corollary~\ref{cor:ch-azumaya}.  A rank-$n$ representation of a degree-$n$ Azumaya algebra is a splitting module, unique up to tensoring by a line bundle; its automorphisms are scalars.  Hence the splitting stack is a $\Gm$-gerbe and becomes $B\Gm\times U_{\bar D}^{\mathrm{irr}}$ after an etale splitting of the Azumaya algebra.  Formula \eqref{eq:azumaya-tangent-splitting} follows by taking the tangent complex of this product.  The local triviality and Morita interpretation are standard for Azumaya algebras and extend to derived Azumaya carriers \cite{ToenAzumaya}.
\end{proof}

\begin{warning}[A universal vector bundle is too strong]\label{warn:brauer-splitting}
A Brauer class may obstruct global splitting.  The intrinsic carrier is therefore the Azumaya algebra, its $\Gm$-gerbe, or its Morita class, not a globally chosen matrix algebra.
\end{warning}

\begin{definition}[Morita-valued arithmetic Cartan carrier]\label{def:morita-carrier}
A Morita-valued arithmetic Cartan carrier over a base $S$ consists of a finite or perfect $S$-linear algebra object $\mathcal A$, considered up to derived Morita equivalence, together with full continuous transport into $\mathcal A^{\times}$ and a perfect cotangent/tangent carrier for its representation stack.  A vector-bundle realization is a local splitting chart, not part of the global definition.
\end{definition}

This Morita formulation also organizes tensor operations.  For algebra carriers $\mathcal A$ and $\mathcal B$,
\[
\mathcal A^{\vee}=\mathcal A^{\mathrm{op}},
\qquad
\underline{\operatorname{Hom}}(\mathcal A,\mathcal B)
=\mathcal A^{\mathrm{op}}\otimes\mathcal B,
\]
and tensor transport takes values in $\mathcal A\otimes\mathcal B$.

\begin{proposition}[Lax monoidal Cayley--Hamilton comparison]\label{prop:ch-lax-monoidal}
Let $\bar D_1$ and $\bar D_2$ be residual determinant sectors of dimensions $n_1$ and $n_2$.  The tensor product of the universal transports defines a deformation of the residual tensor determinant over
$R_{\bar D_1}\widehat\otimes R_{\bar D_2}$ and hence canonical comparison maps
\[
R_{\overline{D_1\otimes D_2}}
\longrightarrow
R_{\bar D_1}\widehat\otimes R_{\bar D_2}
\]
and
\[
\mathcal E^{\mathrm{CH}}_{\overline{D_1\otimes D_2}}
\widehat\otimes_{R_{\overline{D_1\otimes D_2}}}
(R_{\bar D_1}\widehat\otimes R_{\bar D_2})
\longrightarrow
\mathcal E^{\mathrm{CH}}_{\bar D_1}
\widehat\otimes
\mathcal E^{\mathrm{CH}}_{\bar D_2}.
\]
Dual transport similarly gives a comparison with
$(\mathcal E^{\mathrm{CH}}_{\bar D_1})^{\mathrm{op}}$.  These maps make the Cayley--Hamilton carriers lax symmetric monoidal in the Morita category.  They need not be isomorphisms or Morita equivalences on every stratum.
\end{proposition}

\begin{proof}
Tensor and dual are functorial polynomial laws on determinants, hence induce the ring maps.  Their universal representations satisfy the relevant Cayley--Hamilton identities, giving the algebra maps; coherence comes from tensor products of representations.  Surjectivity onto an entire target sector is not expected.
\end{proof}

\begin{corollary}[Cayley--Hamilton tensor transport is the finite monoidal replacement]\label{cor:ch-finite-monoidal-replacement}
The module-finite Cayley--Hamilton algebras and the lax monoidal comparison maps of Proposition~\ref{prop:ch-lax-monoidal} provide finite transport carriers for tensor products, duals, and internal Homs over the corresponding finite-type bases.  In view of Theorem~\ref{thm:infinite-diagonal-transport-rank}, this is the correct finite transport-side replacement for the infinitely generated $A[[G_K]]$ tensor target: combine the lax Morita-monoidal Cayley--Hamilton atlas with finite coefficient gauges and homotopy transfer.  One should not require the universal tensor target itself to be a bounded finite-projective $A[[G_K]]$-complex as a general axiom.
\end{corollary}


\begin{theorem}[Finite Cayley--Hamilton evaluation compression of the universal tensor target]\label{thm:ch-tensor-evaluation-compression}
Let $A$ be a complete Noetherian $\mathcal O_E$-algebra or a $\Qp$-affinoid algebra, let
$\bar D_1,\ldots,\bar D_r$ be residual determinant sectors of dimensions
$n_1,\ldots,n_r$, and suppose that $A$ is equipped with continuous maps
$R_{\bar D_i}\to A$.  Put
\[
 \Lambda:=A[[G_K]],
 \qquad
 \mathcal E_{i,A}:=
 \mathcal E^{\mathrm{CH}}_{\bar D_i}
 \widehat\otimes_{R_{\bar D_i}}A,
\]
and let
\[
 q_i:\Lambda\longrightarrow \mathcal E_{i,A}
\]
be the transport-algebra map induced by the universal Cayley--Hamilton
representation.  Then, for every fixed $r\ge2$:
\begin{enumerate}[label=\textnormal{(\roman*)}]
\item The completed tensor of the transport maps
\begin{equation}\label{eq:ch-finite-tensor-evaluation-map}
 q^{(r)}:=q_1\widehat\otimes_A\cdots\widehat\otimes_A q_r:
 \widehat{\bigotimes}_A^{\,r}\Lambda
 \longrightarrow
 \mathcal E_{1,A}\widehat\otimes_A\cdots\widehat\otimes_A\mathcal E_{r,A}
\end{equation}
is continuous and $\Lambda$-linear for the diagonal transport action.  Its
target is a finite $A$-module.  On a Fitting stratum on which each
$\mathcal E_{i,A}$ is finite projective of rank $e_i$, the target is finite
projective of rank $\prod_i e_i$.

\item Let
\[
 \Delta_\Lambda^{(r)}:\Lambda\longrightarrow
 \widehat{\bigotimes}_A^{\,r}\Lambda,
 \qquad
 g\longmapsto g\otimes\cdots\otimes g,
\]
be the completed group-algebra diagonal.  The composite
$q^{(r)}\Delta_\Lambda^{(r)}$ sends $g$ to
\[
 \rho^{\mathrm{CH}}_{\bar D_1}(g)\otimes\cdots\otimes
 \rho^{\mathrm{CH}}_{\bar D_r}(g).
\]
It factors through the base change of the Cayley--Hamilton algebra of the
residual tensor determinant by the iterated lax monoidal maps of
Proposition~\ref{prop:ch-lax-monoidal}.  Thus the infinitely generated
universal tensor transport has a canonical finite evaluation after the
residual sectors are fixed.

\item Let $M_i$ be finite projective $A$-modules of ranks $n_i$, equipped with
$\mathcal E_{i,A}$-actions $\theta_i$.  Equation~\eqref{eq:ch-finite-tensor-evaluation-map}
further factors through the finite transport algebra
\begin{equation}\label{eq:finite-ch-tensor-image-algebra}
 \mathcal T_{M_1,\ldots,M_r}
 :=\operatorname{im}\!\left(
 \mathcal E_{1,A}\widehat\otimes_A\cdots\widehat\otimes_A\mathcal E_{r,A}
 \longrightarrow
 \operatorname{End}_A(M_1\otimes_A\cdots\otimes_A M_r)
 \right).
\end{equation}
After a finite determinantal/Fitting refinement on which this image has constant
rank, it is finite projective, and its rank is bounded by
\begin{equation}\label{eq:ch-tensor-evaluation-rank-bound}
 \operatorname{rk}_A\mathcal T_{M_1,\ldots,M_r}
 \le
 \min\!\left\{
 \prod_i e_i,
 \left(\prod_i n_i\right)^2
 \right\}.
\end{equation}
In particular this bound is independent of the orders of finite quotients
$G_K/U$, in contrast with the growth $|G_K/U|^{r-1}$ in
Theorem~\ref{thm:infinite-diagonal-transport-rank}.

\item Let $P_\bullet$ be a finite full-transport projective resolution and choose
finite free or idempotent presentations of its terms.  For any chosen
$r$-fold diagonal approximation
\[
 \Delta_P^{(r)}:P_\bullet\longrightarrow
 P_\bullet^{\widehat\otimes_A r},
\]
the target coordinates lie in finite direct summands of finite powers of
$\widehat\otimes_A^{\,r}\Lambda$.  Applying $q^{(r)}$ entrywise and then the
$\theta_i$ therefore turns every coefficient-evaluated component of
$\Delta_P^{(r)}$ into a finite $A$-matrix.  Chain identities, homotopies, and
any fixed coherence identities satisfied before evaluation remain valid after
evaluation.  Hence the infinite transport rank of the universal tensor target
is not an obstruction to finite coefficient-level cup, tensor, or transferred
higher-operation matrices; it is an obstruction only to treating the
universal tensor target itself as a finite $\Lambda$-module.
\end{enumerate}
\end{theorem}

\begin{proof}
Each $\mathcal E^{\mathrm{CH}}_{\bar D_i}$ is module-finite over its
pseudodeformation ring, so after base change the target of
\eqref{eq:ch-finite-tensor-evaluation-map} is finite over $A$.  On a locus
where the factors are finite projective, ranks multiply.  Diagonal
$\Lambda$-linearity follows because $g\in G_K$ acts on the target as the tensor
of the universal transports.  This proves~(i).

The universal tensor representation has determinant equal to the tensor
product determinant.  Proposition~\ref{prop:ch-lax-monoidal} is precisely the
Cayley--Hamilton factorization of this transport, and iteration gives~(ii).
For~(iii), the tensor product of the $\theta_i$ maps the finite algebra in
(i) into the finite projective module
$\operatorname{End}_A(M_1\otimes\cdots\otimes M_r)$.  The image is a finite
$A$-module; after a determinantal/Fitting refinement on which the image rank
is constant it is finite projective, and its rank is bounded both by the
source rank and by the endomorphism rank,
giving \eqref{eq:ch-tensor-evaluation-rank-bound}.

Finally, a finite projective $\Lambda$-module is a direct summand of a finite
free one.  Therefore, after choosing idempotent presentations, every target
term of $P_\bullet^{\widehat\otimes_A r}$ is a direct summand of a finite power
of $\widehat\otimes_A^{\,r}\Lambda$.  Applying
\eqref{eq:ch-finite-tensor-evaluation-map} and the coefficient actions to its
coordinates produces finite matrices.  Algebra maps preserve all displayed
chain and coherence equations, proving~(iv).
\end{proof}

\begin{corollary}[Generator-level Kronecker coordinates for finite tensor transport]\label{cor:ch-kronecker-tensor-coordinates}
On a framed generic-matrix chart, let $X_{i,j}$ denote the universal matrix of
a fixed topological generator $g_j\in G_K$ in the sector $\bar D_i$.  Then the
finite Cayley--Hamilton evaluation of tensor, dual, and internal-Hom transport
is given by the explicit matrices
\begin{equation}\label{eq:ch-kronecker-generator-formulas}
 X_{1\otimes2,j}=X_{1,j}\otimes X_{2,j},
 \qquad
 X_{1^\vee,j}=(X_{1,j}^{-1})^{\mathsf t},
 \qquad
 X_{\underline{\mathrm{Hom}}(1,2),j}
 =(X_{1,j}^{-1})^{\mathsf t}\otimes X_{2,j}.
\end{equation}
These matrices satisfy the group relations and are compatible with the lax
Morita-monoidal maps.  Thus transport-level tensor, dual, and internal-Hom coordinates are already finite and explicit.  The finite cochain diagonals, higher comparison maps, and their homotopy-coherent refinements are supplied later by the abstract diagonal, common-source transfer, and recursively explicit native-globalization theorems; only a demand for one preferred closed tower-native formula would be stronger.
\end{corollary}

\begin{warning}[Cayley--Hamilton evaluation compression is sectorwise]\label{warn:ch-evaluation-sectorwise}
Theorem~\ref{thm:ch-tensor-evaluation-compression} does not construct one
finite Hopf quotient of $A[[G_K]]$ or of
$\widehat\otimes_A^{\,r}A[[G_K]]$ that works for all representations.  The
finite target depends on the chosen residual determinant sectors and
coefficient base, and different sectors are related only through the lax
Morita-monoidal comparison maps.  This sector dependence is essential and is
compatible with Theorem~\ref{thm:infinite-diagonal-transport-rank}: the
universal tensor target remains infinitely generated even though every fixed
Cayley--Hamilton evaluation is finite.
\end{warning}

\subsection{Strong comparison boundary}

The preceding results give a relative finite-type full-transport algebraization and, at every closed framed point, intrinsic and classical finite-perfect carriers on the correctly based strict-gauge formal neighborhood.  Their comparison is known in finite residue and on the nilcomplete $p$-nilpotent locus; characteristic-zero local-field points remain analytic.  A chartwise strong deformation retract from raw continuous or period cochains to a chosen bounded carrier would be a stronger strict gauge, not a prerequisite for the intrinsic comparison.  The later stratified globalization theorem obtains coherent finite charts without any globally strict projection; the raw-Herr section isolates the residual exact-representative issue as an optional phantom-sensitive normalization.


\section{Residual atlases and wildness}\label{sec:residual-atlas}

The determinant atlas has a finite framed residual refinement.  Stable lattices explain the refinement, while the matrix-pair embedding marks the limit of global normal forms.

\subsection{Stable lattices and residual sectors}

Let $E/\Qp$ be finite, with ring of integers $\mathcal O_E$, uniformizer $\varpi_E$, and residue field $k_E$.

\begin{lemma}[Stable lattice]\label{lem:stable-lattice}
Every continuous representation
\[
\rho:G_K\longrightarrow \mathrm{GL}_n(E)
\]
admits a $G_K$-stable $\mathcal O_E$-lattice $T\subset E^n$.
\end{lemma}

\begin{proof}
The image $\rho(G_K)$ is compact.  A compact subgroup of $\mathrm{GL}_n(E)$ is bounded and therefore stabilizes an $\mathcal O_E$-lattice after conjugation.  Conjugating the standard lattice back gives $T$.
\end{proof}

The reduction $T/\varpi_ET$ is a continuous residual representation
\[
\bar\rho_T:G_K\longrightarrow \mathrm{GL}_n(k_E).
\]
It depends on the lattice, but the collection of possible residual sectors is finite.

\begin{lemma}[Finiteness of residual sectors]\label{lem:finite-residual-sectors}
For fixed $K$, $E$, and $n$, there are only finitely many continuous homomorphisms
\[
G_K\longrightarrow \mathrm{GL}_n(k_E),
\]
and hence only finitely many residual isomorphism classes.
\end{lemma}

\begin{proof}
Choose finitely many topological generators $g_1,\ldots,g_e$ of $G_K$.  A continuous homomorphism is determined by the finite tuple of images of these generators.  Since $\mathrm{GL}_n(k_E)$ is finite, there are at most $|\mathrm{GL}_n(k_E)|^e$ such homomorphisms.
\end{proof}

For each residual representation $\bar\rho$, let
\[
\mathfrak X_{\bar\rho}^{\mathrm{fr}}
\]
denote its framed deformation space equipped with its intrinsic derived formal completion.  A finite profinite presentation gives bounded classical generator-and-relation coordinates.  The derived completion is defined intrinsically by continuous cochains; it is represented by a finite Fox tangent complex only in the cohomologically Fox-complete sectors of Theorem~\ref{thm:intrinsic-fox-chart}.

\begin{theorem}[Finite framed refinement of the determinant atlas]\label{thm:finite-residual-atlas}
Fix $K$, $E$, and $n$.  Each canonical residual determinant chart of Corollary~\ref{cor:ch-atlas} is covered by the generic fibers of finitely many framed residual deformation spaces
\[
\left\{\mathfrak X_{\bar\rho_{ij}}^{\mathrm{fr}}\right\}_{j=1}^{N_i},
\qquad
\det(\bar\rho_{ij})=\bar D_i.
\]
Each framed chart has finite classical transport coordinates: finitely many bounded generator matrices satisfying finitely many relation word equations.  Its intrinsic formal derived completion is governed by
\[
C^\bullet_{\mathrm{cont}}(G_K,\operatorname{ad}\rho)[1].
\]
After passage to a free or Demu\v{s}kin pro-$p$ chart, this tangent completion is represented by the finite Fox--Spencer model.  In a general residual representation sector, Theorem~\ref{thm:local-lci} supplies a finite complete-intersection formal carrier, although not an explicit Fox transport presentation.
\end{theorem}

\begin{proof}
Stable lattices give integral models; their determinant is lattice-independent, while only finitely many residual representations occur.  Finite transport coordinates and the intrinsic derived completion then give the stated refinement.  Different lattices may retain different residual extensions, so the cover refines rather than decomposes the determinant sector.
\end{proof}

\subsection{Lattice refinement}\label{subsec:lattice-refinement}

For a fixed representation $V$, let $\mathcal L_\rho$ be the poset of $G_K$-stable $\mathcal O_E$-lattices in $V$, ordered by inclusion.

\begin{proposition}[Filtered stable-lattice nerve]\label{prop:lattice-nerve}
The poset $\mathcal L_\rho$ is nonempty and filtered.  Consequently its simplicial nerve is contractible.
\end{proposition}

\begin{proof}
Nonemptiness is Lemma~\ref{lem:stable-lattice}.  If $T_1,\ldots,T_s$ are stable lattices, then
\[
T_1+\cdots+T_s
\]
is again a finitely generated full $\mathcal O_E$-lattice and is $G_K$-stable.  It is a common upper bound.  A nonempty filtered category has contractible nerve.
\end{proof}

\begin{corollary}[Objectwise lattice-change descent after inversion]\label{cor:lattice-descent}
Let $F(T)$ be any functorial integral bar, Herr, or Fox cochain model attached to a stable lattice, with transition maps induced by inclusions.  After inverting $\varpi_E$, all $F(T)[1/\varpi_E]$ compute the same $E$-linear transport object, and the homotopy colimit over $\mathcal L_\rho$ is canonically equivalent to it.  In particular, two lattice-dependent models admit a common refinement and become canonically comparable in the derived $E$-linear category.
\end{corollary}

\begin{proof}
A common upper lattice compares any finite collection; after inverting $\varpi_E$ all models identify with $V$.  Contractibility removes the residual choice.
\end{proof}

The filtered construction is compatible with the elementary tensor operations:
\[
(T,S)\longmapsto T\otimes_{\mathcal O_E}S,
\qquad
T\longmapsto T^\vee,
\qquad
(T,S)\longmapsto\operatorname{Hom}_{\mathcal O_E}(T,S).
\]
Thus lattice choice disappears objectwise after inversion.  Integral family variation and the coherent comparison of the native Cayley--Hamilton/Fox/lci core are handled by the residual and stratified atlas constructions below.  Herr remains a downstream cohomological gauge; its direct Tate-dual linear comparison and intrinsic multiplicative comparison are both closed later, with only literal normalization on one unreduced Robba representative left as an optional strict gauge.

\begin{remark}[Finite base change and stable lattices]\label{rem:base-change-lattice}
For finite Galois $L/K$, the sum $T^{\mathrm{Gal}}=\sum_{\sigma}\sigma T$ supplies a descended stable lattice.  This removes lattice existence, not the need for coherent actions on chosen Fox or lci presentations.
\end{remark}

\subsection{A wildness theorem}

Finite atlases do not imply simple normal forms; the obstruction is already present in maximal pro-$p$ transport.

Let $F_2^{(p)}$ denote the free pro-$p$ group on two generators $a,b$.

\begin{proposition}[Free two-directional quotient after finite base change]\label{prop:free-two-quotient}
Assume $p$ is odd.  For every $p$-adic local field $K$, there exists a finite extension $L/K$ such that the maximal pro-$p$ quotient $G_L(p)$ admits a continuous quotient
\[
G_L(p)\twoheadrightarrow F_2^{(p)}.
\]
\end{proposition}

\begin{proof}
If $\mu_p\not\subset K$, take $L=K$.  Then $G_L(p)$ is free pro-$p$ of rank $[L:\Qp]+1\geq2$, so it surjects onto $F_2^{(p)}$.  If $\mu_p\subset K$, choose $L/K$ so that the Demu\v{s}kin rank
\[
d=[L:\Qp]+2
\]
is at least six.  For odd $p$, a standard minimal presentation has relation \cite{NSW}
\[
r=x_1^q[x_1,x_2][x_3,x_4]\cdots[x_{d-1},x_d].
\]
Send $x_1,x_2$ to $1$, send $x_3,x_4,x_5,x_6$ to $a,b,b,a$, respectively, and send all remaining generators to $1$.  The relation maps to
\[
[a,b][b,a]=1,
\]
while the image is generated by $a,b$.  Hence the map factors through a surjection onto $F_2^{(p)}$.
\end{proof}

\begin{remark}
The odd-prime hypothesis only avoids the dyadic Demu\v{s}kin case analysis.
\end{remark}

\begin{theorem}[Matrix-pair embedding]\label{thm:matrix-pair-embedding}
Let $L$ be as in Proposition~\ref{prop:free-two-quotient}, and fix $m\geq1$.  Define
\[
(A,B)\longmapsto \rho_{A,B},
\qquad
\rho_{A,B}(a)=I+\varpi_E^mA,
\qquad
\rho_{A,B}(b)=I+\varpi_E^mB.
\]
This assignment gives a fully faithful functor from pairs of matrices in $M_n(\mathcal O_E)$, with simultaneous conjugacies, to continuous $n$-dimensional $E$-representations of $G_L$ that factor through $F_2^{(p)}$.
\end{theorem}

\begin{proof}
The congruence subgroup is pro-$p$, so the universal property of $F_2^{(p)}$ gives the representation.  Intertwiners satisfy the two displayed conjugacy equations, equivalently simultaneous conjugacy of $(A,B)$.
\end{proof}

\begin{warning}[No global Jordan-type classification]\label{warn:wild-normal-form}
Finite-base-change-compatible classification contains simultaneous similarity of matrix pairs, a standard benchmark for wild matrix problems \cite{BelitskiiSergeichuk}.  Hence no uniform finite scalar invariant list or global Jordan form can classify all local $p$-adic representations.  Faithful finite atlases remain possible, but normal forms must be sectorwise.
\end{warning}


\section{Homotopy-coherent perfectification}\label{sec:strict-no-go}

The local carriers of the preceding section are perfect, but they do not determine a functorial chain-level retraction of the continuous bar complex.  The obstruction already occurs for finite-dimensional complexes; the intrinsic replacement is homotopy-coherent perfectification.

\subsection{Pointwise contractions}

\begin{proposition}[Algebraic pointwise contraction]\label{prop:pointwise-contraction}
Let $C^\bullet$ be a cochain complex of vector spaces over a field $E$, and assume that $H^n(C^\bullet)$ is finite-dimensional and vanishes outside a finite interval.  Then there are maps
\[
H^\bullet(C)\xrightarrow{\ i\ }C^\bullet
\xrightarrow{\ \pi\ }H^\bullet(C),
\qquad
\pi i=\mathrm{id},
\]
and a degree $-1$ homotopy $h$ such that
\begin{equation}\label{eq:pointwise-sdr}
\mathrm{id}_{C}-i\pi=dh+hd.
\end{equation}
Thus $C^\bullet$ admits a finite-dimensional strong deformation retract after forgetting topology and all multiplicative or transport structure.
\end{proposition}

\begin{proof}
Write $Z^n=\ker(d:C^n\to C^{n+1})$ and $B^n=\operatorname{im}(d:C^{n-1}\to C^n)$.  Choose decompositions
\[
Z^n=B^n\oplus H^n,
\qquad
C^n=Z^n\oplus L^n.
\]
Then $d:L^n\to B^{n+1}$ is an isomorphism.  Let $i$ be the inclusion of the chosen representatives $H^n$, let $\pi$ be projection onto $H^n$, and define $h:B^{n+1}\to L^n$ as the inverse of $d$, with $h=0$ on $H^{n+1}\oplus L^{n+1}$.  A direct check on the three summands $B^n$, $H^n$, and $L^n$ gives \eqref{eq:pointwise-sdr}.
\end{proof}

After forgetting topology, pointwise contraction is automatic but noncanonical.  Homological perturbation transfers products and brackets from a chosen contraction; it supplies neither boundedness nor the choice itself \cite{ManettiPerturbation,VokrinekHPT}.

\subsection{No functorial strict projection}

\begin{proposition}[No functorial strict projection]\label{prop:no-functorial-projection}
There is no rule assigning to every finite-dimensional cochain complex $C^\bullet$ over $E$ a strong deformation retract onto its cohomology that is natural with respect to all chain isomorphisms and induces the identity on cohomology.
\end{proposition}

\begin{proof}
Consider the two-term complex
\[
C^0=E^2\xrightarrow{\ d\ }C^1=E,
\qquad
d(x,y)=y.
\]
Its only cohomology is $H^0(C)=E e_1$.  Every projection $\pi^0:E^2\to Ee_1$ restricting to the identity on $Ee_1$ has the form
\[
\pi_a(x,y)=(x+ay)e_1
\]
for some $a\in E$.  For every $b\in E$, the maps
\[
u_b^0(x,y)=(x+by,y),
\qquad
u_b^1=\mathrm{id}_E
\]
form a chain automorphism inducing the identity on $H^0(C)$.  Naturality would require
\[
\pi_a u_b=\pi_a
\]
for every $b$, but the left side is $\pi_{a+b}$.  This is impossible.  Any natural strong deformation retract would in particular provide such a natural projection, giving the contradiction.
\end{proof}

\begin{corollary}[The literal strong strictification target is false]\label{cor:strict-target-false}
A globally functorial strict finite-atlas projection
\[
\cC^\bullet_{\mathrm{bar}}
\substack{\xrightarrow{\ \pi\ }\\[-5pt]\xleftarrow[\ i\ ]{}}
\mathcal K^\bullet,
\qquad
\pi i=\mathrm{id},
\]
natural under all coordinate changes, cannot be part of an intrinsic theory.  Strict contractions may be chosen on local charts, but their transition data must be encoded by higher homotopies.
\end{corollary}

The $p$-adic problem is therefore to construct bounded local models and coherent transition homotopies, not to prove abstract existence of a splitting.

\subsection{Global transport and formal perfectification}

Let
\[
q_{\bar D}:\mathfrak X_{\bar D}\to\operatorname{Spec}R_{\bar D},
\qquad
\widetilde{\mathcal E}^{\mathrm{CH}}_{\bar D}:=q_{\bar D}^{*}\mathcal E^{\mathrm{CH}}_{\bar D}.
\]
The pullback places global transport and the pointwise formal carriers on the same stack.

\begin{definition}[Homotopy-coherent arithmetic finite-atlas carrier]\label{def:homotopy-finite-atlas}
A homotopy-coherent arithmetic finite-atlas carrier of a residual determinant sector consists of:
\begin{enumerate}
\item a finite or perfect full-transport algebra carrier, retained up to derived Morita equivalence;
\item intrinsic perfect adjoint cochain carriers on the strict-gauge coefficient-expanded formal quotients, with closed-fiber and bare-completion comparisons;
\item the classical lci tangent carriers and their comparison with the intrinsic carriers---known on the $p$-nilpotent formal locus and additional data at characteristic-zero local-field points;
\item local strictly perfect gauges, contractions, transition maps, and all higher coherences;
\item reconstruction and homotopy-coherent compatibility with products, brackets, tensor products, duals, and internal Hom.
\end{enumerate}
A global strict projection is not intrinsic; coefficient expansion precedes the relative tangent construction, and unframed objects are strict change-of-frame quotients.
\end{definition}

\begin{theorem}[Full transport and formal carriers]\label{thm:canonical-derived-finite-atlas}
For every residual determinant sector $\bar D$, the pullback
\[
\widetilde{\mathcal E}^{\mathrm{CH}}_{\bar D}
=q_{\bar D}^{*}\mathcal E^{\mathrm{CH}}_{\bar D}
\]
carries the universal continuous $G_K$-transport.  At every closed framed point $x$, the strict-gauge formal neighborhood of Theorem~\ref{thm:finite-lci-ch-atlas} carries the perfect cochain carrier
\begin{equation}\label{eq:canonical-derived-carrier}
\mathcal C^{\mathrm{for}}_{\bar D,x}
=R\Gamma_{\mathrm{cont}}(G_K,\operatorname{ad}\rho_x^{\mathrm u})
\quad\text{after descent},
\end{equation}
and the perfect lci carrier $\mathcal K^{\mathrm{lci}}_{\bar D,x}$; both have amplitude $[0,2]$.
For the closed-point map $i_x$,
\begin{equation}\label{eq:canonical-carrier-closed-fiber}
i_x^{*}\mathcal C^{\mathrm{for}}_{\bar D,x}
\simeq C^\bullet_{\mathrm{cont}}(G_K,\operatorname{ad}\rho_x).
\end{equation}
The two carriers agree in finite-residue sectors and, after passage to the Emerton--Gee/Laurent $F$-crystal chart, throughout the nilcomplete $p$-nilpotent formal locus.  No such comparison is asserted at characteristic-zero local-field points without additional analytic derived input.  On the irreducible locus the completed transport and cochain carriers descend Morita-valuedly on the Azumaya splitting gerbe.

Thus full transport and pointwise perfect formal carriers are canonical.  Their stratified coherent globalization and recursively explicit bounded monoidal charts are supplied later by Theorem~\ref{thm:recursive-native-globalization}.  A direct characteristic-zero lci/analytic identification, when desired, remains a separate comparison problem and is not needed for the native core atlas.
\end{theorem}

\begin{proof}
The transport statement is Theorem~\ref{thm:ch-full-transport} after pullback.  Perfectness, strict-gauge descent, closed-fiber base change, the classicality comparisons, the characteristic-$p$ smooth thickening, and Morita descent are precisely the corresponding assertions of Theorem~\ref{thm:finite-lci-ch-atlas} and Proposition~\ref{prop:azumaya-gerbe}.
\end{proof}

\subsection{Analytic coefficient perfectification}\label{subsec:coefficient-perfectification}

Perfect family-level coefficient complexes on analytic charts are already known, not conjectural.

\begin{theorem}[Perfect coefficient closure in analytic families]\label{thm:kpx-coefficient-perfectness}
Let $A$ be a $\Qp$-affinoid algebra and let $\mathcal W$ be a finite projective $A$-module equipped with a continuous $A$-linear action of $G_K$.  Then
\[
\mathcal K_A^{\mathrm{coeff}}(\mathcal W)
:=R\Gamma_{\mathrm{cont}}(G_K,\mathcal W)
\]
is a perfect $A$-complex of Tor-amplitude contained in $[0,2]$.  For every morphism of $\Qp$-affinoid algebras $A\to B$, the canonical map
\begin{equation}\label{eq:coefficient-base-change}
R\Gamma_{\mathrm{cont}}(G_K,\mathcal W)
\otimes_A^{\mathbf L}B
\xrightarrow{\ \sim\ }
R\Gamma_{\mathrm{cont}}
\bigl(G_K,\mathcal W\otimes_A B\bigr)
\end{equation}
is an equivalence.  The construction is compatible with local Tate duality and is independent, as an object of $\Perf(A)$, of the cyclotomic generator used to write a Herr complex.
\end{theorem}

\begin{proof}
Attach to $\mathcal W$ its relative etale $(\varphi,\Gamma_K)$-module.  The relative Herr complex computes continuous Galois cohomology.  The finiteness theorem of Kedlaya--Pottharst--Xiao~\cite[Theorem~4.4.2 and its base-change corollaries]{KPXFamilies} places this Herr complex in $D_{\mathrm{perf}}(A)$ and proves derived base-change compatibility; the local cohomological dimension of $G_K$ bounds the amplitude by $[0,2]$.  Independence from the chosen generator of the open procyclic part of $\Gamma_K$ is given by the canonical comparison isomorphism between the corresponding Herr complexes.  The Herr chart is used here as a proof and a finite coordinate model; the intrinsic object is $R\Gamma_{\mathrm{cont}}(G_K,\mathcal W)$.
\end{proof}

\begin{corollary}[Coefficient carrier on the representation stack]\label{cor:stack-coefficient-carrier}
Let $U\to X_{\bar D}^{\mathrm{gen}}$ be an affinoid chart and let $\mathcal V_U^{\mathrm u}$ be the pullback of the universal representation.  For every perfect coefficient construction $\mathcal F$ generated from $\mathcal V_U^{\mathrm u}$ by finite sums, tensor products, duals, cones, and derived internal Hom, the complex
\[
\mathcal K_U^{\mathrm{coeff}}(\mathcal F)
:=R\Gamma\bigl(G_K,\mathcal F(\mathcal V_U^{\mathrm u})\bigr)
\]
is perfect and compatible with arbitrary affinoid base change.  These complexes are $\mathrm{GL}_n$-equivariant and descend to perfect objects on the quotient stack $\mathfrak X_{\bar D}$.
\end{corollary}

\begin{proof}
Theorem~\ref{thm:kpx-coefficient-perfectness}, extended from modules to perfect coefficient complexes by finite devissage, applies on every affinoid chart.  Equation~\eqref{eq:coefficient-base-change} supplies the overlap equivalences and their higher cocycle compatibilities in the stack $\Perf$.  Equivariance of the universal representation under change of frame gives descent to the quotient stack.
\end{proof}

\begin{proposition}[Abstract homotopy-coherent coefficient atlas]\label{prop:abstract-coeff-atlas}
On the analytic generic fiber of every residual determinant sector, the coefficient objects of Corollary~\ref{cor:stack-coefficient-carrier} admit, after an affinoid refinement, bounded finite-projective representatives.  On intersections these representatives are related by quasi-isomorphisms with all higher coherences supplied by descent in the stable $\infty$-category of perfect complexes.
\end{proposition}

\begin{proof}
A perfect complex is locally represented by a bounded complex of finite projective modules.  Perfect complexes satisfy descent, so the transition quasi-isomorphisms and their higher coherences are intrinsic to the object of $\Perf(\mathfrak X_{\bar D})$.  No globally natural strict projection is produced, in agreement with Proposition~\ref{prop:no-functorial-projection}.
\end{proof}

\begin{warning}[Perfect coefficient closure is not an explicit full-transport model]\label{warn:coeff-perfect-not-explicit}
Analytic perfectness and descent alone do not provide the integral Cayley--Hamilton/Fox native core or a direct tower-free bar/Herr comparison.  The integral native core and its stratified recursive atlas are constructed independently below; this warning only excludes using analytic perfectness by itself as that construction.
\end{warning}

\begin{warning}[Strict models are gauges, not invariants]\label{warn:strict-model-gauge}
Fox, Herr, and minimal $L_\infty$ complexes are gauges for the intrinsic cochain carrier; Koszul complexes gauge the classical lci carrier.  Comparison means quasi-isomorphisms with coherent higher transitions, not equality of strict projections.  The finite native comparison and globalization are established below.  A direct termwise identification with a particular characteristic-zero analytic/lci realization is a separate optional realization comparison and is not required by the native atlas.
\end{warning}

\section{Finite covariant atlases}

Finite Fox, lci/Kuranishi, Cayley--Hamilton, and analytic perfect carriers now exist in their respective regimes.  We call a compatible family of such finite charts, together with its transition, higher-coherence, and descent data, a \emph{finite covariant atlas} (short: \emph{finite atlas}).  On certified constant-rank residual charts, the preceding tower, multiplicative, monoidal, presentation, and syzygy results already assemble to a homotopy-coherent local finite covariant atlas; Theorem~\ref{thm:certified-local-finite-atlas} records this local closure explicitly.  The residual, Fitting, flag, and descent-exponent results then assemble these local gauges globally in the stratified recursive sense of Theorem~\ref{thm:recursive-native-globalization}.  Period and prismatic realizations remain downstream, and direct closed formulas for every comparison gauge are a stronger optional target.

\subsection{Closure precedes realization}

Closure is sought first in the arithmetic transport category, with its Spencer differential, full transport algebra $\Lambda_{K,E}$, holonomy, and contact relations.  Period, prismatic, and de Rham objects are later realizations.

\begin{remark}[Status of the multiplicative full-transport resolution]\label{rem:transport-resolution-status}
The existence and recursive finite effectivity of the full-transport branch are no longer an open problem.  Theorems~\ref{thm:abstract-full-transport-resolution}, \ref{thm:syzygy-augmented-fox}, and \ref{thm:universal-transport-resolution-comparison} provide integral finite projective resolutions and finite comparison equations; Theorem~\ref{thm:abstract-diagonal} supplies integral diagonals and transferred higher operations; Theorem~\ref{thm:ch-tensor-evaluation-compression} makes every fixed coefficient-evaluated tensor or higher-operation component finite; and the residual/Fitting/Morita constructions below globalize these data recursively.  Raw cyclotomic Herr extraction is direct in every cohomological degree by Theorem~\ref{thm:raw-herr-all-degree-perfectification}, Corollary~\ref{cor:raw-herr-all-degree-fox} gives the direct finite-native linear chain map, and Theorem~\ref{thm:raw-herr-abstract-ainfty-existence} closes representative-level multiplicative existence.  Theorem~\ref{thm:raw-herr-phantom-reduced-normalization} identifies its linear term canonically with the direct Tate-dual class at exactly the finite level retained by the atlas.  A literal equality on one unreduced Robba representative is only the optional strict gauge of Remark~\ref{rem:raw-herr-exact-gauge-boundary}.
\end{remark}

Accordingly, period, prismatic, and de Rham constructions enter only through comparison functors applied after the arithmetic transport carrier has been fixed.

\subsection{Arithmetic finite-atlas data}

Let
\[
\cC^\bullet_{\arith,\mathrm{Cart}}
=(\cC^\bullet,D_{\arith},\mathcal T_{K,E},\cI_{\mathrm{con}})
\]
be a completed arithmetic Cartan complex over $E$ (or over a complete coefficient algebra $\mathcal B$ in families), equivariant for the canonical full transport algebra $\mathcal T_{K,E}=\Lambda_{K,E}$.  Here $D_{\arith}$ is the total Spencer differential and $\cI_{\mathrm{con}}$ is the differential ideal generated by the contact relations.

\begin{definition}[Arithmetic finite-atlas datum]\label{def:arithmetic-finite-atlas}
An arithmetic finite-atlas datum on $\cC^\bullet_{\arith,\mathrm{Cart}}$ consists of
\[
\bigl(V^\bullet_{h,\arith},D_{\arith,h},\mathcal T_{K,E,h},\cI_{\mathrm{con},h}\bigr),
\qquad
V^\bullet_{h,\arith}
\xrightarrow{\,i_h\,}
\cC^\bullet_{\arith,\mathrm{Cart}}
\xrightarrow{\,\pi_h\,}
V^\bullet_{h,\arith},
\qquad
\pi_h i_h=\mathrm{id},
\]
The components have the following roles.  The complex $V^\bullet_{h,\arith}$ is finite projective degreewise over $\mathcal B$, and $D_{\arith,h}$ is its degree-one differential.  The symbol $\mathcal T_{K,E,h}$ denotes the induced continuous representation of the full transport algebra; for $T\in\mathcal T_{K,E}$, write $T_h$ for its image.  Finally, $\cI_{\mathrm{con},h}$ is a differential contact ideal.  These data satisfy
\begin{align}
\pi_hD_{\arith}&=D_{\arith,h}\pi_h,
& D_{\arith}i_h&=i_hD_{\arith,h},\label{eq:finite-atlas-d}\\
\pi_hT&=T_h\pi_h,
& Ti_h&=i_hT_h
\qquad(T\in\mathcal T_{K,E}),\label{eq:finite-atlas-transport}\\
\pi_h(\cI_{\mathrm{con}})&\subseteq \cI_{\mathrm{con},h},
& i_h(\cI_{\mathrm{con},h})&\subseteq \cI_{\mathrm{con}}.\label{eq:finite-atlas-contact}
\end{align}
Equivalently, once $i_h$ and $\pi_h$ are fixed one may define
\[
D_{\arith,h}:=\pi_hD_{\arith}i_h,
\qquad
T_h:=\pi_hTi_h,
\qquad
\cI_{\mathrm{con},h}:=\bigl\langle\pi_h(\cI_{\mathrm{con}})\bigr\rangle_{D_{\arith,h},\mathcal T_{K,E,h}},
\]
The last term is the smallest contact subobject generated by the projected contact relations and stable under both $D_{\arith,h}$ and $\mathcal T_{K,E,h}$.  When the finite datum carries the projected product, it also contains the differential ideal generated by those relations.  These formulas apply whenever the stated intertwining and ideal conditions hold.

Here ``finite'' means finite projective over the arithmetic coefficient system.  It does not mean finite-dimensional over $\Qp$ after completion and holonomy have been forgotten.
\end{definition}

These are the strict local finite-atlas conditions: a finite projective subcomplex together with a bounded commuting cochain projection.  They are used only as chart data; the intrinsic atlas is homotopy-coherent and does not require a globally natural projection.  Boundedness is taken in the chosen $p$-adic, Banach, condensed, bornological, or derived-complete topology.

\begin{proposition}[Weak closure by a commuting projection]\label{prop:weak-closure}
An arithmetic finite-atlas datum induces a finite arithmetic Cartan complex
\[
(V^\bullet_{h,\arith},D_{\arith,h},\mathcal T_{K,E,h},\cI_{\mathrm{con},h})
\]
with
\[
D_{\arith,h}^2=0.
\]
Every polynomial or homotopy-coherent identity in $D_{\arith}$ and the transport operators preserved by $i_h$ and $\pi_h$ descends to the finite complex.  In particular, for a chartwise mixed defect
\[
\Xi_T=D_{\arith}T-TD_{\arith},
\]
one has
\[
(\Xi_T)_h=\pi_h\Xi_T i_h.
\]
Hence $\Xi_T=0$ implies $(\Xi_T)_h=0$.
\end{proposition}

\begin{proof}
The cochain identity uses the intertwining relations, not merely the splitting.  From
\[
D_{\arith}i_h=i_hD_{\arith,h}
\]
one obtains
\[
i_hD_{\arith,h}^2
=D_{\arith}i_hD_{\arith,h}
=D_{\arith}^2i_h=0,
\]
and applying $\pi_h$ gives $D_{\arith,h}^2=0$.  Equivalently,
\[
D_{\arith,h}^2
=\pi_hD_{\arith}i_hD_{\arith,h}
=\pi_hD_{\arith}^2i_h=0.
\]
The mixed-defect identity follows from the two intertwining relations:
\[
D_{\arith,h}T_h-T_hD_{\arith,h}
=\pi_h(D_{\arith}T-TD_{\arith})i_h.
\]
The transport and contact statements are proved in the same way.
\end{proof}

Weak closure preserves identities but need not preserve the representation.

\subsection{Weak projection and faithful closure}

\begin{definition}[Weak and faithful closure]\label{def:faithful-closure}
An arithmetic finite-atlas datum is a \emph{weak closure} if it satisfies Definition \ref{def:arithmetic-finite-atlas}.  It is a \emph{faithful closure} if there is a functorial completed prolongation or reconstruction functor
\[
\widehat{\operatorname{Prol}}_{\arith}
\]
and an equivalence
\begin{equation}\label{eq:faithful-reconstruction}
\cC^\bullet_{\arith,\mathrm{Cart}}
\simeq
\widehat{\operatorname{Prol}}_{\arith}
(V^\bullet_{h,\arith})
\end{equation}
in the relevant derived-complete category, compatible with the complete transport algebra, profinite holonomy, the contact ideal, and the higher $L_\infty$ operations of the deformation complex.
\end{definition}

For a strict chart, faithful closure means that the finite datum reconstructs the completed object; otherwise wild inertia, extensions, and higher operations may be lost.  Proposition~\ref{prop:no-functorial-projection} makes strict closures local gauges, while Definition~\ref{def:homotopy-finite-atlas} gives the intrinsic homotopy-coherent form.  Equivalently, one may require conservative descent and a quasi-isomorphism
\[
\operatorname{RHom}_{\mathrm{Cart}}(\cC_1,\cC_2)
\longrightarrow
\operatorname{RHom}_{\mathrm{FinAtl}}(V_{h,1},V_{h,2})
\]
on the admissible image.


\subsection{Certified local closure and realization descent}\label{subsec:certified-local-finite-atlas}

The recursive native globalization theorem below should not obscure the closure statement that is
already a theorem.  The point is to retain the Cayley--Hamilton module as the
faithful transport layer and to regard Fox, Herr, non-cyclotomic, and Koszul
complexes as finite cochain gauges of the same intrinsic object.  In that form
the previously separate comparison results assemble without introducing a new
strict projection.

\begin{theorem}[Certified constant-rank local finite covariant atlas]\label{thm:certified-local-finite-atlas}
Fix a residual determinant sector and a finite affinoid cover by charts on which
the following data are simultaneously available:
\begin{enumerate}[label=\textnormal{(\alph*)}]
\item a finite projective full-transport resolution as in
Theorem~\ref{thm:abstract-full-transport-resolution}, together with the
Cayley--Hamilton transport reconstruction of
Theorem~\ref{thm:reconstruction-factorization};
\item bounded finite-projective Fox, Herr, or lci/Koszul
cochain gauges whose cohomology ranks are constant, so that the determinantal
contractions used in Theorems~\ref{thm:common-source-chart-change} and
\ref{thm:common-source-multiplicative-tower} are defined;
\item for every presentation overlap that is used, finite marked presentation
certificates as in Theorem~\ref{thm:certified-native-ordered-transition}; when
two endpoint-equivalent certificates are compared, a finite marked
two-certificate as in Theorem~\ref{thm:native-ordered-path-homotopy}; and,
whenever their next coherence is invoked, the exact-rank syzygy hypotheses of
Theorem~\ref{thm:native-ordered-third-coherence};
\item on nonsplit irreducible strata, passage to the Cayley--Hamilton splitting
gerbe whenever a module representative is needed, with descent retained
Morita-valuedly.
\end{enumerate}
Then these charts form a homotopy-coherent local arithmetic finite covariant atlas
on the selected chart nerve with the following properties.
\begin{enumerate}[label=\textnormal{(\roman*)}]
\item The faithful layer is the finite Cayley--Hamilton/full-transport carrier:
it reconstructs the continuous representation and its completed transport.
Each finite Fox, Herr, or Koszul gauge represents the
corresponding intrinsic perfect cochain object on its domain.

\item Tower change is strictly transitive on the chosen finite chain gauges up
to the identity contraction homotopies, and changes of contraction give the
explicit homotopies of Equation~\eqref{eq:chart-change-explicit-homotopy}.
All such maps commute with coefficient base change on the chosen minor charts.

\item The adjoint dga and dg-Lie deformation structures admit compatible
homotopy-coherent $A_\infty/L_\infty$ transitions.  For a finite tensor-stable
coefficient family, external cup products, evaluation, duality, and internal
Hom give a homotopy-coherent lax symmetric-monoidal atlas.  On Azumaya strata
these structures descend to the derived Morita category.

\item Certified changes of finite-word presentation act directly on the
noncommutative relation ideals by the ordered formulas of
Theorem~\ref{thm:certified-native-ordered-transition}.  Competing
presentation certificates are joined by the native path homotopies of
Theorem~\ref{thm:native-ordered-path-homotopy}, and on the stated exact-rank
syzygy charts their first nontrivial
higher ambiguity is filled by the degree-$+2$ coherence of
Theorem~\ref{thm:native-ordered-third-coherence}.

\item Consequently the atlas is independent, up to coherent equivalence, of
the selected finite cochain gauge, chart coordinate, contraction, splitting
module, and certified presentation data within this constant-rank sector.
What is intrinsic is the resulting covariant homotopy-coherent object together
with its retained full-transport Cayley--Hamilton layer, not any one strict
Fox, Herr, or Koszul projection.
\end{enumerate}
No assertion is made here that the projective Fox row-reduction idempotent is
already computable in universal transport coordinates, that all presentation
coherences possess closed native formulas in every arity, or that the
characteristic-zero lci/analytic comparison is solved.  Those are precisely
parts of the stronger global problem below.
\end{theorem}

\begin{proof}
The faithful transport statement is
Theorem~\ref{thm:reconstruction-factorization}, while existence of a length-two
finite projective full-transport resolution is
Theorem~\ref{thm:abstract-full-transport-resolution}.  The common-source
linear transitions, their strict transitivity, base-change compatibility, and
choice homotopies are Theorem~\ref{thm:common-source-chart-change}.  The
multiplicative and Lie enhancements are
the corresponding homotopy-transfer comparison, with tensor, dual, internal-Hom, and Morita compatibilities understood on the certified finite charts.

For finite-word presentation changes, Theorem~\ref{thm:certified-native-ordered-transition}
provides native ordered transition maps on the noncommutative relation
algebras and their finite-chain lifts.  The corresponding two-certificate
homotopies are Theorem~\ref{thm:native-ordered-path-homotopy}, and
Theorem~\ref{thm:native-ordered-third-coherence} supplies the displayed next
coherence on exact-rank syzygy charts.  These two families of change---tower
change and presentation change---act through the same intrinsic continuous
cochain object, while the presentation--tower Kuranishi product atlas gives
the strict interchange on the common finite-jet gauges.  Hence their product
indexing diagram admits the stated homotopy-coherent refinement.  Morita
descent removes the auxiliary splitting module on Azumaya overlaps.  The last
paragraph merely records the explicit boundaries already isolated by the
cited results and therefore adds no further existence claim.
\end{proof}


\begin{theorem}[Projector-free residual finite atlas for coefficient realizations]\label{thm:projector-free-realization-atlas}
Fix any prime $p$, a finite extension $K/\mathbf Q_p$, a coefficient field $E$, and a representation rank $r$.
For every residual representation class
\[
 \bar\rho:G_K\longrightarrow \mathrm{GL}_r(k_E)
\]
let $R_{\bar\rho}^{\square}$ be the framed deformation ring,
$\mathcal T_{\bar\rho}^{\square}$ its universal coefficient module, and let
\[
 \bigl(\mathcal C_{\bar\rho}^{\mathrm{allSyl}},
 e_{\bar\rho}^{\mathrm{allSyl}}\bigr)
\]
be the canonical all-Sylow carrier and transfer idempotent of
Theorem~\ref{thm:all-sylow-transfer}.  Put
\[
 \mathcal K_{\bar\rho}^{\mathrm{coef}}
 :=\operatorname{Im}(e_{\bar\rho}^{\mathrm{allSyl}})
 \in\operatorname{Perf}(R_{\bar\rho}^{\square}).
\]
Then:
\begin{enumerate}[label=\textnormal{(\roman*)}]
\item $\mathcal C_{\bar\rho}^{\mathrm{allSyl}}$ is represented by a finite
direct sum of free or Demu\v{s}kin Fox complexes.  Their differentials are the
universal formal matrices of
Proposition~\ref{prop:universal-formal-fox-matrices}; hence the coefficient
charts exist over the whole framed residual deformation space and require no
constant-rank determinantal refinement.

\item There is a canonical equivalence
\begin{equation}\label{eq:projector-free-coef-rgamma}
 \mathcal K_{\bar\rho}^{\mathrm{coef}}
 \simeq
 R\Gamma(G_K,\mathcal T_{\bar\rho}^{\square}),
\end{equation}
compatible with derived coefficient base change and equivariant for the
residual representation groupoid.  The object in
\eqref{eq:projector-free-coef-rgamma} is canonical in the idempotent-complete
derived category even though a chosen strict matrix representative of its
retract need not be canonical.

\item Pair $\mathcal K_{\bar\rho}^{\mathrm{coef}}$ with the pullback of the
Cayley--Hamilton transport module on the residual determinant sector.  The
Cayley--Hamilton factor reconstructs the continuous representation and all
completed transport, while
$\mathcal K_{\bar\rho}^{\mathrm{coef}}$ reconstructs its intrinsic coefficient
cochain object.  Thus this two-layer pair is faithful for representation plus
cohomology data on the residual chart, with Morita-valued descent on Azumaya
strata.

\item For fixed $K$, $E$, and $r$ there are only finitely many residual
representation classes.  Their two-layer pairs therefore form a finite
residual derived atlas for coefficient realizations.  Its construction uses
neither a full-group Fox row-reduction idempotent
$e_{\mathrm{Fox}}\in M_m(\mathcal O_E[[G_K]])$, nor a global frame for the
full-group projective syzygy, nor a choice of constant-rank contraction.

\item The derived faithful finite closure is already obtained by pairing this
canonical coefficient carrier with the Cayley--Hamilton transport layer.  What is
not thereby supplied is a single global \emph{explicit native presentation} of that
closure in completed-group-algebra coordinates.  A full-group projector or an
equivalent syzygy frame can still be useful for universal chain maps, native higher
diagonals, and closed transition formulas before coefficient evaluation; these are
strictly stronger explicitization requirements.
\end{enumerate}
\end{theorem}

\begin{proof}
Part~(i) is Definition~\ref{def:all-sylow-carrier},
Theorem~\ref{thm:family-sylow-fox}, and
Proposition~\ref{prop:universal-formal-fox-matrices}.  The key point is that
after residual Sylow restriction the universal coefficient action factors
through a maximal pro-$p$ quotient, whose local Galois group is free or
Demu\v{s}kin; the resulting Fox complex is finite free over the entire framed
deformation ring.

Part~(ii) is Theorem~\ref{thm:all-sylow-transfer} together with its family
base-change statement.  Idempotent completeness of perfect complexes makes
the derived image canonical; no choice of a splitting of the full completed
Fox relation matrix enters.  Part~(iii) follows from
Theorem~\ref{thm:reconstruction-factorization}: the Cayley--Hamilton module is
the faithful transport layer, whereas the all-Sylow image computes precisely
the intrinsic coefficient cochains.  The Morita qualification is
Proposition~\ref{prop:all-sylow-morita}.  Finiteness in~(iv) is
Corollary~\ref{cor:finite-residual-sylow-atlas}.  Finally,
Theorems~\ref{thm:full-fox-projector} and
\ref{thm:universal-transport-resolution-comparison} show exactly where a
chosen full-transport resolution remains useful, proving the boundary in~(v).
\end{proof}


\subsection{Derived faithful finite closure theorem}

\begin{theorem}[Derived faithful finite closure]\label{thm:derived-faithful-finite-closure}
Fix a finite extension $K/\Qp$, a coefficient field $E$, and a representation rank $r$.
Then the category of continuous rank-$r$ $E$-representations of $G_K$ admits a finite residual derived/Morita atlas with the following properties.
\begin{enumerate}[label=\textnormal{(\roman*)}]
\item On every residual representation sector, the faithful transport layer is the corresponding Cayley--Hamilton module, which reconstructs the representation and its completed profinite transport.
\item The intrinsic coefficient cochain object is represented canonically in the idempotent-complete perfect derived category by the all-Sylow carrier of Theorem~\ref{thm:projector-free-realization-atlas}; this representation is compatible with derived coefficient base change.
\item Pairing the two layers gives a finite residual derived chart faithful for representation plus cohomology data, with Morita-valued descent on Azumaya strata.  Since only finitely many residual representation classes occur for fixed $K,E,r$, these charts form a finite residual derived atlas.
\item On every certified constant-rank sector satisfying the hypotheses of Theorem~\ref{thm:certified-local-finite-atlas}, the derived charts refine to a homotopy-coherent finite covariant atlas carrying the compatible multiplicative, $A_\infty/L_\infty$, tensor, duality, internal-Hom, and certified presentation-transition structures proved there.
\end{enumerate}
Thus finite faithful closure exists at the derived/Morita level.  What is not asserted is the existence of one globally strict, presentation-independent, computable matrix model for all charts and all higher coherences simultaneously.
\end{theorem}

\begin{proof}
The Cayley--Hamilton reconstruction is Theorem~\ref{thm:reconstruction-factorization}.  The canonical coefficient carrier, its base-change compatibility, its pairing with the transport layer, and finiteness of the residual atlas are Theorem~\ref{thm:projector-free-realization-atlas}.  The homotopy-coherent refinement on certified constant-rank sectors is Theorem~\ref{thm:certified-local-finite-atlas}.  Morita descent is supplied by Proposition~\ref{prop:all-sylow-morita} together with the Cayley--Hamilton splitting-gerbe descent used in those theorems.  No additional existence input is required.
\end{proof}

\subsection{Stratified recursively explicit native globalization}

\begin{definition}[Recursive native explicitness]\label{def:recursive-native-explicitness}
A finite arithmetic atlas is \emph{recursively native-explicit} if, after fixing one of its finite residual/transport strata and a requested finite coherence arity, every displayed differential, transition, homotopy, higher operation, and descent idempotent is obtained from a finite sequence of the following operations:
\begin{enumerate}[label=\textnormal{(\alph*)}]
\item evaluation of finitely many transport words, Fox derivatives, or Cayley--Hamilton generator matrices;
\item inversion of explicitly declared nonvanishing minors on a principal chart;
\item solution of finitely presented linear module equations on a finite matrix complex;
\item finite rooted-, planar-, or leveled-tree formulas, or an equivalent fixed finite stage of a chosen cofibrant/Cech indexing resolution.
\end{enumerate}
The definition requires termination at every prescribed finite stage and coherent equivalence under change of the auxiliary choices.  It does not require one canonical matrix frame, one global minimal contraction, or one closed formula valid simultaneously at all coherence arities.
\end{definition}

\begin{theorem}[Stratified recursively explicit native globalization]\label{thm:recursive-native-globalization}
Fix a finite extension $K/\Qp$, a coefficient field $E$, and a representation rank $r$.  The derived faithful finite closure of Theorem~\ref{thm:derived-faithful-finite-closure} admits an integral, stratified, Morita-covariant globalization which is recursively native-explicit in the sense of Definition~\ref{def:recursive-native-explicitness}.  More precisely:
\begin{enumerate}[label=\textnormal{(\roman*)}]
\item \emph{Native core charts.}  There are finitely many residual representation sectors.  On each one, the faithful transport layer is the finite Cayley--Hamilton carrier and the intrinsic coefficient layer is the canonical all-Sylow Karoubi image.  After Sylow restriction the coefficient differential is represented by universal finite free Fox matrices on the whole framed residual deformation ring, and these matrices algebraize on the finite-type Cayley--Hamilton pro-$p$ factorization loci.  No full-$G_K$ free frame or functorial Fox projector is part of the global datum.

\item \emph{Finite stratified strictification.}  The finite-type transport/lci charts admit finite Fitting and conormal-minor stratifications on which the Fox and Koszul differentials, contractions, and row reductions are explicit finite matrices.  Non-rigid flag strata descend through the derived $h$-cover and admit finite affine strict models; after the finite descent-exponent refinement, a finite Cech--Koszul stage and its retract idempotent are obtained from explicit finite linear systems.  On effectively presented coefficient charts for which membership in finitely presented submodules is decidable, the search for a working descent stage terminates.

\item \emph{Rank-jump compatibility.}  Across a Fitting-rank boundary the unreduced bounded perfect complex is retained.  Its differential matrices specialize entrywise by derived base change, while the hyper-Tor spectral sequence records the possible cohomology jump.  Thus rank jumps require a change of stratum or minor gauge, not a new global differential or a global minimal contraction.

\item \emph{Transitions and all finite coherence levels.}  Comparisons between finite projective full-transport resolutions reduce to finite lifting systems and are unique up to finite chain homotopy.  On constant-rank charts the common-source and Cech--Stasheff constructions give explicit chain, $A_\infty$, and $L_\infty$ transitions.  After fixing a coherent cofibrant/Boardman--Vogt convention, the full Cech nerve is homotopy coherent; every prescribed Taylor component or finite coherence level is a finite rooted- or leveled-tree expression in the chart matrices and inverse minors.  Certified finite presentation changes may instead be evaluated by the direct ordered Fox--Magnus formulas and their native two- and three-coherences.  Theorem~\ref{thm:presentation-linear-comparison-contractible} shows that these stages exhaust the independent linear presentation coherence; higher linear fillers introduce no new obstruction class.  Different choices determine coherently equivalent atlases.

\item \emph{Multiplicative and Morita structure.}  The atlas carries the transported integral $E_1/A_\infty$ coefficient algebra and, after passage to $E$, the adjoint $L_\infty$ deformation structure, together with homotopy-coherent lax symmetric-monoidal tensor, duality, evaluation, and internal-Hom maps.  Cayley--Hamilton evaluation compresses every fixed tensor or higher-operation component to a finite coefficient matrix on the chosen residual sectors.  On Azumaya strata the construction descends in the derived Morita category.

\item \emph{Presentation and gauge independence.}  The invariant output is independent up to coherent equivalence of Sylow subgroup, finite presentation, minor, contraction, flag, splitting module, and descent retraction.  A chosen presentation supplies a strict native representative; the canonical syzygy-augmented Fox object and the intrinsic all-Sylow image remove any need for a functorial projector.  Herr complexes remain downstream finite cohomological gauges: Theorem~\ref{thm:fox-herr-zigzag} supplies the functorial derived comparison, Theorem~\ref{thm:raw-herr-tate-extraction-core} and Corollary~\ref{cor:raw-herr-direct-fox-core} give the direct raw $H^1$-to-Fox formula, and Theorem~\ref{thm:raw-herr-all-degree-perfectification} with Corollary~\ref{cor:raw-herr-all-degree-fox} gives the all-degree finite-output chain comparison.  Theorem~\ref{thm:raw-herr-abstract-ainfty-existence} supplies the representative-level multiplicative equivalence, while Theorem~\ref{thm:raw-herr-phantom-reduced-normalization} fixes its finite-native linear class.  No preferred strict formula on the entire unreduced Robba representative is required.
\end{enumerate}
Consequently the global finite-atlas problem is closed at the stratified recursively explicit level.  No further existence or globalization problem remains.  Remark~\ref{rem:optional-exact-raw-herr-normalization} records only the stronger choice of an exact raw-period representative.  This conclusion is compatible with Warning~\ref{warn:wild-normal-form}: no universal Jordan form, finite orbit-invariant list, globally functorial strict projection, globally free full-$G_K$ syzygy frame, or global minimal model is constructed or required.
\end{theorem}

\begin{proof}
Finiteness of the residual cover and the projector-free two-layer faithful charts are Theorem~\ref{thm:projector-free-realization-atlas}; the universal formal Fox matrices and their finite-type Cayley--Hamilton algebraization are Proposition~\ref{prop:universal-formal-fox-matrices} and Theorem~\ref{thm:ch-fox-effectivity}.  The canonical projector-free full-group replacement is Theorem~\ref{thm:syzygy-augmented-fox}, while Corollary~\ref{cor:sylow-fox-free-frames} supplies free frames on the pro-$p$ local charts.

Finite Fitting/conormal strictification and rank-jump specialization are Theorem~\ref{thm:transport-explicit-lci-row-reduction}, Proposition~\ref{prop:determinantal-contraction}, and Proposition~\ref{prop:rank-jump-specialization}.  Derived flag descent, finite strict local effectivity, finite-stage Cech--Koszul extraction, the finite Schur--Koszul Morita presentation, and the terminating retraction search are Theorems~\ref{thm:nonrigid-flag-hdescent}, \ref{thm:nonrigid-strict-local}, \ref{thm:nonrigid-finite-stage-retract}, \ref{thm:schur-koszul-flag-dg}, and \ref{thm:linear-descent-retraction}; Theorem~\ref{thm:descent-exponent-stratification} makes the working descent level a finite constructible stratification on every quasi-compact affine chart.

The universal projective-resolution comparison is Theorem~\ref{thm:universal-transport-resolution-comparison}.  The common-source chain and multiplicative transitions are Theorems~\ref{thm:common-source-chart-change} and \ref{thm:common-source-multiplicative-tower}; all higher Cech fillers at any prescribed finite level are supplied by Theorem~\ref{thm:cech-stasheff-gauges}.  The direct presentation-native branch is given by Theorems~\ref{thm:certified-native-ordered-transition}, \ref{thm:native-ordered-path-homotopy}, and \ref{thm:native-ordered-third-coherence}; Theorem~\ref{thm:presentation-linear-comparison-contractible} shows that the identity component of the derived linear comparison space is contractible; the displayed native path and third-coherence formulas are explicit strict witnesses rather than additional invariants.  Proposition~\ref{prop:decomposable-native-coherence-contraction} and Theorem~\ref{thm:relative-decomposable-ainfty-lifting} record the stronger direct-word relative lifting mechanism, but the global atlas uses the already established common-source Cech--Stasheff coherence and does not assume the extra quotient-level word condition.  Multiplicative all-Sylow transport, lax monoidal comparison, and Morita descent are Theorem~\ref{thm:all-sylow-ainfty}, Theorem~\ref{thm:common-source-lax-monoidal-tower}, Proposition~\ref{prop:ch-lax-monoidal}, and Proposition~\ref{prop:all-sylow-morita}; finite coefficient evaluation of tensor targets and fixed higher components is Theorem~\ref{thm:ch-tensor-evaluation-compression}.

Every operation appearing in these constructions is one of the finite operations listed in Definition~\ref{def:recursive-native-explicitness}.  The no-functorial-projection, no-functorial-Fox-projector, rank-jump, Brauer-gerbe, and wildness results show simultaneously why the conclusion is naturally stratified, Morita-valued, and homotopy coherent rather than a single strict normal form.  This proves the theorem.
\end{proof}

\subsection{Optional strict raw-Herr normalization}

\begin{remark}[Optional exact raw-Herr normalization]\label{rem:optional-exact-raw-herr-normalization}
The finite native atlas requires only the automatic phantom-reduced
normalization of Theorem~\ref{thm:raw-herr-phantom-reduced-normalization}.
One may ask in addition for an $A_\infty$ representative whose first Taylor
coefficient is literally the chosen map
\[
 \Theta_{HF}^{\mathrm{raw}}:
 C^\bullet_{\varphi,\gamma}(D(\mathcal A))\longrightarrow F^\bullet
\]
on the unreduced Robba complex, compatibly with family base change and change of
cyclotomic generator.  This asks for extra continuous chain-homotopy splittings
in the raw period topology.  It is a strict-gauge refinement only; the paper
neither assumes nor needs such a splitting, and its existence or failure does
not change the intrinsic multiplicative equivalence or any finite-native datum.
\end{remark}

\begin{warning}[Closed formulas are stronger than globalization]\label{warn:no-tautological-chart}
Remark~\ref{rem:optional-exact-raw-herr-normalization} is not solved by taking a $(\varphi,\Gamma)$-module as the foundational object, placing $V$ in degree zero and renaming the bar construction as prolongation, or retaining only the pseudodeformation ring.  Conversely, failure to produce one canonical strict formula does not reopen Theorem~\ref{thm:recursive-native-globalization}: strict projector choices and global normal forms are excluded on structural grounds, while the atlas is defined by coherent descent of finite gauges.
\end{warning}

\begin{remark}[Methodological tower and presentation independence]\label{rem:methodological-tower-free}
No cyclotomic, Lubin--Tate, or chosen profinite presentation is intrinsic.  Theorem~\ref{thm:recursive-native-globalization} obtains invariance by descent through the common intrinsic object.  Remark~\ref{rem:optional-exact-raw-herr-normalization} records only the optional demand for literal equality of a chosen raw-Herr chain representative; Theorem~\ref{thm:raw-herr-phantom-reduced-normalization} already supplies the canonical finite-native normalization modulo phantoms.  Presentation invariance itself is closed by the native atlas and its syzygy coherences.
\end{remark}

\subsection{Closed low-rank tests and Herr comparison}

The former low-rank reconstruction tests are theorem-level consequences of the native atlas.  We record them here as consistency checks and then summarize the final Herr comparison boundary; no additional explicitization problem is introduced in this subsection.

\begin{theorem}[Rank-one native presentation independence]\label{thm:rankone-native-independence}
Let $\chi:G_K\to E^{\times}$ be a continuous rank-one character.  Its finite native closure is independent, up to the coherent equivalence of the atlas, of the stable lattice, residual-killing/Sylow chart, finite pro-$p$ presentation, and contraction/minor gauge.  More precisely:
\begin{enumerate}[label=\textnormal{(\roman*)}]
\item every stable $\mathcal O_E$-lattice in the one-dimensional $E$-space is homothetic to every other one, and the scalar transport matrices $\chi(g)$ are unchanged by that homothety;
\item the all-Sylow carrier canonically removes the residual-killing and Sylow choices, while Theorem~\ref{thm:presentation-linear-comparison-contractible} and the native presentation-transition theorems remove the presentation choice up to coherent equivalence;
\item the Cayley--Hamilton transport layer reconstructs the full character $\chi$, hence in particular its finite torsion part and every cyclotomic or Lubin--Tate logarithmic depth invariant defined from its restriction to an open pro-$p$ subgroup;
\item the Lubin--Tate level readout is the connection-depth invariant of Proposition~\ref{prop:LT-depth}, and changing the finite cohomological gauge does not change it because the transport character itself is retained.
\end{enumerate}
Thus the former rank-one presentation-independence test is closed; what is noncanonical is only a chosen strict matrix gauge, not the rank-one finite native object.
\end{theorem}

\begin{proof}
Over the discrete valuation ring $\mathcal O_E$, two full lattices in a one-dimensional $E$-space differ by multiplication by a power of a uniformizer and a unit, so conjugation does not change the scalar matrices $\chi(g)$.  The all-Sylow independence is Theorems~\ref{thm:all-sylow-transfer} and \ref{thm:projector-free-realization-atlas}; presentation coherence is Theorem~\ref{thm:presentation-linear-comparison-contractible} together with Theorems~\ref{thm:certified-native-ordered-transition}--\ref{thm:native-ordered-third-coherence}.  Full character reconstruction is the rank-one case of Theorem~\ref{thm:reconstruction-factorization}.  Since torsion and logarithmic depth are functions of the reconstructed character, they are invariant under all finite-chart changes; Proposition~\ref{prop:LT-depth} gives the stated Lubin--Tate identification.
\end{proof}

\begin{theorem}[Trianguline finite-atlas reconstruction]\label{thm:trianguline-finite-atlas-reconstruction}
Let
\[
 0\longrightarrow D_1\longrightarrow D\longrightarrow D_2\longrightarrow0
\]
be a representation-induced trianguline extension on a residual chart, with notation as in Equations~\eqref{eq:phi-matrix}--\eqref{eq:herr-normalization}.  Then the finite native atlas reconstructs the extension class and its gauge action as follows.
\begin{enumerate}[label=\textnormal{(\roman*)}]
\item the normalized off-diagonal pair
\[
 (y,z)=\left(\frac{a_\varphi}{\alpha_2}f,
              \frac{a_\gamma}{\beta_2}f\right)
\]
is a Herr $1$-cocycle exactly when the semilinear extension matrices satisfy Equation~\eqref{eq:extension-compatibility};
\item changing the splitting $e_2\mapsto e_2+be_1$ changes $(y,z)$ by the Herr coboundary of $x=bf$, so the class in $H^1_{\varphi,\gamma}(D_1D_2^{\vee})$ is precisely the nonsplit extension datum modulo gauge;
\item Theorem~\ref{thm:raw-herr-all-degree-perfectification} and Corollary~\ref{cor:raw-herr-all-degree-fox} send the raw Herr representative directly to the finite Fox/native cochain core, while the Cayley--Hamilton layer reconstructs the full upper-triangular profinite transport and therefore does not forget the nonsplit extension;
\item changes of cyclotomic generator, Fox presentation, Sylow chart, splitting module, and finite contraction change only the strict gauge representative and preserve the resulting extension object up to coherent equivalence.
\end{enumerate}
Hence the degree-one contact equation and the derived extension reconstruction required by the former trianguline test are already theorem-level consequences of the finite atlas.
\end{theorem}

\begin{proof}
The calculation preceding Proposition~\ref{prop:herr-spencer} shows directly that the Herr cocycle equation is Equation~\eqref{eq:extension-compatibility} and that Equations~\eqref{eq:gauge-phi}--\eqref{eq:gauge-gamma} are exactly the Herr coboundary action.  Herr cohomology therefore classifies the extension modulo splitting gauge.  The raw-to-finite chain comparison is Corollary~\ref{cor:raw-herr-all-degree-fox}; faithful retention of the nonsplit transport is Theorem~\ref{thm:reconstruction-factorization}.  Generator and presentation independence are Theorem~\ref{thm:raw-herr-tate-extraction-core}, Theorem~\ref{thm:presentation-linear-comparison-contractible}, and the coherent atlas of Theorem~\ref{thm:recursive-native-globalization}.
\end{proof}

\begin{corollary}[Herr is not a recursive-globalization obstruction]\label{cor:recursive-herr-comparison}
Let $U=\operatorname{Sp}A$ be an affinoid residual chart on which the Herr perfect object and the all-Sylow/Fox coefficient object are represented by bounded complexes of finite projective $A$-modules, and refine $U$ by the common Fitting/minor stratification so that their cohomology modules are finite locally free.  Then, after choosing the standard cohomology identifications with $H^*(G_K,V)$, the Herr--Fox chain transition is recursively explicit: it may be written as
\[
 \Theta_{\mathrm{FH}}=i_{\mathrm{Fox}}\,
 \Psi_{\mathrm{FH}}\,p_{\mathrm{Herr}},
\]
with $i_{\mathrm{Fox}}$ and $p_{\mathrm{Herr}}$ obtained from inverse minors.  Its entries are rational in the finite differential matrices and the selected minors; a change of contractions changes the transition by the explicit finite chain homotopy of Theorem~\ref{thm:common-source-chart-change}.  The multiplicative, $A_\infty/L_\infty$, and lax monoidal refinements at every prescribed finite arity are the corresponding finite common-source transfer formulas.

Consequently Herr contributes no additional existence, descent, or higher-coherence obstruction to Theorem~\ref{thm:recursive-native-globalization}.  Moreover, Theorem~\ref{thm:raw-herr-tate-extraction-core} closes direct extraction of arbitrary raw cyclotomic $H^1$ classes, Corollary~\ref{cor:raw-herr-direct-fox-core} sends those classes to finite Fox coordinates, and Theorem~\ref{thm:raw-herr-all-degree-perfectification} with Corollary~\ref{cor:raw-herr-all-degree-fox} closes the all-degree finite-output chain comparison.  Theorem~\ref{thm:raw-herr-phantom-reduced-normalization} identifies the linear term of the existing multiplicative equivalence with the displayed Tate-dual native class.  Literal equality on the unreduced Robba representative is only the optional gauge of Remark~\ref{rem:raw-herr-exact-gauge-boundary}.
\end{corollary}

\begin{proof}
Perfectness and affinoid strict representatives are Theorem~\ref{thm:kpx-coefficient-perfectness} and Proposition~\ref{prop:abstract-coeff-atlas}.  After Fitting refinement, Proposition~\ref{prop:determinantal-contraction} gives contractions by inverse minors.  The displayed formula and its choice homotopy are Theorem~\ref{thm:common-source-chart-change}; Theorems~\ref{thm:common-source-multiplicative-tower} and \ref{thm:common-source-lax-monoidal-tower} give the higher and monoidal refinements.  The intrinsic derived identification is Theorem~\ref{thm:fox-herr-zigzag}.
\end{proof}

\begin{remark}[Raw-Herr forward higher-comparison status]\label{rem:herr-finite-atlas-status}
The linear all-degree raw Herr--to--finite-native comparison is closed by Theorem~\ref{thm:raw-herr-all-degree-perfectification} and Corollary~\ref{cor:raw-herr-all-degree-fox}.  Representative-level multiplicative comparison is closed by Theorem~\ref{thm:raw-herr-abstract-ainfty-existence}, and finite-native normalization is automatic by Theorem~\ref{thm:raw-herr-phantom-reduced-normalization}.  Passing through the finite Tate/Fox core and homotopy transfer supplies the recursively explicit higher operations.  Thus no separate Herr finite-atlas problem remains; Remarks~\ref{rem:raw-herr-exact-gauge-boundary} and \ref{rem:optional-exact-raw-herr-normalization} record only the optional literal raw-period gauge.
\end{remark}

\section{Conclusion}

A continuous $p$-adic representation is naturally a flat profinite transport system.  Its continuous cochains form the intrinsic arithmetic Spencer complex, its deformation theory is Maurer--Cartan, Herr complexes provide finite cyclotomic cohomological gauges, completed Fox calculus gives explicit transport resolutions in free and Demu\v{s}kin pro-$p$ sectors, and Cayley--Hamilton algebras retain the full transport information that cohomology alone forgets.

The finite-atlas problem is therefore not a new classification of representations by period invariants.  At the derived/Morita level it is closed by Theorem~\ref{thm:derived-faithful-finite-closure}; at the stronger effective level, Theorem~\ref{thm:recursive-native-globalization} gives a finite stratified native atlas whose charts, transitions, descent retracts, monoidal maps, and every prescribed finite coherence level are obtained recursively from finite transport words, minors, matrix equations, and finite homotopy-transfer formulas.  Rank jumps are absorbed by the unreduced perfect complexes and derived specialization, while non-rigid flags are handled by finite-stage Schur--Koszul descent.

No stronger direct-gauge problem remains at the level claimed by the paper.  Theorem~\ref{thm:raw-herr-abstract-ainfty-existence} closes representative-level multiplicative existence, and Theorem~\ref{thm:raw-herr-phantom-reduced-normalization} identifies its linear term with the direct Tate-dual native class at exactly the finite level retained by the atlas.  Literal equality on one unreduced Robba representative is only the optional strict gauge recorded in Remarks~\ref{rem:raw-herr-exact-gauge-boundary} and \ref{rem:optional-exact-raw-herr-normalization}.  Linear presentation coherence is closed by Theorem~\ref{thm:presentation-linear-comparison-contractible}, while the required multiplicative/simplicial higher coherence is supplied by the common-source Cech--Stasheff atlas; Theorem~\ref{thm:relative-decomposable-ainfty-lifting} concerns only the stronger direct ordered-word refinement.  Non-rigid descent is algorithmically finite chartwise, while Theorem~\ref{thm:no-uniform-descent-exponent-bound} rules out a universal combinatorial exponent bound under arbitrary nonreduced base change.  Wildness explains why none of these structures should be replaced by a universal normal form; it does not obstruct the stratified Morita/homotopy-coherent globalization already constructed.

\appendix

\section{Trianguline matrix verification}

We give the computation leading to \eqref{eq:extension-compatibility}.  From \eqref{eq:gamma-matrix},
\[
\gamma(e_2)=a_\gamma e_1+\beta_2e_2.
\]
Apply $\varphi$ semilinearly:
\begin{align*}
\varphi\gamma(e_2)
&=\varphi(a_\gamma)\varphi(e_1)+\varphi(\beta_2)\varphi(e_2)\\
&=\bigl(\alpha_1\varphi(a_\gamma)+\varphi(\beta_2)a_\varphi\bigr)e_1
+\varphi(\beta_2)\alpha_2e_2.
\end{align*}
Similarly,
\begin{align*}
\gamma\varphi(e_2)
&=\gamma(a_\varphi)\gamma(e_1)+\gamma(\alpha_2)\gamma(e_2)\\
&=\bigl(\beta_1\gamma(a_\varphi)+\gamma(\alpha_2)a_\gamma\bigr)e_1
+\gamma(\alpha_2)\beta_2e_2.
\end{align*}
Equality of the $e_1$ coefficients gives \eqref{eq:extension-compatibility}; equality of the $e_2$ coefficients is the diagonal rank-one compatibility.

For the gauge change $e_2'=e_2+be_1$,
\begin{align*}
\varphi(e_2')
&=a_\varphi e_1+\alpha_2e_2+\varphi(b)\alpha_1e_1\\
&=\bigl(a_\varphi+\alpha_1\varphi(b)-\alpha_2b\bigr)e_1+\alpha_2e_2',
\end{align*}
which proves \eqref{eq:gauge-phi}; the $\gamma$ formula is identical.

\section{The nonlinear cocycle calculation}

Let $\rho_A(g)=(I+a(g))\rho(g)$.  Then
\begin{align*}
\rho_A(g)\rho_A(h)
&=(I+a(g))\rho(g)(I+a(h))\rho(h)\\
&=(I+a(g))(I+g\cdot a(h))\rho(gh)\\
&=\bigl(I+a(g)+g\cdot a(h)+a(g)g\cdot a(h)\bigr)\rho(gh).
\end{align*}
Equating with $(I+a(gh))\rho(gh)$ gives \eqref{eq:nonlinear-cocycle}.  This calculation is the exact source of the Maurer--Cartan quadratic term.

\begin{thebibliography}{99}\bibitem{JannsenWingberg}
U.~Jannsen and K.~Wingberg,
\emph{Die Struktur der absoluten Galoisgruppe $p$-adischer Zahlk\"orper},
Invent. Math. \textbf{70} (1982), 71--98.

\bibitem{JardenShusterman}
M.~Jarden and M.~Shusterman,
\emph{The absolute Galois group of a $p$-adic field},
Math. Ann. \textbf{386} (2023), 1595--1603.

\bibitem{NeukirchLocal}
J.~Neukirch,
\emph{The absolute Galois group of a $p$-adic number field},
Ast\'erisque \textbf{94} (1982), 153--164.

\bibitem{Brumer}
A.~Brumer,
\emph{Pseudocompact algebras, profinite groups and class formations},
J. Algebra \textbf{4} (1966), 442--470.

\bibitem{NSW}
J.~Neukirch, A.~Schmidt, and K.~Wingberg,
\emph{Cohomology of Number Fields}, 2nd ed.,
Grundlehren Math. Wiss. \textbf{323}, Springer, 2008.

\bibitem{Fox}
R.~H.~Fox,
\emph{Free differential calculus. I. Derivation in the free group ring},
Ann. of Math. (2) \textbf{57} (1953), 547--560.

\bibitem{FontaineLocal}
J.-M.~Fontaine,
\emph{Repr\'esentations $p$-adiques des corps locaux. I},
in The Grothendieck Festschrift, Vol. II, Progr. Math. \textbf{87},
Birkh\"auser Boston, 1990, 249--309.

\bibitem{CherbonnierColmez}
F.~Cherbonnier and P.~Colmez,
\emph{Repr\'esentations $p$-adiques surconvergentes},
Invent. Math. \textbf{133} (1998), 581--611.

\bibitem{ColmezTrianguline}
P.~Colmez,
\emph{Repr\'esentations triangulines de dimension $2$},
Ast\'erisque \textbf{319} (2008), 213--258.


\bibitem{BhattScholzePFC}
B.~Bhatt and P.~Scholze,
\emph{Prismatic $F$-crystals and crystalline Galois representations},
Camb. J. Math. \textbf{11} (2023), 507--562; arXiv:2106.14735.

\bibitem{DuLiuPrismatic}
H.~Du and T.~Liu,
\emph{A prismatic approach to $(\varphi,\widehat G)$-modules and $F$-crystals},
arXiv:2107.12240.

\bibitem{HerrSurvey}
L.~Herr,
\emph{$\Phi$-$\Gamma$-modules and Galois cohomology},
in Invitation to Higher Local Fields, Part II, 2000; arXiv:math/0012156.

\bibitem{LiuCohomology}
R.~Liu,
\emph{Cohomology and duality for $(\varphi,\Gamma)$-modules over the Robba ring},
Int. Math. Res. Not. (2008); arXiv:0711.4346.

\bibitem{BergerTrianguline}
L.~Berger,
\emph{Trianguline representations},
Bull. Lond. Math. Soc. \textbf{43} (2011), 619--635; arXiv:1011.3447.

\bibitem{KPXFamilies}
K.~Kedlaya, J.~Pottharst, and L.~Xiao,
\emph{Cohomology of arithmetic families of $(\varphi,\Gamma)$-modules},
J. Amer. Math. Soc. \textbf{27} (2014), 1043--1115; arXiv:1203.5718.

\bibitem{AntonioPadic}
J.~Ant\'onio,
\emph{Moduli of $p$-adic representations of a profinite group},
arXiv:1709.04275.

\bibitem{AntonioProEtale}
J.~Ant\'onio,
\emph{Moduli of $\ell$-adic pro-\'etale local systems for smooth non-proper schemes},
arXiv:1904.08001.

\bibitem{BIPLocal}
G.~B\"ockle, A.~Iyengar, and V.~Pa\v{s}k\=unas,
\emph{On local Galois deformation rings},
Forum Math. Pi \textbf{11} (2023), Paper No. e30, 54 pp.; arXiv:2110.01638.

\bibitem{BIPCorrigendum}
G.~B\"ockle, A.~Iyengar, and V.~Pa\v{s}k\=unas,
\emph{On local Galois deformation rings -- Corrigendum},
Forum Math. Pi \textbf{12} (2024), Paper No. e5; DOI: 10.1017/fmp.2024.3.

\bibitem{MinDerivedEG}
Y.~Min,
\emph{Classicality of derived Emerton--Gee stack},
Math. Ann. \textbf{393} (2025), no.~1, 439--494; DOI: 10.1007/s00208-025-03259-7.

\bibitem{GalatiusVenkatesh}
S.~Galatius and A.~Venkatesh,
\emph{Derived Galois deformation rings},
Adv. Math. \textbf{327} (2018), 470--623; arXiv:1608.07236.

\bibitem{PalQuickOdd}
A.~P\'al and G.~Quick,
\emph{$A_3$-formality for Demu\v{s}kin groups at odd primes},
arXiv:2601.07551.

\bibitem{PalQuickTwo}
A.~P\'al and G.~Quick,
\emph{$A_3$-formality for pro-$2$ Demu\v{s}kin groups},
arXiv:2607.01028.

\bibitem{WangErickson}
C.~Wang-Erickson,
\emph{Deformations of residually reducible Galois representations via $A_\infty$-algebra structure on Galois cohomology},
J. Reine Angew. Math. \textbf{769} (2020), 1--55; arXiv:1809.02484.

\bibitem{ChenevierDeterminants}
G.~Chenevier,
\emph{The $p$-adic analytic space of pseudocharacters of a profinite group and pseudorepresentations over arbitrary rings},
in \emph{Automorphic Forms and Galois Representations}, Vol.~1,
London Math. Soc. Lecture Note Ser. \textbf{414}, Cambridge Univ. Press, 2014, 221--285; arXiv:0809.0415.

\bibitem{WangEricksonFamilies}
C.~Wang-Erickson,
\emph{Algebraic families of Galois representations and potentially semi-stable pseudodeformation rings},
Math. Ann. \textbf{371} (2018), 1615--1681; arXiv:1501.05629.

\bibitem{ToenAzumaya}
B.~To\"en,
\emph{Derived Azumaya algebras and generators for twisted derived categories},
Invent. Math. \textbf{189} (2012), 581--652; arXiv:1002.2599.

\bibitem{HalpernLeistnerPreygel}
D.~Halpern--Leistner and A.~Preygel,
\emph{Mapping stacks and categorical notions of properness},
Compos. Math. \textbf{159} (2023), no.~3, 530--589; arXiv:1402.3204.


\bibitem{BhattScholzeWitt}
B.~Bhatt and P.~Scholze,
\emph{Projectivity of the Witt vector affine Grassmannian},
Invent. Math. \textbf{209} (2017), 329--423; arXiv:1507.06490.

\bibitem{MathewGalois}
A.~Mathew,
\emph{The Galois group of a stable homotopy theory},
Adv. Math. \textbf{291} (2016), 403--541; arXiv:1404.2156.

\bibitem{EfimovGrassmannians}
A.~I.~Efimov,
\emph{Derived categories of Grassmannians over integers and modular representation theory},
Adv. Math. \textbf{304} (2017), 179--226; arXiv:1410.7462.

\bibitem{BuchweitzLeuschkeVdBGrass}
R.-O.~Buchweitz, G.~J.~Leuschke, and M.~Van den Bergh,
\emph{On the derived category of Grassmannians in arbitrary characteristic},
Compos. Math. \textbf{151} (2015), no.~7, 1242--1264; arXiv:1006.1633.

\bibitem{BKRDerived}
Y.~Berest, G.~Khachatryan, and A.~Ramadoss,
\emph{A simple construction of derived representation schemes},
J. Noncommut. Geom. \textbf{6} (2012), 1--20; arXiv:1010.4901.

\bibitem{BKRCyclic}
Y.~Berest, G.~Khachatryan, and A.~Ramadoss,
\emph{Derived representation schemes and cyclic homology},
Adv. Math. \textbf{245} (2013), 625--689; arXiv:1112.1449.

\bibitem{BrownRazakCrossed}
R.~Brown and A.~Razak Salleh,
\emph{Free crossed resolutions of groups and presentations of modules of identities among relations},
LMS J. Comput. Math. \textbf{2} (1999), 28--61; arXiv:math/9812132.

\bibitem{HuebschmannCrossed}
J.~Huebschmann,
\emph{Crossed modules},
arXiv:2403.15900.

\bibitem{BrownCohomology}
K.~S. Brown,
\emph{Cohomology of Groups},
Graduate Texts in Mathematics 87, Springer, New York, 1982.

\bibitem{BurkeTransfer}
J.~Burke,
\emph{Transfer of $A_\infty$-structures to projective resolutions},
arXiv:1801.08933.

\bibitem{ManettiPerturbation}
M.~Manetti,
\emph{A relative version of the ordinary perturbation lemma},
Rend. Mat. Appl. (7) \textbf{30} (2010), 221--238; arXiv:1002.0683.

\bibitem{VokrinekHPT}
L.~Vok\v{r}inek,
\emph{Enriched categorical aspects of homological perturbation theory},
arXiv:2412.21182.

\bibitem{RoeTurtureanGQ2}
D.~Roe and D.~Turturean,
\emph{A presentation of the absolute Galois group of $\mathbf Q_2$},
online manuscript (2026),
\url{https://roed314.github.io/gq2/paper/paper.html}.

\bibitem{LabuteClassification}
J.~P. Labute,
\emph{Classification of Demushkin groups},
Canadian J. Math. \textbf{19} (1967), 106--132.

\bibitem{GartnerMassey}
J.~G\"artner,
\emph{Higher Massey products in the cohomology of mild pro-$p$-groups},
J. Algebra \textbf{422} (2015), 788--820; arXiv:1212.2118.

\bibitem{SuciuWangCup}
A.~I. Suciu and H.~Wang,
\emph{Cup products, lower central series, and holonomy Lie algebras},
J. Pure Appl. Algebra \textbf{223} (2019), no.~8, 3359--3385.

\bibitem{CookFP}
G.~Corob Cook,
\emph{On profinite groups of type $\operatorname{FP}_\infty$},
arXiv:1412.1876.

\bibitem{StacksPerfect}
The Stacks Project Authors,
\emph{Perfect complexes},
Stacks Project, Tag 0656, especially Lemma 0658,
\url{https://stacks.math.columbia.edu/tag/0656}.

\bibitem{StacksLCI}
The Stacks Project Authors,
\emph{Local complete intersection rings},
Tag 09PY,
\url{https://stacks.math.columbia.edu/tag/09PY}.

\bibitem{StacksKoszul}
The Stacks Project Authors,
\emph{Koszul regular sequences},
Tag 062D,
\url{https://stacks.math.columbia.edu/tag/062D}.

\bibitem{StacksFitting}
The Stacks Project Authors,
\emph{Fitting ideals and flattening stratifications},
Tags 0C3C and 052F,
\url{https://stacks.math.columbia.edu/tag/0C3C}.


\bibitem{StacksEtale}
The Stacks Project Authors,
\emph{Finite locally free and etale morphisms},
Tags 02KB, 025J, and 0476,
\url{https://stacks.math.columbia.edu/tag/025J}.

\bibitem{BelitskiiSergeichuk}
G.~R. Belitskii and V.~V. Sergeichuk,
\emph{Complexity of matrix problems},
Linear Algebra Appl. \textbf{361} (2003), 203--222; arXiv:0709.2488.


\end{thebibliography}

\end{document}